%------------------------------------------------------------------------------%
% FUNCTIONAL ANALYSIS                                                          %
%------------------------------------------------------------------------------%
% File  : Functional_Analysis.tex                                              %
% Author: Klaus Hoeflinger, Beethovenstrasse 11, D-73117 Wangen                %
% Email : khoeflinger@t-online.de                                              %
%------------------------------------------------------------------------------%

%------------------------------------------------------------------------------%
%                       DOCUMENT CLASS and INCLUDE FILES                       %
%------------------------------------------------------------------------------%
\documentclass[11pt,twoside]{article}
\usepackage{geometry}
\geometry{paperheight=241mm,
          paperwidth=152mm,
          top=22mm,
          bottom=17mm,
          left=20mm,
          right=20mm,
          textwidth=112mm,
          textheight=202mm,
          headheight=10mm,
          headsep=3mm,
          footskip=9mm,
          includehead,
          includefoot,
          marginparsep=0mm,
          marginparwidth=0mm}

\renewcommand{\baselinestretch}{0.945}\normalsize

\AtBeginDocument{\setlength\abovedisplayskip{2mm}}
\AtBeginDocument{\setlength\belowdisplayskip{2mm}}

%------------------------------------------------------------------------------%
% define title formats                                                         %
%------------------------------------------------------------------------------%
\usepackage{titlesec}

\renewcommand{\thesection}{\arabic{section}}
\renewcommand{\sfdefault}{FiraSans}
\renewcommand{\ttdefault}{pcr}
\renewcommand{\rmdefault}{antiqua}

\titleformat{\part}[display]
{\normalfont \large \rmfamily \bfseries \centering}
{\small{\thepart.}}{2pt}{}

\titleformat{\section}[runin]
{\normalfont \rmfamily \bfseries}
{\thesection}{1pt}{.\;}[.]

%\titleformat{\section}
%{\normalfont \bfseries \small \filcenter}
%{\thesection.\;}{0mm}{}

\titleformat{\subsection}[runin]
{\normalfont \itshape}
{}{0mm}{}

\titleformat{\subsubsection}[runin]
{\normalfont \itshape}
{}{}{}{}

\titleformat{\paragraph}[runin]
{\normalfont}
{}{}{}{}

\titleformat{\subparagraph}[runin]
{\normalfont}
{}{}{}{}

\titlespacing*{\part}          { 0pt}{ 0pt}{ 8pt}
\titlespacing*{\section}       { 0pt}{ 7pt}{ 4pt}
\titlespacing*{\subsection}    { 0pt}{14pt}{ 6pt}
\titlespacing*{\subsubsection} { 0pt}{ 6pt}{12pt}
\titlespacing*{\paragraph}     { 0pt}{ 6pt}{12pt}
\titlespacing*{\subparagraph}  { 0pt}{ 0pt}{12pt}

%------------------------------------------------------------------------------%
% define enumerate parameters                                                  %
%------------------------------------------------------------------------------%
\usepackage{enumitem}
\setlength{\parindent}{3pt}
\setlength{\parskip}  {0pt}

%\setlist{noitemsep}

\setlist[enumerate]{
         leftmargin=12pt,
         rightmargin=0pt,
         itemsep=1pt,
         itemindent=3pt,
         topsep=3pt,
         labelsep=4pt,
         partopsep=0pt,
         labelwidth=0pt,
         listparindent=0pt,
         parsep=0pt}

\setlist[itemize]{
         leftmargin=12pt,
         rightmargin=0pt,
         itemsep=0pt,
         itemindent=1pt,
         topsep=1pt,
         labelsep=4pt,
         partopsep=1pt,
         labelwidth=0pt,
         listparindent=0pt,
         parsep=0pt}

%------------------------------------------------------------------------------%
% define packages                                                              %
%------------------------------------------------------------------------------%
\usepackage{upref}
\usepackage{amssymb}
\usepackage{slantsc}
\usepackage{amsmath}
\usepackage{amsthm}
\usepackage{thmtools}
\usepackage{amscd}
\usepackage{verbatim}
\usepackage{latexsym}
\usepackage{multicol}
\usepackage{graphicx}
\usepackage{setspace}
\usepackage[ngerman,english]{babel}
\usepackage[protrusion=true,expansion=true]{microtype}
\usepackage{tikz}
\usetikzlibrary{arrows,arrows.meta,decorations.markings}
\usepackage{mathrsfs}
\usepackage{amsfonts}
\usepackage{datetime2}
\usepackage{textgreek}
\usepackage{hyperref}
\usepackage{pifont}
\usepackage{relsize}
%\usepackage{biblatex}
\relsize{-0.05}
\usepackage[skip=0mm]{parskip}
\usetikzlibrary{decorations.pathmorphing}

\usepackage{mlmodern}
\usepackage[T5,T1]{fontenc}
\usepackage{ragged2e}

%\usepackage{mathabx}
%\usepackage{times}

%\usepackage{fontspec}
%\setmainfont{Nimbus Roman No9 L}
%\usepackage[T1]{fontenc}

%------------------------------------------------------------------------------%
% define page layout                                                           %
%------------------------------------------------------------------------------%
\usepackage{fancyhdr}
\pagestyle{fancy}
\setlength{\headheight}{30.0pt}
\renewcommand{\headrulewidth}{0pt}

\AtBeginDocument{
  \addtolength\abovedisplayskip{-0.0\baselineskip}
  \addtolength\belowdisplayskip{-0.0\baselineskip}
}

\fancyhead{}
\fancyfoot{}
\addtolength{\headwidth}{0.02\marginparwidth}

\makeatletter
\let\ps@plain\ps@fancy
\makeatother

\renewcommand\sectionmark[1]
{
 \def\sectionname{#1}
 \markright{\thesection #1}
}

\makeatletter
\@addtoreset{section}{part}
\makeatother

%------------------------------------------------------------------------------%
% define footnote without number                                               %
%------------------------------------------------------------------------------%
\newcommand{\myfootnote}[1]{%
\renewcommand{\thefootnote}{}%
\footnotetext{#1}%
\renewcommand{\thefootnote}{\arabic{footnote}}%
}

%------------------------------------------------------------------------------%
% define numbering in equations                                                %
%------------------------------------------------------------------------------%
\numberwithin{equation}{section}

%------------------------------------------------------------------------------%
% define newtheorem                                                            %
%------------------------------------------------------------------------------%
\declaretheoremstyle[
headfont=\scshape, 
headindent = 0mm, %\parindent,
%postheadhook = {\hspace{0mm}\newline},
%postheadspace = \newline, 
spaceabove = 3mm, 
spacebelow = 3mm,
numberwithin=section,
name=Theorem]{khMATH}
\declaretheorem[style=khMATH]{theorem}

\declaretheoremstyle[
headfont=\bfseries, 
headindent = 0mm, %\parindent,
%postheadhook = {\hspace{0mm}\newline},
%postheadspace = \newline, 
spaceabove = 3mm, 
spacebelow = 3mm,
numberwithin=part,
name=Theorem]{khMATH1}
\declaretheorem[style=khMATH1]{theorem1}

\declaretheoremstyle[
headfont=\scshape, 
headindent = 0mm, %\parindent,
%postheadhook = {\hspace{0mm}\newline},
%postheadspace = \newline, 
spaceabove = 3mm, 
spacebelow = 3mm,
numberwithin=section,
name=Corollary]{khMATH1}
\declaretheorem[style=khMATH1]{corollary}

%            name         counter     label              numeration-mode
\theoremstyle{plain}
\newtheorem*{introduction}            {\sc Introduction}
\newtheorem {standard}                {}                 [section]
\newtheorem*{definition}              {\sc Definition}
\newtheorem{satz1}                    {\sc Satz}
\newtheorem{lemma}                    {\sc Lemma}
\newtheorem*{proposition}             {\sc Proposition}
%\newtheorem{corollary}               {\sc Corollary}
\newtheorem*{corollar}                {\sc Korollar}
\newtheorem*{remark}                  {\sc Remark}
\newtheorem*{remarks}                 {\sc Remarks}
\newtheorem*{example}                 {Example}
\newtheorem*{examples}                {Examples}
\newtheorem*{theo}                    {\sc Theorem}

%------------------------------------------------------------------------------%
% define tabular                                                               %
%------------------------------------------------------------------------------%
\usepackage{tabularx}
\newcolumntype{L}[1]{>{\raggedright\arraybackslash}p{#1}}
\newcolumntype{C}[1]{>{\centering\arraybackslash}  p{#1}} 
\newcolumntype{R}[1]{>{\raggedleft\arraybackslash} p{#1}} 

%------------------------------------------------------------------------------%
% define \sh, \zB                                                              %
%------------------------------------------------------------------------------%
\def\dsh{{d.\,h.}}
\def\ie{{i.\,e.}}
\def\zB{z.\,B.}
\renewcommand{\abstractname}{ }

%------------------------------------------------------------------------------%
% change default spacing for cases                                             %
%------------------------------------------------------------------------------%
\makeatletter
\renewcommand*\env@cases[1][1.6]{%
  \let\@ifnextchar\new@ifnextchar
  \left\lbrace
  \def\arraystretch{#1}%
  \array{@{}l@{\quad}l@{}}%
}
\makeatother

\newenvironment{rcases}
  {\left.\begin{aligned}}
  {\end{aligned}\right\rbrace}

%------------------------------------------------------------------------------%
% change default for \predisplaypenalty                                        %
%------------------------------------------------------------------------------%
\predisplaypenalty=990
\allowdisplaybreaks \numberwithin{equation}{section}

%------------------------------------------------------------------------------%
% general notations                                                            %
%------------------------------------------------------------------------------%
\def\*{\star}
\def\={\,=\,}
\def\+{\,+\,}
\def\xsubset{{\,\subset\,}}
\def\xtimes{{\,\times\,}}
\def\xeof{{\,\in\,}}
\def\eq{{\,=\,}}
\def\xto{{\,\to\,}}
\def\eqdef{\,:=\,}
\def\:{{\,:\,}}
\def\ele{{\,\in\,}}
\def\exists{\mbox{ exists }}
\def\and{\mbox{ and }}
\def\for{\mbox{ for }}
\def\fa{\mbox{ for all }}
\def\fe{\mbox{ for each }}
\def\qfa{\quad \mbox{ for all }}
\def\df{\,\dotfill}
\def\iff{if and only if }
\def\wrt{with respect to}

\makeatletter
\newcommand \Df {\leavevmode \cleaders \hb@xt@ .70em{\hss .\hss }\hfill \kern \z@}
\makeatother

\def\sql{\mathfrak{a}}               % sesquilinear form a
\def\DF{B}                           % Dirichlet form a
%------------------------------------------------------------------------------%
% number sets and special elements                                             %
%------------------------------------------------------------------------------%
\def\P{\mathbb{H}}                   % half plane of $\C$
\def\J{I}                            % interval I
\def\N{{\mathbb{N}}}                 % unsigned integers
\def\Q{{\mathbb{Q}}}                 % rational integers
\def\Z{{\mathbb{Z}}}                 % integers
\def\R{{\mathbb{R}}}                 % real numbers
\def\Rn{{\mathbb{R}}^n}              % real numbers in Rn
\def\Rd{{\mathbb{R}}^d}              % real numbers in Rn
\def\C{{\mathbb{C}}}                 % complex numbers
\def\F{{\mathbb{K}}}                 % arbitrary field
\def\Torus{{\mathbb{T}}}             % complex circle line
\def\T{\tau}                         % upper bound of interval (0,t)
\def\sign#1{\operatorname{sign}(#1)} % sign of a number

%------------------------------------------------------------------------------%
% set theory                                                                   %
%------------------------------------------------------------------------------%
\def\G{G}                            % domain $G$
\def\clG{\overline{\G}}              % closure of $G$
\def\bndG{\partial \G}               % boundary of $G$
\def\setA{\mathcal{A}}               % set A
\def\setB{\mathcal{B}}               % set B
\def\setC{\mathcal{C}}               % set C
\def\setM{M}                         % set M
\def\setN{N}                         % set N
\def\setS{\mathcal{S}}               % set S
\def\famF{\mathfrak{F}}              % family of sets
\def\indI{\mathcal{I}}               % index set I
\def\pot#1{\mathfrak{P}(#1)}         % power set

\def\relR{\mathbf{R}}                % relation R
\def\mapf{f}                         % mapping f
\def\mapg{g}                         % mapping g

\def\set#1#2{\{ \, #1 \,: \, #2 \, \}}
\def\tensor#1#2{#1 \otimes #2}
\def\Tens{\oplus}

\def\cal#1{{\mathcal{#1}}}
\def\aut#1{\operatorname{Aut}(#1)}

\def\cross#1#2{#1 {\,\times\,} #2}
\def\multicross#1#2{#1 \times  \ldots \times  #2}
\def\multitens#1#2 {#1 \otimes \ldots \otimes #2}

%------------------------------------------------------------------------------%
% group theory                                                                 %
%------------------------------------------------------------------------------%
\def\1{{\mathsf{1}}}                % unit element of a monoid
\def\0{{\mathsf{0}}}                % unit element of an Abelian monoid
\def\kernel#1{\mathfrak{N}(#1)}     % kernel of a homomorphism
\def\cokernel#1{\mathfrak{C}{(#1)}} % cokernel of a homomorphism
\def\Hom#1{\operatorname{Hom}(#1)}

%------------------------------------------------------------------------------%
% linear algebra                                                               %
%------------------------------------------------------------------------------%
\def\vecV{\mathit{V}}               % vector space V
\def\vecH{\mathit{H}}               % vector space H
\def\vecW{\mathit{W}}               % vector space W
\def\vecU{\mathit{U}}               % subspace U
\def\convC{\mathcal{C}}             % convex set C
\def\basH{\mathfrak{B}}             % Hamel base B

\def\LO#1{\mathscr{L}(#1) }         % algebra of linear transformations

\def\scalar#1#2{ \langle #1,#2 \rangle }
\def\trace#1{\operatorname{tr}(#1)}
\def\dim#1{{\operatorname{dim} \, #1 }}
\def\codim#1{{\operatorname{codim} \, #1 }}
\def\dim{\operatorname{dim}}
\def\codim{\operatorname{codim}}
\def\ind{\operatorname{ind}}

\def\genG{\cal{G}}

\def\algA{\cal{A}}
\def\algB{\cal{B}}
\def\algU{\cal{U}}

\def\idI{\cal{I}}

%------------------------------------------------------------------------------%
% topology                                                                     %
%------------------------------------------------------------------------------%
\def\compK{{K}}          % compact set C
\def\openO{\mathcal{O}}  % open set O
\def\openU{\mathcal{U}}  % open set U
\def\U{\mathcal{U}}      % neighbourhood U
\def\V{\mathcal{V}}      % neighbourhood V
\def\d{\mathit{d}}       % metric function d

\def\interior#1{\mathring{#1}}
\def\closure#1{{\overline{#1}}}

\def\topT{\mathcal{T}}   % topology T
\def\topX{X}             % topological space X
\def\topY{Y}             % topological space Y
\def\topZ{Z}             % topological space Z
\def\metrS{M}            % metric space M

\def\grpG{G}             % topological group G
\def\grpA{A}             % topological group A
\def\grpB{B}             % topological group B
\def\grpC{C}             % topological group C
\def\grpH{H}             % topological group H
\def\grpU{U}             % topological group U
\def\grpV{V}             % topological group V
\def\topG{G}             % topological group G
\def\topH{H}             % topological group H
\def\topU{U}             % topological subgroup U

\def\subS{\cal{S}}
\def\riR{R}              % ring R

%------------------------------------------------------------------------------%
% calculus                                                                     %
%------------------------------------------------------------------------------%
\def\supp#1{\operatorname{supp}(#1)}
\def\singsupp#1{\operatorname{sing \, supp}(#1)}

\def\re{\operatorname{Re}}
\def\im{\operatorname{Im}}
\def\grad{\operatorname{grad}}

%\def\abs#1 {\left|#1\right|}
\def\abs#1 {\lvert #1 \rvert}
%\def\abs#1 {\left|\,#1\,\right|}

%------------------------------------------------------------------------------%
% measure theory                                                               %
%------------------------------------------------------------------------------%
\def\salgA{\Sigma}               % sigma-algebra S
\def\salgB{\cal{B}}              % sigma-algebra B
\def\Borel#1{\cal{B}_{#1}}       % Borel-algebra 
\def\LB{\lambda}                 % Lebesgue-Borel measure

%------------------------------------------------------------------------------%
% function spaces                                                              %
%------------------------------------------------------------------------------%
\def\Lp{L^p(\G)}                 % spaces L_p
\def\Lq{L^q(\G)}                 % spaces L_q
\def\Lh{L^2(\G)}                 % spaces L_2
\def\Li{L^{\infty}(\G)}          % spaces L_infty
\def\Lc{L_{loc}^1(\G)}           % spaces L_1^{loc}
\def\Ck{C^k(\G)}                 % spaces C_k

\def\BV#1{BV(#1) }               % set of functions of bounded variation
\def\AC#1{AC(#1) }               % set of absolutely continuous functions

\def\MP#1{\mathscr{M}^+(#1) }    % set of positive measures
\def\MS#1{\mathscr{M}^-(#1) }    % vector space of finite signed measures
\def\MC#1{\mathscr{M}(#1) }      % vector space of complex measures
\def\MCR#1{\mathscr{M}_{r}(#1) } % vector space of regular complex measures

\def\SobW{W^{k,p}(\G)}                % Sobolev space W
\def\SobO{\mathring{W}^{k,p}(\G)}     % Sobolev space W
%\def\WN#1#2{\mathring{W}^{#1}(#2) }   % Sobolev space W
\def\WN#1#2{W_0^{#1}(#2) }   % Sobolev space W
%\def\HN#1#2{\mathring{H}^{#1}(#2) }   % Sobolev space H
\def\HN#1#2{H_0^{#1}(#2) }   % Sobolev space H
%------------------------------------------------------------------------------%
% functional analysis                                                          %
%------------------------------------------------------------------------------%
\def\snorm#1{\lVert #1 \rVert}
\newcommand{\norm}[1]{\left\lVert #1 \right\rVert}

\def\topV{V}                       % topological vector space V
\def\topW{W}                       % topological vector space W
\def\normS{X}                      % normed space

\def\banX{\mathit{X}}              % Banach space X
\def\banY{\mathit{Y}}              % Banach space Y
\def\banZ{\mathit{Z}}              % Banach space Z
\def\dbanX{\mathit{X^*}}           % dual space of Banach space X
\def\dbanY{\mathit{Y^*}}           % dual space of Banach space Y
\def\dbanZ{\mathit{Z^*}}           % dual space of Banach space Z

\def\banA{\mathit{A}}              % Banach algebra A
\def\banB{\mathit{B}}              % Banach algebra B
\def\banC{\mathit{C}}              % C*-algebra C


\def\domain#1{\mathfrak{D}(#1)}    % domain
\def\range#1{\mathfrak{R}(#1)}     % range

\def\isoJ{\mathfrak{J}}            % isometry J
\def\repA{{\cal{A}_{\sql}}}        % representation operator A
\def\linS{{S}}                     % linear operator S
\def\linG{{G}}                     % linear operator G
\def\linL{{L}}                     % linear operator L
\def\linT{{T}}                     % linear operator A
\def\linA{\mathit{A}}              % linear operator A
\def\linB{{B}}                     % linear operator B
\def\linM{{M}}                     % linear operator M
\def\linU{{U}}                     % linear operator U
\def\linI{{I}}                     % linear operator I
\def\A{\cal{A}}                    % linear differential operator A
\def\BDO{\cal{B}}                  % linear boundary differential operator B
\def\linU{{U}}                     % unitary operator U
\def\linP{{P}}                     % projection P
\def\compA{k}                      % compact operator k
\def\linFR{f}                      % finite rank operator f
\def\fredA{A}                      % Fredholm operator F

\def\HAM{\cal{H}}                  % Hamilton operator
\def\PA{\partial^{\alpha}}         % elementary differential operator
\def\DO{\cal{A}}                   % linear differential operator
\def\BC{\cal{B}}                   % boundary operator
\def\SLO{\cal{L}_{p,q}}            % Sturm-Liouville operator
\def\MO{\cal{L}_{M}}               % momentum operator

\def\HS#1{S(#1)           }        % algebra of Hilbert-Schmidt operators
\def\FR#1{F(#1)           }        % algebra of finite rank operators
\def\TC#1{N(#1)           }        % algebra of trace class operators
\def\BU#1{\mathscr{U}(#1) }        % algebra of unitary linear operators
\def\BO#1{\mathscr{L}(#1) }        % algebra of bounded linear operators
\def\CO#1{\mathscr{C}(#1) }        % algebra of compact linear operators
\def\CL#1{\mathscr{C}(#1) }        % algebra of compact linear operators
\def\FO#1{\Phi(#1) }               % set of Fredholm operators
\def\FK#1#2{\Phi_{#2}(#1) }        % set of Fredholm operators with index k

\def\hilH{\mathit{H}}              % Hilbert space H
\def\clU{\mathit{U}}               % closed subspace U

\def\orthO{\mathfrak{O}}           % orthonormal system O
\def\orthS{\mathfrak{S}}           % orthonormal system S
\def\basB{\mathfrak{B}}            % basis B of a Hilbert space

\def\GL#1 {\mathbf{GL}(#1)}        % linear group
\def\OG#1 {\mathbf{O}(#1)}         % linear group
\def\SOG#1 {\mathbf{SO}(#1)}       % linear group

%------------------------------------------------------------------------------%
% distribution theory                                                          %
%------------------------------------------------------------------------------%
\def\disF{F}                  % distribution F
\def\disG{G}                  % distribution G
\def\disH{H}                  % distribution H
\def\disU{U}                  % distribution U
\def\disT{T}                  % distribution T
\def\SF{C^{\infty}(\G)}       % space of smooth functions
\def\TF{C_c^{\infty}(\G)}     % space of compactly supported smooth functions
\def\TFRd{C_c^{\infty}(\R^d)} % space of test functions
\def\SRd{\mathscr{S}(\R^d)}   % Schwartz space
\def\SR1{\mathscr{S}(\R)}     % Schwartz space
\def\B1#1{C^1(#1) }           % space of continuously differentiable functions
\def\Bm#1{C^m(#1) }           % space of continuously differentiable functions
\def\Ck{C^k(\G)}
\def\TD{\mathscr{S}'(\R^d)}   % space of temperated distributions
\def\E{\cal{E}}               % space of distributions with compact support
\def\D{\mathscr{D}}           % D
\def\FT{\mathcal{F}}          % Fourier transformation F
\def\FT{\mathscr{F}}          % Fourier transformation F

%------------------------------------------------------------------------------%
% other definitions                                                            %
%------------------------------------------------------------------------------%
\def\Proof{\textsc{Proof}}
\def\FRECHET{Fr\'{e}chet}
\def\ARZELA{Arzel\'{a}}
\def\GARDING{G\r{a}rding}
\def\HOELDER{H\"older}
\def\SCHROEDINGER{Schr\"odinger}
\def\JOERGENS{J\"orgens}
\def\WUEST{W\"uest}
\def\KAEHLER{K\"aehler}
\def\HOEFLINGER{H\"oflinger}
\def\POINCARE{Poincar\'{e}}
\def\FA{Functional Analysis}

\def\Author   {K.\,{\HOEFLINGER}}
\def\AuthorUC {K.\,HOEFLINGER}
\def\Version  {1.0}
\def\Email    {khoeflinger@t-online.de}


\usepackage{newunicodechar}
\newunicodechar{²}{\ifmmode{}^2\else\texttwosuperior\fi}

\renewcommand\thesection{\arabic{section}}

\def\Title {Functional Analysis}

\titleformat{\part}[display]
{\normalfont \large \rmfamily \bfseries \centering}
{\small{\thepart.}}{2pt}{}

\titleformat{\section}%[runin]
{\normalfont \bfseries \filcenter}
{\S \thesection.\;}{0mm}{}

\titlespacing*{\section}   {0pt}{16pt}{ 4pt}
\titlespacing*{\subsection}{0pt}{10pt}{12pt}

\def\G{G}

\newlength{\bibitemsep}\setlength{\bibitemsep}{.5\baselineskip plus .05\baselineskip minus .05\baselineskip}
\newlength{\bibparskip}\setlength{\bibparskip}{0pt}
\let\oldthebibliography\thebibliography
\renewcommand\thebibliography[1]{%
  \oldthebibliography{#1}%
  \setlength{\parskip}{\bibitemsep}%
  \setlength{\itemsep}{\bibparskip}%
}

%------------------------------------------------------------------------------%
%                                BEGIN OF DOCUMENT                             %
%------------------------------------------------------------------------------%
\begin{document}
\thispagestyle{empty}

\centerline{\Large{\textbf{Functional Analysis}}}

\vspace{5pt}
\centerline{\small{\textsc{\Author}}}
\vspace{20pt}

\myfootnote
{
  Version {\Version} from \today;
  Email:  \textit{khoeflinger@\,t-online.de}
}

%------------------------------------------------------------------------------%
%                                INTRODUCTION                                  %
%------------------------------------------------------------------------------%
This text consists of a series of articles
on the mathematical area of {functional analysis}
and of a final article concerning with the applications of the abstract theory
to linear partial differential equations,
which arise as models for problems
in mathematical physics and engineering science.

\subparagraph{}
Functional analysis was mainly developed
in the first half of the 20$^{\textrm{th}}$ century,
and valuable contributions have been done by
{D.\;Hilbert},
{S.\;Banach},
{F.\;Riesz},
{J.\;von Neumann},
{F.\,E.\;Browder},
{L.\;Schwartz},
{J.-L.\;Lions} and others.

\subparagraph{}
The topics within functional analysis are mainly divided into two parts:
the linear and the non-linear theory.
In our treatise we shall exclusively deal with the linear one,
hence the fundamental concept of {"linear operator"}
will play a central role in our considerations.
Especially
what historically is meant by the concept of unbounded operator
is of great importance,
since the theory of this specific class of mappings can be successfully applied
to the study of partial differential operators
and their related boundary and initial-value problems.

\paragraph{}
It is assumed that
the reader is familiar with elementary concepts in linear algebra,
general topology and measure theory.
In fact,
functional analysis can be seen as the synthesis of the both first subjects,
especially {\wrt} infinite-dimensional linear spaces.
Further preliminaries
will be introduced as soon as there is the need to use them.

\subparagraph{}
It is not the idea of this text to give all the proofs of the stated theorems,
except in the case where these are elementary.
Instead,
the main purpose is a transparent and overall presentation of the topics,
especially for readers, which have learned functional analysis in the past
and want to have a refresh of the main results now.

\paragraph{} %- honoration to Prof. Dr. Werner
I wrote the treatise in honour of Prof.\;Dr.\;P.\;Werner (1932-2018),
who conveyed my fascination of these beautiful mathematical topics
by his lessons hold in 1984 and 1985 at
Mathematical Institute A of University Stuttgart.

%------------------------------------------------------------------------------%
%                              TABLE OF CONTENTS                               %
%------------------------------------------------------------------------------%
\fancyhead[OL,ER]{\small{(\thepage)}}
\fancyhead[C] {\small{\textsc{Table of Contents}}}
\pagenumbering{roman}

\setcounter{page}{0}
\newpage
\thispagestyle{empty}
\section*{Table of Contents}
%------------------------------------------------------------------------------%
\begin{small}
\vspace{7pt}
%\centerline{{\footnotesize{Part I}}}
%\vspace{1pt}
\centerline{\textbf{I. Topological and Normed Linear Spaces}}
\vspace{2pt}

\centerline
{
\begin{tabular}{R{5mm}L{94mm}R{5mm}}
%------------------------------------------------------------------------------%
 1. & Generalities on Topological Linear Spaces                      \Df &  1 \\
 2. & Seminorms and Locally Convex Spaces                            \Df &  3 \\
 3. & Normed Linear Spaces                                           \Df &  4 \\
 4. & Strictly and Uniformly Convex Spaces                           \Df &  5 \\
 5. & Banach and Hilbert Spaces                                      \Df &  6 \\
 6. & Spaces of Scalar-Valued Functions and Measures                 \Df &  9 \\
 7. & Hierarchy of Topological and Normed Linear Spaces              \Df & 18 \\
%------------------------------------------------------------------------------%
\end{tabular}
}

\vspace{7pt}
%\centerline{{\footnotesize{Part II}}}
%\vspace{1pt}
\centerline{\textbf{II. Linear Operators in Banach Spaces}}
\vspace{2pt}

\centerline
{
\begin{tabular}{R{5mm}L{94mm}R{5mm}}
%------------------------------------------------------------------------------%
 1. & Fundamentals of Operator Theory                                \Df & 19 \\
 2. & Continuous Linear Operators                                    \Df & 20 \\
 3. & Compact Linear Operators                                       \Df & 21 \\
 4. & Closed and Closable Linear Operators                           \Df & 23 \\
 5. & The Closed Graph Theorem                                       \Df & 25 \\
%------------------------------------------------------------------------------%
\end{tabular}
}

\vspace{7pt}
%\centerline{{\footnotesize{Part III}}}
%\vspace{1pt}
\centerline{\textbf{III. Adjoint Spaces and Operators}}
\vspace{2pt}

\centerline
{
\begin{tabular}{R{5mm}L{94mm}R{5mm}}
%------------------------------------------------------------------------------%
 1. & Continuous Linear Functionals                                  \Df & 27 \\
 2. & Weak Topologies, Bidual Spaces and Reflexivity                 \Df & 27 \\
 3. & Representation Theorems for Adjoint Spaces                     \Df & 29 \\
 4. & Spaces of Distributions                                        \Df & 32 \\
 5. & Adjoint Operators in Locally Convex Spaces                     \Df & 35 \\
 6. & Adjoint Operators in Banach Spaces                             \Df & 36 \\
 7. & Hilbert Adjoints and Normal Operators                          \Df & 38 \\
%------------------------------------------------------------------------------%
\end{tabular}
}

\vspace{7pt}
%\centerline{{\footnotesize{Part IV}}}
%\vspace{1pt}
\centerline{\textbf{IV. Normally Solvable Operators}}
\vspace{2pt}

\centerline
{
\begin{tabular}{R{5mm}L{94mm}R{5mm}}
%------------------------------------------------------------------------------%
 1. & Linear Operators with Closed Range                             \Df & 42 \\
 2. & The Fredholm Index                                             \Df & 43 \\
 3. & Perturbations of Semi-Fredholm Operators                       \Df & 44 \\
 4. & Basic Theory on Fredholm Operators                             \Df & 45 \\
 5. & \textsc{Fredholm}'s Alternative Theorem                        \Df & 46 \\
 6. & Concepts of Solvability for Linear Equations                   \Df & 46 \\
%------------------------------------------------------------------------------%
\end{tabular}
}

\vspace{7pt}
%\centerline{{\footnotesize{Part V}}}
%\vspace{1pt}
\centerline{\textbf{V. Calculus of Vector-Valued Functions}}
\vspace{2pt}

\centerline
{
\begin{tabular}{R{5mm}L{94mm}R{5mm}}
%------------------------------------------------------------------------------%
 1. & Spaces of Classical Differentiable Functions                   \Df & 49 \\
 2. & Riemann Integration                                            \Df & 50 \\
 3. & Linear Evolution Equations                                     \Df & 52 \\
 4. & Bochner Integration                                            \Df & 54 \\
 5. & Inhomogeneous Cauchy Problems                                  \Df & 55 \\
 6. & Weak Derivatives and Bochner Sobolev-Spaces                    \Df & 56 \\
 7. & Hilbert-Bochner-Sobolev Spaces                                 \Df & 57 \\
 8. & Holomorphic Vector-Valued Functions                            \Df & 58 \\
%------------------------------------------------------------------------------%
\end{tabular}
}

\newpage
%\thispagestyle{empty}
%\vspace{7pt}
%\centerline{{\footnotesize{Part VI}}}
%\vspace{1pt}
\centerline{\textbf{VI. General Spectral Theory}}
\vspace{2pt}

\centerline
{
\begin{tabular}{R{5mm}L{94mm}R{5mm}}
%------------------------------------------------------------------------------%
 1. & Resolvent Sets and Resolvent Functions                         \Df & 60 \\
 2. & Different Kinds of Spectra                                     \Df & 62 \\
 3. & Riesz-Schauder Spectral Theory                                 \Df & 64 \\
 4. & Linear Operators with Pure Point Spectrum                      \Df & 66 \\
 5. & The Riesz-Dunford Functional Calculus                          \Df & 66 \\
 6. & Essential and Discrete Spectrum                                \Df & 69 \\
%------------------------------------------------------------------------------%
\end{tabular}
}

\vspace{7pt}
%\centerline{{\footnotesize{Part VII}}}
%\vspace{1pt}
\centerline{\textbf{VII. Commutative Banach Algebras}}
\vspace{2pt}

\centerline
{
\begin{tabular}{R{5mm}L{94mm}R{5mm}}
%------------------------------------------------------------------------------%
 1. & Normed Algebras                                                \Df & 71 \\
 2. & Resolvent and Spectrum of Banach Algebra Elements              \Df & 73 \\
 3. & Gelfand Representation                                         \Df & 74 \\
 4. & Involutional Algebras and C*-Algebras                          \Df & 75 \\
 5. & The Gelfand-Naimark Theorem                                    \Df & 77 \\
 6. & Locally Compact Abelian Groups                                 \Df & 79 \\
 7. & Abstract Fourier Transformation                                \Df & 82 \\
 8. & Examples of Fourier Transformation                             \Df & 85 \\
 9. & Fourier-Plancherel Transformation                              \Df & 87 \\
10. & The Fourier Transform in Distributional Spaces                 \Df & 89 \\
%------------------------------------------------------------------------------%
\end{tabular}
}

\vspace{7pt}
%\centerline{{\footnotesize{Part VIII}}}
%\vspace{1pt}
\centerline{\textbf{VIII. Linear Operators in Hilbert Space}}
\vspace{2pt}

\centerline
{
\begin{tabular}{R{5mm}L{94mm}R{5mm}}
%------------------------------------------------------------------------------%
 1. & Numerical Range and External Field of Regularity               \Df & 91 \\
 2. & Unitary Operators and Isometries                               \Df & 92 \\
 3. & Symmetric Linear Operators                                     \Df & 92 \\
 4. & Self-Adjoint Operators                                         \Df & 94 \\
 5. & Essentially Self-Adjointness                                   \Df & 95 \\
 6. & The von\,Neumann Extension Theory                              \Df & 95 \\
 7. & Perturbations of Self-Adjoint Operators                        \Df & 96 \\
%------------------------------------------------------------------------------%
\end{tabular}
}

\vspace{7pt}
%\centerline{{\footnotesize{Part IX}}}
%\vspace{1pt}
\centerline{\textbf{IX. Spectral Theory of Self-Adjoint Operators}}
\vspace{2pt}

\centerline
{
\begin{tabular}{R{5mm}L{94mm}R{5mm}}
%------------------------------------------------------------------------------%
 1. & Spectral Properties of Normal Operators                        \Df & 99 \\
 2. & Properties of the Resolvent and the Spectrum                   \Df &100 \\
 3. & The Discrete Spectral Theorem                                  \Df &102 \\
 4. & Projections and Projection-Valued Measures                     \Df &103 \\
 5. & Functional Calculus and Spectral Operators                     \Df &105 \\
 6. & The Spectral Theorem for Self-Adjoint Operators                \Df &107 \\
 7. & Positive Self-Adjoint Operators and Cauchy Problems            \Df &110 \\
%------------------------------------------------------------------------------%
\end{tabular}
}

\vspace{7pt}
%\centerline{{\footnotesize{Part X}}}
%\vspace{1pt}
\centerline{\textbf{X. Semigroups of Bounded Linear Operators}}
\vspace{2pt}

\centerline
{
\begin{tabular}{R{5mm}L{94mm}R{5mm}}
%------------------------------------------------------------------------------%
 1. & Strongly Continuous Semigroups                                 \Df &114 \\
 2. & Infinitesimal Generators                                       \Df &115 \\
 3. & Accretive Operators and Contractive Semigroups                 \Df &116 \\
 4. & Sectorial Operators and Analytic Semigroups                    \Df &118 \\
 5. & Skew-Adjoint Operators and Unitary Groups                      \Df &120 \\
 6. & Well-Posedness of Second-Order Cauchy Problems                 \Df &121 \\
%------------------------------------------------------------------------------%
\end{tabular}
}

\newpage
%\thispagestyle{empty}
%\vspace{7pt}
%\centerline{{\footnotesize{Part XI}}}
%\vspace{1pt}
\centerline{\textbf{XI. Coercive Sesquilinear Forms}}
\vspace{2pt}

\centerline
{
\begin{tabular}{R{5mm}L{94mm}R{5mm}}
%------------------------------------------------------------------------------%
 1. & Generalities on Sesquilinear Forms                             \Df &123 \\
 2. & Coercive Forms and Variational Equations                       \Df &125 \\
 3. & Triples of Hilbert Spaces                                      \Df &126 \\
 4. & Linear Operators Associated with a H-Elliptic Form             \Df &126 \\
 5. & Sectorial Forms and m-Sectorial Operators                      \Df &129 \\
 6. & The Friedrichs Extension Theorem                               \Df &132 \\
 7. & Variational Evolution Equations                                \Df &133 \\
%------------------------------------------------------------------------------%
\end{tabular}
}

\vspace{7pt}
%\centerline{{\footnotesize{Part XII}}}
%\vspace{1pt}
\centerline{\textbf{XII. Applications to Differential Operators}}
\vspace{2pt}

\centerline
{
\begin{tabular}{R{5mm}L{94mm}R{5mm}}
%------------------------------------------------------------------------------%
 1. & Free Space and Boundary Value Problems                         \Df &136 \\
 2. & Operators of Elliptic, Parabolic and Hyperbolic Type           \Df &138 \\
 3. & Functional Analytic Methods for Differential Operators         \Df &140 \\
 4. & Realizations in {\HOELDER} and Lebesgue Spaces                 \Df &142 \\
 5. & Closeness and Closability of Differential Operators            \Df &143 \\
 6. & Formally Adjoint Differential Expressions                      \Df &144 \\
 7. & {\GARDING}'s Inequality                                        \Df &146 \\
 8. & Weak Solutions of Elliptic Boundary Value Problems             \Df &149 \\
 9. & The Fredholm Property of Properly Elliptic Operators           \Df &149 \\
10. & Spectral Theory of Strongly Elliptic Operators                 \Df &150 \\
11. & Self-Adjoint Extensions of Half-Bounded Operators              \Df &152 \\
12. & Eigenvalue Problems in Bounded Domains                         \Df &157 \\
13. & Linear Differential Operators in Distributional Spaces         \Df &159 \\
14. & Free Space Problems                                            \Df &162 \\
15. & Fourier Transform Methods                                      \Df &165 \\
%------------------------------------------------------------------------------%
\end{tabular}
}

\vspace{7pt}
\centerline
{
\begin{tabular}{R{5mm}L{94mm}R{5mm}}
%------------------------------------------------------------------------------%
    & Bibliography                                                   \Df &    \\
%------------------------------------------------------------------------------%
\end{tabular}
}
\end{small}

%------------------------------------------------------------------------------%
%                                    CONTENT                                   %
%------------------------------------------------------------------------------%

%------------------------------------------------------------------------------%
%   I. TOPOLOGICAL AND NORMED LINEAR SPACES ---
%------------------------------------------------------------------------------%
%  §1. Generalities on Topological Linear Spaces ---
%  §2. Seminorms and Locally Convex Spaces ---
%  §3. Normed Linear Spaces ---
%  §4. Strictly and Uniformly Convex Spaces ---
%  §5. Banach and Hilbert Spaces ---
%  §6. Spaces of Scalar-Valued Functions and Measures ---
%  §7. The Hierarchy of Topological and Normed Linear Spaces ---
%  ---

\newpage
\pagenumbering{arabic}
\setcounter{page}{1}
\thispagestyle{empty}
\fancyhead[OL,ER]{\small{\thepage}}
\def\TitleI{Topological and Normed Linear Spaces}
\fancyhead[C] {\small{\textsc{I. \TitleI}}}
\part{\TitleI}
%------------------------------------------------------------------------------%
Our aim in this first section
is to review those fundamental classes of topological linear spaces,
in which the analysis of so-called \textit{linear operators} is usually done.
These classes include,
besides some generic classes like
{locally convex spaces} and {completely metrizable spaces},
especially the class of {normed linear spaces} and related subclasses,
where we want to emphasize
{Banach} and {Hilbert spaces}.
Finally,
there are a lot of normed linear spaces appearing in higher analysis,
whose elements are either arbitrary functions or measures.

\section{Generalities on Topological Linear Spaces}
%------------------------------------------------------------------------------%
Basically,
the scalar field $\F$ of the linear spaces under consideration in this text
will either be the field $\R$ of real numbers
or the field $\C$ of complex numbers.
Notice that these fields are furnished
with the so-called {standard topology},
whose subbase is given by the totality of open intervals in the real line
and by the totality of open disks in the complex plane,
respectively.
In general,
if results on topological or normed linear spaces
prove to be independent of the choice of the scalar field,
we shall use the symbol $\F$ instead of $\R$ or $\C$. 

\subparagraph{} %- topological linear spaces
By a \textit{topological linear space} is meant an ordinary vector space $\vecV$
over a topological field $\F$ (usually that of real or complex numbers),
together with a topology $\topT_{\,\vecV}$ 
such that addition $(x,y) \mapsto x\,{+}\,y$
and outer multiplication $(\lambda,x) \mapsto \lambda x$ are continuous.
Hence a multiplicity of topological concepts,
like
the \textit{closure} of a set,
the \textit{neighbourhood} of a point,
the \textit{convergence} of a sequence,
the \textit{continuity} of a mapping and others
are meaningful in the context of topological linear spaces.

\paragraph{} %- Hausdorff's separation axiom
We continue with the introduction of a basic property
that our topological linear spaces under consideration have to own:
all singletons $\{x\}$ with $x \xeof \vecV$ are closed sets,
which is equivalent to say
that the topology $\topT_{\,\vecV}$
fulfills \textsc{Hausdorff}'s separation axiom.
In other words,
each pair of points $x,y \xeof \banX$ can be separated
by open sets $\openO_x, \openO_y$,
that is,
there are $\openO_x \subset \vecV$ and $\openO_y \subset \vecV$ with
$x \xeof \openO_x$,
$y \xeof \openO_y$ and
$\openO_x \cap \openO_y \eq \emptyset$.
The validness of this separation axiom implies not only
that every singleton $\{x\}$ is closed,
but also
that the limit of a convergent sequence of elements in $\vecV$ is unique.

\paragraph{} %- metrizable linear spaces
Among a lot of properties,
a topological linear space may have,
there are two special ones
that are of importance for our purposes:

\begin{itemize}[leftmargin=12pt,itemsep=5pt]
\item[1.]
A topological linear space $\vecV$ is said to be \textit{metrizable},
if there is a function
$\d\:\vecV \xtimes \vecV \xto [0,\infty)$
fulfilling the properties of a \textit{metric}\footnote{
\,Given a non-empty set $\metrS$,
a mapping $\d\:\cross{\metrS}{\metrS} \xto \R$ is called a \textit{metric}
on $\metrS$
iff the following four conditions hold for all $x,y,z \in \metrS$:

\begin{itemize}[leftmargin=40pt,itemsep=0pt]
\item[(M-1)]
$d(x,y) \geq 0$,
\item[(M-2)]
$d(x,y) \eq 0$ if and only if $x=y$,
\item[(M-3)]
$d(x,y) \eq d(y,x)$;
\item[(M-4)]
$d(x,z) \leq d(x,y) + d(y,z)$ \hspace{3mm}(\textit{triangle inequality}).
\end{itemize}

Such a function $\d$ gives rise to a topology in $\metrS$,
called the \textit{metric topology} associated with $\d$,
and a subbase is given by all open balls
$B_r(x) \eqdef \{y \xeof \metrS: \d(x,y) < r\}$
with center $x \xeof \metrS$ and radius $r > 0$.
}
such that the metric topology associated with $\d$
resembles the given topology $\topT_{\,\vecV}$.

\item[2.]
A sequence $\{x_n\}_{n \in \N}$ with values
in a topological linear space $\vecV$ is said to be Cauchy
if for each neighbourhood $\U$ of $\0 \xeof \vecV$
there is an integer $N_{\U}$
such that $x_n - x_m \in \U$ for all $m,n > N_{\U}$.
Then $\vecV$ is called \textit{sequentially complete},
if every Cauchy sequence converges to a unique element $x \xeof \vecV$.
In the setting of metric spaces,
however,
the definition of Cauchy sequence may be simplified by means of the metric:
$\{x_n\}_{n \in \N}$ is Cauchy,
if for each $\epsilon > 0$ there is an integer $N_{\epsilon}$ 
such that
$x_m - x_n \xeof B_{\epsilon}(\0)$ for all $m,n > N_{\epsilon}$,
where $B_{\epsilon}(\0) \eqdef \{x \xeof \vecV: \d(x,\0) < \epsilon\}$
is the open ball with radius $\epsilon$ and center $\0 \xeof \vecV$.
\end{itemize}

\subparagraph{}
For a topological linear space being metrizable and sequentially complete
at the same time,
both concepts of completeness are compatible to each other,
that is,
such a space is complete as a metric space
\iff
it is sequentially complete as a topological linear space.
This equivalence is founded by the fact
that a metric space fulfills the \textit{first countability axiom}.

\paragraph{} %- continuous embeddings
In the sequel we also use the notion of \textit{continuous embedding}
of a topological linear space $\banX$ into a topological linear space $\banY$,
which is nothing but a continuous and \textit{injective} linear mapping
$\iota\:\banX \xto \banY$\footnote{
\,The linearity of $\iota$ means that
$\iota(\alpha x + \beta y) \eq \alpha \iota(x) + \beta \iota(y)$ holds
for all $\alpha, \beta \xeof \F$ and $x,y \xeof \banX$.}.

\subparagraph{} %- linear homeomorphisms
If $\iota$ is surjective and $\iota^{-1}$ is continuous as well,
then both mappings are called \textit{linear homeomorphisms}
between $\banX$ and $\banY$.
By $\banX \cong \banY$ we mean
that $\banX$ and $\banY$ are linear-homeomorphic to each another.

\section{Seminorms and Locally Convex Spaces}
%------------------------------------------------------------------------------%
The class of \textit{locally convex spaces} can be seen as the most important
one within all kinds of topological linear spaces.
The topology $\topT_{\,\vecV}$ of a locally convex space $\vecV$ distinguishes
by the fact
that its origin $\0$
(and thus any other element in $\vecV$) owns a neighbourhood base
consisting of \textit{convex}, \textit{balanced} and \textit{absorbing} sets.

\subparagraph{} %- seminorms
It turns out
that each such pair $(\vecV,\topT_{\,\vecV})$ gives rise to
a family of \textit{seminorms} on the space $\vecV$.
By the concept of seminorm
we mean a mapping $\abs{\,\cdot\,} \:\vecV \xto \R$ having the properties

\begin{itemize}[leftmargin=55pt,itemsep=2pt]
\item[(SN-1)] \quad
$\abs{\lambda x} \eq \abs{\lambda} \cdot \abs{x} \for \lambda \xeof \F$
\textit{(homogenity)},
\item[(SN-2)] \quad
$\abs{x+y} \, \leq \, \abs{x} + \abs{x} \for x,y \xeof \vecV$
\textit{(subadditivity)},
\end{itemize}

which especially implies $\abs{x} \leq 0$ for $x \xeof \vecV$
and $\abs{\0} \eq 0$.

\subparagraph{} %- Minkowski functional
Moreover,
if $\U \xsubset \vecV$ forms a convex, balanced and absorbing neighbourhood
of the origin $\0$,
then the \textit{Minkowski functional}
\[
\begin{cases}[1.3]
\quad \abs{\,\cdot\,} _{\,\U}\:\vecV \xto \R_0^+,\\
\quad x \mapsto \abs{x} _{\,\U} \, \eqdef \,
\mbox{inf} \, \{\lambda \geq 0: x \in \lambda \, \U\} \\
\end{cases}
\]
fulfills all the properties of a seminorm.
On the other hand,
a family $\famF$ of seminorms defined on a linear space $\vecV$
leads in a natural way to a locally convex topology $\topT_{\,\vecV}$:
it is the coarsest topology on $\vecV$
such that all seminorms in $\famF$ remain continuous.

\subparagraph{}
In summary,
there is a one-to-one correspondence
between linear spaces endowed with locally convex topologies
and those endowed with families of seminorms.
Notice that a locally convex topology is Hausdorff
iff
the family of seminorms is \textit{separating},
{\ie} for each element $x \neq 0$
there is at least one seminorm $\abs{\,\cdot\,} _{\alpha}$
such that $\abs{x} _{\alpha} \neq 0$.

\paragraph{} %- F-spaces and Frechet spaces
There are a two subclasses of locally convex spaces,
we want to review in the following:

\begin{itemize}[leftmargin=12pt,itemsep=5pt]
\item[1.]
A topological linear space $\vecV$ is called \textit{F-space},
if it is also a \textit{complete} metric space with translation-invariant
metric $\d$,
{\ie}
\[
\d(x,y) \= \d(x+z,y+z) \; \fa x,y,z \in \vecV,
\]
and the metric topology coincides with the given topology on $\vecV$.

\item[2.]
If the topology of the F-space $\vecV$ is locally convex,
then $\vecV$ is said to be a \textit{{\FRECHET} space}\footnote{
A metrizable locally convex space that is not necessarily complete,
is called a \textit{pre-{\FRECHET} space}.}.
Alternatively,
a locally convex space is {\FRECHET}
\iff
it is complete and its family of seminorms is mostly countable.
In this case
we may define a translation-invariant metric on a {\FRECHET} space $\vecV$
according to
\[
\d(x,y) \, \eqdef \,
\sum_{n=1}^{\infty} p_n \cdot \frac{\abs{x-y} _n}{1+ \abs{x-y} _n}
\quad (x,y \in \topV)
\]
where $\{p_n\}_{n=1}^{\infty}$ is any sequence of positive real numbers,
whose infinite sum converges.
\end{itemize}

\section{Normed Linear Spaces}
%------------------------------------------------------------------------------%
Based on the preceded results,
a normed linear space is nothing but a pre-{\FRECHET} space,
whose family of separating seminorms consists of only one element
$\norm{\,\cdot\,} \eqdef \abs{\,\cdot\,} _0$,
which of course must be a separating one,
that is,
$\norm{x} \eq 0$ implies $x \eq \0$.

\subparagraph{}
It is possible to introduce the concepts of norm and normed linear space
in a direct way, 
{\ie} without using the theory on families of seminorms
and locally convex spaces.
Following this approach,
a \textit{normed space} $\banX$ is a $\F$-linear space
that differs from other abstract linear spaces
by the presence of a particular mapping $\norm{\,\cdot\,}: \banX \xto \R$,
denoted the \textit{norm} on $\banX$,
which fulfills for all $x,y \xeof \banX$ and $\lambda \xeof \F$

\begin{itemize}[leftmargin=55pt,itemsep=2pt]
\item[(N-1)] \quad
$\norm{\lambda \, x} \= \abs{\lambda} \cdot \norm{x}$,

\item[(N-2)] \quad
$\norm{x+y} \, \leq \, \norm{x} + \norm{y}$
\; \textit{(triangle inequality)},

\item[(N-3)] \quad
$\norm{x} \geq 0$ and $\norm{x} \eq 0 \Leftrightarrow x \eq \0$.
\end{itemize}

\paragraph{}
Linear spaces can be furnished with different norms.
Then two arbitrary norms $\norm{\,\cdot\,}_{\alpha}$
and $\norm{\,\cdot\,}_{\beta}$ defined on a linear space $\banX$
are called \textit{equivalent} to each another,
if there are constants $c_1,c_2 > 0$ such that
$
c_1 \norm{x}_{\alpha} \leq \norm{x}_{\beta}  \leq c_2 \norm{x}_{\alpha} 
$
for all $x \xeof \banX$.

\subparagraph{} %- norm topologies
Each normed linear space $\banX$ can be equipped
with a very natural topology $\topT$,
called the \textit{norm topology} on $\banX$.
It can be defined either by the assertion
that $\topT \eqdef \topT_{\norm{\cdot}}$ is the smallest topology on $\banX$
such that the norm mapping remains continuous,
or it is defined to be the metric topology $\topT_d$,
when $\banX$ is considered as a metric space with metric function
$\d(x,y) \eqdef \norm{x-y}$ for all $x,y \xeof \banX$.
In fact,
both definitions lead to the same collection of open subsets in $\banX$,
so $\topT \eqdef \topT_{\norm{\cdot}} \eq \topT_d$,
which justifies to speak of \textit{the} topology of a normed linear space.
Clearly,
equivalent norms on $\banX$ own one and the same topology.
For example,
if $\banX$ is a finite-dimensional linear space,
then all norms on $\banX$ are equivalent and each of them
generates the so-called \textit{standard topology} on $\banX$.

\subparagraph{} %- norm-isomorphisms
if $\banX$ and $\banY$ are normed linear spaces
and $\iota$ is a \textit{norm-preserving} linear homeomorphism
between $\banX$ and $\banY$,
that is,
$\norm{\iota x}_{\banY} \eq \norm{x}_{\banX}$ for all $x \xeof \banX$,
we denote this mapping a \textit{norm-isomorphism} between $\banX$ and $\banY$,
and in this case both spaces are called \textit{norm-isomorphic} to each other
(we again use the notion $\banX \cong \banY$ in this case).

\subparagraph{} %- separable spaces
A normed linear space $\banX$ is called \textit{separable}
if there is a mostly countable subset lying dense in $\banX$.

\paragraph{} %- norm convergence
We finally introduce the concept of convergence
for sequences in a \textit{normed} linear space $\banX$.
Any sequence $\{x_n\}_{n \in \N}$ in $\banX$
is called \textit{strongly convergent}
or simply \textit{norm convergent} to $x \xeof \banX$,
if it converges to $x$ {\wrt} the norm topology,
that is,
if $\lim_{\,n \xto \infty} \norm{x_n-x} \eq 0$.

\section{Strictly and Uniformly Convex Spaces}
%------------------------------------------------------------------------------%
In this paragraph we consider certain normed linear spaces,
in which the best approximation of arbitrary elements
by other elements contained in a convex subset is always possible.

\paragraph{} %- uniformly convex spaces
A \textit{uniformly convex space} is a \textit{real} normed linear space $\banX$
having the following property:
for every $\epsilon > 0$ there is $\delta > 0$ 
such that
for any elements $x,y$ of the closed unit ball $B_1(\0)$ of $\banX$
the validness of the inequality $\norm{(x+y)/2} > 1- \delta$
implies $\norm{x-y} < \epsilon$.
In other words,
if the center $(x+y)/2$ of the line segment $[x,y]$
between elements $x,y$ belonging to the \textit{unit sphere}
$S_1(\0) \eqdef \partial B_1(\0)$ tends to $S_1(\0)$,
then the $x$ and $y$ tend together.
This property is possibly lost
when switching to another but equivalent norm on $\banX$.

\subparagraph{} %- strictly convex spaces
A normed linear space $\banX$ is called \textit{strictly convex},
if the conditions
$\norm{x} \leq 1$, $\norm{y} \leq 1$ and $x \neq y$
imply
$\norm{x+y}/2< 1$.
This means,
that the \textit{inner} line segment $(x,y)$
between $x,y \xeof S_1(\0)$ lies completely in the interior of $B_1(\0)$,
that is
$\norm{\,\alpha x + (1-\alpha)\,y\,} < 1$ for $0 < \alpha < 1$.
Each linear subspace of a uniformly convex Banach space owns this feature.
In the light of the forthcoming theorem \ref{I_APPROXIMATION_THEOREM},
it is an important fact
that uniformly convex spaces are also strictly convex.

\paragraph{} %- approximation theorem
Recall that a subset $\convC$ of a linear space is called \textit{convex},
if the line segment $[x,y]$ of any pair of points $x,y \xeof \convC$
belongs to $\convC$ again.
Then many approximation problems may be formulated as follows:
it is to find an element $y_0$ in a convex set $\convC \subset \banX$
such that $y_0$ has minimal distance to a given element $y \notin \convC$.
The following statement,
known as the \textit{approximation theorem},
provides sufficient conditions regarding existence and uniqueness
of a solution for that problem:

\begin{theorem1}\label{I_APPROXIMATION_THEOREM}
Let $(\banX,\norm{\,\cdot\,})$ be a real normed linear space,
let $y \xeof \banX$ and
$\convC \subset \banX$ be a \textit{closed} and \textit{convex} subset.
Then we have:

\begin{itemize}[leftmargin=17pt,itemsep=5pt]
\item[(a)] \textit{Uniqueness}:
if $\banX$ is \textit{strictly convex},
then there is mostly one element $y_0 \xeof \convC$
such that $\norm{y_0-y} \eq \inf_{x \xeof \convC} \norm{y-x}$.

\item[(b)] \textit{Existence}:
if $\banX$ is a \textit{uniformly convex Banach space}
then there is a unique element $y_0 \xeof \vecV$
such that $\norm{y_0-y} \eq \inf_{x \xeof \convC} \norm{y-x}$.
\end{itemize}
\end{theorem1}

\section{Banach and Hilbert Spaces}
%------------------------------------------------------------------------------%
A normed linear space $\banX$ is said to be \textit{sequentially complete}
(or briefly \textit{complete}),
if for any Cauchy sequence $\{x_n\}_{n \in \N}$ in $\banX$
there is $x \xeof \banX$
such that $\lim_{\,n \xto \infty} \norm{x_n-x} \eq 0$.
Equivalently,
the space $\banX$ is complete
\iff
for every absolutely convergent series
$\sum_{n=1}^{\infty} \norm{x_n} < \infty$ there is $x \xeof \banX$
such that $\norm{x-x_n} \xto 0$ for $n \xto \infty$.
Each complete normed linear space is denoted a \textit{Banach space}.
Like in the non-complete case,
a Banach space $\banX$ with norm $\norm{\,\cdot\,}$
is nothing but a {\FRECHET} space,
whose separating family of seminorms contains only this norm.

A complex (real) \textit{Hilbert space} differs
from an arbitrary Banach space by the presence of an \textit{inner product},
which is nothing but a symmetric and positive sesquilinear (bilinear) form
on such a space.
This form defines the norm and the metric topology of the Hilbert space,
but also determines the inner geometry to be close
to that of an Euclidean space.

\subparagraph{} %-- inner geometry of Hilbert spaces
Let us make these assertions more precise.
Given a real or complex linear space $\hilH$,
an \textit{inner product} in $\hilH$ is defined to be a mapping
$\scalar{\,\cdot\,}{\,\cdot\,}\: \hilH \xtimes \hilH \xto \F$,
which is

\begin{itemize}[leftmargin=35pt]
\item[(I-1)]
linear in the first argument,

\item[(I-2)]
either linear or antilinear in the second argument (conditional upon,
whether $\F$ is the real or the complex number field),
\end{itemize}

together with the additional properties

\begin{itemize}[leftmargin=35pt]
\item[(I-3)]
$\scalar{y}{x} \eq           \scalar{x}{y}$  ($\F \eq \R$) rsp.
$\scalar{y}{x} \eq \overline{\scalar{x}{y}}$ ($\F \eq \C$),
\item[(I-4)]
$\scalar{x}{x} \geq 0$ for all $x \neq 0$,
\item[(I-5)]
$\scalar{x}{x} \eq 0$ implies $x \eq 0$.
\end{itemize}

\subparagraph{}
Property (I-3) is denoted the \textit{Hermitean symmetry}
of $\scalar{\,\cdot\,}{\,\cdot\,}$,
while (I-4) and (I-5) together mean \textit{positive definiteness}.
The presence of the inner product gives rise to define a norm in $\hilH$
according to
\begin{equation}\label{I_DERIVED_NORM}
\norm{x} \, \eqdef \, \scalar{x}{x} ^{1/2} \quad \quad (x \in \hilH),
\end{equation}
so each linear space furnished with an inner product is normed
in a natural manner.
The inner product and the norm are coupled by
the \textit{Cauchy-Schwarz} \textit{inequality}
\[
\abs{\scalar{x}{y}} \, \leq \,
\norm{x} \cdot \norm{y} \quad \quad (x,y \in \hilH),
\]
which implies that the mapping
$\scalar{\,\cdot\,}{\,\cdot\,}\:\hilH \xtimes \hilH \xto \F$ is continuous.
An inner product space $\hilH$ is called a \textit{Hilbert space}
or shortly \textit{Hilbertian},
if it is complete {\wrt} the norm defined by \eqref{I_DERIVED_NORM}.

\subparagraph{} %- parallelogram equation and polarization formula
The pleasant geometric properties of inner product spaces
may be attributed to the validness of the \textit{parallelogram equation}
\begin{equation}\label{I_PARALLELOGRAM_EQUATION}
2 \norm{x}^2 + 2 \norm{y}^2 \= \norm{x+y}^2 + \norm{x-y}^2
\qfa x,y \xeof \hilH.
\end{equation}
In general,
the norm $\norm{\,\cdot\,}$ on a linear space $\banX$
can be induced by an inner product
\iff \eqref{I_PARALLELOGRAM_EQUATION} holds.
In this case and if $\banX$ is a complex linear space,
the inner product can be recovered from the norm
by the \textit{polarization formula}
\[
\scalar{x}{y}
\=
\frac{1}{4} \,
(\,
\norm{x + y}^2 - \norm{x-y}^2 +i \norm{x+iy}^2 - i \norm{x-iy}^2
\,).
\]
If $\banX$ is a real space,
however,
the polarization formula simplifies to
\[
\scalar{x}{y} \=
\frac{1}{4} \,
(\,
\norm{x + y}^2 - \norm{x-y}^2
\,).
\]

\paragraph{} %- properties of Hilbert spaces
In the following we want to point out
that Hilbert spaces differ from general Banach spaces
by a couple of desirable properties:

\begin{itemize}[leftmargin=12pt,itemsep=5pt]
\item[1.] \textit{Orthogonality}.
Two elements $x,y \xeof \hilH$ are called \textit{orthogonal} to each other,
if $\scalar{x}{y} \eq 0$ (notation: $x \bot y$).
Any subset $\orthO \xsubset \hilH$ is said to be \textit{orthogonal}
or an \textit{orthogonal system},
if all elements in $\orthO$ are orthogonal to each other.
Finally,
the \textit{orthogonal space} $\setA^{\bot}$ of a subset $\setA \xsubset \hilH$
contains exactly the elements $y \xeof \hilH$,
which are orthogonal to all $x \xeof \setA$,
and forms a closed linear subspace of $\hilH$.

\item[2.] \textit{Projection Theorem}.
Let $\setA \xsubset \hilH$ be a \textit{closed} linear subspace.
Then for each $x \xeof \hilH$ there is exactly one $x_0 \xeof \setA$
such that
\[
\norm{ x - x_0 } \= \inf \, \{ \, \norm {x-y}: y \in \setA \,\}.
\]
The element $x_0$ is called the \textit{projection} of $x$ onto $\setA$
and minimizes the distance between $\setA$ and $x$.
The associated linear mapping
$\pi_{\setA} \: x \xeof \hilH \mapsto x_0 \xeof \setA$
is also denoted the \textit{projection} of $\hilH$ onto $\setA$.
In general,
by a \textit{projection} is meant
any linear mapping $\pi\:\hilH \xto \hilH$ satisfying
$\pi \circ \pi \eq \pi$ and
$\scalar{\pi x}{y} \eq \scalar{x}{\pi y}$ for all $x,y \xeof \hilH$.

\item[3.] \textit{Decomposition Theorem}.
Using the notation of orthogonal space $\setA^{\bot}$,
the projection of an element $x$ onto $\setA$ is characterized as follows:
$x_0 \eq \pi_{\setA} x$ iff $x-x_0 \xeof \setA^{\bot}$.
Hence each $x \xeof \hilH$ can be uniquely represented by the formula
$x \eq x_0 + x^{\bot}$,
where $x_0 \xeof \setA$ and $x^{\bot} \xeof \setA^{\bot}$,
so the entire space $\hilH$ is isometrically isomorphic
to the direct sum $\setA \, \oplus \, \setA^{\bot}$
of the closed subspaces $\setA$ and $\setA^{\bot}$.
This fact is known as \textit{decomposition theorem} in Hilbert spaces.
In other words,
each closed linear subspace $\setA$ of $\hilH$
is \textit{topologically complementable},
that is,
to each closed subspace $\setA$
there is a second closed subspace $\setB \xsubset \hilH$
such that $\setA \cap \setB \eq \{\0\}$ and $[\setA \cup \setB] \eq \hilH$.
Of course,
a closed subspace $\A$ is topologically complemented
by its orthogonal space. %$\setA^{\bot}$.

\item[4.] \textit{Existence of Orthonormal Bases}.
An orthogonal system $\orthO \xsubset \hilH$ is denoted
an \textit{orthonormal system},
if $\norm{x} \eq 1$ for all $x \xeof \orthO$.
However,
if $\orthO$ contains only \textit{finitely} many elements,
then the span $[\orthO]$ forms a \textit{closed} subspace in $\hilH$,
and the projection $\pi_{[\orthO]}$ of any element $x$ onto $[\orthO]$
can be expressed by the formula
\[
\pi_{[\orthO]} \, x \= \sum_{i=1}^n \, \scalar{x}{e_i} \, e_i.
\]
However,
if $\orthO$ consists of arbitrarily many elements,
the inequality $\scalar{x}{e_{\alpha}} \neq 0$ can only be valid
for mostly countable many elements $e_{\alpha} \xeof \orthO$,
which in conclusion enables to add the coefficients $\scalar{x}{e_{\alpha}}$
together about all elements of $\orthO$.

\subparagraph{} %- orthonormal bases
An orthonormal system $\orthO$ is called
an \textit{orthonormal base} in $\hilH$,
if its span $[ \orthO ]$ lies dense in $\hilH$.
In this case we may represent each $x \xeof \hilH$ according to
\begin{equation}\label{I_FOURIER_COEFFICIENTS}
x \= \sum_{e_{\alpha} \in \, \orthO} \, \scalar{x}{e_{\alpha}} \, e_{\alpha},
\end{equation}
which is primarily a formal sum,
but consists of mostly countable many summands.
The inner products $\scalar{x}{e_{\alpha}}$ are called
the \textit{Fourier coefficients} of $x$ {\wrt} $\orthO$.
The right hand side of \eqref{I_FOURIER_COEFFICIENTS}
consists of \textit{mostly countable many} non-zero summands,
and this series is invariant
by realignments of the elements $e_{\alpha} \xeof \orthO$.

\subparagraph{}
There is the need to determine,
whether a given orthonormal system $\orthO$ forms an orthonormal base.
In fact,
it can be shown that
the following statements are equivalent:

\begin{itemize}[leftmargin=20pt]
\item[(a)]
$\orthO$ is an orthonormal base in $\hilH$.

\item[(b)]
$\orthO$ is \textit{complete},
that is,
$\scalar{x}{e} \eq 0$ for all $e \xeof \orthO$ implies $x \eq \0$.

\item[(c)]
$\sum_{e_{\alpha} \in \,\orthO} \;
\abs{ \scalar{x}{e_{\alpha}} } ^2 \eq \norm{x}^2$ for all $x \xeof \hilH$
(\textit{Parseval}\textit{'s equation}).
\end{itemize}

\paragraph{} %- existence of orthonormal bases
It is also nearby to ask for the existence of orthonormal bases.
First of all,
the collection of all orthonormal systems in a Hilbert space $\hilH$
forms a partial ordered set,
and under the assumption that \textsc{Zorn}{'s lemma}\,\footnote{\,
\,The \textit{axion of choice} claims that
if $\cal{F}$ is a disjoint family $F_{\alpha}$ ($\alpha \in \J$)
of non-empty sets,
then there is a set $A$ that shares exactly one element
with each other set $F_{\alpha} \in \cal{F}$.}
is valid,
it contains a \textit{maximal} element $\orthO$
that cannot be increased,
hence each Hilbert space $\hilH \neq \{\0\}$ owns an orthonormal base.

\subparagraph{} %- Hilbert space dimension
Furthermore,
if $\orthO_1$ and $\orthO_2$ are orthonormal bases in $\hilH$,
then there is a bijective mapping between them,
thus they own identical cardinal numbers.
This legitimates us to introduce the notation
of \textit{Hilbert space dimension} $\abs{\hilH} $
being the cardinal number of any orthonormal base in $\hilH$.

\subparagraph{} %- separable Hilbert spaces
The space $\hilH$ is separable \iff $\abs{ \hilH} \leq \abs{\N} $.
Indeed,
each orthonormal base in $\hilH$ forms a mostly countable set in this case,
and $\hilH$ proves to be norm-isomorphic to the Hilbert space $\ell^2$
of square summarizable sequences $f\:\N \xto \F$.
\end{itemize}

\section{Spaces of Scalar-Valued Functions and Measures}
%------------------------------------------------------------------------------%
This paragraph forms an enumeration of well-known and important
topological linear spaces of functions or measures.
It is assumed that the involved functions are either real or complex-valued.

\begin{itemize}[leftmargin=12pt,itemsep=5pt]
\item[1.] \textit{Spaces of continuous functions.}\\
For any topological space $\topX$
the function spaces $C_b(\topX)$ and $C_0(\topX)$,
both endowed with the \textit{supremum norm}
\[
\norm{f} \eqdef \sup_x \, \abs{f(x)} \,,
\]
where $x$ runs through the elements of the underlying space $\topX$.
Both are complete and hence examples of Banach spaces.
If the space $\topX$ is compact,
then $C(\topX)$,
the space of continuous functions $f\:\topX \xto \C$,
forms a Banach space as well.

\item[2.] \textit{The space $\MS{X}$.}\\
Recall from measure theory
that a \textit{signed measure},
defined on a measurable space $(X,\salgB)$,
is a $\sigma$-\textit{additive} set function
\[
\mu: \salgB \to (-\infty,+ \infty] \quad \mbox{ or } \quad
\mu: \salgB \to [-\infty,+ \infty).
\]
If $\mu$ is a \textit{finite signed measure}
(the values $+\infty$ und $-\infty$ are excluded in this case),
it can be represented as the difference of positive measures $\mu_+$
(\textit{Hahn-Jordan decomposition theorem}).
In this case
the positive measure $\abs{\mu} := \mu_+ + \mu_-$ is called the
\textit{variation}\ of $\mu$,
and the mapping $\norm{\mu} := \abs{\mu} (X)$ is said to be
the \textit{total variation} of $\mu$.
It can be shown 
that the total variation constitutes a norm on the space $\MS{X}$
of finite signed measures,
which is even complete and hence forms a Banach space.

\item[3.] \textit{The space $\MC{X}$.}\\
Similarly to finite signed measures,
a \textit{complex measure} on a measurable space $(X,\salgB)$
is a $\sigma$-additive set function $\mu:\salgB \to \C$.
The set $\MC{X}$ of complex measures on $(X,\salgB)$
actually forms a $\C$-linear space.
It is possible to attach to each complex measure $\mu$
an arbitrary positive measure $\abs{\mu} $,
called the \textit{total variation} of $\mu$.
It is the mapping
$\abs{\mu} :\salgB \to [0,\infty]$,
which assigns to each measurable set $A \xeof \salgB$ the value
\[
\abs{\mu} (A)
\, := \,
\sup \, \sum_{n=1}^{\infty} \, \norm{\,\mu(A_n)},
\]
where $(A_n)$ may be any mostly countable and disjoint decomposition
of $A$ into measurable subsets of $A$.
Indeed,
the mapping $\abs{\mu} $ is a positive measure,
and a complex measure $\mu$ is called of \textit{bounded variation},
if $\abs{\mu} (X) < \infty$.
The mapping 
\[
\mu \mapsto \norm{\mu}
\, := \,
\abs{\mu} (X)
\]
satisfies all conditions of a norm,
hence the pair $(\MC{X} ,\norm{\cdot})$ forms a normed linear space.
Like the space $\MS{X}$ of finite signed measures,
the space $\MC{X}$ is complete and thus a Banach space.

\item[4.] \textit{The spaces $BV[a,b]$ and $AC[a,b]$.}\\
Let $BV[a,b]$ be the linear space consisting
of all functions of bounded variation.
This space can be equipped with the norm
\[
\norm{f}_{bv} \, \eqdef \,
\abs{ f(a) } \, + \,
\sup_{P} \sum_{i=0}^{n_P-1} \abs{f(x_{i+1}-f(x_i)} ,
\]
where the supremum is taken over all finite partitions of $[a,b]$.
The space $BV(a,b)$ forms a Banach space in this norm.

\subparagraph{} %- the space AC[a,b]
A measurable function $f\:[a,b] \xto \R$ is called 
\textit{absolutely continuous},
if for each $\epsilon > 0$ there is $\delta > 0$ 
such that for a finite decomposition
\[
x_0=a < x_1 < \ldots < x_n=b \mbox{ mit }
\sum_{k=1}^n \abs{x_k-x_{k-1}} < \delta
\]
of the interval $[a,b]$
the condition $\sum_{k=1}^n \abs{f(x_k)-f(x_{k-1})} < \delta$ is valid.
Every absolutely continuous function is of bounded variation.
The etenge $AC[a,b]$ of absolutely continuous functions on $[a,b]$
endowed with the norm $\norm{\,\cdot\,}_{BV}$
forms a closed linear subspace of $BV[a,b]$ and hence 
is a Banach space in its own right. 
\end{itemize}

\paragraph{} %- spaces of functions defined on a domain G
Let us continue with the consideration of function spaces,
where the underling set $\G$ is a \textit{domain} in $\Rd$,
that is,
an open and connected set within an Euclidean space $\Rd$,
which may be bounded, unbounded or even the space $\Rd$ itself.

\subparagraph{} %- multi-indices
We will also work in the sequel with \textit{multi-indices},
which are nothing but $d$-tuples
$\alpha \eq (\alpha_1,\ldots,\alpha_d)$
with values $\alpha_i \xeof \N_0$
and modulus $\abs{\alpha} \eqdef \alpha_1 + \ldots + \alpha_d$.
They are mostly used to shorten multiple partial derivatives
$
\PA f \eqdef \partial_1^{\alpha_1}
         \partial_2^{\alpha_2} \ldots
         \partial_d^{\alpha_d} f
$
of a sufficiently often differentiable and complex-valued function $f$,
but also to denote the multiple powers
$\xi^{\alpha} \eqdef \xi_1^{\alpha_1}
\cdot \xi_2^{\alpha_2} \cdot \ldots \cdot \xi_d^{\alpha_d}$
of a vector $\,\xi \eq (\xi_1,\xi_2,\ldots,\xi_d) \xeof \C^d$.

\subparagraph{}
\begin{itemize}[leftmargin=12pt,itemsep=5pt]
\item[5.] \textit{Spaces of continuously differentiable functions}.\\
The linear subspace $C^k(\clG) \subset C(\clG)$
consists of those continuously differentiable functions $f\:\G \xto \F$,
whose partial derivatives $\partial^{\alpha} f$ are bounded on $\clG$
for every multi-index $\alpha$ with $\abs{\alpha} \leq k$.
It is complete if endowed with the norm
\[
\norm{f}_k
\, \eqdef \,
\sum_{\abs{\alpha} \leq k} \,
\sup_{x \in \clG } \; \abs{\partial^{\alpha} f(x)} .
\]

\item[6.] \textit{{\HOELDER} spaces}.\\
It has been shown
that the spaces $C^k(\clG)$ are not very convenient
for the investigation of partial differential equations,
but the so-called \textit{{\HOELDER} spaces} $C^{k,\gamma}(\clG)$ are so.
Preliminary,
a function $f\:\G \xto \F$ is called \textit{{\HOELDER}-continuous}
{\wrt} the parameter $\gamma \xeof (0,1]$,
if its \textit{{\HOELDER}-norm} is finite,
that is,
\[
\abs{f} _{\,\gamma} \, \eqdef \,
\sup_{x,y \, \in \, \G, \, x \neq y} \,
\frac{\abs{f(x) - f(y)} } {\abs{x-y} ^{\gamma}} \, < \, \infty.
\]
Let $C^{k,\gamma}(\clG)$ be the linear space
of functions $f \xeof C^k(\clG)$,
whose partial derivatives $\PA f$
are {\HOELDER}-continuous for all $\abs{\alpha} \eq k$.
Then the \textit{{\HOELDER} norm} $\abs{f} _{k,\gamma}$
of a function $f \xeof C^{k,\gamma}(\clG)$ is defined
by the sum of the norm $\norm{f}_k$
and the {\HOELDER} norms of the partial derivatives $\PA f$
for $\abs{\alpha} \eq k$:
\[
\abs{f} _{k,\gamma}
\, \eqdef \,
\norm{f}_k
\, + \,
\sum_{\abs{\alpha} \eq k} \abs{\,\PA f} _\gamma.
\]
The space $C^{k,\gamma}(\clG)$ is complete {\wrt} this norm.
Notice that for $k \eq 0$ and $\gamma \eq 1$ we obtain the space
$C^{0,1}(\clG)$ of \textit{Lipschitz-continuous} functions.

\subparagraph{}
If $\G$ is a bounded domain, we have for $0 < \alpha < \beta \leq 1$
\[
C^1(\clG) \subsetneq
C^{0,\beta}(\clG) \subsetneq
C^{0,\alpha}(\clG) \subsetneq
C(\clG).
\]

\item[7.] \textit{Spaces of smooth functions}.\\
Next we consider three function spaces
consisting of indefinitely differentiable functions $f\:\G \xto \F$
that frequently occure in applied analysis
and constitute {\FRECHET} spaces respectively
complete locally convex linear spaces.

\begin{itemize}[leftmargin=17pt]
\item[(a)]
The space $\SF$ of functions $f\:\G \xto \F$ being indefinitely differentiable.
Unfortunately,
there is no norm that makes $\SF$ into a complete normed space,
but after all it is possible to define a countable family $\famF$
of seminorms to make $\SF$ into a {\FRECHET} space.
For example,
$\famF$ can be given by the mappings
\[
p_{N,n} \, \eqdef \,
\max_{x \in \compK_n} \;
\{ \, \abs{\, (\PA \phi)(x) \,} : \abs{\alpha} \leq N \}
\quad \quad (N,n \in \N)\,,
\]
where $(\compK_n)$ is a sequence of compact subsets of $\G$
such that
\[
\compK_1 \subset \compK_2 \ldots \subset \G \, \mbox{ for } \, n \in \N,
\quad
\bigcup_{n \in \N} \, \compK_n \= \G,
\]
and each $\compK_n$ lies in the interior of $\compK_{n+1}$
(such a sequence is called a \textit{compact exhaustion} of $\G$).
The space $\SF$ is often denoted $\E(\G)$,
when considered in the locally convex topology induced by the totality
of seminorms $p_{N,n}$.
It is complete {\wrt} this topology,
so it forms a {\FRECHET} space.

\item[(b)]
The subspace $\TF \subset \SF$ of smooth functions
having \textit{compact} support
\[
\supp{\,\phi} \, \eqdef \, \overline{\{x \in \G: \phi(x) \neq 0\}},
\]
usually denoted the \textit{space of test functions}.
This function space plays an outstanding role in analysis,
since it lies dense in a multiplicity of other function spaces.
Unfortunately,
there is neither a countable family of seminorms
making $\TF$ into a {\FRECHET} space nor is $\TF$ completely metrizable.

\subparagraph{}
But it is still possible to define a locally convex topology $\topT$ on $\TF$
such that $\TF$ is complete in the sense of a topological linear space.
Instead of saying how the open sets of this topology $\topT$ look like,
let us define $\topT$ by means of a suitable notion of convergence,
known as $\D$-convergence in $\TF$:\,
a sequence $(\phi_n)$ in $\TF$ is convergent to $\phi \xeof \TF$,
if there is a compact set $\compK \subset \G$
such that
(i)  $\supp{\phi_n} \subset \compK$ for all $n \xeof \N$ and
(ii) for each multi-index $\alpha$ the sequence $(\PA \phi_n)$
     converges to $\PA \phi$ \textit{uniformly} on $\compK$.

The space $\TF$ equipped with the topology
induced by $\D$-convergence is usually denoted $\D(\G)$.
Notice
that we have $\D(\G) \hookrightarrow \E(\G)$,
that is,
the space $\D(\G)$ is continuously embedded into the space $\E(\G)$
{\wrt} the locally convex topologies of these spaces.

\subparagraph{}
We have $\D(\G) \hookrightarrow \E(\G)$,
that is,
the space $\D(\G)$ is continuously embedded into the space $\E(\G)$
{\wrt} their locally convex topologies.

\item[(c)]
The \textit{Schwartz space} $\SRd$,
which is the space of all smooth functions $\phi \xeof C^{\infty}(\Rd)$
having the additional property
\[
\abs{ \PA \phi(x) \cdot x^{\beta}} \xto 0 \,
\for \norm{x}_d \xto \infty
\]
and arbitrary multi-indices $\alpha,\beta$,
where $\norm{x}_d$ is the Euclidean norm of $x \xeof \Rd$.
The elements of $\SRd$ are called \textit{rapidly decreasing}
or briefly \textit{Schwartz functions}.
Hence a Schwartz function and all its derivatives vanish at infinity
faster then the reciprocal of any polynomial.
The most famous example of a Schwartz function is the Hermite function
\[
\cal{H}_0(x) \, \eqdef \, e^{- \frac{\norm{x}^2}{2}} \quad \quad (x \in \Rd).
\]

\subparagraph{}
Rather similar to the space $\SF$,
it is possible to find for the Schwartz space $\SRd$
a family of seminorms as follows:
define for $\phi \xeof C^{\infty}(\Rd)$ and multi-indices $\alpha, \beta$
\begin{equation}\label{I_SRN_FAMILY_OF_SEMINORMS}
p_{\alpha,\beta}(\phi) \, \eqdef \,
\sup_{x \in \Rd} \abs{ \, x^{\alpha} (D^{\beta} \phi) (x) \, } .
\end{equation}

\subparagraph{}
Then $\phi$ is a Schwartz function
iff $p_{\alpha,\beta}(\phi) < \infty$
holds for all indices $\alpha, \beta$.
The collection of mappings $p_{\alpha,\beta}\:\SRd \xto \R$
defines a countable family of seminorms,
which gives rise to a locally convex topology $\topT_{\mathscr{S}}$ on $\SRd$.
Moreover,
since the Schwartz space is complete {\wrt} $\topT_{\mathscr{S}}$,
it actually forms a \textit{{\FRECHET} space}.

\subparagraph{}
Comparing the Schwartz space $\SRd$ with the function spaces
$\D(\Rd)$ and $\E(\Rd)$,
it is easy to see that there is the chain of inclusions
\[
\D(\Rd) \, \hookrightarrow \, \SRd \, \hookrightarrow \, \cal{E}(\Rd).
\]
Moreover,
by comparing the locally convex topologies of these spaces,
it turns out that these inclusions are continuous.
\end{itemize}

\item[8.] \textit{Lebesgue spaces}.\\
To introduce the spaces $L^p(\G)$ for $p \xeof [1,\infty)$,
let $\cal{L}^p(\G)$ be the $\F$-linear space of measurable functions
$f\:\G \xto \F$
such that $\abs{f} ^p$ is Lebesgue-integrable.

\subparagraph{}
We define for $f,g \xeof \cal{L}^p(\G)$ the equivalence relation
$f \sim g$ according to $f \eq g$ almost everywhere.
Then the elements of the space $L^p(\G)$
are the equivalence classes in $\cal{L}^p(\G)$ {\wrt} $\sim$,
that is,
$L^p(\G) \eqdef \cal{L}^p(\G) / \sim$.
The \textit{Lebesgue norm} in $\Lp$ is given by
\[
\norm{f}_p
\, \eqdef \,
\big( \int_{\,\G} \abs{f} ^p \, d\mu \,\big) \, ^{\frac{1}{p}}.
\]
Finally,
the \textsc{Fischer-Riesz} \textit{theorem} claims
that the spaces $\Lp$ are complete {\wrt} the Lebesgue norm.
It turns out that $\Lp$ is separable \iff $1 \leq p < \infty$.

\subparagraph{} %- Lebesgue space L2
The case $p \eq 2$ gives rise to a special attention:
here it is possible to derive the norm $\norm{\,\cdot\,}_2$
from the inner product
\[
\scalar{f}{g} \, \eqdef \, \int_{\,\G} f \, \overline{g} \; d\mu,
\]
thus $\norm{f}_2 \eq \scalar{f}{f} ^{1/2}$
and $\Lh$ becomes into a Hilbert space.

\paragraph{}
For $p \eq \infty$ and any measurable function $f\:\G \xto \F$
we define the mapping
\begin{equation}\label{I_SUPREMUM_NORM}
\norm{f}_{\infty} \, \eqdef \, \inf
\{\,a \geq 0: \abs{f} \leq a \,\, \mu \mbox{-almost everywhere}\,\},
\end{equation}
taking values in the set $\R^+ \cup \{\infty\}$.
Let $\cal{L}^{\infty}(\G)$ be the linear space
of all measurable functions $f\:\G \xto \F$
with $\norm{f}_{\infty} < \infty$
and $\Li$ be the linear space of all equivalence classes
in $\cal{L}^{\infty}(\G)$,
where $f \sim g$ means $f=g$ $\mu$-almost everywhere.
Members of these spaces are called \textit{essentially bounded}.
Similar to the case $p < \infty$,
the space $\Li$ is complete {\wrt} \eqref{I_SUPREMUM_NORM},
but not separable.

\item[9.] \textit{Sobolev spaces}.\\
A function $f\:\G \xto \C$ is said to belong to the space $\Lc$,
if $f$ is Lebesgue-integrable on each compact subset of $\G$.
This space is useful to define the concept of \textit{weak differentiability}.
A function $f \xeof \Lc$ is said to be \textit{weakly differentiable}
{\wrt} a given multi-index $\alpha$,
if there is $g \xeof \Lc$
such that 
\[
\int_{\,\G} g(x) \, \phi(x) \, dx
=
(-1)^{ \abs{\alpha} } \,
\int_{\,\G} f(x) \, \PA \phi(x) \, dx
\for \phi \in \TF.
\]
We denote $\PA f \eqdef g$ in this case,
and weak differentiability arises as a linear mapping
$\PA\:D_{\alpha} \subset \Lc \xto \Lc$,
where each $D_{\alpha}$ is a linear subspace of $\Lc$.
Notice that
the concept of weak derivative is compliant
with that of the classical derivative:
if $f$ has classical derivative $\PA f$,
then $f$ is also weakly differentiable,
and the weak derivative is equal to $\PA f$.

\begin{itemize}[leftmargin=17pt]
\item[(a)] The next function spaces under consideration are those,
whose elements are $p$-integrable and weakly differentiable up to order $k$
at the same time.
More precisely,
for any integer $k$ and real $p \geq 1$,
the \textit{Sobolev space} $\SobW$ is the set
of functions $f$ in the Lebesgue space $\Lp$,
whose weak derivatives $\PA f$ exist and also belong to $\Lp$
for all multi-indices $\alpha$ with $ \abs{\alpha} \leq k$,
briefly written
\[
\SobW
\, \eqdef \,
\{\,f \in \Lp: \PA f \in \Lp, \; \abs{\alpha} \leq k \, \}.
\]
Sobolev spaces are of course linear spaces,
and especially for $k \eq 0$
the spaces $W^{0,p}(\G)$ are nothing but the Lebesgue spaces $\Lp$ themselves.

\subparagraph{} %- Sobolev norms
For functions $f \xeof \SobW$ let the mapping $f \mapsto \abs{f} _{k,p}$
be defined according to
\begin{equation}\label{I_SOBOLEV_NORM}
\abs{f} _{k,p} \, \eqdef \,
\big(
\sum_{\abs{\alpha } \leq k}
\int_{\,\G} \, \abs{\PA f} ^p \, dx
\big) ^{1/p}.
\end{equation}
It fulfills all the properties of a norm,
thus the space $\SobW$ is actually a normed space,
and the mapping defined by \eqref{I_SOBOLEV_NORM}
is usually denoted a \textit{Sobolev norm}.
With respect to the question of completeness,
it is not hard to show that
the spaces $\SobW$ are complete
and therefore constitute Banach spaces as well.

\subparagraph{} %- alternative definition of Sobolev space
There is a second way to define Sobolev space.
Consider the subspace
\[
C_{k,p}^{\infty}(\G) \, \eqdef \, \{\, f \in \SF: \abs{f} _{k,p} < \infty \,\},
\]
and since this space is \textit{not} complete
{\wrt} the norm $\abs{\,\cdot\,} _{k,p}$,
we set
\[
\cal{W}^{k,p}(\G) \, \eqdef \, \overline{C_{k,p}^{\infty}(\G)}.
\]
The space $\cal{W}^{k,p}(\G)$ is defined by a purely topological operation,
that is,
by taking the closure of a smaller space.
Thus for elements in this space there is no instrinsic notion of derivative,
and one introduces the notion of \textit{strong derivative} as follows:
given $u \xeof \cal{W}^{k,p}(\G)$
and a multi-index $\alpha$ ($\abs{\alpha} \leq k$),
a further function $\PA u$ also belonging to $\cal{W}^{k,p}(\G)$
is said to be the $\alpha$-strong derivative of $u$,
if there is a sequence $\{\phi_n\} \xeof C_{k,p}^{\infty}(\G)$
converging to $u$ such that
\[
\abs{\PA u - \PA \phi_n)} _{k,p} \xto 0 \quad \quad (n \xto \infty).
\]
Sobolev spaces are reflexive for $1 < p < \infty$,
but not for $p \eq 1$.

\paragraph{} %- Meyers-Serrin theorem
The \textsc{Meyers-Serrin} \textit{theorem} claims
$\cal{W}^{k,p}(\G) \cong \SobW$ in the sense of Banach space isomorphisms.
Moreover,
the notions of strong and weak derivative are equivalent
when identifying functions in $\cal{W}^{k,p}(\G)$ and $\SobW$.

\paragraph{} %- embedding theorems for Sobolev spaces
One of the most important properties of Sobolev spaces is given by the fact
that under certain hypotheses they can continuously be embedded
into further spaces and those of type $C^k(\clG)$.
We define for fixed $1 \leq p < \infty$
the \textit{Sobolev numbers} $\theta_{\,k,d,p}$ by
\[
\theta_{\,k,d,p} \eqdef k - d/p
\quad \quad (k,d \in \N, p \geq 1),
\]
which play a crucial role in \textsc{Sobolev}\textit{'s embedding theorem}:

\begin{theorem1}\label{I_SOBOLEV_EMBEDDING_THEOREM_1}
Let $\G \subset \Rd$ be a domain having Lipschitz-continuous boundary.
If $k,m$ are chosen such that $\theta_{\,k,d,p} > m$,
then there are continuous embeddings
$\SobW \hookrightarrow C^m(\clG)$ and
$\SobW \hookrightarrow C^{m,\gamma}(\clG)$ for $0 < \gamma < 1$.
\end{theorem1}

Hence each class $f \xeof \SobW$ has a representative in these spaces,
supposed the Sobolev number is sufficiently large.

\paragraph{} %- embeddings of Sobolev spaces into other Sobolev spaces
In case of $\theta_{\,k,d,p} \leq 0$,
however,
one cannot embed $\SobW$ into $C(\clG)$.
But it is at least possible to find an embedding of $\SobW$
into some lower-dimensional Lebesgue spaces $L^q(\G)$:

\begin{theorem1}\label{I_SOBOLEV_EMBEDDING_THEOREM_2}
For $\G \subset \Rd$,
$k \xeof \N$ and $p,q \geq 1$ such that
\[
1 < q < \frac{dp}{d-kp},
\]
there are continuous embeddings $\SobW \hookrightarrow L^q(\G)$.
\end{theorem1}

\subparagraph{} %- embeddings for fractional Sobolev spaces
Finally,
there are also embeddings of Sobolev spaces into further Sobolev spaces,
where the parameter $p$ is fixed:

\begin{theorem1}\label{I_SOBOLEV_EMBEDDING_THEOREM_3}
If $\G \subset \Rd$ has Lipschitz-continuous boundary,
$s,s' \xeof \N$ and $1 < s < s' < \infty$,
then $W^{s',p}(\G) \hookrightarrow W^{s,p}(\G)$.
\end{theorem1}

\item[(b)] Let us consider in the family $\SobW$ of Sobolev spaces
the case $p \eq 2$ and denote the space $H^k(\G) \eqdef W^{k,2}(\G)$
for $k \xeof \N_0$,
which is often called a \textit{Hilbert-Sobolev space}.

\subparagraph{} %- inner product and norm of Hilbert-Sobolev spaces
Each space $H^k(\G)$ can be endowed with the inner product
\[
\scalar{f}{g}_k \, \eqdef \,
\sum_{\abs{\alpha} \leq k} \int_{\,\G} \;
\PA f \; \PA \overline{g} \, dx\,,
\quad \quad f,g \xeof H^k(\G)
\]
and with norm
$\abs{f} _k \eq \scalar{f}{f}^{1/2}$ for $f \xeof H^k(\G)$,
so $H^k(\G)$ forms a Hilbert space in its own right.

\item[(c)] The space $\TF$ of test functions can surely be endowed
with the norms $\abs{\cdot} _{k,p}$,
but this leads to uncomplete normed spaces.
Hence it is nearby to introduce the spaces $\WN{k,p}{\G}$ for $k \xeof \N$
by completion of the space $\TF$ {\wrt} the Sobolev norms:
\[
\WN{k,p}{\G} \, \eqdef \, \overline{\TF_{\abs{\cdot} _{k,p}}}.
\]
The elements of these spaces can be identified with functions in $\SobW$.
Furthermore,
the spaces $\WN{k,p}{\G}$ form closed linear subspaces
of the Sobolev spaces $\SobW$.
Clearly,
the spaces $\HN{k}{\G} \eqdef \HN{k,2}{\G}$ are Hilbert spaces
with inner products and norms as defined for $H^k(\G)$.
\end{itemize}
\end{itemize}

\begin{comment}
\paragraph{}
Further normed linear spaces like Bochner spaces
and Sobolev spaces will be introduced later
after having available the notion of \textit{adjoint space} (section III)
and methods for integration and differentiation of vector-valued functions
(section V).
\end{comment}

\section{Hierarchy of Topological and Normed Linear Spaces}
%------------------------------------------------------------------------------%
\paragraph{} %- hierarchy of spaces
The following figure shows the inheritance tree
of topological and normed linear spaces,
where Banach and Hilbert spaces are positioned at the very end
of this tree.
For the notion of \textit{reflexive space} we refer to section III.

\begin{center}
\scalebox{0.67}{ % optimum = 0.69

\begin{tikzpicture}[node distance=4.0cm,auto]

\node[align=center,rectangle,draw](TLS) at (0,-1.0)
{\;topological linear space\; \\ $(\topV,\topT_V)$};

\node[align=center,rectangle,draw](LCTLS) at (-5,-3.5)
{locally convex \\ \;topological linear space\; \\ $(\topV,\{ p_{\alpha} \}_{\alpha \in \cal{I}})$};

\node[align=center,rectangle,draw](SCTLS) at (5,-3.5)
{\textbf{sequentially complete}\\ \; \textbf{topological linear space} \; \\ $(V,\topT_V)$};

\node[align=center,rectangle,draw](MS) at (-3,-5.9)
{metric space \\ $(\metrS,\topT_d)$};

\node[align=center,thick,rectangle,draw](CMS) at (1,-5.9)
{\textbf{complete} \\ \; \textbf{metric Space} \; \\ $(\metrS,\topT_d)$};

\node[align=center,thick,rectangle,draw](F-S) at (5,-5.9)
{\; \textbf{F-space}\; \\ $(F,\topT_V=\topT_d)$};

\node[align=center,rectangle,draw](PFS) at (-5,-8.0)
{\;pre-{\FRECHET} space\; \\ $(F,\{p_n\}_{n \in \N})$};

\node[align=center,thick,rectangle,draw](FS) at (5,-8.0)
{\;\textbf{{\FRECHET} space}\; \\ $(F,\{p_n\}_{n \in \N})$};

\node[align=center,rectangle,draw](NS) at (-5,-10.5)
{\;{normed linear space}\; \\$(\banX,\norm{\cdot})$};

\node[align=center,thick,rectangle,draw](BS) at (5,-10.5)
{\;\textbf{Banach space}\; \\$(\banX,\norm{\cdot})$};

\node[align=center,thick,rectangle,draw](RBS) at (5,-12.0)
{\;\textbf{reflexive space}\; \\$(\banX,\norm{\cdot})$};

\node[align=center,thick,rectangle,draw](UCS) at (5,-13.5)
{\;\textbf{uniformly convex space}\; \\$(\banX,\norm{\cdot})$};

\node[align=center,rectangle,draw](IPS) at (-5,-15)
{\;{inner product space}\; \\$(\vecH,\scalar{\cdot}{\cdot})$};

\node[align=center,thick,rectangle,draw](HS) at (5,-15)
{\;\textbf{Hilbert space}\; \\$(\hilH,\scalar{\cdot}{\cdot})$};

\draw[->,>={Stealth}, thick] (MS) to (CMS);
\draw[->,>={Stealth}, thick] (MS) to (PFS);
\draw[->,>={Stealth}, thick] (CMS) to (F-S);
\draw[->,>={Stealth}, thick] (TLS) to (SCTLS);
\draw[->,>={Stealth}, thick] (TLS) to (LCTLS);
\draw[->,>={Stealth}, thick] (SCTLS) to (F-S);
\draw[->,>={Stealth}, thick] (F-S) to (FS);
\draw[->,>={Stealth}, thick] (LCTLS) to (PFS);
\draw[->,>={Stealth}, thick] (PFS) to (NS);
\draw[->,>={Stealth}, thick] (NS) to (IPS);
\draw[->,>={Stealth}, thick] (PFS) to (FS);
\draw[->,>={Stealth}, thick] (NS) to (BS);
\draw[->,>={Stealth}, thick] (IPS) to (HS);
\draw[->,>={Stealth}, thick] (FS) to (BS);
\draw[->,>={Stealth}, thick] (BS) to (RBS);
\draw[->,>={Stealth}, thick] (RBS) to (UCS);
\draw[->,>={Stealth}, thick] (UCS) to (HS);

\node[align=center](I) at (0,0.4) {\textit{metric and topological linear spaces}};
\node[align=center](II) at (   0,-16.4) {\textit{normed linear spaces}};
\draw[dashed,thick]        (-7.8,  0.0) rectangle (8.0,-9.0);
\draw[dashed,thick]        (-7.8, -9.5) rectangle (8.0,-16.0);
\end{tikzpicture}
}

\footnotesize{Figure \thesection.1:
              The hierarchy of metric and topological linear spaces}
\end{center}

% BIBLIOGRAPHY
%------------------------------------------------------------------------------%
\vspace{5mm}
\begin{thebibliography}{99}
\begin{small}

\bibitem{I_DS} N.\,Dunford, J.\,T.\,Schwartz:
\textit{Linear Operators I}.
Wiley \& Sons, 1958.

\bibitem{I_RS1} M.\,Reed, B.\,Simon:
\textit{Modern Methods of Mathematical Physics I: {\FA}},
Academic Press Inc., 1980.

%\bibitem{I_Yo} K.\,Yosida:
%\textit{Functional Analysis}.
%Springer, 1980.

\end{small}
\end{thebibliography}

%------------------------------------------------------------------------------%
%  II. LINEAR OPERATORS IN BANACH SPACES ---
%------------------------------------------------------------------------------%
%  §1. Fundamentals of Operator Theory ---
%  §2. Continuous Linear Operators ---
%  §3. Compact Linear Operators ---
%  §4. Closed and Closable Linear Operators ---
%  §5. The Closed Graph Theorem ---
%  ---

\newpage
\thispagestyle{empty}
\def\TitleII{Linear Operators in Banach Spaces}
\fancyhead[C] {\small{\textsc{II. \TitleII}}}
\part{\TitleII}
%------------------------------------------------------------------------------%
The aim of this second part
is to present some fundamental facts on linear operators
acting between Banach spaces.
Furthermore,
we want to introduce various classes of linear operators
by means of the validness of certain topological properties,
which to the end lead to the concept of
\textit{continuous},
\textit{compact},
\textit{closed} and
\textit{closable} operator.
It is also our intention to point out certain relationships
between these notions and to establish some equivalences.

\section{Fundamentals of Operator Theory}
%------------------------------------------------------------------------------%
Let $\banX$ and $\banY$ be normed linear spaces
(or more general topological linear spaces).
By a \textit{linear operator} between $\banX$ and $\banY$
is meant a mapping
$\linA \: \domain{\linA} \xsubset \banX \xto \banY$,
where the \textit{domain of definition} $\domain{\linA}$
forms a linear subspace of $\banX$ and
$\linA$ is \textit{linear},
that is,
\[
\linA(\lambda x + \mu y) \= \lambda \, \linA x + \mu \, \linA y
\]
holds for all $x,y \xeof \domain{\linA}$ and $\lambda,\mu \xeof \F$.
The space $\banX$ is denoted the \textit{domain space} of $\linA$,
whereas $\banY$ is said to be the \textit{range space} of $\linA$.

\subparagraph{} %- choices of the domain space
The domain of definition $\domain{\linA}$ may of course
coincide with the entire space $\banX$,
but it is often sufficient to require that 
$\domain{\linA}$ lies dense in $\banX$,
so each $x \xeof \banX$ can be approximated
by a sequence of elements in $\domain{\linA}$.
It is usual to say that $\linA$ is \textit{densely defined} in $\banX$
in this case.
However,
if the set $\domain{\linA}$ is \textit{not} dense in $\banX$,
it is nearby to consider the closure
$\banX_0 {\,:=\,} \overline{\domain{\linA} } \xsubset \banX$ 
and to pay attention to the slightly modified but densely defined operator
$\linA_0\:\domain{\linA} \xsubset \banX_0 \xto \banY$.

\paragraph{} %-- sums and compositions of linear operators
Rather the domains of definition for operators $\linA_1$ and $\linA_2$
with same domain space $\banX$ and range space $\banY$
may be proper linear sub\-spaces of $\banX$,
the \textit{sum} $\linA_1 + \linA_2$ and
the \textit{composition} $\linA_1 \circ \linA_2$ are well-defined,
if their domains of definition are given by
\[
\begin{cases}[1.2]
\quad 
\domain{\linA_1 + \linA_2} \, \eqdef \,
\domain{\linA_1} \, \cap \, \domain{\linA_2},\\
\quad
\domain{\linA_1 \circ \linA_2} \, \eqdef \,
\{\, x \in \domain{\linA_2}: \linA_2 \, x \in \domain{\linA_1}\,\}.\\
\end{cases}
\]
Furthermore,
the operator $\lambda \linA$ is well-defined by
$\domain{\lambda \linA} \eqdef \domain{\linA}$ and
$(\lambda \linA) x \eqdef \lambda (\linA x)$
for $x \xeof \domain{\lambda \linA}$,
where $\lambda \xeof \F$ is arbitrary.

\subparagraph{} %-- extensions of linear operators
If the domain of definition of a linear operator
is a proper subset of the domain space,
the concept of \textit{extension} is meaningful.
Let $\linA_1$ and $\linA_2$ be linear operators
with domain of definitions $\domain{\linA_1} \xsubset \banX$
and $\domain{\linA_2} \xsubset \banX$,
respectively.
We introduce the partial ordered relation
\[
\linA_1 \prec \linA_2
: \Longleftrightarrow 
\domain{\linA_1} \xsubset \domain{\linA_2}
\, \and \,
\linA_1 \, x = \linA_2 \, x
\fa x \in \domain{\linA_1}
\]
and denote $\linA_2$ an \textit{extension} of $\linA_1$
if $\linA_1 \prec \linA_2$.
In addition,
the operators $\linA_1$ and $\linA_2$ are said to be \textit{equal}
(briefly $\linA_1 \eq \linA_2$)
iff $\linA_1 \prec \linA_2$ and $\linA_2 \prec \linA_1$,
or equivalently,
if $\domain{\linA_1} \eq \domain{\linA_2}$ and
$\linA_1\,x \eq \linA_2\,x$ for all $x \xeof \domain{\linA_1}$.

\subparagraph{} %-- kernels, ranges and cokernels
Moreover,
to every linear operator $\linA\:\domain{\linA} \xsubset \banX \xto \banY$
there are basically assigned three linear subspaces,
namely

\begin{itemize}[leftmargin=12pt]
\item[1.]
the \textit{kernel} $\kernel{\linA} \eq
\{\,x \in \domain{\linA}: \linA x \eq \0\,\} \subseteq \banX$;

\item[2.]
the \textit{range} $\range{\linA} \eq
\{\, \linA x: x \in \domain{\linA}\,\} \subseteq \banY$;

\item[3.]
the \textit{cokernel} $\cokernel{\linA} \subseteq \banY$,
{\ie} the linear-algebraic complement of $\range{\linA}$
within the space $\banY$,
hence
$\range{\linA} \oplus \cokernel{\linA} \eq \banY$.
\end{itemize}

\paragraph{} %- classification by topological properties
A first classification of linear operators
between a domain space $\banX$ and a range space $\banY$
can be performed by checking certain properties,
which are defined by means of the norm topologies in these spaces.
It leads in conclusion to the concepts of \textit{continuous}, \textit{compact},
\textit{closed} and \textit{closable} operator.

\section{Continuous Linear Operators}
%------------------------------------------------------------------------------%
A linear operator $\linA \: \domain{\linA} \xsubset \banX \xto \banY$
is called \textit{continuous},
if the preimage $\linA^{-1}(\openO)$ of each open set $\openO \xsubset \banY$
is open in the domain space $\banX$,
which of course is the well-known definition of continuity
for mappings between arbitrary topological spaces.

\subparagraph{} %- continuity in normed linear spaces
For the mappings between normed linear spaces,
however,
we may replace continuity by \textit{sequential continuity},
which means that a mapping $f$ is continuous in $x \xeof \banX$
if for each convergent series $(x_n)$ with limit element $x \xeof \banX$
we have $\lim_{n \to \infty} f(x_n) \eq f(x)$.
The equivalence is especially true
if $f$ forms a linear operator $\linA$
and the domain of definition $\domain{\linA}$ is a proper subset of $\banX$.

\paragraph{} %- B.L.T. theorem
If $\banX,\,\banY$ are normed,
the domain is dense in $\banX$ and $\linA$ is continuous,
then $\linA$ can uniquely be extended to a linear mapping
defined on the closure $\overline{\domain{\linA}} \eq \banX$,
which proves to be continuous as well.
This is the assertion of the \textit{B.L.T.\,theorem}
(written out the \textit{bounded linear transformation theorem},
see \cite{II_RS1}, page 6).
It is given by the fact that
any continuous linear operator $\linA$ is even \textit{uniformly} continuous,
that is,
for each $\epsilon > 0$ there exists $\delta > 0$,
such that
$\norm{x-y}_X < \delta$ implies $\norm{\linA x - \linA y}_Y < \epsilon$
for \textit{all} $x,y \in \domain{\linA}$,
hence $\delta$ can be chosen independently of $x$ and $y$.

\subparagraph{}
Hence we may basically assume
that if $\linA$ is continuous,
then $\domain{\linA} \eq \banX$ by applying this canonical extension,
so $\linA$ appears as a mapping $\linA\:\banX \xto \banY$.

\paragraph{} %- space of continuous linear operators
Let us turn to the next concept of importance.
A linear operator $\linA$
acting between normed linear spaces $\banX$ and $\banY$
is called \textit{bounded}
\iff it is \textit{norm-bounded},
that is,
if there is $c \geq 0$
such that $\norm{\linA x} \leq c \cdot \norm{x}$ for each $x \xeof \banX$.
In this case we define the \textit{operator norm} of $\linA$ according to
\begin{equation}\label{II_OPERATOR_NORM}
\norm{\linA} \eqdef \sup_{\norm{x}=1} \norm{\linA x}.
\end{equation}

\paragraph{}
The following result on continuous operators is well-known:

\begin{theorem1}\label{II_MAIN_THEOREM_ON_CONTINUOUS_OPERATORS}
For $\linA\:\banX \xto \banY$ being a linear operator,
the subsequent assertions are equivalent to each other:

\begin{itemize}[leftmargin=9mm]
\item[(a)] $\linA$ is continuous;
\item[(b)] $\linA$ is uniformly continuous;
\item[(c)] $\linA$ is sequentially continuous in every $x \xeof \banX$;
\item[(d)] $\linA$ is sequentially continuous at the origin $\0 \xeof \banX$;
\item[(e)] $\linA$ is bounded.
\end{itemize}
\end{theorem1}

\paragraph{} %- the operator algebra B(X)
Let $\BO{\banX,\banY}$ be the linear space
of all continuous linear operators $\linA\:\banX \xto \banY$,
which forms a subspace of $L(\banX,\banY)$,
the space of \textit{all} linear but not necessarily continuous mappings.
The space $\BO{\banX,\banY}$ from a normed linear space $\banX$
to a Banach space $\banY$
and endowed with the norm given by \eqref{II_OPERATOR_NORM} is complete,
since $\banY$ is supposed to be complete,
hence $\BO{\banX,\banY}$ forms a Banach space in its own right.

\subparagraph{}
In addition,
in case of $\banY \eq \banX$
the space $\BO{\banX} := \BO{\banX,\banX}$
constitutes a \textit{normed algebra}
with composition $(\linA_1,\linA_2) \mapsto \linA_2 \circ \linA_1$
as its multiplicative operation
and owns the property
\[
\norm{\linA_1 \circ \linA_2} \, \leq \, \norm{\linA_1} \cdot \norm{\linA_2}.
\]
It is called the \textit{operator algebra}
associated with $\banX$.

\section{Compact Linear Operators}
%------------------------------------------------------------------------------%
A linear operator $\compA \: \banX \xto \banY$
acting between normed linear spaces $\banX,\,\banY$
is called \textit{compact},
if the image of each bounded set under $\compA$ is relative compact
{\wrt} the norm topology in $\banY$,
or equivalently,
if $\{x_n\}_{n \in \N}$ is a bounded sequence in $\banX$,
then the sequence $\{\compA \, x_n\}_{n \in \N}$ in $\banY$
owns a convergent subsequence.

\subparagraph{}
For example,
a linear operator $\compA \: \banX \xto \banY$ is compact,
if at least one of the spaces $\banX$ and $\banY$ is finite-dimensional.
In case of $\banY \eq \banX$,
the identity operator $\1_X\:\banX \xto \banX$ is compact
iff $\banX$ is finite-dimensional.

\paragraph{} %- properties of compact operators
Further properties of compact linear operators are:

\begin{itemize}[leftmargin=17pt]
\item[(a)]
every compact operator is continuous
(the opposite assumption leads to a contradiction);

\item[(b)]
the composition $\linA_1 \circ \linA_2$ of continuous operators is compact,
if one of them is so;

\item[(c)]
finite linear combinations of compact linear operators are compact, too;

\item[(d)]
the limit of a norm convergent sequence of compact operators is compact.

\item[(e)]
the first and second property together imply that
the set $\CO{\banX}$ of compact operators $\compA\:\banX \xto \banX$
is a \textit{closed} ideal within the operator algebra $\BO{\banX} $,
hence it is complete and forms a Banach algebra in its own right;
\end{itemize}

\paragraph{} %- finite-rank operators
A linear operator $\linFR \: \banX \xto \banY$ is said
to have \textit{finite rank},
if the range $\range{\linFR}$ is finite-dimensional.
It is obvious that $\linFR$ is compact.

\subparagraph{}
Suppose that $(\linFR_n)$ is a sequence of finite-rank operators
converging to the operator $\linA$ {\wrt} the norm of $\BO{\banX,\banY}$.
Then $\linA$ is compact
and the set $\CO{\banX,\banY}$ of compact linear operators
proves to be a subset of the closure of the set $\FR{\banX,\banY}$
of finite-rank operators,
{\ie}
$\CO{\banX,\banY} \subset \overline{\FR{\banX,\banY}}$.

\subparagraph{}
Reversely,
in the setting of Hilbert spaces 
each compact linear operator $\compA$ appears in this way,
that is,
$\compA$ is the norm-limit of a sequence of finite-rank operators
(but this is not necessarily true in the general setting of Banach spaces).

\paragraph{} %- Hilbert Schmidt operators
Let $\basB \eq \{e_1,e_2,\ldots\}$ be any orthonormal base
of a \textit{separable} Hilbert space $\hilH$ and let $\linA \xeof \BO{\hilH}$.
We define the number $S(\linA)$ according to
\[
S(\linA) \eqdef \sum_{i \in \basB} \, \norm{ \linA \, e_i}^2.
\]
Then it turns out that this definition for $S(\linA)$
does not depend on the choice of the orthonormal base.
In case of $S(\linA) < \infty$
the operator $\linA$ is called a \textsc{Hilbert-Schmidt} \textit{operator},
and the number
\[
\norm{\linA}_{HS} \eqdef \sqrt{S(\linA)}
\]
is denoted the \textit{Hilbert-Schmidt} \textit{norm} of $\linA$,
since the mapping $\linA \mapsto \norm{\linA}_{HS} $
owns all the properties of a norm.

\subparagraph{}
The set $\HS{\hilH}$ of {Hilbert-Schmidt} operators in $\hilH$
is topologically and algebraically closed in $\BO{\hilH}$
{\wrt} the operations $\linA_1 + \linA_2$ and $\lambda \linA$,
hence $\HS{\hilH}$ endowed with norm
$\norm{\,\cdot\,}_{HS}$ forms a Banach space in its own right.
Each {Hilbert-Schmidt} operator proves to be compact,
so $\HS{\hilH}$ forms a two-sided ideal within the space $\CO{\hilH}$
of compact operators acting in $\hilH$.

\section{Closed and Closable Linear Operators}
%------------------------------------------------------------------------------%
A linear operator $\linA \: \domain{\linA} \xsubset \banX \xto \banY$
is called \textit{graph-closed} (or shortly \textit{closed}) if the set
\[
\Gamma_{\linA} \eqdef
\{\,(x,\linA x): x \xeof \domain{\linA}\,\} \; \xsubset \; \banX \xtimes \banY,
\]
denoted the \textit{graph} of $\linA$,
is closed in the space $\banX \xtimes \banY$
{\wrt} the product topology $\topT_{\,\banX \xtimes \banY}$.
This property especially implies that
the kernel $\kernel{\linA}$ of $\linA$ forms a closed set in $\banX$.

\subparagraph{} %- sequentially closed operators
It is a matter of fact
that an operator $\linA$ is closed iff it is \textit{sequentially closed},
that is,
$x_n \xto x$ and $\linA x_n \xto y$ implies
$x \xeof \domain{\linA}$ and $\linA x \eq y$,
briefly written
\[
\begin{rcases}
 x_n       \xto x \quad \\
 \linA x_n \xto y \quad \\
\end{rcases}
\quad \Longrightarrow \quad
x \xeof \domain{\linA} \, \and \, \linA x \eq y.
\]
This can easily proved by the fact
that a subset of a normed linear space is closed
\iff it is sequentially closed.
Checking the sequentially closeness is the usual method
to verify that a given linear operator is graph-closed.
But notice that the equivalence of these both notations is in general not valid,
if $\banX$ and $\banY$ are non-metrizable topological linear spaces.

\subparagraph{} %- graph norm
It is often helpful to make the domain of definition $\domain{\linA}$
of a linear operator $\linA$ into a normed linear space by its own.
This can be done by introducing the so-called \textit{graph norm}
\[
\norm{x}_{\linA} \eqdef \norm{x}_X \+ \norm{\linA x}_Y
\]
for $x \xeof \domain{\linA}$.
In fact,
the pair $(\domain{\linA},\norm{\,\cdot\,}_{\linA})$
forms a normed linear space,
which constitutes a Banach space \iff 
$\linA$ is sequentially closed.
This equivalence provides a further method
to verify closeness of a given linear operator.

\paragraph{} %- closeness of operators based on Schauder estimates
The following theorem
uses the possible validness of an abstract \textsc{Schauder} \textit{estimate}
to ensure sequentially closenss:

\begin{theorem1}\label{II_BROWDER_THEOREM}
Let $\linA_0 \: \banX \xto \banY$ be a continuous linear operator,
where the Banach space $\banX$ is continuously and densely embedded
into the Banach space $\banY$ by a mapping $\iota \: \banX \xto \banY$.
If the estimate
\begin{equation}\label{II_SCHAUDER_ESTIMATE}
\norm{x}_{\banX} \, \leq \,
c \cdot
\{ \,\norm{\iota x}_{\banY} + \norm{\linA_0 \, x}_{\banY} \, \}
\end{equation}
is fulfilled for some $c {\,>\,} 0$ and all $x \xeof \banX$,
then the linear operator $\linA$ acting in $\banY$ with
$\domain{\linA} := \iota(\banX)$ and
$\linA\, \iota x := \linA_0 \, x$ for $x \xeof \banX$ is closed.
\end{theorem1}

\subparagraph{} %- proof
In fact,
let $(x_n)_{n \in \N}$ be a sequence in $\domain{\linA}$
with $\norm{x_n-x}_{\banX} \to 0$
and $\norm{\linA \,x_n-y}_{\banY} \to 0$.
It is to verify that $x \xeof \domain{\linA}$ and $\linA x \eq y$.
Inequality \eqref{II_SCHAUDER_ESTIMATE} implies that
$(x_n)_{n \in \N}$ is a Cauchy sequence in $\banX$
with limit element $x \xeof \banX$.
Since $\linA_0$ is continuous,
the element $y$ is also the limit
of the sequence $(\linA \, \iota x_n)_{n \in \N}$,
hence $\linA$ is closed.

\paragraph{} %-- pre-closed operators
Many linear operators acting between Banach or Hilbert spaces are not closed,
but can be extended to operators being closed.
We shall call such operators \textit{closable} or \textit{pre-closed}.
If a linear operator $\linA$ is closable,
then there is a \textit{minimal closed extension} $\linA_c$ of $\linA$,
called the \textit{closure} of $\linA$.
This mapping $\linA_c$ is given by the fact
that each other closed extension of $\linA$ extends $\linA_c$,
and $\linA_c$ is nothing but the intersection
of all closed extensions of $\linA$.
Clearly,
each closed operator is closable,
since it constitutes a closed extension of itself.
Furthermore,
given a closed linear operator $\linA$
with domain of definition $\domain{\linA}$,
any linear subspace $\cal{L} \xsubset \domain{\linA}$
is called a \textit{core} of $\linA$,
if the closure of $\linA$ restricted to $\cal{L}$ resembles $\linA$.

\paragraph{} %- criterion for closability
There is a well-known criterion for closability:

\begin{theorem1}\label{II_CLOSABILITY_THEOREM}
A linear operator $\linA$ with domain of definition $\domain{\linA}$
is closable,
if each sequence $(x_n,\linA x_n)_{n \in \N}$
with $x_n \xeof \domain{\linA}$ for $n \xeof \N$ and converging to $(\0,y)$
implies $y \eq \0$,
briefly written
\[
\begin{rcases}
 x_n       \xto \0 \quad \\
 \linA x_n \xto  y \quad \\
\end{rcases}
\quad \Longrightarrow \quad
y \= \0.
\]
\end{theorem1}

\begin{comment}
\subparagraph{}
\Proof:
If $\linA$ is closable,
then $\overline{\Gamma_{\linA}}$ is the graph of the closure $\overline{\linA}$
and forms a closed subspace in $\banX \xtimes \banY$.
Let $x_n \to 0$ and $\linA x_n \to y$,
so $(x_n,\linA x_n)$ converges to $(\0,y)$,
which implies $y \eq \0$.

\subparagraph{}
On the other hand,
let $\linA$ own the following property:
$x_n \to \0$ with $x_n \xeof \domain{\linA}$
and $\linA x_n \to y$ implies $y \eq \0$.
We introduce the subspace $D$ consisting of all
$x \xeof \banX$,
for which there is $x_n \xeof \domain{\linA}$ and $y \xeof \banX$
such that $x_n \to x$ and $\linA x_n \to y$.
Define $\overline{\linA} x \eqdef y$ for $x \xeof D$,
then $\domain{\linA} \xsubset D$ and $\linA \prec \overline{\linA}$.
It is still to show that

\begin{itemize}[leftmargin=7mm]
\item[(1)] $\overline{\linA}$ is well defined on $D$ and
\item[(2)] $\overline{\linA}$ is closed, hence $\linA$ is closable.
\end{itemize}

\subparagraph{}
To verify (1),
let $(\tilde{x}_n)$ be a further sequence converging to $x$ with
$\linA \tilde{x}_n \to \tilde{y}$.
But since $x_n - \tilde{x}_n \to \0$
and $\linA (x_n - \tilde{x}_n) \to y-\tilde{y}$,
we have $y \eq \tilde{y}$.
Moreover,
it is easy to verify that $\overline{\linA}$ is linear on $D$.

\subparagraph{}
To verify (2),
let $\{x_n\}_{n \in \N}$ be a sequence in $D$ such that $x_n \to x$
and $\overline{\linA} \to y$.
We have to verify that $x \xeof D$ and $\overline{\linA} x \eq y$.
Since for arbitrary $\epsilon > 0$ and $x_n \in D$
there is $z_n \xeof \domain{\linA}$ such that
$\norm{z_n - x_n} < \frac{\epsilon}{2}$ and
$\lVert \overline{\linA} x_n - \linA z_n \rVert < \frac{\epsilon}{2}$,
we obtain
\[
\norm{x_n - z_n} \, + \,
\lVert \overline{\linA} x_n - \linA z_n \rVert
\, < \, \epsilon.
\]
Hence $z_n \to x$ and $\linA z_n \to y$,
that is,
$x \in D$ and $\overline{\linA} x \eq y$.
\qed
\end{comment}

\paragraph{} %- comparison of continuous and closed operators
Let us emphasize that in general 
the concepts of closeness and continuity of a linear operator $\linA$
are not comparable.
However,
if $\banX$ and $\banY$ are Banach or Hilbert spaces
and $\domain{\linA}$ is a \textit{closed} set in $\banX$,
then closeness of $\linA$ is an easy consequence of continuity of $\linA$.

\subparagraph{}
This assertion is especially true in the following situation:
if $\linA$ is densely defined and continuous,
then it can be extended to a continuous linear operator
$\linA_{ext} \: \banX \xto \banY$,
and since $\banX$ forms a closed linear space,
the extension $\linA_{ext}$ is closed,
too.
Hence continuous linear operators of the form $\linA \: \banX \xto \banY$
fit in a natural way
into the class of densely defined and closed linear operators.

\section{The Closed Graph Theorem}
%------------------------------------------------------------------------------%
Let us first emphasize that in general 
the concepts of closeness and continuity of a linear operator $\linA$
are not comparable.
However,
if $\banX$ and $\banY$ are Banach or Hilbert spaces
and $\domain{\linA}$ is a \textit{closed} set in $\banX$,
then closeness of $\linA$ is an easy consequence of continuity of $\linA$.
This fact can be applied to the following situation:
if $\linA$ is densely defined and continuous,
then it can be extended to a continuous linear operator
$\linA_{ext} \: \banX \xto \banY$,
and since $\banX$ forms a closed linear space,
the extension $\linA_{ext}$ is closed,
too.
Hence continuous linear operators of the form $\linA \: \banX \xto \banY$
fit in a natural way
into the class of densely defined and closed linear operators.

\paragraph{} %- closed graph theorem
For the rest of this paragraph it is assumed that the spaces $\banX,\,\banY$
are Banach.
The \textit{closed graph theorem} claims that
under a certain condition the reverse statement is also true:

\begin{theorem1}\label{CLOSED_GRAPH_THEOREM}
If the linear subspace $\domain{\linA}$ is closed in $\banX$,
then closeness of $\linA$ implies continuity,
so both concepts are equivalent to each other in this case.
\end{theorem1}

\paragraph{} %- bounded inverse theorem
The following application of the closed graph theorem
is of outstanding meaning:

\begin{theorem1}\label{BOUNDED_INVERSE_THEOREM}
If a linear operator $\linA$ is closed and bijective,
then its inverse $\linA^{-1}$ is also linear and closed,
and since the domain of definition $\domain{\linA^{-1}}$
is the entire space $\banY$,
the closed graph theorem implies that $\linA^{-1} \: \banY \xto \banX$
is continuous\footnote{
\,One also says that $\linA$ is \textit{invertible}
if it is bijective and $\linA^{-1}$ is continuous.}.
\end{theorem1}

\subparagraph{} %- open mapping theorem
The closed graph theorem is a consequence of
\textsc{Banach}'s \textit{open mapping theorem}:

\begin{theorem1}\label{OPEN_MAPPING_THEOREM}
If $\linA$ is continuous and surjective,
then the image $\linA(\openO)$ of an open set $\openO \subset \banX$
is open in $\banY$.
Moreover,
if $\linA$ is injective,
then $\linA^{-1}$ is continuous
and constitutes a linear homeomorphism between the spaces $\banX$ and $\banY$
(this is the \textit{bounded inverse theorem} for continuous operators).
\end{theorem1}

The closed graph theorem,
the open mapping theorem and
the bounded inverse theorem
are equivalent to each another
and can be derived from \textit{\textsc{Baire}'s category theorem},
which is a fundamental result in general topology.

\subparagraph{} %- summary
In conclusion,
if the domain of definition $\domain{\linA}$ of a linear operator $\linA$
is a \textit{closed} subspace of the domain space $\banX$,
then the following chain of implications is valid:
\[
\linA \mbox{ compact }     \Longrightarrow
\linA \mbox{ bounded }     \Longleftrightarrow
\linA \mbox{ continuous }  \Longleftrightarrow
\linA \mbox{ closed}.     %\Longrightarrow \linA \mbox{ closable}.
\]

\paragraph{} %- solvability of linear equations
Regarding the solvability of $\linA \, x \eq y$
where $\linA$ is acting between Banach spaces $\banX$ and $\banY$,
the following concepts are available:
the equation $\linA \, x \eq y$ is called

\begin{itemize}[leftmargin=12mm]
\item[(a)]
\textit{densely solvable} (\cite{II_Kr}),
if $\overline{\range{\linA}} \eq \banY$;

\item[(b)]
\textit{uniquely solvable} (\cite{II_Kr}),
if $\linA$ is injective;

\item[(c)]
\textit{everywhere solvable} (\cite{II_Kr}),
if $\linA$ is surjective;

\item[(d)]
\textit{solvable} (\cite{II_Br}),
if it is uniquely and everywhere solvable;

\item[(e)]
\textit{well-posed},
if it is solvable and $\linA^{-1}$ is continuous.
\end{itemize}

\subparagraph{}
Hence well-posedness means that
solutions of the linear equations depend continuously on the right hand side.
We want to emphasize
that if the linear operator $\linA$ is continuous or closed,
then property (d) implies property (e) by the bounded inverse theorem.

% BIBLIOGRAPHY
%------------------------------------------------------------------------------%
\vspace{5mm}
\begin{thebibliography}{99}
\begin{small}

\bibitem{II_Br} F.\,Browder:
\textit{Functional Analysis and Partial Differential Equations I}.
Math. Annalen 138, pp.\,55-79, 1959.

\bibitem{II_RS1} M.\,Reed, B.\,Simon:
\textit{Modern Methods of Mathematical Physics I: {\FA}},
Academic Press Inc., 1980.

\bibitem{II_Kr} S.\,Krein:
\textit{Linear Equations in Banach Spaces}.
Birk\-h\"{a}user, 1982.

\bibitem{II_Yo} K.\,Yosida:
\textit{Functional Analysis}.
Springer, 1980.

\end{small}
\end{thebibliography}

%------------------------------------------------------------------------------%
% III. ADJOINT SPACES AND OPERATORS ---
%------------------------------------------------------------------------------%
%  §1. Continuous Linear Functionals ---
%  §2. Bidual Spaces and Reflexivity ---
%  §3. Representation Theorems for Adjoint Spaces ---
%  §4. Spaces of Distributions ---
%  §5. Adjoint Operators in Locally Convex Spaces ---
%  §6. Adjointness in Banach Spaces ---
%  §7. Hilbert Adjoints and Normal Operators ---
%  ---

\newpage
\thispagestyle{empty}
\def\TitleIII{Adjoint Spaces and Operators}
\fancyhead[C] {\small{\textsc{III. \TitleIII}}}
\part{\TitleIII}
%------------------------------------------------------------------------------%
It is often useful to combine the study of a topological or normed linear space
$\banX$ with that of its \textit{dual space} $\banX^*$,
which is nothing but the set of continuous and $\F$-valued linear mappings
endowed with a linear-algebraic structure.
Dual spaces are sometimes of interest in its own right,
since they may constitute extensions of the original space $\banX$,
which has many helpful applications in higher analysis.

\section{Continuous Linear Functionals}
%------------------------------------------------------------------------------%
Let $\banX$ be an arbitrary topological $\F$-linear space.
The underlying fields $\R$ or $\C$ surely form topological
$\F$-linear spaces again,
where the topologies are given by the so-called \textit{standard topologies}.

\subparagraph{}
The main objects in this paragraph are \textit{continuous linear functionals}
on $\banX$,
that is,
continuous mappings of the form $\phi \: \banX \xto \F$,
where $\phi$ is linear at the same time,
so
\[
\phi(\alpha x + \beta y) \= \alpha \, \phi(x) + \beta \, \phi(y)
\]
holds for all $\alpha,\beta \xeof \F$ and $x,y \xeof \banX$.
We define the mappings $\phi_1+\phi_2$ and $\alpha \phi$
pointwise for $x \xeof \banX$,
so the set $\banX^*$ of all continuous linear functionals on $\banX$
forms a $\F$-linear space,
too,
called the \textit{adjoint space} or the \textit{dual space} of $\banX$.
If the topology on $\banX$ is locally convex,
then a specific version of the celebrated \textit{Hahn-Banach} \textit{theorem}
ensures that
there are non-trivial continuous linear mappings on $\banX$,
thus we may assume $\banX^* \neq \{\0\}$.

\subparagraph{} %- conjugate linear functionals
In case of $\F \eq \C$ there are often reasons
to consider continuous mappings $\phi\:\banX \xto \F$,
which are \textit{antilinear} in the sense that 
\[
\phi(\alpha x + \beta y) \=
\overline{\alpha} \, \phi(x) + \overline{\beta} \, \phi(y)
\]
for all $\alpha,\beta \xeof \C$ and $x,y \xeof \banX$.
Let us denote $\banX_a^*$ the linear space of continuous antilinear functionals 
on $\banX$,
usually called the \textit{conjugate linear space}
or the \textit{antidual space} of $\banX$.

\section{Weak Topologies, Bidual Spaces and Reflexivity}
%------------------------------------------------------------------------------%
Let us first remember that if a linear space $\banX$ is normed,
then the adjoint space $\banX^*$ is also normed,
where the norm of an element $\phi \xeof \banX^*$ is defined by
\[
\norm{\phi} \, \eqdef \, \sup_{x \in \banX,\,\norm{x} \leq 1} \, \abs{\phi(x)} .
\]
Since the fields $\R$ and $\C$ are complete,
the {adjoint space} $\banX^*$ is complete, too,
which can easily be proved.

\paragraph{} %- weak topologies
Recall also from section I
that a normed linear space $\banX$ owns a specific topology,
called the \textit{norm topology},
which is the coarsest one to make the norm into a continuous mapping.
But there is a further and useful topology on $\banX$,
denoted the \textit{weak topology} $\topT_w$.
It is the coarsest topology on $\banX$
such that all continuous linear functionals
$\phi \xeof \banX^*$ remain continuous.
The notion of convergence {\wrt} this topology
is known as \textit{weak convergence}.
More precisely,
a sequence $\{x_n\}_{n \in \N}$ of elements in $\banX$
is called \textit{weakly convergent} to $x \xeof \banX$
(notation: $x_n \rightharpoonup x$),
if $\lim_{n \xto \infty} \phi(x_n) \eq \phi(x)$ for each $\phi \xeof \banX^*$.

\subparagraph{} %- properties of the weak topology
The weak topology is also Hausdorff.
But it does not contain more open sets than the norm topology,
so $\topT_w$ is weaker than $\topT_{\norm{\cdot}}$.
In other words,
norm convergence of any sequence in $\banX$ implies weak convergence.
If $\banX$ is finite-dimensional,
however,
both topologies coincide,
but this is in general not true in arbitrary normed linear spaces.

\paragraph{} %- bidual spaces
Let for a given normed linear space $\banX$ the \textit{bidual space}
defined by $\banX^{**} := (\banX^*)^*$,
which forms a complete normed linear space again.
Then there is a natural embedding
\[
\begin{cases}[1.2]
\quad \iota \: \banX \hookrightarrow \banX^{**},\\
\quad \iota(x) := \phi_x \for \phi_x \xeof \banX^*,\\
\end{cases}
\]
where $\phi_x(f) := f(x)$ for $f \xeof \banX^*$.
The mapping $\iota$ is linear, continuous and injective.
In other words,
each element $x \xeof \banX$ gives rise to a mapping $\phi_x \xeof \banX^{**}$,
where $\norm{x} \eq \norm{\iota_x}$ holds,
so $\banX$ can be seen as a linear subspace of $\banX^{**}$.

\subparagraph{} %- reflexive spaces
It can in general nothing be said whether the mapping $\iota$ is onto.
But if this is the case,
then $\iota$ constitutes a norm-isomorphism between $\banX$ and $\banX^*$.
Many normed linear spaces considered in functional analysis
own this remarkable property,
which is indicated by the concept of \textit{reflexivity}:
the space $\banX$ is called \textit{reflexive},
if $\iota$ is onto and consequently $\banX \cong \banX^{**}$ holds.

\paragraph{} %- properties of reflexive spaces
We collect some propositions concerning reflexivity:

\begin{itemize}[leftmargin=12pt,itemsep=5pt]
\item[1.]
Since adjoint spaces of normed spaces are throughout Banach spaces,
every reflexive space must be complete.

\item[2.]
If a Banach space $\banX$ is norm-isomorphic to a reflexive space,
then $\banX$ must be reflexive, too.

%\item[]
%All closed linear subspaces $\clU$ of a reflexive Banach space $\banX$
%and their quotient spaces $\banX / \clU$ prove to be reflexive.

\item[3.]
Dual spaces of reflexive spaces are reflexive.

\item[4.]
Reflexivity of a Banach space $\banX$ is strongly related to reflexivity
of \textit{all} its adjoint spaces,
that is,
$\banX$ is reflexive
iff
one of the spaces
$
\banX^{*}, \banX^{**}, \banX^{***}, \ldots
$
is reflexive.
On the other hand,
if $\banX$ is not reflexive,
then all these spaces prove to be pairwise unisomorphic to each other.
\end{itemize}

\subparagraph{} %- Milman-Pettis theorem
Important examples of reflexive Banach spaces are given
by the \textit{Milman-Pettis} \textit{theorem},
which claims that
\textit{each uniformly convex Banach space is reflexive}.

\paragraph{} %- weak* topologies
We finally introduce the \textit{weak*}-topology on the adjoint $\banX^*$
of a locally convex topological linear space $\banX$:
it is the smallest topology in $\banX^*$
such that
for each $x \xeof \banX$ the linear functional $\phi_x$ remains continuous.

\paragraph{} %- Alaoglu-Bourbaki theorem
The weak*-topology on the adjoint space $\banX^*$ may possibly contain
fewer open sets than the weak topology and is therefore coarser.
But $\banX^*$ is at least a Hausdorff space {\wrt} the weak*-topology.
The main result concerning weak*-topologies in adjoint spaces in the setting
of normed linear spaces is given by the \textit{Alaoglu theorem}:

\begin{theorem1}\label{IV_ALAOGLU_BOURBAKI_THEOREM}
The closed unit ball
$B_1(\0) : = \{ f \in \banX^*: \norm{f} \leq 1 \}$ is compact
{\wrt} the weak*-topology on $\banX^*$.
\end{theorem1}

\section{Representation Theorems for Adjoint Spaces}
%------------------------------------------------------------------------------%
In this paragraph we describe the adjoint spaces
for a couple of normed linear spaces $\banX$ arising in applied analysis
by quoting a linear homeomorphism $\banX^* \cong \banY$,
where the representation space $\banY$ has to be chosen in a suitable way.
Indeed,
it is one of the basic tasks in functional analysis to identify
the adjoint space $\banX^*$ of a given normed linear space $\banX$
with another known Banach space $\banY$
by finding a suitable norm-isomorphism $\Phi \: \banX^* \xto \banY$
in order to obtain $\banX^* \cong \banY$.

\begin{itemize}[leftmargin=12pt,itemsep=5pt]
\item[1.] \textit{The adjoint space of $C_c(\G)$.}\\
It is known from measure theory
that a continuous function $f\:\G \xto \R$ is measurable,
but in general not Lebesgue-integrable.
If $f \xeof C_c(\G)$,
however,
then $f$ is bounded and clearly Lebesgue-integrable.

\subparagraph{} %- Riesz-Markov theorem I
Every Borel measure $\mu$ gives rise to
a positive and continuous linear functional $\phi_{\mu}$
on the space $C_c(\G)$ endowed with the supremum norm.
Conversely,
the question arises
whether each continuous linear functional $\phi \xeof C_c(\G)^*$
appears in this way.
The answer is given by the \textsc{Riesz-Markov} \textit{theorem},
quoting
that the representation of positive linear functionals on $\G$
can just be done by the \textit{Radon measures} on $\G$:

\begin{theorem1}\label{IV_RIESZ-MARKOV_THEOREM_1}
Let $\phi$ be a positive and continuous linear functional
on the normed linear space $C_c(\G)$.
Then there is exactly one Radon measure
$\mu\:\Borel{\G} \xto [0,\infty]$ such that
\[
\phi(f) \eqdef \int_{\,\G} f \, d\mu, \quad \quad f \in C_c(\G).
\]
\end{theorem1}

\item[2.] \textit{The adjoint space of $C_0(\G)$.}\\
A \textit{complex} measure $\mu \: \Borel{\G} \xto \C$
is said to be \textit{Radon},
if the so-called \textit{variation} $\abs{\mu} $
is Radon in the sense of a positive measure on $\G$.

\subparagraph{}
Notice that $\abs{\mu} $ is always finite.
The space of complex Radon measures $\MCR{\G}$
proves to be complete {\wrt} the norm $\norm{\mu} := \abs{\mu} $
for $\mu \xeof \MCR{\G}$.

\paragraph{} %- Riez-Markov theorem II
The following theorem of \textsc{Riesz-Markov} type
describes the adjoint space of $C_0(\G)$:

\begin{theorem1}\label{IV_RIESZ_MARKOV-THEOREM_2}
Let $\phi$ be a complex-valued and continuous linear functional
defined on the space $C_0(\G)$.
Then there is exactly one complex Radon measure $\mu$ such that
\[
\phi(f) \, \eqdef \, \int_{\,\G} f \, d\mu, \quad \quad f \in C_0(\G).
\]
\end{theorem1}
In conclusion,
$C_0(\G)^* \cong \MCR{\G}$,
so $C_0(\G)^*$ is norm-isomorphic to the Banach space $\MCR{\G}$
of complex Radon measures.

\item[3.] \textit{The adjoint space of $C[a,b]$.}\\
Especially in the case $\G \eq [a,b]$ we may quote
that \textit{every} continuous linear functional $G$ on $C[a,b]$ arises
as Riemann-Stieltjes integral
{\wrt} a certain function $g$ of bounded variation,
that is,
we have $C[a,b]^* \cong BV(a,b)$.
In fact,
functions of bounded variation correspond to Radon measures
on the interval $[a,b]$,
hence the norm-isomorphism is an immediate consequence
of theorem \ref{IV_RIESZ-MARKOV_THEOREM_1}.

\item[4.] \textit{The adjoint space of $L^p(X)$.}\\
If $1 < p <\infty$ and $(X,\salgB,\mu)$ is an ordinary measure space,
we have $L^p(X)^* \cong L^q(X)$,
where $q$ is conjugated to $p$:
$p^{-1} + q^{-1} \eq 1$.
For $p \eq 1$ and $(X,\salgB,\mu)$ being a $\sigma$-finite measure space,
however,
$L^1(X)^* \cong L^{\infty}(X)$ holds.

\item[5.] \textit{The adjoint space of $L^p(a,b;\banX)$.}\\
If $1 < p <\infty$ and the Banach space $\banX$ is reflexive and separable,
we have $L^p(a,b;\banX)^* \cong L^q(a,b,\banX^*)$,
where $q$ is conjugated to $p$.

\item[6.] \textit{The conjugate space of a complex Hilbert space.}\\
Due to the properties of the inner product,
each element $x \xeof \hilH$ gives rise to
a continuous and conjugate linear functional $\phi_x \xeof \hilH_a^*$
according to
\[
\phi_x(y) \, \eqdef \, \scalar{x}{y} \qfa y \in \hilH.
\]
The mapping $\Phi\:\hilH \xto \hilH_a^*$ defined by $x \mapsto \phi_x$
is linear,
called the \textsc{Riesz} \textit{map} of $\hilH$.
Since the inner product is positive definite,
$\Phi$ turns out to be injective.

\subparagraph{} %- Riesz's representation theorem
\textit{Riesz}\textit{'s representation theorem} claims
that the mapping $\Phi$ is actually surjective,
so there is a canonical correspondence
between the elements of $\hilH$ and those of $\hilH_a^*$:

\begin{theorem1}\label{IV_RIESZ_REPRESENTATION_THEOREM}
If $\hilH$ is a complex Hilbert space,
then $\Phi\: \hilH \xto \hilH_a^*$ is bijective and norm-preserving,
{\ie}
$\norm{\Phi(x)} \eq \norm{x}$ for $x \xeof \hilH$.
\end{theorem1}

\subparagraph{}
This theorem implies that
to each $\phi \xeof \hilH_a^*$
there is a unique element $x_{\phi} \xeof \hilH$,
which fulfills the equation
$\scalar{x_{\phi}}{y} \eq \phi(y)$ for every $y \xeof \hilH$.
Moreover,
the \textsc{Riesz} map gives rise to introduce an inner product
on the conjugate space $\hilH_a^*$ by
\[
\scalar{\phi}{\psi} \, \eqdef \, \scalar{\Phi^{-1}(\phi)}{\Phi^{-1}(\psi)}
\quad \quad (\phi,\psi \in \hilH_a^*),
\]
hence $\Phi$ is not only norm-preserving,
but also \textit{isometric},
that is,
it preserves the inner products $\scalar{x}{y}$ for $x,y \xeof \hilH$.

\item[7.] \textit{The adjoint space of a real Hilbert space.}\\
On the other hand,
if $\hilH$ is a \textit{real} Hilbert space,
we may replace $\hilH_a^*$ by $\hilH^*$ and verify
that the \textsc{Riesz} map $\Phi$ is a norm-isomorphism
between $\hilH$ and $\hilH^*$,
since in this case
the inner product in $\hilH$ is also linear in the second argument.
\end{itemize}

\paragraph{}
The adjoint spaces of smooth function spaces will be treated
in the next paragraph and lead us to so-called \textit{distributional spaces}.

\section{Spaces of Distributions}
%------------------------------------------------------------------------------%
The notion of \textit{distribution} forms a generalization
of an ordinary function $f\:\G \xto \F$.
The corresponding subject in functional analysis,
called the \textit{theory of distributions},
was developed by \textsc{L.\,Schwartz} in the midth of the 20th century
to describe ideas and concepts of theoretical physics
by an exact mathematical formalism.

\subparagraph{}
Distribution theory can be seen as a direct application
of the abstract theory of locally convex linear spaces and their dual spaces.
In the following let us enumerate the most important spaces of distributions:

\begin{itemize}[leftmargin=12pt,itemsep=5pt]
\item[1.] \textit{The adjoint space of $\D(\G)$.}
Remember the space $\D(\G)$,
consisting of test functions and furnished with a locally convex topology,
in which this space is complete but not metrizable.
Then every continuous linear functional $\disF\:\D(\G) \xto \C$ 
is called a \textit{Schwartz distribution}.
We denote $\D'(\G) \eqdef \D(\G)^*$ the adjoint space of $\D(\G)$
respectively the space of Schwartz distributions on $\G$.
It is also possible to characterize a Schwartz distribution
{\wrt} $\D$-convergence,
when applied to test functions:
a linear functional $\disF\:\TF \xto \C$
is a Schwartz distribution
iff $\disF(\phi_n) \xto 0$ for any $\D$-convergent sequence $(\phi_n)$
converges to the zero function.

\subparagraph{} %- regular distributions
Notice that the linear functional $\disF_f$
derived from a locally integrable function $f$ by
\[
\disF_f(\phi) \, \eqdef \, \int_{\,\G} f \, \phi \, dx,
\quad \quad \phi \in \TF
\]
is continuous,
and this is also true for the linear functional
induced by the Radon measure $\mu$ according to
\[
\disF_{\mu}(\phi) \, \eqdef \, \int_{\,\G} \phi \, d\mu,
\quad \quad \, \phi \in \TF.
\]
Thus locally integrable functions as well as Radon measures
may be regarded as Schwartz distributions.
Notice that the space $\D'(\G)$ is indeed extremely extensive
and contains nearly all functions being important in applications,
amongst others the functions belonging to the well-known spaces
$\Lp$, $\Li$, $C(\G)$, $C(\clG)$ and $C_0(\G)$,
but also measures on $\G$,
which are regular enough in some sense.

\subparagraph{}
A Schwartz distribution is said to be a \textit{regular distribution},
if it is the continuous linear functional $\disF_f$
induced by a locally integrable function $f \xeof \Lc$,
otherwise it is called \textit{singular}.
Well-known examples of singular distributions are
the \textit{Dirac distributions}
\[
\delta_a(\phi) \, \eqdef \, \phi(a)
\qfa \phi \in \TF
\quad \quad  (a \in \G).
\]

\item[2.] \textit{The adjoint space of $\E(\G)$.}
Regarding the {\FRECHET} space $\E(\G)$,
each element $\disF \xeof \E'(\G) \eqdef \E(\G)^*$ is called
a \textit{distribution with compact support}.
This notion is founded by the fact
that for such $\disF$ there is a compact set $K_{\disF} \subset \G$,
called the \textit{support} of $\disF$
such that
$\disF \phi \eq 0$ for all $\phi \xeof \E(\G)$
having support in $\G \setminus K_{\disF}$.
On the other hand,
each measurable function $f$ with compact support
gives rise to a compactly supported distribution $\disF_f \xeof \E'(\G)$.

\subparagraph{}
Since it is always possible to increase an adjoint space
by decreasing the underlying base space,
there is a continuous inclusion $\E'(\G) \, \hookrightarrow \, \D'(\G)$,
which is dual to embedding
$\D(\G) \, \hookrightarrow \, \E(\G)$.
So we basically may regard
each compactly supported distribution as a Schwartz distribution.

\item[3.] \textit{The adjoint space of $\SRd$.}
Similar to the spaces $\D(\G)$,
we introduce the notion of \textit{tempered distribution}
as a continuous linear functional $\disT \: \SRd \xto \C$,
which generalizes that of bounded and slow-growing locally integrable functions.
The space of tempered distributions is usually denoted $\TD$.

\subparagraph{} %- comparison of distributional spaces
Comparing the known kinds of distributional spaces,
we obtain a chain of inclusions,
which is dual to that of the base spaces:
\[
\cal{E}'(\Rd) \, \subset \, \TD \, \subset \, \D'(\Rd).
\]
IF we endow each of these spaces with the \textit{weak*-topology},
then these embeddings prove to be continuous.
Hence a distribution having compact support can also be regarded
as a tempered distribution,
which in turn gives rise to a Schwartz distribution.
\end{itemize}

\paragraph{} %- convolutions and tensor products
The usage of distributions for the solution theory of linear partial
differential equations can only be exhausted in full strength
if the notions of \textit{convolution} and \textit{tensor product}
between elements of the spaces $\D'(\Rd)$ and $\E'(\Rd)$ has been introduced
as far as possible.

\begin{itemize}[leftmargin=12pt,itemsep=5pt]
\item[1.] \textit{Convolutions.}
The convolution $\phi \ast \psi$ of functions $\phi, \psi \xeof \SRd$
is defined according to
\[
(\phi \ast \psi) (x) \, := \, \int_{\,\Rd} \phi(y) \, \psi(x-y) \, dy
\]
and always exists,
hence the space of Schwartz functions forms an algebra
{\wrt} convolution as multiplicative operation.
Furthermore,
convolution between a regular distribution
$\disF_f \xeof \D'(\Rd)$ and a test function $\phi \xeof \D(\Rd)$
is similarly defined:
\[
(\disF_f \ast \phi) (x) \, := \, \int_{\,\Rd} f(y) \, \phi(x-y) \, dy.
\]

\subparagraph{}
More general,
the convolution $\disF \ast \disG$ $\disF, \disG \xeof \D'(\Rd)$ exists,
if the \textit{support condition} is fulfilled,
meaning that for each compact subset $K \xsubset \Rd$ the set
\[
K_{\disF,\disG} :=
\{(x,y): x \xeof \supp{F}, y \xeof \supp{G}, x+y \xeof K \}
\xsubset \Rd \xtimes \Rd
\]
is compact, too.
In this case we may define the convolution $\disF \ast \disG$ according to
\[
(\disF \ast \disG) \, \phi \, := \,
\iint \limits_{\compK_{\disF,\disG}}
\disF(x)\,\disG(y)\,\phi(x+y) \, dx \, dy, \quad \phi \xeof C_c^{\infty}(\Rd).
\]
The support condition is sufficient
but not necessary for the existence of the convolutional distribution
$\disF \ast \disG$.
Nevertheless,
the condition is always fulfilled
if at least one of the distributions $\disF, \disG$
is compactly supported.
Notice also that
convolution on the space $\D'(\Rd)$ is commutative, if exists.

\subparagraph{}
Especially the convolution between $\disF \xeof \D'(\G)$
and the Dirac distribution $\delta_0$ always exists and
satisfies the equation $\disF \ast \delta_0 \eq \disF$,
thus $\delta_0$ serves as identity element of this operation.
But convolution is also associative,
hence for distributions $\disF, \disG, \disH \xeof \D'(\Rd)$ the equality
\[
(\disF \ast \disG) \ast \disH \= \disF \ast (\disG \ast \disH)
\]
holds,
provided that the involved convolutions exist.

\item[2.] \textit{Tensor products.}
Besides the notation of convolution one also requires
that of \textit{tensor product} between distributions.
Tensor products are useful to transform an inital value
of a distributional differential equation
into an inhomogeneous part of this equation.
Based on the fact
that for $\phi \xeof \D(\R^{d_1})$ and $\psi \xeof \D(\R^{d_2})$ 
the tensor product $\phi \otimes \psi$ defined by
\[
(\phi \otimes \psi)(x,y) \= \phi(x) \cdot \psi(y)\,,
\quad (x,y) \xeof \R^{d_1+d_2}
\]
belongs to the space $\D(\R^{d_1+d_2})$,
there is a unique distribution
$\disF \otimes \disG \xeof \D'(\R^{d_1+d_2})$
for $\disF \xeof \D'(\R^{d_1})$ and $\disG \xeof \D'(\R^{d_2})$,
which for all test functions
$\phi \xeof \D(\R^{d_1})$ and $\psi \xeof \D(\R^{d_2})$
fulfills 
\[
\disF \otimes \disG \, (\phi \otimes \psi)
\= \disF(\phi) \cdot \disG(\psi).
\]
It is called the \textit{tensor product} between $\disF$ and $\disG$.
Performing the tensor product is a commutative and associative operation,
{\ie}
\[
\disF \otimes \disG = \disG \otimes \disF
\quad \and \quad
(\disF \otimes \disG) \otimes \disH \=
 \disF \otimes (\disG \otimes \disH).
\]
\end{itemize}

\section{Adjoint Operators in Locally Convex Spaces}
%------------------------------------------------------------------------------%
The concept of {adjoint operator} is helpful to generalize
the domains of linear partial differential operators and Fourier transformation
to the very extensive spaces of distributions,
but one first has to introduce this notion
in the setting of locally convex spaces.

\subparagraph{}
Let $\banX,\banY$ be locally convex topological linear spaces
and $\banX^*,\,\banY^*$ their adjoint spaces.
Given a continuous linear operator $\linA\:\banX \xto \banY$,
each linear functional $f \xeof \banX^*$ gives rise to
a linear functional $g \xeof \banY^*$ by $g:= f \circ \linA$.
The assignment $f \mapsto g$ induces a continuous mapping
$\linA^*\:\banY^* \xto \banX^*$,
called the \textit{adjoint operator} of $\linA$.
One often says that
the functional $g$ is the \textit{pullback} of $f$ along the operator $\linA$.
The adjoint operator $\linA^*$ is also linear.

\subparagraph{}
The operator $\linA^*$ is also linear.
If $\linA \: \vecV \xto \vecW$ is continuous,
then $\linA^*$ is injective
iff if $\range{\linA} $ is dense in $\vecW$.
Moreover,
if $\linA$ is bijective,
then $\linA^*$ is bijective, too,
and its inverse $(\linA^*)^{-1}$ is the adjoint operator of $\linA^{-1}$.

\subparagraph{}
Finally,
if $\banX$ and $\banY$ are {\FRECHET} spaces,
then the solvability of the linear equation $\linA x \eq y$ 
is related to certain properties of the adjoint operator:
the operator $\linA \: \banX \xto \banY$ is surjective
iff
$\linA^*$ is injective and $\range{\linA^*} $ is closed in $\banY^*$.

\paragraph{} %- distributional derivatives
We use the notion of adjoint operator to define derivatives
of Schwartz distributions in the space $\D'(\G)$.
Let $\PA$ be the elementary differential expression
given by the multi-index $\alpha$.
Then $\PA$ is a continuous linear operator $\PA \:\D(\G) \xto \D(\G)$,
due to the fact
that the space $\D(\G)$ is endowed with a very fine topology.

\subparagraph{}
The action $\PA_{\D}$ of the expression $\PA$
on the distributional space $\D'(\G)$
is defined by the \textit{formally adjoint} expression
$(\PA)^* := (-1)^{\alpha} \PA$,
or equivalently
\[
(\PA_{\D} \disF) \, \phi \, := \, \disF((\PA)^* \phi)
\= (-1)^{\abs{\alpha} } \, \disF(\PA \phi)
\]
for all $\disF \xeof \D'(\G)$ and $\phi \xeof \D(\G)$.

\subparagraph{}
This definition is chosen in order to obtain compatibility
with possible classical derivatives of \textit{regular} distributions.
In fact,
if $F_f \xeof \D'(\G)$ is the distribution generated
by a differentiable function $f \:\G \xto \C$,
then $\PA_{\D} F_f \eq F_{\PA f}$ holds.

\paragraph{} %- compatibility between convolution and tensor product
The mapping $\disF \mapsto \PA_{\D} \disF$ proves to be continuous
on the Schwartz space $\D'(\G)$.
We also have the following compatibility rules
between convolutions and tensor products:
\[
\PA_{\D} (\disF \ast \disG)  \=
(\PA_{\D} \disF) \ast \disG \=
\disF \ast (\PA_{\D} \disG) \,,
\]
\[
 \PA_{\D} (\disF \otimes \disG) =
(\PA_{\D} \disF) \otimes \disG  =
 \disF \otimes (\PA_{\D} \disG).
\]

\section{Adjointness in Banach Spaces}
%------------------------------------------------------------------------------%
Adjoint operators are also useful to study properties
of a densely defined linear operator acting between Banach spaces.
Indeed,
it is profitable to combine the study of a linear operator
with that of its \textit{adjoint operator}.

\subparagraph{} %- definition of adjoint operator
To introduce this concept,
let $\banX,\,\banY$ be Banach spaces,
$\banX^*\,\banY^*$ be the adjoint spaces and
$\linA\:\domain{\linA} \subset \banX \xto \banY$
be a linear operator with domain of definition $\domain{\linA}$
lying dense in $\banX$.
Then the domain of definition $\domain{\linA^*}$
of the \textit{adjoint operator} $\linA^*$ is defined by
\[
\domain{\linA^*} \, \eqdef \,
\{\, g \in \banY^*:
g \circ \linA\:\domain{\linA} \xto \F \mbox{ is continuous}\,\}.
\]
For $g \xeof \domain{\linA^*} $ let
$\linA^* g \eqdef \overline{g \circ \linA}$,
the unique extension of the continuous linear functional $g \circ \linA$
to the whole of $\banX$, 
which implies $\linA^* g \xeof \banX^*$.
Thus the linear operator $\linA^*$
has domain space $\banY^*$ and range space $\banX^*$,
and the domain of definition $\domain{\linA^*}$
of the adjoint operator $\linA^*$ forms a linear subspace of $\banY^*$,
on which $\linA^*$ is linear.

\paragraph{} %- pairing between the spaces X and X*
In the following we denote $\scalar{\cdot}{\cdot}$
the natural pairing between the linear spaces $\banX$ and $\banX^*$,
{\ie}
$\scalar{x}{x^*} \eqdef x^*(x)$ for $x \xeof \banX$ and $x^* \xeof \banX^*$.
Then it is possible to define the adjoint operator $\linA^*$
by the validness of the equations
\[
\scalar{\linA \,x}{y^*} \= \scalar{x}{\linA^* y^*}
\quad \quad (x \in \banX, \, y^* \in \banY^*).
\]
Hence $y^* \xeof \domain{\linA^*}$,
if there is $x^* \xeof \banX^*$ such that
$\scalar{\linA \,x}{y^*} \eq \scalar{x}{x^*}$ for all $x \xeof \banX$,
which gives rise to define $\linA^* y^* \eqdef x^*$.

\paragraph{} %- properties of the adjoint 
Adjoint operators enjoy a couple of nice properties:

\begin{theorem1}\label{IV_MAIN_THEOREM_ON_ADJOINT_OPERATORS}
Let $\linA\:\domain{\linA} \subset \banX \xto \banY$ be densely defined.
\begin{itemize}[leftmargin=9mm]
\item[(a)] The adjoint operator $\linA^*$ is closed.
\item[(b)] If $\banY \eq \banX$ and $\linA$ is closable,
           then $\linA^*$ is densely defined
           and the bi-adjoint $\linA^{**} \eqdef (\linA^*)^*$ is well-defined.
\item[(c)] If $\banY \eq \banX$ is reflexive and $\linA$ is closable,
           then $\linA^{**} \eq \overline{\linA}$,
           which implies that
           if $\linA$ is closed, then $\linA^{**} \eq \linA$.
\end{itemize}
\end{theorem1}

\paragraph{} %- annihilators and orthogonality relations
Let the \textit{annihilator}     $M^{\bot}$ and
    the \textit{pre-annihilator} $^{\bot}N$
of $M \subset \banX$ and $N \subset \banX^*$ be defined by
\[
\begin{cases}[1.2]
\quad M^{\bot}:= \{\,x^*\in \banX^*: \scalar{x}{x^*} = 0 \fa x   \in M \,\},\\
\quad ^{\bot}N:= \{\,x  \in \banX:   \scalar{x}{x^*} = 0 \fa x^* \in N \,\}.\\
\end{cases}
\]
We then obtain the both fundamental relations
\begin{equation}\label{IV_ORTHOGONALITY_RELATIONS}
\overline{\range{\linA^*} } \= \kernel{\linA}^{\bot}
\quad \and \quad
\overline{\range{\linA} } \= ^{\bot}\kernel{\linA^*}
\end{equation}
between the kernels and ranges of the linear operators $\linA$ and $\linA^*$.

\subparagraph{} %- surjectivity of linear operators
With respect to solvability of the equation $\linA \,x \eq y$,
the relationships \eqref{IV_ORTHOGONALITY_RELATIONS} imply
that $\linA$ has dense range in $\banY$ iff $\linA^*$ is injective.

\paragraph{} %- adjoints of continuous linear operators
In the case,
where $\linA$ is \textit{continuous} and $\domain{\linA} \eq \banX$,
the adjoint of $\linA$ is a continuous linear operator
$\linA^*\:\banY^* \xto \banX^*$ and the following assertions can be done:

\begin{itemize}[leftmargin=12pt,itemsep=5pt]
\item[1.]
If $\linA$ is bijective,
then $\linA^*$ proves to be continuous and bijective,
too,
and $(\linA^*)^{-1} \eq (\linA^{-1})^*$ holds.

\item[2.]
Taking the adjoint is a linear and injective operation
in the normed linear space $\BO{\banX,\banY}$,
that is,
\[
(\lambda \linA_1 + \mu \linA_2)^* = \lambda \, {\linA_1}^* + \mu {\linA_2}^*
\]
for $\lambda,\mu \xeof \F$,
where the norm is preserved, 
{\ie} $\norm{\linA^*} \eq \norm{\linA}$ for $\linA \xeof \BO{\banX,\banY}$.

\item[3.]
If $\banZ$ is a further Banach space,
then for
$\linA_1 \xeof \BO{\banY,\banZ}$ and $\linA_2 \xeof \BO{\banX,\banY}$
the equality
$(\linA_1 \circ \linA_2)^* \eq {\linA_2}^* \circ {\linA_1}^*$ holds.

\item[4.]
With respect to compactness of $\linA$, \textsc{Schauder}'s theorem is valid:
the operator $\linA$ is compact \iff $\linA^*$ is compact.
\end{itemize}

\paragraph{} %- formally adjoint pairs
To the end
we introduce the concept of \textit{adjoint pair} of linear operators.
The densely defined linear operators
$\linA\:\domain{\linA} \subset \banX   \xto \banY$ and
$\linB\:\domain{\linB} \subset \banY^* \xto \banX^*$
are said to be \textit{formally adjoint} to each other
or to be a \textit{formally adjoint pair} $(\linA,\linB)$ if
\[
\scalar{\linA \,x}{f}_{\banY \times \banY^*} \=
\scalar{x}{\linB f}_{\banX \times \banX^*}
\]
holds for all $x \xeof \domain{\linA}$ and $f \xeof \domain{\linB}$.
In addition,
we define the \textit{restricted adjoint} $\linB_r$
of the linear operator $\linB$ according to
$\domain{\linB_r} :=
\{ x \in \banX: \mbox{there is } y \in \banY \mbox{ such that }
(\linB f)\,x \eq f(y) \fa f \in \domain{\linB} \}$
and $\linB_r \, x \eqdef y$.
The mapping $\linB_r$ acts between the Banach spaces $\banX$ and $\banY$
and proves to be linear.
We may reformulate the notation of {formally adjoint pair} by saying that 
$\linA$ and $\linB$ are formally adjoint to one another
if $\linB_r$ is an extension of $\linA$.

\subparagraph{} %- realization of a formally adjoint pair
Finally,
by a \textit{realization} of a formally adjoint pair $(\linA,\linB)$ 
we mean any densely defined and closed linear operator $R$
lying between $\linA$ and $\linB_r$,
so $\linA \prec R \prec \linB_r$ holds.

\section{Hilbert Adjoints and Normal Operators}
%------------------------------------------------------------------------------%
We finally give a more direct definition of the adjoint operator
in the Hilbert space setting, 
that in conclusion leads to the concept of \textit{normal operator}.

\subparagraph{}
Remember that
the \textsc{Riesz} $\tilde{\Phi}$ map
is a norm-isomorphism between a complex Hilbert space $\hilH$
and its conjugated space $\hilH^{\times}$ consisting of continuous and
antilinear functionals $\phi \:\hilH \xto \C$.
Then,
given a linear operator $\linA$ with dense domain $\domain{\linA}$ in $\hilH$,
the operator $\linA_H^*$,
often called the \textit{Hilbert adjoint} of $\linA$,
is nothing but the adjoint $\linA^*$ conjugated by the Riesz map,
that is,
\begin{equation}\label{IV_ADJOINT_OPERATOR}
\linA_H^* \, \eqdef \, \tilde{\Phi}^{-1} \circ \linA^* \circ \tilde{\Phi}.
\end{equation}
Hence $\linA_H^*$ constitutes a linear mapping in $\hilH$.

\subparagraph{}
In case of a real Hilbert space,
however,
we may replace $\tilde{\Phi}$ in \eqref{IV_ADJOINT_OPERATOR}
by the \textit{linear} \textsc{Riesz} map $\Phi\:\hilH \xto \hilH^*$,
so $\linA_H^* := \Phi{-1} \circ \linA^* \circ \Phi$.
As a rule for the sequel,
whenever we consider an adjoint operator in a Hilbert space,
we will write $\linA^*$ instead of $\linA_H^*$,
so $\linA^*$ has to be understood as an operator acting in $\hilH$.

\paragraph{} %- another definition of the Hilbert adjoint
In the following
we use a direct approach to define the operator $\linA^*$
that does not rely on the concept of adjoint operator
in the Banach space setting.

\subparagraph{}
If $\linA$ is linear and densely defined in $\hilH$,
the domain of definition $\domain{\linA^*}$ of the adjoint operator $\linA^*$
consists of those elements $y \xeof \hilH$,
for which there is $y^* \xeof \hilH$ fulfilling
$\scalar{\linA \,x}{y} \eq \scalar{x}{y^*}$ for all $x \xeof \domain{\linA}$.
We then define $\linA^* y \eqdef y^*$ in this case,
and $\linA^*$ proves to be linear.

\subparagraph{}
The set $\domain{\linA^*}$ coincides with the set of elements $y \xeof \hilH$,
for which the linear functional $f_y\:\domain{\linA} \xto \F$
given by $f_y(x) \eqdef \scalar{\linA \,x}{y}$ is continuous.
By the B.\,L.\,T.\,theorem,
we may extend $f_y$ to a linear mapping $\tilde{f}$
defined on the whole of $\hilH$,
and due to \textsc{Riesz}'s representation theorem
there is a unique element $y^* \xeof \hilH$
such that $\tilde{f}_y(x) \eq \scalar{x}{y^*}$ for all $x \xeof \hilH$.
Thus we may define $\linA^* y \eqdef y^*$ for $y \xeof \domain{\linA^*}$.

\subparagraph{}
Notice that
the operators $\linA$ and $\linA^*$ are coupled by the formula
\begin{equation}\label{IV_DEFINITION_OF_ADJOINT_OPERATOR}
\scalar{\linA \,x}{y} \= \scalar{x}{\linA^* y},
\end{equation}
where $x \xeof \domain{\linA}$ and $y \xeof \domain{\linA^*}$.
Moreover,
equation \eqref{IV_DEFINITION_OF_ADJOINT_OPERATOR} gives rise
to the following relationships
between the kernels and ranges of $\linA$ and $\linA^*$,
which are variants of \eqref{IV_ORTHOGONALITY_RELATIONS}
for operators in Hilbert space:
\[
\overline{\range{\linA^*} } \= \kernel{\linA}^{\bot}
\quad \mbox{and} \quad
\overline{\range{\linA} } \= \kernel{\linA^*}^{\bot}.
\]

\subparagraph{}
They are useful
to investigate the possible property of surjectivity of the operator $\linA$.
More precisely,
if one can prove that $\linA^*$ is injective,
then the range of $\linA$ is dense in $\hilH$,
and if in addition $\linA$ has closed range
(for example if $\linA$ is bounded below),
then we can conclude that $\linA$ is even surjective.

\paragraph{} %- properties of the adjoint operator
We collect the main properties of adjoint operators to the following

\begin{theorem1}
Let $\linA$ be a linear and densely defined operator
in Hilbert space $\hilH$.

\begin{itemize}[leftmargin=25pt]
\item[(a)]
$\linA^*$ is closed.
But $\linA^*$ is not necessarily densely defined.
In fact,
it can happen that $\domain{\linA^*} \eq \{\0\}$.

\item[(b)]
If $\linB$ is an extension of $\linA$,
then $\linA^*$ is an extension of $\linB^*$,
that is,
$\linA \prec \linB$ implies $\linB^* \prec \linA^*$.
Indeed,
let $y \xeof \domain{\linB^*}$,
then there is $y^* \xeof \hilH$ such that
\[
\scalar{\linB x}{y} \= \scalar{\linA \,x}{y} \= \scalar{x}{y^*}
\qfa x \in \domain{\linB}\,,
\]
and it is easy to see
that $y \xeof \domain{\linA^*}$ with $\linA^* y \eq \linB^* y$.

\item[(c)]
We especially have
$(\linA - \lambda \1_H)^* \eq \linA^* - \overline{\lambda} \1_H$
for $\lambda \xeof \C$,
where $\1_H$ is the identity operator in $\hilH$.
The domain of definition of the operator $(\linA - \lambda \1_H)^*$ is given by
\[
\domain{(\linA  - \lambda \1_H)^*} \=
 \domain{\linA^* - \overline{\lambda} \1_H} \=
 \domain{\linA^*}.
\]

\item[(d)]
$(\lambda \linA)^* \eq \overline{\lambda} \linA^*$
for $\lambda \xeof \C \setminus \{0\}$ 
(but this is not true in case of $\lambda \eq 0$).

\item[(e)]
If $\linA$ is closable with closure $\overline{\linA}$,
then taking the adjoint of both operators leads to the same result,
that is,
$\linA^* \eq \overline{\linA}\,^*$.
\begin{comment}
Indeed,
since $\overline{\linA}$ is an extension of $\linA$,
we immediately obtain that $\overline{\linA}^* \prec \linA^*$.
To show that $\linA^* \prec \overline{\linA}^*$,
let $y \xeof \domain{\linA^*}$.
If $x \xeof \domain{\overline{\linA}}$,
then there is a sequence $\{x_n\}_{n \in \N}$ in $\domain{\linA}$ such that
$x_n \to x$ and $\linA \,x_n \to \overline{\linA} x$.
Since $\scalar{\linA \,x_n}{y} \eq \scalar{x_n}{\linA y}$ for all $n \xeof \N$,
we obtain
$\scalar{\overline{\linA} x}{y} \eq \scalar{x}{\linA^* y}$
by letting $n \to \infty$,
which shows that $y \xeof \domain{\overline{\linA} ^*}$ and
$\overline{\linA}^* y \eq \linA^* y$.
\end{comment}

\item[(f)]
$\linA$ is closable
\iff
$\; \overline{\domain{\linA^*}} \eq \hilH$.
In this case the operator $\linA^{**} \eqdef (\linA^*)^*$
is well-defined and forms the closure of $\linA$,
that is,
$\linA^{**} \eq \overline{\linA}$.
Thus the closure of a closable linear operator can be obtained
by taking the adjoint twice.

\begin{comment}
\Proof:
An essential tool to prove this theorem is the mapping
\[
\Psi\:\hilH \times \hilH \xto \hilH \times \hilH,
\quad \quad
(x,y) \mapsto (-y,x),
\]
having the property $\Psi^2 \eq \1_H$.
This mapping is also linear and unitary,
since
\[
\norm{\Psi(x,y)} \= \norm{(x,y)} \= (\abs{x} ^2 + \abs{y} ^2)^{1/2}
\quad \mbox{ for all } (x,y) \in \hilH \times \hilH.
\]
This implies that for any subspace $U \subset \hilH$ the equality
$\Psi(U^{\bot}) \eq \Psi(U)^{\bot}$ holds.
Now,
if $\Gamma_{\linA}$ is the graph of $\linA$, we have
\[
(\Gamma_{\linA})^{\bot} \= \{\,(y,z) \xeof \hilH \times \hilH:
\scalar{x}{y} + \scalar{z}{\linA \,x} \= 0
\; \mbox{ for all } x \xeof \domain{\linA} \,\}.
\]
On the other hand,
the image of the set $(\Gamma_{\linA})^{\bot}$ under $\Psi$ is given by
\[
\Psi\, (\Gamma_{\linA})^{\bot} \=
\{\,(-z,y) \in \hilH \times \hilH:
\scalar{x}{y} + \scalar{\linA \,x}{z} \= 0
\; \mbox{ for all } x \in \domain{\linA} \,\}
\]
\[
\=
\{\, (z,y) \in \hilH \times \hilH:
z \in \domain{\linA^*} \mbox{ and } y = \linA^*z\,\}
\= \Gamma_{\linA^*},
\]
so after all
we obtain $\Psi \,(\Gamma_{\linA})^{\bot} \eq \Gamma_{\linA^*}$.

Since $\Gamma_{\linA}$ forms a linear subspace of the Hilbert space
$\hilH \xtimes \hilH$,
we have
\begin{equation}\label{IV_CRITERION_FOR_CLOSABILITY}
\overline{\Gamma_{\linA}} \=
(\Gamma_{\linA})^{\bot \bot} \=
\Psi^2 \, (\Gamma_{\linA})^{\bot \bot} \=
(\Psi \, (\Psi (\Gamma_{\linA})^{\bot}))^{\bot} \=
(\Psi \, \Gamma_{\linA^*})^{\bot}.
\end{equation}
After these preliminaries we are able to prove the claim of the theorem.
First, if $\linA^*$ is densely defined,
then $\linA^{**}$ exists and
$\Psi \,(\Gamma_{\linA})^{\bot} \eq \Gamma_{\linA^*}$
together with \eqref{IV_CRITERION_FOR_CLOSABILITY} implies
$\Gamma_{\overline{A}} \eq \overline{\Gamma_{\linA}} \eq \Gamma_{\linA^{**}}$.

\subparagraph{}
Reversely,
suppose that $\domain{\linA^*}$ is not a dense subset in $\hilH$.
Then there exists $\0 \neq z \xeof \domain{\linA^*} ^{\bot}$
fulfilling the equations $\scalar{y}{z} \eq 0$
for all $y \xeof \domain{\linA^*}$,
which means nothing but $(z,0) \xeof (\Gamma_{\linA^*})^{\bot}$.
But then the set $ \Psi \, (\Gamma_{\linA^*}^{\bot})$
contains $(0,z)$ and cannot be the graph of any operator.
If $\linA$ would be closable,
we necessarily would have
$\Gamma_{\overline{\linA}} \eq \psi \, (\Gamma_{\linA^*})^{\bot}$
in opposite to the conclusion above.
Thus $\linA$ is not closable.

\subparagraph{}
To the end,
if $\linA^*$ is densely defined and $\linA$ has a closure,
the operator $\linA^{**}$ is well-defined and
$\Gamma_{\overline{\linA}} \eq \psi \, (\Gamma_{\linA^*})^{\bot} =
\Gamma_{A^{**}}$ holds, so $\overline{\linA} \eq \linA^{**}$.
\qed
\end{comment}
\end{itemize}
\end{theorem1}

\subparagraph{}
Statement (f) is also true in the Banach space setting
under the condition that the involved Banach spaces are reflexive.

\paragraph{} %- formally adjoint pairs in Hilbert space
Furthermore,
in the Hilbert space setting the densely defined and closed linear operators
$\linA,\,\linB$ constitute a formally adjoint pair,
if $\scalar{\linA x}{y} \eq \scalar{x}{\linB y}$
for all $x \xeof \domain{\linA}$
and $y \xeof \domain{\linB}$.
Equivalently,
the pair of operators $(\linA,\linB)$ is formally adjoint,
if $\linA \prec \linB^*$ and $\linB \prec \linA^*$.

\subparagraph{}
For example,
if $\linA$ is densely defined and closed,
the pairs $(\linA,\linA^*)$ and $(\linA^*,\linA)$ are formally adoint to each
another,
but also the operators $\linA+\lambda$ and $\linA^*+\overline{\lambda}$ are so.

\paragraph{} %- adjoints of continuous linear operators
Let us consider the Hilbert adjoints of \textit{continuous}
and everywhere defined operators more in detail.
If $\linA$ is continuous and densely defined,
then the adjoint $\linA^*$ is continuous and everywhere defined,
too.
Hence both operators belong to the Banach algebra $\BO{\hilH}$,
which proves to be closed {\wrt} taking the adjoint.
But the mapping $*\:\BO{\hilH} \xto \BO{\hilH}$
has some more salient properties:

\begin{itemize}[leftmargin=17pt]
\item[(a)]
    $\linA \mapsto \linA^*$ is antilinear,
    that is,
    $(\alpha \linA_1 + \beta \linA_2)^* \=
    \overline{\alpha} \, \linA_1^* + \overline{\beta} \, \linA_2^*$.
\item[(b)]
    $\norm{ \linA} \= \norm{ \linA^*}$,
    so $\linA \mapsto \linA^*$ is norm-preserving in $\BO{\hilH}$.

\item[(c)]
    $(\linA_1 \circ \linA_2)^* \= \linA_2^* \circ \linA_1^*$.

\item[(d)]
    $\linA^{**} \= \linA$,
    {\ie} taking the adjoint is an involutional operation.

\item[(e)]
$\norm{ \linA \circ \linA^*} \= \norm{\linA} \, \norm{\linA^*}$,
called the \textit{C*-\,property} of $\BO{\hilH}$,
which makes $\BO{\hilH}$ into a unital \textit{C*-\,algebra}
(see section VII).
%\footnote{
%\,The C*-\,property of the Banach algebra $\BO{\hilH}$
%is one of the reasons to investigate so-called abstract C*-\,algebras,
%that leads to a multiplicity of useful results.}
%.
\end{itemize}

\paragraph{} %- normal operators
The concept of Hilbert adjoint gives rise
to introduce the class of \textit{normal} linear operators:
a densely defined linear operator $\linA$ is denoted \textit{normal},
if
\[
\linA \circ \linA^* \= \linA^* \circ \linA.
\]
This definition of course implies the identity
of the sets $\domain{\linA}$ and $\domain{\linA^*}$.
Notice that a normal operator is closed (since $\linA^*$ is so),
the normed space $(\domain{\linA^*},\norm{\,\cdot\,}_{\linA^*})$
is complete and the norm $\norm{\,\cdot\,}_{\linA}$ is equivalent to the norm
$\norm{\,\cdot\,}_{\linA^*}$.

\paragraph{} %- characterization of normal operators
The next theorem characterizes normality of a linear operator $\linA$
by other conditions:

\begin{theorem1}
The subsequent statements are equivalent:

\begin{itemize}[leftmargin=30pt]
\item[(a)] $\linA$ is normal;
\item[(b)] $\linA^*$ is normal;
\item[(c)] $\domain{\linA} \eq \domain{\linA^*}$ and
           $\norm{\linA \,x} \eq \norm{\linA^* x}$
           for each $x \xeof \domain{\linA}$.
\end{itemize}
\end{theorem1}

\subparagraph{}
Let us also mention that
the class of \textit{normal} and \textit{bounded} linear operators
in Hilbert space includes a couple of important subclasses:

\begin{itemize}[leftmargin=12pt]
\item[1.] 
A \textit{unitary} operator $\linU$ distinguishes 
by the property $\linU \circ \linU^* \eq 1$ respectively
$\linU^{-1} \eq \linU^*$;
\item[2.] 
a \textit{symmetric operator} $\linA$ is defined by the property
$\linA \prec \linA^*$,
and
\item[3.] 
by a \textit{projection} $P$ is meant a linear operator,
which is symmetric and fulfills $P \circ P \eq P$.
\end{itemize}

% BIBLIOGRAPHY
%------------------------------------------------------------------------------%
\vspace{5mm}
\begin{thebibliography}{99}
\begin{small}

\bibitem{III_Ka} T.\,Kato:
\textit{{Perturbation theory of linear operators}}.
Springer, Grund\-lehren der mathematischen Wissenschaften, Band 132, 1966.

\bibitem{III_RN} F.\,Riesz, B.\,Sz.-Nagy:
\textit{{Functional Analysis}}.
Springer, 1956.

\bibitem{III_Yo} K.\,Yosida:
\textit{Functional Analysis}.
Springer, 1980.

\end{small}
\end{thebibliography}

%------------------------------------------------------------------------------%
%  IV. NORMALLY SOLVABLE OPERATORS ---
%------------------------------------------------------------------------------%
%  §1. Linear Operators with Closed Range ---
%  §2. The Fredholm Index ---
%  §3. Perturbations of Semi-Fredholm Operators ---
%  §4. Basic Theory on Fredholm Operators ---
%  §5. Fredholm's Alternative Theorem ---
%  §6. Concepts of Solvability for Linear Equations ---
%  ---

\newpage
\thispagestyle{empty}
\def\TitleIV{Normally Solvable Operators}
\fancyhead[C] {\small{\textsc{IV. \TitleIV}}}
\part{\TitleIV}
%------------------------------------------------------------------------------%
Given a densely defined linear operator $\linA$ acting between Banach spaces
$\banX$ and\,$\banY$,
the orthogonality relations in section IV can be improved
if the ranges of $\linA$ and $\linA^*$ are \textit{closed} subspaces
of $\banY$ and\,$\banX^*$,
so in these cases we actually would have
\begin{equation}\label{V_ORTHOGONALITY_RELATIONS}
{\range{\linA} } \= ^{\bot}\kernel{\linA^*}
\quad \and \quad
{\range{\linA^*} } \= \kernel{\linA}^{\bot}.
\end{equation}
One also says that
a linear operator $\linA$ is \textit{normally solvable}
if it is densely defined and
$\range{\linA} \eq \, ^{\bot}\kernel{\linA^*}$ holds,
which means that $\linA$ has closed range.

\section{Linear Operators with Closed Range}
%------------------------------------------------------------------------------%
In this paragraph we shall mention two important theorems,
which give us sufficient conditions for a linear operator to own closed range.

\paragraph{} %- closed range theorem
The \textit{closed range theorem}
is one of the classical results in functional analysis
and provides equivalent conditions
for a densely defined and closed linear operator to have closed range:

\begin{theorem1}\label{CLOSED_RANGE_THEOREM}
Let $\banX,\,\banY$ be Banach spaces
and $\linA \: \domain{\linA} \xsubset \banX \xto \banY$
be densely defined and closed.
Then it is equivalent to quote:

\begin{itemize}[leftmargin=9mm]
\item[(a)] $\range{\linA}$ is closed.
\item[(b)] $\range{\linA^*}$ is closed.
\item[(c)] $\range{\linA} $ is the pre-annihilator of $\kernel{\linA^*}$,
           {\ie} $\range{\linA} \= ^{\bot}\kernel{\linA^*} $.
\item[(d)] $\range{\linA^*}$ is the annihilator of $\kernel{\linA}$,
           {\ie} $\range{\linA^*} \= \kernel{\linA} ^{\bot}$.
\end{itemize}
\end{theorem1}

\subparagraph{}
For a proof of this theorem see for example \cite{V_Yo}, pp. 205-207.
The statement was originally proved by \textsc{S.\,Banach}
for continuous linear operators and
have later been extended to closed linear operators,
see \cite{IV_Br}, theorem\,1.1 and theorem\,1.2.

\paragraph{} %- operators bounded below
Besides theorem \ref{CLOSED_RANGE_THEOREM}
it is natural to search for further conditions
that imply closeness of the range of an operator.
One of them is given as follows:
a linear operator $\linA$ is called \textit{bounded below}
if there is $\theta {\,>\,} 0$ such that
\begin{equation}\label{V_BOUNDED_BELOW_CONDITION}
\norm{\linA x}_Y \, \geq \, \theta \norm{x}_X \fa x \xeof \domain{\linA}.
\end{equation}

\subparagraph{} %- bounded below theorem
The \textit{bounded-below theorem} is a useful tool
to verify the closeness of $\range{\linA}$:

\begin{theorem1}\label{V_BOUNDED_BELOW_THEOREM}
If $\linA$ fulfills inequality \eqref{V_BOUNDED_BELOW_CONDITION},
then $\linA$ is injective and the operator
$\linA^{-1} \: \range{\linA} \xsubset \banY \xto \banX$
is bounded with operator norm $\norm{\linA^{-1}} \leq \theta^{-1}$.
Furthermore,
the range $\range{\linA}$ is closed \iff $\linA$ is closed.
\end{theorem1}

Indeed,
injectivity of the operator $\linA$ can easily proved,
supposed \eqref{V_BOUNDED_BELOW_CONDITION} holds,
and closeness of $\range{\linA}$ is a consequence of sequentially closeness
of $\linA$.

\section{The Fredholm Index}
%------------------------------------------------------------------------------%
Given the normed linear spaces $\banX,\,\banY$ and
a linear operator
\[
\linA \: \domain{\linA} \subset \banX \xto \banY,
\]
let $\alpha(\linA)$ (the \textit{nullity index} of $\linA$)
and $\beta(\linA)$ (the \textit{deficiency index} of $\linA$)
be the Hamel dimensions of the linear subspaces
$\kernel{\linA} \xsubset \banX$ and $\cokernel{\linA} \xsubset \banY$.

\subparagraph{} %- Fredholm index
Both integers together can be used
to measure the deviation of $\linA$ from bijectivity.
Moreover,
if at least one of them is finite,
then the so-called \textit{Fredholm index}
\[
\chi(\linA) \eqdef
\begin{cases}[1.2]
\quad + \infty, & \alpha(\linA) = \infty\\
\quad \alpha(\linA) \, - \, \beta(\linA), &
\alpha(\linA) < \infty, \, \beta(\linA) < \infty\\
\quad - \infty, & \beta(\linA) = \infty\\
\end{cases}
\]
is well-defined and takes values in the set
$\{- \infty\} \cup \N \cup \{+ \infty\}$.
In the special case,
where $\alpha(\linA) {\,<\,} \infty$,
$\beta(\linA) {\,<\,} \infty$
and consequently $\chi(\linA) \xeof \Z$,
the operator $\linA$ is said to have \textit{finite index}.

\paragraph{} %- definition of semi-Fredholm operator
The linear operator $\linA \: \domain{\linA} \xsubset \banX \xto \banY$
is called \textit{semi-Fredholm},
if
\begin{itemize}[leftmargin=12mm]
\item[(F-1)] the range $\range{\linA}$ is closed,
\item[(F-2)] the kernel $\kernel{\linA} $ or the cokernel $\cokernel{\linA} $
             is finite-dimensional
             (rsp. $\alpha(\linA) {\,<\,} \infty$ or
                   $\beta(\linA)  {\,<\,} \infty$).
\end{itemize}

Notice the remarkable fact
that if $\linA$ is continuous with $\domain{\linA} \eq \banX$
and $\beta(\linA) {\,<\,} \infty$,
then the range $\range{\linA}$ is necessarily closed,
since it constitutes the topological complement
of the finite-dimensional linear subspace $\cokernel{\linA} \xsubset \banY$,
that is,
$\banX \eq \range{\linA} \oplus \cokernel{\linA}$.

\subparagraph{} %- lower and upper semi-Fredholm operators
We refine the concept of semi-Fredholm operator as follows:
a linear operator $\linA \: \domain{\linA} \xsubset \banX \xto \banY$
with closed range $\range{\linA}$ is called

\begin{itemize}[leftmargin=12mm]
\item[(F-3)]
\textit{upper semi-Fredholm},
if $\alpha(\linA) {\,<\,} \infty$ rsp. $\chi(\linA) {\,<\,} \infty$,

\item[(F-4)]
\textit{lower semi-Fredholm},
if $\beta(\linA) {\,<\,} \infty$ rsp. $\chi(\linA) {\,>\,} {-}\infty$.
\end{itemize}

\subparagraph{} %- example of upper semi-Fredholm operator
Examples of upper semi-Fredholm operators are the bounded-below ones.
In fact,
theorem \ref{V_BOUNDED_BELOW_THEOREM} implies that
these operators are upper semi-Fredholm with vanishing nullity index.

\paragraph{} %- Schauder estimates and Peetre's lemma
If a continuous linear operator $\linA$ satisfies
an abstract \textit{Schauder estimate},
one may quote a sufficient condition for $\linA$ to be upper semi-Fredholm,
which is known as \textsc{Peetre}\textit{'s lemma}:

\begin{theorem1}\label{V_PEETRE_THEOREM}
Let $\banX,\,\banY,\,\banZ$ be three Banach spaces,
where $\banX$ is compactly embedded into $\banZ$,
that is,
there is an injective and compact linear mapping $\iota \: \banX \xto \banZ$.
Then a continuous operator $\linA \: \banX \xto \banY$ is upper semi-Fredholm
iff there is $c {\,>\,} 0$ such that
\[
\norm{x}_{\banX} \, \leq \,
c \, (\norm{\iota x}_{\banZ} \+ \norm{\linA \,x}_{\banY})
\quad \quad (x \in \banX).
\]
\end{theorem1}

\subparagraph{}
For a proof of this theorem, see \cite{IV_Wl}, theorem 12.12, page 180.

\section{Perturbations of Semi-Fredholm Operators}
%------------------------------------------------------------------------------%
Nowadays there is a well-elaborated \textit{perturbation theory}
for semi-Fredholm operators available,
where conditions are studied,
under which a linear operator of the form $\linA + \linB$ is semi-Fredholm,
assumed $\linA$ is semi-Fredholm and $\linB$ fulfills certain properties.
In this case the operator $\linB$ is regarded as a perturbation of $\linA$.

\subparagraph{} %- Yood's theorem
We next quote two theorems {\wrt} conditions on $\linB$,
which ensure the semi-Fredholm property of the operator $\linA + \linB$.
We basically assume that $\banX$ and $\banY$ are Banach spaces.

\begin{theorem1}\label{YOOD_THEOREM}
Let $\linA \: \banX \xto \banY$ be semi-Fredholm
and $\linB \: \banX \xto \banY$ be linear and bounded
with sufficiently small norm $\norm{\linB}$.
Then the operator $\linA + \linB$ remains semi-Fredholm
with $\chi(\linA+\linB) \eq \chi(\linA)$.
\end{theorem1}

\subparagraph{} %- punctured neighbourhood theorem
In other words,
this theorem asserts
that the set of semi-Fredholm operators is open in the set $\BO{\banX,\banY}$.
Especially linear operators of the form $\linA + \lambda$ for $\abs{\lambda} $ 
small enough are semi-Fredholm again.
But more can be said on the values $\alpha(\linA+\lambda)$,
$\beta(\linA+\lambda)$ and $\chi(\linA+\lambda)$
by \textsc{Gohberg}\textit{'s punctured neighbourhood theorem}:

\begin{theorem1}\label{PUNCTURED_NEIGHBOURHOOD_THEOREM}
Let the continuous linear operator $\linA \: \banX \xto \banX$ be semi-Fredholm.
Then there is $\lambda_0 {\,>\,} 0$
such that for all $\lambda \xeof (0,\lambda_0)$ we have
$\alpha(\linA + \lambda) \leq \alpha(\linA)$,
$\beta (\linA + \lambda) \leq \beta (\linA)$
and
$\chi  (\linA + \lambda) \eq  \chi  (\linA)$.
\end{theorem1}

\section{Basic Theory on Fredholm Operators}
%------------------------------------------------------------------------------%
Let us refine the definition of semi-Fredholm operator
to a stronger concept:
in addition to the properties (F-1) to (F-4),
a continuous linear operator $\linA \: \banX \xto \banY$
acting between normed linear spaces $\banX$ and $\banY$ is called
\begin{itemize}[leftmargin=12mm]
\item[(F-5)]
\textit{Fredholm},
if $\chi(\linA)$ is finite,
{\ie} $\alpha(\linA) {\,<\,} \infty$ and $\beta(\linA) {\,<\,} \infty$.
\end{itemize}

\subparagraph{} %- Fredholm property of the adjoint
The relationships \eqref{V_ORTHOGONALITY_RELATIONS} imply
that iff $\linA$ is densely defined and Fredholm,
then the adjoint operator $\linA^*$ is also Fredholm and
$\chi(\linA^*) \eq {-} \chi(\linA)$ holds.
Surely,
if $\linA$ is self-adjoint and Fredholm at the same time,
then $\chi(\linA) \eq 0$ must hold.

\paragraph{} %- properties of Fredholm operators
Next we collect a couple of properties of Fredholm operators:

\begin{itemize}[leftmargin=4mm]
\item[1.]
$\linA^*$ is Fredholm \iff $\linA$ is so.

\item[2.]
If one of $\linA$ or $\linA^*$ is Fredholm,
then the dimensions of the kernels and ranges of $\linA$ and $\linA^*$
are related to each other by
$\alpha(\linA^*) \eq \beta(\linA)$ and
$\beta(\linA^*) \eq \alpha(\linA)$,
which implies $\chi(\linA^*) \eq - \chi(\linA)$.

\item[3.]
Fredholm operators acting in $\banX$ and
belonging to one and the same connected component
within the algebra $\BO{\banX}$ have identical index.
Thus,
if $\linA$ is Fredholm and the linear operator $\linB$ can be reached
by a continuous path in $\BO{\banX}$ with starting point $\linA$,
then $\linB$ is Fredholm,
too,
and $\chi(\linB) \eq \chi(\linA)$ holds.
\end{itemize}

\paragraph{} %- left and right parametrix
Our next subject is the representation of Fredholm operators
using compact linear operators.
More precisely,
a Fredholm operator can be represented by means of
a \textit{left} and \textit{right parametrix},
respectively.

\subparagraph{}
Let $\linA\:\banX \xto \banY$ be a continuous linear operator.
Then a linear mapping $P_l\:\banY \xto \banX$
is called a \textit{left parametrix} of $\linA$
if $\linA \circ P_l \eq \1_X + \compA_l$
for some compact operator $\compA_l \xeof \BO{\banX}$.
In a similar way,
a linear mapping $P_r\:\banY \xto \banY$
is said to be a \textit{right parametrix} of $\linA$,
if $P_r \circ \linA \eq \1_Y + \compA_r$
for some compact operator $\compA_r \xeof \BO{\banY}$.
Every Fredholm operator owns a left and right parametrix,
where the involved compact operators $\compA_l,\,\compA_r$
even are of \textit{finite-rank},
that is,
they have finite-dimensional ranges.

\paragraph{} %- Atkinson's theorem
\textit{Atkinson's theorem} quotes that a Fredholm operator does not lose
the Fredholm property if perturbed by a compact linear operator:

\begin{theorem1}\label{ATKINSON_THEOREM}
Let $\linA$ be Fredholm and $\compA$ be compact.
Then $\linA + \compA$ is Fredholm again, and
$\chi{\linA+\compA} \eq \chi{\linA}$ holds.
\end{theorem1}

\paragraph{} %-- the Fredholm property of closed linear operators
We have introduced the Fredholm property of linear operators
under the assumption that
these are continuous and everywhere defined in their domain spaces.
If a linear operator $\linA$ is closed instead of continuous,
however,
we define the possible Fredholm property by regarding
$\linA$ as a continuous linear operator $\linA \: \domain{\linA} \xto \banY$,
where the domain of definition $\domain{\linA}$ is endowed
with the \textit{graph norm}
\[
\norm{x}_{\linA} \eqdef \norm{x}_X \+ \norm{\linA x}_Y
\]
for $x \xeof \domain{\linA}$.
The pair $(\domain{\linA},\norm{\,\cdot\,}_{\linA})$
forms a Banach space due to the fact that
the operator $\linA$ is assumed to be closed.

\section{Fredholm's Alternative Theorem}
%------------------------------------------------------------------------------%
In the following we cite the \textsc{Fredholm} \textit{alternative theorem}
as the outstanding result in Fredholm theory
{\wrt} the solvability of the linear equation $\linA x \eq y$.
Remember the \textit{rank-nullity theorem} known from linear algebra,
claiming that if a linear space $\vecV$ is finite-dimensional,
then a linear mapping $\linA \: \vecV \xto \vecV$ is surjective
\iff it is injective.

\subparagraph{}
The alternative theorem is a considerable generalization
of the rank-nullity theorem to arbitrary Banach spaces and Fredholm operators
with vanishing index:

\begin{theorem1}\label{FREDHOLM_ALTERNATIVE_THEOREM}
Let $\banX,\,\banY$ be Banach spaces and
the continuous linear operator $\linA \: \banX \xto \banY$ be Fredholm
with index $\chi(\linA) \eq 0$.

\subparagraph{}
According to the solvability of the linear equation $\linA \, x \eq y$,
exactly one of the following cases occurs:

\begin{itemize}[leftmargin=12pt]
\item[1.]
$\linA \, x \eq y$ is uniquely solvable for all $y \xeof \banX$,
so $\linA$ is bijective,

\item[2.]
$\linA \, x \eq y$ is not for all $y \xeof \banX$ solvable,
and $\linA$ is neither injective nor surjective in this case.
However,
if the equation is solvable for an element $y_0 \xeof \banX$,
then the set of solutions forms a finite-dimensional subspace of $\banX$.
Moreover,
the equation $\linA \, x \eq y$ owns a solution $x \xeof \banX$
\iff
$y \xeof \,^{\bot}\kernel{\linA^*}$.
\end{itemize}
\end{theorem1}

\section{Concepts of Solvability for Linear Equations}
%------------------------------------------------------------------------------%
Let us finally illustrate the interplay between Fredholm theory
and the theory of bijective linear operators having compact inverse.

\subparagraph{} %- the operator 1+k is Fredholm
Remember that if $\banX$ is a Banach space and
$k \: \banX \xto \banX$ is compact,
then the operator $\1_{\banX}{+}k$ has finite nullity index and closed range,
hence it is upper semi-Fredholm.
This is also true for the operator $\1_{\banX^*}{+}k^*$,
since $k^*$ is compact, too,
and both properties together imply that
$\1_{\banX}{+}k$ is Fredholm with index $\chi(\1_{\banX}{+}k) \eq 0$.

\subparagraph{} %- equivalence of equations
The following computation is of central meaning 
and for example useful in solving the \textit{Dirichlet problem}
for linear-elliptic differential operators.
Consider the linear equation $\linA x \eq y$,
where $\banX$ is a Banach space and
$\linA \: \domain{\linA} \xsubset \banX \xto \banX$
is bijective linear operator with compact inverse for some $\zeta \in \C$.
Then for $\zeta \xeof \rho(\linA)$ we surely have
\[
\linA \,x \= (\linA {-}\zeta)\,x + \zeta \,x \= y,
\]
and applying the mapping $(\linA {-}\zeta)^{-1}$ from the left yields
\[
x \+ \zeta (\linA {-} \zeta)^{-1}\,x
\=
(\linA {-}\zeta)^{-1}\,y \, =: \, \tilde{y}
\]
or equivalently $x + \tilde{\compA} \,x \eq \tilde{y}$,
where $\tilde{\compA} := \zeta (\linA {-} \zeta)^{-1}$ is compact.
The equations $\linA \,x \eq y$ and
$(\1_{\banX} + \tilde{\compA})\,x \eq \tilde{y}$
are called \textit{equivalent to each other} in the sense
that the element $x$ is a solution of the first one
\iff
it is a solution of the second one.

\paragraph{}
We summarize the result of the previous computation to the following

\begin{theorem1}\label{FREDHOLM_ALTERNATIVE_THEOREM_II}
Let $\linA \: \domain{\linA} \xsubset \banX \xto \banX$
be an operator such that
$(\linA-\zeta)^{-1}$ exists for some $\zeta \in \C$ and is compact.
Then the equation $\linA x \eq y$ is equivalent to an equation
of the form $\fredA x \eq (\1_{\banX} + \tilde{\compA})x \eq \tilde{y}$,
where
the operator $\fredA$ is Fredholm rsp. the operator $\tilde{\compA}$ is compact.
\end{theorem1}

\subparagraph{}
Clearly,
\textsc{Fredholm}'s alternative theorem \ref{FREDHOLM_ALTERNATIVE_THEOREM}
holds for the operator $\fredA$ and therefore provides assertions
also on the solutions of the original equation $\linA x \eq y$.

\paragraph{} %- solvability of linear equations
We want to finalize our elaborations on linear equations
with a complete collection of concepts
regarding the solvability of $\linA \, x \eq y$
where again $\linA$ is assumed to be closed and
acting between Banach spaces $\banX$ and $\banY$.
The equation $\linA \, x \eq y$ is called

\begin{itemize}[leftmargin=15pt]
\item[(a)]
\textit{densely solvable} (\cite{II_Kr}),
if $\overline{\range{\linA}} \eq \banY$;

\item[(b)]
\textit{normally solvable} (\cite{II_Kr}),
if $\overline{\range{\linA}} \eq \range{\linA}$;

\item[(c)]
\textit{correctly solvable} (\cite{II_Kr}),
if $\linA$ is bounded below;

\item[(d)]
\textit{semi-solvable} (\cite{II_Br}), if $\linA$ is Fredholm;

\item[(e)]
\textit{almost solvable} or \textit{regularly solvable} (\cite{II_Br}),
if $\linA$ is semi-solvable and $\chi(\linA) \eq 0$;

\item[(f)]
\textit{uniquely solvable} (\cite{II_Kr}),
if $\linA$ is injective;

\item[(g)]
\textit{everywhere solvable} (\cite{II_Kr}),
if $\linA$ is surjective;

\item[(h)]
\textit{solvable} (\cite{II_Br}),
if it is uniquely and everywhere solvable;

\item[(i)]
\textit{well-posed},
if it is solvable and $\linA^{-1}$ is continuous,
or equivalently,
if $0 \xeof \rho(\linA)$;

\item[(j)]
\textit{completely solvable} (\cite{II_Br}) or \textit{compactly solvable}
if it is solvable and $\linA^{-1}$ is compact.
\end{itemize}

% BIBLIOGRAPHY
%------------------------------------------------------------------------------%
\vspace{5mm}
\begin{thebibliography}{99}
\begin{small}

\bibitem{IV_Br} F.\,Browder:
\textit{Functional Analysis and Partial Differential Equations I}.
Math. Annalen 138, pp.\,55-79, 1959.

\bibitem{IV_Kr} S.\,Krein:
\textit{Linear Equations in Banach Spaces}.
Birk\-h\"{a}user, 1982.

\bibitem{IV_Wl} J.\,Wloka:
\textit{Partial Differential Equations}.
Cambridge University Press, 1987.

\bibitem{IV_Yo} K.\,Yosida:
\textit{Functional Analysis}.
Springer, 1980.

\end{small}
\end{thebibliography}

%------------------------------------------------------------------------------%
%   V. CALCULUS OF VECTOR-VALUED FUNCTIONS ---
%------------------------------------------------------------------------------%
%  §1. Spaces of Classical Differentiable Functions ---
%  §2. Riemann Integration ---
%  §3. Linear Evolution Equations ---
%  §4. Bochner Integration ---
%  §5. Inhomogeneous Cauchy Problems ---
%  §6. Weak Derivatives and Bochner-Sobolev Spaces ---
%  §7. Hilbert-Bochner-Sobolev Spaces ---
%  §8. Holomorphic Vector-Valued Functions ---
%  ---

\newpage
\thispagestyle{empty}
\def\TitleV{Calculus of Vector-Valued Functions}
\fancyhead[C] {\small{\textsc{V. \TitleV}}}
\part{\TitleV}
%------------------------------------------------------------------------------%
In section I we have considered topological and normed linear spaces,
whose elements are real or complex-valued functions.
There are good reasons to extend the definitions of some of these spaces
to contain so-called \textit{vector-valued} functions,
which are nothing but mappings $f\:(a,b) \xto \banX$,
where $(a,b) \subset \R$ is a finite interval
and $\banX$ is a real or complex Banach space.

\subparagraph{}
The most important spaces containing such functions are
the vector-valued {\HOELDER} spaces,
the Bochner spaces and
the Bochner-Sobolev spaces.
For the latter one the notion of weak differentiability of a vector-valued
function is required.

\section{Spaces of Classical Differentiable Functions}
%------------------------------------------------------------------------------%
If $\banX$ is an arbitrary Banach space and $(a,b) \xsubset \R$,
then a vector-valued function $x \: [a,b] \xto \banX$
is called \textit{classical differentiable} in $t_0 \xeof (a,b)$,
if there is an element $x'(t_0) \xeof \banX$ such that
\[
\lim_{\,t \xto t_0} \,
\norm{ \, \frac{x(t)-x(t_0)}{t-t_0}\, - x'(t_0) \, } \= 0.
\]
Moreover,
the function $x$ is called {classical differentiable in} $(a,b)$,
if $x$ is so in each $t \xeof (a,b)$.
The linear space consisting of functions having \textit{continuous}
classical derivatives in $(a,b)$ is denoted $C^1(a,b;\banX)$.

\subparagraph{} %- spaces of k-times differentiable functions
The spaces $C^k(a,b;\banX)$ ($k \xeof \N$) are defined recursively as usual,
and the space $C^{\infty}(a,b;\banX)$ is defined
by taking the intersection of all spaces $C^k(a,b;\banX)$,
$k \xeof \N_0$.

\paragraph{} %-- vector-valued Hölder spaces
We continue with the introduction of vector-valued {\HOELDER} spaces.
First of all,
let $0 < \alpha \leq 1$ and $C^{\alpha}(a,b;\banX)$ be the linear space of
\textit{{\HOELDER}-continuous} functions $f:[a,b] \to \banX$,
defined by the property
\begin{equation}\label{II_HOELDER_NORM_2}
\norm{f} _{\alpha} \, := \,
\sup_{x,y \in [a,b], \, x \neq y} \,
\frac{\norm{\, f(x) - f(y) \,} }
     {\norm{x-y} ^{\alpha}}
\, < \, \infty.
\end{equation}
The mapping $f \xeof C^{\alpha}(a,b;\banX) \mapsto \norm{f}_{\alpha}$ is a norm,
in which the space $C^{\alpha}(a,b;\banX)$ is complete.
Then for $k \xeof \N$ and $0 < \alpha \leq 1$ we define
$C^{k,\alpha}(a,b;\banX)$ as the space of functions $f:[a,b] \to \banX$
such that

\begin{itemize}[leftmargin=15pt]
\item[(a)]
$f \xeof C^k(a,b;\banX)$;

\item[(b)]
the derivatives
$f^{(1)},f^{(2)},\ldots,f^{(k)}$
can continuously be extended from $(a,b)$ to the closure $[a,b]$;

\item[(c)]
The highest derivative $f^{(k)}$ belongs to $C^{\alpha}(a,b;\banX)$.
\end{itemize}

\subparagraph{}
In fact, the mapping $\norm{\,\cdot\,}_{k,\alpha}$ defined by
\[
\norm{f}_{k,\alpha} \, := \,
\sum_{i=0}^k \, \sup_{x \in [a,b]} \snorm{f^{(i)}(x)} \, + \,
\snorm{f^{(k)}}_{\alpha}
\]
constitutes a norm on $C^{k,\alpha}(a,b;\banX)$,
in which this space is complete.

\section{Riemann Integration}
%------------------------------------------------------------------------------%
It is the purpose in the current paragraph
to sketch the setup of the \textit{Riemann integral}
for vector-valued functions $f \: [a,b] \xto \banX$.
Like in the case $\banX \eq \R$,
one usually starts with a partition
\[
\cal{Z} \=
\{x_0:=a,x_1,\ldots,x_{n-1},x_n:=b\},
\quad x_{i-1} < x_i \quad (i=1,\ldots,n),
\]
where the length $\snorm{\cal{Z}} $ of the largest subinterval
is called the \textit{grid size} of $\cal{Z}$:
\[
\snorm{\cal{Z}} \eqdef \max_{1 \leq i \leq n} \, (x_i-x_{i-1}).
\]
Furthermore,
a \textit{tagged partition} $\cal{Z}(t_1,\ldots,t_n)$
of $[a,b]$ is a partition $\cal{Z}$
together with a sequence of intermediate points $t_1,\ldots,t_n$,
where each $t_i$ belongs to the sub\-interval $[x_{i-1},x_i]$.
Then the \textit{Riemannian sum} of $f \: [a,b] \xto \banX$
{\wrt} the tagged partition $\cal{Z}(t_1,\ldots,t_n)$ is given by
\[
R(f,\cal{Z}) \eqdef \sum_{i=1}^n f(t_i) \cdot (x_i - x_{i-1}).
\]

\subparagraph{}
The \textit{Riemann integral} of $f$ on $[a,b]$
is defined as the possible limit of Riemannian sums $R(f,\cal{Z})$
for $\snorm{\cal{Z}} \xto 0$.
Alternatively,
a function $f$ is called \textit{Riemann-integrable} on $[a,b]$
to the value $R_f \xeof \banX$,
if to each $\epsilon > 0$ there is a number $\delta > 0$ such that
\[
\norm{ R_f - R(f,\cal{Z})} < \epsilon
\]
for each tagged partition $\cal{Z}(t_1,\ldots,t_n)$
with $\abs{\cal{Z} } < \delta$,
and in this case the vector $R_f$ is said to be
the \textit{Riemann integral} of $f$,
denoted
\[
\int_a^b f \, dx \eqdef R_f.
\]

\subparagraph{}
Continuity of a vector-valued function $f$ is a sufficient condition
for the existence of the Riemann integral for $f$.
Indeed,
every continuous function $f \: [a,b] \xto \banX$ is Riemann-integrable.
The Riemann integral constitues a linear map on $C([a,b],\banX)$,
satisfying the estimate
\[
\norm{\, \int_a^b f(t) \, dt \,} \, \leq \,  \int_a^b \norm{f(t)} \, dt.
\]
Moreover, 
any Riemann-integrable function fulfills the \textit{mean value theorem}
\[
\int_a^b f(t) \, dt \= (b-a)\,\overline{f}\,,
\]
where the function $\overline{f}$ belongs to the closed convex hull
of the set $\{\,f(t): t \xeof [a,b]\,\}$.

\subparagraph{}
As in the scalar-valued case,
in order to commute integration and taking the limit of a sequence
$\{f_n\}_{n \in \N}$ of continuous vector-valued functions
$f_n \: [a,b] \xto \banX$,
it suffices that
$\{f_n\}$ is \textit{uniformly convergent}
to a continuous function $f \: [a,b] \xto \banX$,
and under this condition
\[
\lim_{n \to \infty} \, \int_a^b f_n(t)\,dt \= \int_a^b f(t) \, dt.
\]

\subparagraph{} %- improper Riemann integral
Like in classical calculus,
we also introduce the notion of the \textit{improper} Riemann integral.
If $f$ is defined on the real line and is R-integrable
on each compact interval,
then by the unproper Riemann integral of $f$ on $\R$
is meant the possible limit
\[
\int_{- \infty}^{+ \infty} f(t) \, dt
\eqdef
\lim_{a \to \infty} \, \int_{-a}^{+a} f(t) \, dt.
\]

\paragraph{} %- complex path integrals
If $f \: \openU \xto \banX$ is defined on a domain $\openU$
within the complex plane
and $\gamma \: [a,b] \xto \openO$ is a continuous \textit{path},
then we define
\[
\int_{\gamma} f(z)\, dz \eqdef \int_a^b (f \circ \gamma) dt,
\]
and if $\gamma$ is continuously differentiable,
the path integral may be computed by 
\[
\int_a^b f(\gamma(t)) \, \gamma(t) \, dt.
\]

\subparagraph{}
The following statement claims
that Riemann integration of vector-valued functions is commutable
with the action of closed linear operators:

\begin{theorem1}
Let $\linA \: \domain{\linA} \xto \banX$ be a closed linear operator
acting in a Banach space $\banX$
and $f \: [a,b) \xto \banX$ ($a < b \leq \infty$) continuous.
Suppose that

\begin{itemize}[leftmargin=22pt]
\item[1.] $f(t) \in \domain{\linA}$ for all $t \xeof [a,b]$,
\item[2.] the mapping $t \mapsto \linA f(t)$ is continuous in $[a,b)$,
\item[3.] the possibly improper integral values $x$ and $y$ exist, where
\[
x \, := \, \int_a^b f(t) \, dt
\quad \and \quad
y \, := \, \int_a^b \linA f(t) \, dt.
\]
\end{itemize}
Then $x \in \domain{\linA}$ and $\linA x \eq y$ holds.
\end{theorem1}

\paragraph{} %- Riemann-Stieltjes integration for vector-valued functions
We finally want to mention the concept
of \textit{Riemann-Stieltjes integral}
for a vector-valued function $f \: [a,b] \xto \banX$.
If $g\:[a,b] \xto \R$ is a further function and
\[
\cal{Z} := \{t_0=a,t_1,\ldots,t_n=b;\xi_0,\ldots,\xi_{n-1}\}
\]
a tagged partition of $[a,b]$,
then $f$ is called \textit{Riemann-Stieltjes-integrable},
if there is an element $R_f \in \banX$ such that
\[
\sum_{i=0}^{n-1} \, f(\xi_i) \, [g(t_{i+1})-g(t_i)],
\quad
\xi_i \in [t_{i+1},t_i]
\]
converges to $R_f$ for $\abs{\cal{Z}} \to 0$,
and in this case we define
\[
\int_a^b f \,dg \, := \, R_f.
\]
The mapping $g$ plays the role of a weight function.
A sufficient criterion for Riemann-Stieltjes integrability
of a function $f \: [a,b] \xto \banX$ is for example,
that $f$ is piecewise continuous and $g$ is monotonly growing in $[a,b]$.

\subparagraph{}
Improper Riemann-Stieltjes integrals may be defined in a similar manner.
In fact,
such integrals exist if and only if the following limits exist:
\[
\lim_{\genfrac{}{}{0pt}{}{a \to -\infty}{b \to +\infty}}
\int_a^b f \, dg
\]

\section{Linear Evolution Equations}
%------------------------------------------------------------------------------%
By an \textit{evolution equation} is meant
an ordinary or partial differential equation
that describes the behaviour of a dynamical system.
More precisely,
it describes the propagation of one or more physical magnitudes along the time,
where their initial states are given in a starting point $t_0 \xeof \R$.

\subparagraph{} %- example of evolution equation
An elementary example of an evolution equation
is given by the differential equation
$x' + \lambda x \eq \0$ for $x(0) \eq x_0 \xeof \R$ and $\lambda {\,>\,} 0$.
It is well known that
this problem is solved by the exponential function
$x(t,x_0) \eq x_0 \, e^{-\lambda t}$ for $t {\,>\,} 0$.

\paragraph{} %- functional analysis and evolution equations
Functional analysis provides a rich theory
on \textit{homogeneous} and \textit{linear} evolution equations of the form
\begin{equation}\label{V_HOMOGENEOUS_CAUCHY_PROBLEM_1}
\begin{cases}[1.2]
\quad x'(t) \+ \linA \, x(t) \= \0,\\ 
\quad x(0) \= x_0 \xeof \banX,\\
\end{cases}
\end{equation}
where $\linA$ is a linear operator acting in a Banach space $\banX$
and with domain of definition $\domain{\linA}$.
The task is to find \textit{vector-valued} solutions $x \: [0,\T) \xto \banX$
of the equation $x' + \linA x \eq \0$ 
satisfying the initial value condition $x(0) \eq x_0$.
The evolution equation \eqref{V_HOMOGENEOUS_CAUCHY_PROBLEM_1} with prescribed
initial value is denoted the \textit{Cauchy problem of first order}.

\subparagraph{} %- second order equations
The homogeneous Cauchy problem of \textit{second order} consists of solving
the initial value problem
\begin{equation}\label{V_HOMOGENEOUS_CAUCHY_PROBLEM_2}
\begin{cases}[1.2]
\quad x''(t) \+ \linA \, x(t) \= \0,\\
\quad x(0)  \= x_0, \\
\quad x'(0) \= x_1,\\
\end{cases}
\end{equation}
where $\linA$ is densely defined and closed.
We also suppose $x_0 \xeof \domain{\linA}$,
but the element $x_1$ may belong to any other linear subspace of $\banX$.

\paragraph{} %- concepts of solution for homogeneous Cauchy problems
In the following we want to introduce some concepts of solution
for the first order Cauchy problem:

\begin{itemize}[leftmargin=12pt]
\item[1.]
A vector-valued function
$x \xeof C([0,\T),\domain{\linA}) \, \cap \, C^1((0,\T),\banX)$ is called
a \textit{classical solution} of \eqref{V_HOMOGENEOUS_CAUCHY_PROBLEM_1},
if it satisfies the initial condition $x(0) \eq x_0$
and the differential equation in each point $t \xeof (0,\T)$.

\item[2.]
a continuous function $x \: [0,\T) \xto \banX$ is said to be
a \textit{mild solution} of \eqref{V_HOMOGENEOUS_CAUCHY_PROBLEM_1} if
\[
\mbox{(1) } \, \tilde{x}(t) \eqdef
\int_{\,0}^t x(s) \, ds \, \xeof \domain{\linA}
\quad \and \quad
\mbox{(2) } \, \linA \, \tilde{x}(t) \eq x(t)-x_0
\]
holds for all $0 {\,<\,} t {\,<\,} \T$,
where the integral on the right-hand side of (1) 
is meant in the Riemannian sense.
One is led to the notion of mild solution in a natural way
by integrating both sides of the equation $x' + \linA \, x \eq \0$,
which then becomes into
\[
\int_{\,0}^t x' ds \=
x(t)-x(0) \= -
\int_{\,0}^t \linA \, x(s) \, ds.
\]
If $\linA$ is closed,
then the action of $\linA$ and performing the integration commutes
for continuous functions $x \: [0,t] \xto \domain{\linA}$,
that is,
\[
\int_{\,0}^t \linA \, x(s) \, ds  \= \linA \, \int_{\,0}^t x(s) \, ds,
\]
and to the end we obtain the equation
\[
x(t) \= x(0) - \linA \, \int_{\,0}^t x(s) \,ds
\quad \quad (t \xeof (0,\T)).
\]
\end{itemize}

\subparagraph{} %- comparing classical and mild solutions
Notice that each classical solution is a mild solution,
and reversely a mild solution is classical
iff it is continuously differentiable in $(0,\T)$.
Thus classical and mild solutions of homogeneous Cauchy problems
only differ {\wrt} this regularity condition.

\paragraph{} %-- correctly and well-posed Cauchy problems
%------------------------------------------------------------------------------%
After having the notions of classical and mild solution
for the homogeneous Cauchy problem available,
we want to introduce
the concepts of \textit{correct posedness} and \textit{well-posedness}.

\subparagraph{} %- definitions of correctly and well-posed problems
First,
the homogeneous Cauchy problem \eqref{V_HOMOGENEOUS_CAUCHY_PROBLEM_1}
is called \textit{correctly posed},
if there is a unique classical solution $x$
for any initial value $x_0 \xeof \domain{\linA}$.
Second,
it is denoted \textit{well-posed},
if it is correctly posed and
the solution $x$ depends continuously on the initial value $x_0$,
or in other words,
if for each $\T {\,>\,} 0$ there is $c_{\T} {\,>\,} 0$ such that
\[
\norm{x(t)} \leq c_{\T} \norm{x_0}
\]
for all $t \xeof (0,\T)$ and all $x_0 \xeof \banX$.

\subparagraph{}
In \cite{V_EN}
the continuous dependency of the solutions from the initial values
is expressed as follows:
if $\{x_n\}_{n \xeof \N}$ is a converging sequence in $\domain{\linA}$
with limit element $\0 \xeof \banX$,
then the related solutions $t \mapsto x_n(t)$ are uniformly convergent
to the zero function in all compact intervals $[0,t_0]$
for $0 {\,<\,} t_0 < \T$
(similar definitions for correct posedness and well-posedness
can also be quoted for the second-order homogeneous Cauchy problem
\eqref{V_HOMOGENEOUS_CAUCHY_PROBLEM_2}).

\subparagraph{}
Let us finally mention that problem \eqref{V_HOMOGENEOUS_CAUCHY_PROBLEM_1}
governed by a densely defined and closed linear operator $\linA$ is well-posed
iff
it is correctly posed and
in addition $\linA$ is regular,
that is,
$\rho(\linA) \ne \emptyset$ holds.

\section{Bochner Integration}
%------------------------------------------------------------------------------%
To introduce a concept of integration for vector-valued functions,
which generalizes the Lebesgue integral for scalar-valued functions,
we first have to introduce a suitable notation of measurability 
for vector-valued functions.

\subparagraph{} %- measurable functions
A function $f\:\Omega \xto \banX$ is called \textit{weakly measurable},
if $\phi \circ f\:\Omega \xto \F$ is measurable for every $\phi \xeof \banX^*$.
Then a weakly measurable function $f\:\Omega \xto \banX$
is called \textit{Bochner-integrable},
if
\[
\int_{\Omega} \norm{f} \,d\mu \, < \, \infty.
\]
For any $p \geq 1$
we denote $L^p(\Omega,\banX)$ the space of functions $f\:\Omega \xto \banX$
having Bochner-integrable powers $\abs{f} ^{p}$.
The elements of this space have also to be understood
as equivalence classes of functions
being identical almost everywhere in $\Omega$.
Endowed with the norm
\[
f \mapsto \norm{f}_p \eqdef (\int_{\Omega} \, \norm{f} ^p \, d\mu)^{1/p},
\]
the Bochner spaces $L^p(\Omega,\banX)$ are complete
and thus constitute Banach spaces.

\subparagraph{} %- weak and strong measurability
It should be mentioned that
if we suppose the Banach space $\banX$ to be separable,
then a function $f\:\Omega \xto \banX$ is weakly measurable
if and only if $f$ is strongly (or Pettis-) measurable.
The latter notion of measurability uses approximations of $f$ by sequences
$\{e_n\}_{n \in \N}$ of elementary functions $e_n\:\Omega \xto \banX$,
and the Bochner integral of $f$ is obtained by taking the limit of the
Bochner integrals of the functions $e_n$,
that is,
\[
\int_{\Omega} f \,d\mu \eqdef \lim_{n \xto \infty} \int_{\Omega} e_n \, d\mu.
\]
This definition proves to be independent of the specific choices
of the approximating sequences for $f$.

\subparagraph{} %- specific choice of Omega and X for Bochner spaces
With respect to applications in partial differential equations,
Bochner spaces are of interest on the one hand 
for the choices $\Omega \eq \G$ and $\banX \eq \F \xeof \{\R,\C\}$,
where $\G$ is a domain in one of the Euclidean spaces $\Rd$
and $L^p(\G,\F)$ is nothing but the Lebesgue space $\Lp$.
On the other hand,
in the weak solution theory of evolution equations
we chose $\Omega$ to be a finite or infinite interval $\J \subset \R$,
and the Banach space $\banX$ may be any complete function space $\mathbb{F}$
consisting of differentiable mappings $f\:\G \xto \F$,
so we are dealing with the Bochner spaces $L^p(\J,\mathbb{F})$ in this case.

\section{Inhomogeneous Cauchy Problems}
%------------------------------------------------------------------------------%
In this paragraph we deal with the solution theory
of the \textit{inhomogeneous} Cauchy problem
\begin{equation}\label{X_INHOMOGENEOUS_CAUCHY_PROBLEM}
\begin{cases}[1.2]
\quad   x'(t) + \linA \, x(t) \, = \, y(t)\\
\quad \; x(0) \= x_0 \in \banX, \\
\end{cases}
\end{equation}
where $y$ is a function $y \: [0,\T) \xto \banX$ and
the solution $x \: (0,\T) \xto \banX$ of the problem
\eqref{X_INHOMOGENEOUS_CAUCHY_PROBLEM} has to be searched.

\subparagraph{} %- Duhamel's principle
Let us first remember \textsc{Duhamel}\textit{'s principle},
also called the \textit{method of variation of constants}:
if $x$ is a classical solution of \eqref{X_INHOMOGENEOUS_CAUCHY_PROBLEM},
the right hand side $y \: [0,\T) \xto \banX$ is continuous
and the linear operator $\linA$ is bounded,
then 
\begin{equation}\label{DUHAMEL_SOLUTION}
x(t) \= \linS(t) \, x_0 \, + \, \int_{\,0}^t \linA \, y(s) \, ds,
\quad (0 < t \leq \T),
\end{equation}
where $\linS(t)$ is the solution operator of the homogeneous problem
\eqref{V_HOMOGENEOUS_CAUCHY_PROBLEM_1} at time $t > 0$.

\paragraph{} %- solvability of the inhomogeneous Cauchy problem
The basic problem in solving the inhomogeneous Cauchy problem is to find
conditions on the function $y$,
which ensure existence and uniqueness of a solution.
We first review the both notions of classical and mild solution:

\begin{itemize}[leftmargin=5mm]
\item[1.]
A function $x \: [0,\T] \xto \banX$ is called a \textit{classical solution}
for problem \eqref{X_INHOMOGENEOUS_CAUCHY_PROBLEM},
if $u \xeof C([0,\T)) \cap C^1((0,\T))$,
fulfills the differential equation $x' + \linA x \eq y$ in every $t \in (0,\T)$
and the initial condition $x(0) \eq x_0$.
\item[2.]
A function $x \: [0,\T] \xto \banX$ is called a \textit{mild solution},
if $x$ fulfills equation \eqref{DUHAMEL_SOLUTION}.
\end{itemize}

\paragraph{}
While in the homogeneous case a mild solution is classical
\iff it is continuously differentiable,
the validness of this assertion remains true for the inhomogeneous problem
only
under some additional regularity properties of the right hand side $y$.

\subparagraph{}
To be precise,
we have to check that $y$ belongs to one of the following function spaces:

\begin{itemize}[leftmargin=30pt]
\item[1.]
$y \xeof L^1([0,\T),\banX) \, \cap \, C((0,\T),\banX)$,

\item[2.]
$y(t) \xeof \domain{\linA}$ for each $t \in (0,\T)$,

\item[3.]
$\linA y \xeof L^1((0,\T),\banX)$.
\end{itemize}

\section{Weak Derivatives and Bochner-Sobolev Spaces}
%------------------------------------------------------------------------------%
We next introduce so-called \textit{Bochner-Sobolev spaces},
which differ by the property
that their members are at the same time
Bochner-integrable and \textit{weakly differentiable}.

\subparagraph{}
Let $\hilH$ be a separable Hilbert space and $L^2(a,b;\hilH)$ be
the Bochner space consisting of functions $f\:(a,b) \xto \hilH$,
that is,
we have
\[
\int_a^b \norm{f(t)}^2 \,dt \, < \, \infty
\]
for $f \xeof L^2(a,b,;\hilH)$.
We endow this space with the inner product
\[
\scalar{f}{g} \eqdef \int_a^b \scalar{f(t)}{g(t)} \,dt,
\quad \quad
f,g \xeof L^2(a,b;\hilH)
\]
such that the norm is given as usual by
$\norm{f}_2 \eq \scalar{f}{f}^{1/2}$.
Thus $L^2(a,b;\hilH)$ is Hilbertian,
and if $(a,b)$ is a \textit{finite} interval,
the inclusion $L^2(a,b;\hilH) \subset L^1(a,b;\hilH)$ holds.
We shall denote $L^2(a,b;\hilH)$ a \textit{Hilbert-Bochner space}.

\subparagraph{} %- weak derivatives of vector-valued functions
For $f \xeof L^2(a,b;\hilH)$,
a further function $g\:(a,b) \xto \hilH$ is said to be
the \textit{weak derivative} of $f$,
if
\[
\int_a^b g \phi \,dx \= - \int_a^b f \phi' \,dx
\]
holds for each $\phi \xeof C_c^{\infty}(a,b,\banX)$,
that is,
$\phi$ is an indefinitely differentiable function with compact support
in $(a,b)$.
If such a function $g$ exists,
we denote $f' := g$ in this case,
which should not be confused with the classical derivative of $f$.
But notice that if the weak derivative of $f$ is present,
then $f$ is also classical differentiable.

\section{Hilbert-Bochner-Sobolev Spaces}
%------------------------------------------------------------------------------%
Having the notion of weak derivative available,
we finally introduce the \textit{Hilbert-Bochner-Sobolev space}
\[
H^1(a,b;\hilH) :=
W^{1,2}(a,b,\hilH) :=
\{f \in L^2(a,b;\hilH): f' \in L^2(a,b;\hilH) \},
\]
which is a Hilbert space as well with inner product given by
\[
\scalar{f}{g}_1 \eqdef
\int_a^b \scalar{f}{g} \, dt
\, + \,
\int_a^b \scalar{f'}{g'} \, dt\,,
\quad f,g \in H^1(a,b;\hilH).
\]

\subparagraph{} %- mixed Bochner-Sobolev spaces
Finally,
let $\hilH_1, \hilH_2$ be \textit{real} Hilbert spaces,
where $\hilH_1$ is continuously embedded into $\hilH_2$.
Then there is also a continuous embedding
\[
H^1(a,b;\hilH_1) \, \overset{\kappa}{\hookrightarrow} \, H^1(a,b;\hilH_2).
\]
Consequently,
given a Hilbert space triple 
$
\vecV \, {\overset{\iota_1}{\hookrightarrow}} \,
\hilH \, {\overset{\iota_2}{\hookrightarrow}} \vecV^*,
$
we may regard the space $H^1(a,b;\vecV)$
as a linear subspace of $H^1(a,b;\vecV^*)$,
which enables us to define the \textit{mixed Hilbert-Bochner-Sobolev space}
\[
H^1(a,b;\vecV,\vecV^*) \eqdef
\{ f \in L^2(a,b;\vecV): f' \in L^2(a,b;\vecV^*) \}
\]
or equivalently
\[
H^1(a,b;\vecV,\vecV^*) \eqdef L^2(a,b;\vecV) \, \cap H^1(a,b;\vecV^*).
\]
Then the mapping
\[
(f,g) \mapsto \scalar{f}{g}_1 \eqdef
\int_a^b \, \scalar{f(t)}{g(t)}_{\vecV} \, dt
\, + \,
\int_a^b \, \scalar{f'(t)}{g'(t)}_{V^*} \, dt
\]
fulfills the properties on an inner product
and makes $H^1(a,b;\vecV,\vecV^*)$ into a Hilbert space
having the following properties:

\begin{itemize}[leftmargin=17pt]
\item[(a)]
Each $f \xeof H^1(a,b;\vecV,\vecV^*)$
is $\lambda$-nearly everywhere identical
to a continuous function $f_c\:[a,b] \xto \hilH$,
and by identifying the functions $f$ and $f_c$
the space $H^1(a,b;\vecV,\vecV^*)$ can continuously be embedded
into the space $C([a,b],\hilH)$.

\item[(b)]
The space $C^{\infty}([a,b],\vecV)$ is dense in $H^1(a,b;\vecV,\vecV^*)$.

\item[(c)]
The following product rule is valid for each $t \xeof (a,b)$:
\[
\frac{d}{dt} \norm{u(t)} ^2
\=
2 \re \, \scalar{u'(t)}{u(t)}_1
\fa u \xeof H^1(a,b;\vecV,\vecV^*).
\]
\end{itemize}

\section{Holomorphic Vector-Valued Functions}
%------------------------------------------------------------------------------%
To the end we introduce the concept of differentiability
for functions  $x\:\openO \xto \banX$,
where $\openO$ is an open and connected subset of the complex plane
and $\banX$ is a complex Banach space.
The function $x$ is called \textit{complex differentiable} in a point
$z_0 \xeof \openO$,
if 
\[
\lim_{z \to z_0} \frac{x(z)-x(z_0)}{z-z_0}
\]
exists,
and this limit is usually denoted $x'(z_0)$,
called the \textit{complex derivative} of $x$ in $z_0$.
Furthermore,
$x$ is said to be \textit{holomorphic} or \textit{analytic}
in $z_0 \xeof \openO$,
if there is an open neighbourhood $U$ of $z_0$
such that $x$ has complex derivative in each $z \xeof U$.
Finally,
$x$ is said to be holomorphic or analytic in $\openO$,
if $x$ is so in each $z_0 \xeof \openO$.

\subparagraph{}
The property of $x$ being holomorphic in $\openO$ can be characterized
as follows:
$x$ is holomorphic \iff
the functions $f \circ x \: \openO \xto \C$ are holomorphic
considered as complex-valued functions for each $f \xeof \banX^*$.
As in complex analysis,
a holomorphic function $x\:\openO \xto \banX$ can be locally represented
by a power series with coefficients in $\banX$.
It should also be mentioned that
the concepts just introduced also apply to \textit{operator-valued} functions,
where the Banach space $\banX$ is nothing but the space $\BO{\banY}$
for some further Banach space $\banY$.

% BIBLIOGRAPHY
%------------------------------------------------------------------------------%
\vspace{5mm}
\begin{thebibliography}{99}
\begin{small}

\bibitem{V_DU} J.\,Diestel, J.\,J.\,Uhl:
\textit{Vector Measures}.
American Mathematical Soc., 1977.

\bibitem{V_EN} K.-J.\,Engel, R.\,Nagel:
\textit{A Short Course on Operator Semigroups}.
Universitext, Springer, 2006.

\bibitem{V_Yo} K.\,Yosida:
\textit{Functional Analysis}.
Springer, 1980.

\end{small}
\end{thebibliography}

%------------------------------------------------------------------------------%
%  VI. GENERAL SPECTRAL THEORY ---
%------------------------------------------------------------------------------%
%  §2. Resolvent Sets and Resolvents ---
%  §3. Different Kinds of Spectra ---
%  §4. Riesz-Schauder Spectral Theory ---
%  §5. Linear Operators with Pure Point Spectrum ---
%  §6. The Discrete Spectral Theorem ---
%  §7. The Riesz-Dunford Functional Calculus ---
%  §8. Essential and Discrete Spectrum ---
%  ---

\newpage
\thispagestyle{empty}
\def\TitleVI{General Spectral Theory}
\fancyhead[C] {\small{\textsc{VI. \TitleVI}}}
\part{\TitleVI}
%------------------------------------------------------------------------------%
The aim in this part
is to give an outline on the spectral theory of a closed or continuous
linear operator $\linA$ acting in a complex Banach space $\banX$.

\subparagraph{}
Let us start with the remark
that each real Banach space $\banX$ can be made
into a complex Banach space $\banX_{\C}$,
and the related procedure is denoted the \textit{complexification} of $\banX$.
It can be performed as follows:
let $\banX_{\C} \eqdef \banX + i \banX \eq \{x+iy: x,y \xeof \banX\}$
and furnished with the addition and multiplication
derived from those in $\banX$.
The norm in $\banX_{\C}$ is defined according to
\[
\norm{x+iy} \, \eqdef \,
\sup \, \{\norm{x} \cos \theta + \norm{y} \sin \theta: \theta \in [0,2\pi)\},
\]
in which this space is complete.

\paragraph{}
Let $\1_X \: \banX \xto \banX$ be the identity operator in $\banX$.
In the sequel we shall consistently use the abbreviation $\linA - \zeta$
instead of $\linA -\zeta \1_X$.
Given a linear operator
$\linA \: \domain{\linA} \subset \banX \xto \banX$,
the complex number $\zeta \xeof \C$ is called a

\begin{itemize}[leftmargin=12pt]
\item[1.]
\textit{regular value} of $\linA$,
if $\linA - \zeta \: \domain{\linA} \xto \banX$ is bijective and
the inverse operator $(\linA-\zeta)^{-1} \: \banX \xto \domain{\linA}$
is continuous.
The set $\rho(\linA)$ of regular values is denoted
the \textit{resolvent set} of $\linA$;

\item[2.]
\textit{spectral value} of $\linA$,
if $\zeta \notin \rho(\linA)$.
The set $\sigma(\linA)$ of spectral values of $\linA$ is said to be
the \textit{spectrum} of $\linA$,
thus $\sigma(\linA) \eq \C \setminus \rho(\linA)$.
\end{itemize}

\subparagraph{}
Having these concepts available,
some further questions arise {\wrt} the location and shape
of the resolvent set and the spectrum within the complex plane.
Moreover,
various classes of linear operators acting in $\banX$ can be defined
by requiring specific properties of these both fundamental sets.

\section{Resolvent Sets and Resolvent Functions}
%------------------------------------------------------------------------------%
Let us continue with a discussion of the resolvent set
and the resolvent functions.
If $\linA$ is a linear operator acting in a complex Banach space $\banX$
and $\zeta \xeof \rho(\linA)$,
then the inverse operator
\[
R(\linA,\zeta) \eqdef (\linA - \zeta)^{-1}\: \banX \xto \domain{\linA}
\]
is everywhere defined in $\banX$.
Since we want to assume $\linA$ to be closed,
the operators $\linA - \zeta$ and $R(\linA,\zeta)$ are closed,
too,
and the \textit{closed graph theorem}
implies that $R(\linA,\zeta)$ is continuous,
that is,
$R(\linA,\zeta) \xeof \BO{\banX}$\footnote{
\,If $\linA$ is \textit{not} closed,
one has in addition to require that
for $\zeta \xeof \rho(\linA)$ the inverse operator $R(\linA,\zeta)$
is bounded.}.
The operator $R(\linA,\zeta)$ is called the \textit{resolvent} of $\linA$
{\wrt} $\zeta \xeof \rho(\linA)$,
and the operator-valued mapping
\begin{equation}\label{VI_RESOLVENT_FUNCTION}
\begin{cases}[1.2]
\quad R(\linA,\cdot)\: \rho(\linA) \xto \BO{\banX}\,,\\
\quad \zeta \mapsto (\linA - \zeta)^{-1}\\
\end{cases}
\end{equation}
is denoted the \textit{resolvent function} (or shortly the \textit{resolvent})
of $\linA$.

\subparagraph{}
But the resolvents of a linear operator $\linA$ are not independent
from each other.
In particular,
for all $\zeta_1,\zeta_2 \xeof \rho(\linA)$
the \textit{first resolvent equation}
\begin{equation}\label{VI_RESOLVENT_EQUATION}
R(\linA,\zeta_1) - R(\linA,\zeta_2)
\=
(\zeta_2-\zeta_1) \, R(\linA,\zeta_1) \circ R(\linA,\zeta_2)
\end{equation}
holds,
which implies that
the resolvents $R(\linA,\zeta_1)$ and $R(\linA,\zeta_2)$
commute with each other.

\paragraph{}
Let us turn to the resolvent function of a linear operator $\linA$
with non-empty resolvent set $\rho(\linA)$.
In fact,
if $\zeta_0 \xeof \rho(\linA)$,
we have $\zeta \xeof \rho(\linA)$ for all $\zeta \xeof \C$ satisfying
\[
\abs{\zeta - \zeta_0 } \, < \, r \eqdef \norm{R(\linA,\zeta_0)}^{-1},
\]
thus $\rho(\linA)$ is an open set in $\C$.
For values $\zeta$ in the disk with center $\zeta_0$ and radius $r$
one may represent the resolvent function $R(\linA,\cdot)$ as a power series
with coefficients in the operator algebra $\BO{\banX}$:
\[
R(\linA,\zeta) \=
\sum_{n=0}^{\infty} \, R(\linA,\zeta_0)^{n+1} \, (\zeta_0-\zeta)^n.
\]
This implies that
$R(\linA,\cdot)$ is locally holomorphic in $\rho(\linA)$.

\paragraph{}
The question,
whether a given complex number $\zeta$ belongs to
the resolvent set or the spectrum of $\linA$,
may be answered by the following criterion:
let $\{\zeta_n\}_{n=1}^{\infty}$ be a sequence in $\rho(\linA)$
with $\lim_{\,n \to \infty} \xi_n \eq \zeta$.
Then $\zeta \xeof \rho(\linA)$ \iff
\[
\sup_{n \in \N} \norm{R(\linA,\zeta_n)} < \infty.
\]
On the other hand,
if $\zeta$ is an element of the boundary of $\rho(\linA)$,
then $\norm{R(\linA,\zeta_n)} \xto \infty$ holds
for each sequence $\{\zeta_n\}_{n=1}^{\infty}$ with $\zeta_n \xeof \rho(\linA)$
and $\zeta_n \xto \zeta$.

\paragraph{}
We finally emphasize
that if $\rho(\linA)\,{\neq}\,\emptyset$,
then $\linA$ is closed.
In fact,
let $\zeta \xeof \rho(\linA)$ such that
$\linA-\zeta$ is bijective with domain of definition
$\domain{(\linA-\zeta)^{-1}} \eq \banY$ and range
$\range{(\linA-\zeta)^{-1}} \eq \domain{\linA}$.
Let $x_n \xto x$ with $x_n \xeof \domain{\linA}$ for $n \xeof \N$
and $\linA x_n \xto y$.
Then
\begin{equation}\label{VI_XXX_TOOL}
x = \lim_{n \to \infty} x_n =
\lim_{n \to \infty} (\linA-\zeta)^{-1}(\linA-\zeta) x_n =
(\linA-\zeta)^{-1}(y-\zeta x).
\end{equation}
Thus $x \xeof \range{(\linA-\zeta)^{-1}} \eq \domain{\linA}$.
Applying $(\linA-\zeta)^{-1}$ from left to \eqref{VI_XXX_TOOL},
we obtain $\linA x - \zeta x \eq y -\zeta x$,
hence $\linA x \eq y$,
which is nothing but closeness of $\linA$.

\subparagraph{}
Thus a reasonable spectral theory makes only sense
for linear operators being at least closed.
Operators,
which are closed and in addition own non-empty resolvent sets,
are often called \textit{regular}.

\section{Different Kinds of Spectra}
%------------------------------------------------------------------------------%
We next turn to the spectrum of $\linA$.
First of all,
the set $\sigma(\linA)$ is closed,
since $\rho(\linA)$ is open.
It may happen that $\sigma(\linA)$ degenerates to the empty set.
In the special case,
where $\linA$ is everywhere defined and \textit{bounded},
the spectrum of $\linA$ proves to be non-empty and forms a bounded subset
of the complex plane.

\paragraph{} %- various kinds of spectra
We next will have a look at the spectral values
of a closed linear operator $\linA$ and study possible reasons,
which for $\lambda \xeof \sigma(\linA)$
prevent the operator $\linA - \lambda$ to be bijective.

\begin{itemize}[leftmargin=12pt]
\item[1.]
Recall from linear algebra
that a complex number $\lambda$ is called an \textit{eigenvalue}
of a linear operator $\linA$,
if there is an element $x \xeof \banX$ such that
$x \neq 0$ and $\linA \,x \eq \lambda x$.
In this case the element $x$ is called an \textit{eigenvector} of $\linA$
{\wrt} $\lambda$,
and all these eigenvectors form a linear subspace of $\banX$,
called the \textit{eigenspace} of $\linA$ for $\lambda$.
The \textit{point spectrum} of $\linA$ is defined to be
the set $\sigma_p(\linA)$ of all eigenvalues of $\linA$.
Clearly,
for $\lambda \xeof \sigma_p(\linA)$ the operator
$\linA - \lambda$ is not injective,
since its kernel is non-trivial.

\subparagraph{}
If $\linA - \lambda$ is injective,
but not surjective,
then $(\linA-\lambda)^{-1}$ is defined on the set $\range{\linA - \lambda}$.
We distinguish between two cases:

\item[2.]
if $\range{\linA - \lambda}$ is dense in $\banX$,
then $\lambda$ is said to be an element of the set $\sigma_c(\linA)$,
called the \textit{continuous spectrum} of $\linA$,
otherwise

\item[3.]
the number $\lambda$ belongs to the \textit{residual spectrum}
$\sigma_r(\linA)$.
There is an interesting identity between the spectrum
of a densely defined and closed linear operator $\linA$
and that of its adjoint $\linA^*$, namely
$\sigma_r(\linA) \eq \sigma_p(\linA^*)$ (see \cite{VI_EN}, p.\,161).
\end{itemize}

\subparagraph{}
In summary,
for each closed linear operator
$\linA \: \domain{\linA} \subset \banX \xto \banX$
the set of complex numbers forms a disjoint decomposition
of the resolvent set and the different kinds of spectra, namely
\[
\C \= \rho    (\linA) \, \sqcup
        \, \sigma_p(\linA) \, \sqcup
        \, \sigma_c(\linA) \, \sqcup
        \, \sigma_r(\linA).
\]

\subparagraph{}
The \textit{approximate point spectrum} $\sigma_{ap}(\linA)$
of a closed linear operator $\linA$ is the set of $\lambda \xeof \C$,
for which $\linA-\lambda$ is not injective
or $\range{\linA-\lambda}$ is not closed in $\banX$.
The elements of $\sigma_{ap}(\linA)$ are called
the \textit{approximate eigenvalues} of $\linA$.
Any number $\lambda \xeof \C$ belongs to $\sigma_{ap}(\linA)$
\iff
there is a sequence $\{x_n\}_{n \in \N} \subset \domain{\linA}$ such that
$\norm{x_n} \eq 1$ for all $n \xeof \N$ and
$\lim_{\,n \to \infty} \norm{\linA \,x_n - \lambda x_n} \eq 0$.
Since $\linA$ is assumed to be closed,
the topological boundary $\partial \sigma(\linA)$ of the spectrum 
forms a subset of $\sigma_{ap}(\linA)$.

\subparagraph{}
It is easy to see that
each eigenvalue of $\linA$ belongs to $\sigma_{ap}(\linA)$,
so this set generalizes the point spectrum $\sigma_p(\linA)$.
But $\sigma_{ap}(\linA)$ can be empty only in the case,
where $\sigma(\linA) \eq \emptyset$ or $\sigma(\linA) \eq \C$.

\paragraph{} %- spectral properties of the adjoint
Recall that the adjoint of a closable linear operator in Hilbert space
is closed,
so the basic concepts of spectral theory are applicable as well.
Let $\linA \: \domain{\linA} \subset \hilH \xto \hilH$
be densely defined and closed.
Then the following identities hold:

\begin{itemize}[leftmargin=12mm,itemsep=2pt]
\item[1.] $\rho(\linA^*) =
          \{\, \overline{\zeta} \xeof \C: \zeta \xeof \rho(\linA) \,\}$.
\item[2.] $\sigma(\linA^*) =
          \{\, \overline{\lambda} \xeof \C: \lambda \xeof \sigma(\linA) \,\}$.
\item[3.] $\sigma_c(\linA^*) =
          \{\, \overline{\lambda} \xeof \C: \lambda \xeof \sigma_c(\linA) \,\}$.
\item[4.] If $\lambda \xeof \sigma_p(\linA) \setminus \sigma_c(\linA)$,
          then $\overline{\lambda} \xeof (\sigma_p(\linA^*) \setminus
               \sigma_c(\linA^*)) \cup \sigma_r(\linA^*)$.
\item[5.] If $\lambda \xeof \sigma_r(\linA)$,
          then $\overline{\lambda} \xeof \sigma_p(\linA^*) \setminus
                                       \sigma_c(\linA^*)$.
\end{itemize}

\paragraph{} %-- spectral mapping theorems
There are a couple of relationships between the spectrum of a linear operator
$\linA$ and those of its resolvents,
that are known as
and is known as \textit{spectral mapping theorems} for resolvents,
see for example \cite{VI_EN}, pp.\,161-162:

\begin{theorem1}\label{SPECTRAL_MAPPING_THEOREM}
For fixed $\zeta \xeof \C$,
various kinds of spectra of the operator $R(\linA,\zeta)$ are related
to the spectra of $\linA$ by the identities

\begin{itemize}[leftmargin=12mm,itemsep=2pt]
\item[1.]
$\sigma(R(\linA,\zeta)) \setminus \{0\} =
\{(\lambda-\zeta)^{-1}: \lambda \xeof \sigma(\linA)\}$,
\item[2.]
$\sigma_p(R(\linA,\zeta)) \setminus \{0\} =
\{(\lambda-\zeta)^{-1}: \lambda \xeof \sigma_p(\linA)\}$,
\item[3.]
$\sigma_{ap}(R(\linA,\zeta)) \setminus \{0\} =
\{(\lambda-\zeta)^{-1}: \lambda \xeof \sigma_{ap}(\linA)\}$,
\item[4.]
$\sigma_r(R(\linA,\zeta)) \setminus \{0\} =
\{(\lambda-\zeta)^{-1}: \lambda \xeof \sigma_r(\linA)\}$.
\end{itemize}
\end{theorem1}

\subparagraph{} %- spectral radius
If $\linA$ is bounded,
the \textit{spectral radius} of $\linA$ is defined according to
\[
r(\linA) \eqdef \sup \, \{\abs{\lambda} : \lambda \xeof \sigma(\linA)\}.
\]
It can be shown that $r(\linA) \leq \norm{\linA}$,
and since each resolvent $R(\linA,\zeta)$ is bounded,
we have
\begin{equation}\label{VI_RESOLVENT_ESTIMATE}
\norm{R(\linA,\zeta)} \, \geq \,
r(R(\linA,\zeta))     \, \geq \,
\delta(\zeta,\sigma(\linA))^{-1}
\end{equation}
for any $\zeta \xeof \rho(\linA)$,
where $\delta(\zeta,\sigma(\linA))$ is the minimal distance between
the singleton $\{\zeta\}$ and the spectrum of $\linA$.
Notice that inequality \eqref{VI_RESOLVENT_ESTIMATE} is a consequence
of the spectral mapping theorem \ref{SPECTRAL_MAPPING_THEOREM}.

\section{Riesz-Schauder Spectral Theory}
%------------------------------------------------------------------------------%
We next turn to the spectral theory of compact linear operators
developed by \textsc{F.\,Riesz} and \textsc{J.\,Schauder} in the early
decades of the 20$^{\textrm{th}}$ century.
In advance,
their results are known
as the \textsc{Riesz-Schauder} \textit{spectral theorem},
which resemble those of $(d,d)$-matrices
considered as linear mappings in the finite-dimensional spaces $\C^d$.

\paragraph{} %- geometric and algebraic multiplicity
Before quoting the main theorem of this paragraph,
let us have a short look to this case,
which can completely be covered by the theory of square matrices.
In finite-dimensional linear algebra,
the spectral theory begins with the study of eigenvalues and eigenvectors
of a square matrix $\A \xeof \C^{d \times d}$.
Remember that an eigenvalue $\lambda$ of $\A$ is a complex number
satisfying $\A x \eq \lambda x$ for some nonzero vector $x \xeof \C^d$,
called an eigenvector corresponding to $\lambda$.
The set of all eigenvalues of $\A$, counting multiplicities,
forms the spectrum $\sigma(\A)$.
Moreover,
the eigenvalues of $\A$ are precisely the roots
of its characteristic polynomial $p_{\A}(\lambda) \eq det(\lambda {-} \A)$,
a monic polynomial of degree $d$.
The algebraic multiplicity of an eigenvalue $\lambda$ is
the multiplicity of $\lambda$ as a root of $p_{\A}(\lambda)$.
The geometric multiplicity of $\lambda$,
however,
is the dimension of the eigenspace $\kernel{\A {-} \lambda}$,
which is at most the algebraic multiplicity.

\subparagraph{} %- Jordan canonical form
A matrix $\A$ is diagonalizable \iff
for every eigenvalue geometric multiplicity equals the algebraic multiplicity.
When $\A$ is not diagonalizable,
its structure is captured by the \textit{Jordan canonical form},
which decomposes $\C^d$ into generalized eigenspaces.
For an eigenvalue $\lambda$,
the generalized eigenspace is nothing but $\kernel{\A{-}\lambda}^m$,
where $m$ is the \textit{index} of $\lambda$,
defined as the smallest positive integer
such that $\kernel{\A{-}\lambda}^m \eq \kernel{\A{-}\lambda}^{m+1}$.
This index equals the size of the largest Jordan block
associated with $\lambda$.
The space $\C^n$ then decomposes as a direct sum
of these generalized eigenspaces over all distinct eigenvalues.

\subparagraph{} %- Jordan canonical form
The \textit{Jordan canonical form theorem} states that
every square matrix $\A$ over an algebraically closed field such as $\C$
is similar to a unique (up to permutation of blocks) Jordan matrix $J$,
meaning there exists an invertible matrix $P$ such that $\A \eq PJP^{-1}$.
The matrix $J$ is block-diagonal,
consisting of Jordan blocks $J_k(\lambda)$ for each eigenvalue $\lambda$,
where a $k \xtimes k$ Jordan block has the value $\lambda$ on the main diagonal
and the value $1$ on the superdiagonal.
This form was established by \textsc{C.\,Jordan} in his 1870 treatise.
A simple example of a non-diagonalizable matrix is
\[
\A \eq \begin{pmatrix} 0 & 1 \\ 0 & 0 \end{pmatrix},
\]
which has the eigenvalue $0$ with algebraic multiplicity $2$
but geometric multiplicity $1$.

\subparagraph{}
In summary,
the \textsc{Jordan} decomposition of a $(d,d)$-matrix
propagates the existence of $k$ pairwise different complex eigenvalues
$\lambda_1,\ldots,\lambda_k$ and associated eigenspaces $E_1,\ldots,E_k$
such that
$\C^d \eq E_1 \oplus E_2 \oplus \ldots \oplus E_k$.
Moreover,
the poles of the resolvent function $R(\linA,\cdot)$ 
coincides with the eigenvalues of the operator $\linA$.

\paragraph{} %- Riesz-Schauder spectral theorem
The \textit{spectral theorem} for compact linear operators shows
that the spectral theory of such operators in principle does not differ
from that of linear transformations acting in the spaces $\C^d$: 

\begin{theorem1}\label{RIESZ-SCHAUDER_SPECTRAL_THEOREM}
Let $\banX$ be a complex Banach space and $\compA\:\banX \xto \banX$ be compact.
Then the following assertions hold:

\begin{itemize}[leftmargin=15pt]
\item[(a)]
The spectrum $\sigma(\compA) \subset \C$ is mostly countable.

\item[(b)]
$\lambda \eq 0$ is the only possible limit point of $\sigma(\compA)$.

\item[(c)]
If $\banX$ is infinite-dimensional, then $0 \xeof \sigma(\compA)$.

\item[(d)]
Each $0 \neq \lambda \xeof \sigma(\compA)$ is an eigenvalue of $\compA$
with finite-dimensional eigenspace.

\item[(e)]
Every $0 \neq \lambda \xeof \sigma(\compA)$
is a pole of the resolvent function $R(\compA,\, \cdot \,)$.
\end{itemize}
\end{theorem1}

\subparagraph{}
Notice
that the forth assertion in theorem \ref{RIESZ-SCHAUDER_SPECTRAL_THEOREM}
is nothing but the second part of \textsc{Fredholm}'s {alternative theorem}
(see paragraph IV.5).
Indeed,
this is true
since for each compact operator $\compA$
the operators $\compA - \lambda$ ($\lambda \xeof \C$) are semi-solvable,
that is,
they are Fredholm operators with vanishing index,
see theorem \ref{ATKINSON_THEOREM}.

\section{Linear Operators with Pure Point Spectrum}
%------------------------------------------------------------------------------%
In this paragraph
we consider a closed linear operator acting in a Banach space,
whose spectrum consists of eigenvalues only
and does not have accumulation points in the complex plane.
By this reason such an operator is called to \textit{have pure point spectrum}.

\subparagraph{}
In other words,
we study linear operators $\linA$ having the property
$\sigma(\linA) \eq \sigma_p(\linA)$,
that is,
the continuous and residual spectrum of $\linA$ are both empty
and $\sigma(\linA)$ consists of eigenvalues only.
Such operators are of outstanding meaning in applied analysis,
since they appear as
realizations of linear-ellipic partial differential operators
in Sobolev spaces for \textit{bounded} domains $\G \subset \Rd$.
Recall also from paragraph IV.6
that in this case
the operator equation $\linA x \eq y$ is equivalent to an operator
equation of the form $(\1_X + \compA) x \eq y$ with $\compA$ being compact,
hence \textsc{Fredholm}'s alternative theorem may be applied. 

\subparagraph{} %- operators with compact resolvents
Furthermore,
a closed linear operator $\linA$
with domain of definition $\domain{A} \subset \banX$
is said to \textit{own compact resolvents},
if at least one resolvent $(\linA -\zeta)^{-1}$ with $\zeta \xeof \rho(\linA)$
is compact.
The first resolvent equation \eqref{VI_RESOLVENT_EQUATION} implies
that in this case all other resolvents of $\linA$ are compact.
It is easy to verify
that $\linA$ has pure point spectrum
\iff it has compact resolvents.

\subparagraph{} %- order of the eigenvalues
The eigenvalues of a linear operator $\linA$ with pure point spectrum
are mostly countable and can be ordered by their moduli,
that is,
there is a sequence $\{ \lambda_n \}_{n \xeof \N}$ consisting
of eigenvalues of $\linA$
such that
\[
\abs{\lambda_1} \,\leq\, \abs{\lambda_2} \,\leq\, \abs{\lambda_3} \,\leq\,
\ldots \, .
\]
More can be said in the "symmetric" case,
see the \textit{discrete spectral theorem} (paragraph IX.3).

\paragraph{} %- criterion for operators having pure point spectrum
There is a useful criterion to check
whether a linear operator owns the property of having pure point spectrum:

\begin{theorem1}\label{CRITERION_FOR_OPERATORS_HAVING_PURE_POINT_SPECTRUM}
Let the domain of definition $\domain{\linA} \subset \banX$
of a densely defined operator $\linA$ in $\banX$
be equipped with the graph norm induced by $\linA$.
Then $\linA$ has compact resolvents
\iff the embedding
$\iota\:
( \domain{\linA} ,\norm{\,\cdot\,}_{\linA} )
\hookrightarrow
(\banX,\norm{\,\cdot\,})$
is compact.
\end{theorem1}

\section{The Riesz-Dunford Functional Calculus}
%------------------------------------------------------------------------------%
Our subject in this paragraph is
the exposition of the \textit{holomorphic functional calculus},
sometimes also called the \textsc{Riesz-Dunford} \textit{calculus}
for bounded linear operators acting in a Banach space $\banX$.
It allows to define functions $f(\linA)$ for a given linear operator $\linA$.

To introduce the Riesz-Dunford functional calculus,
first recall
that a function $f \: \G \subset \C \xto \C$ is called \textit{analytic}
or \textit{holomorphic},
if $f(\zeta)$ for $\zeta$ close to $\zeta_0$ 
can locally be represented by a power series
\[
f(\zeta) = \sum_{k=0}^{\infty} a_k(\zeta_0) \, (\zeta - \zeta_0)^k
\quad \for \; \abs{\zeta-\zeta_0} < R(\zeta_0).
\]
Let $H(\G)$ be the space of analytic functions in $\G$.
We denote $H^{\infty}(\G)$ the subspace
of all \textit{bounded} functions in $H(\G)$,
which is a Banach space {\wrt} the norm
\[
\norm{f}_{\infty} \eqdef \sup_{\zeta \in \G} \, \abs{f(\zeta)} .
\]

\subparagraph{} %- contour integrals
Let $\openO \subset \C$ be an open and connected subset of the complex plane,
$g \: [a,b] \xto \openO$ be a rectificable path
and $f \: \openO \xto \C$ be a complex-valued function.
Then
\[
\int_{\,\gamma} \, f(\zeta) \, d\zeta \eqdef \int_a^b f \, dg
\]
is called the \textit{contour integral} of $f$ {\wrt} the function $g$,
where $\gamma$ is the image of $[a,b]$ under $g$.

\paragraph{} %- Cauchy's integral theorem
\textsc{Cauchy}\textit{'s integral theorem} is a corner stone in
complex analysis and quotes
that if $f$ is \textit{analytic}\, in $\openO$,
then
for \textit{every} closed curve $\gamma$ we have
\[
\int_{\,\gamma} \, f(\zeta) \, d\xi \= 0.
\]
Furthermore,
\textsc{Cauchy}\textit{'s integral formula} says
that for each $z_0 \xeof \openO$ belonging to the interior
of a closed curve $\gamma$ the value $f(z_0)$ is nothing but
\begin{equation}\label{VI_CAUCHY_INTEGRAL_FORMULA_1}
f(z_0) \= \frac{1}{2 \pi i} \int_{\,\gamma} \, \frac{f(z)}{z-z_0} \, dz\,,
\end{equation}
where the \textit{winding number} of $\gamma$ has value $1$.
This representation of $f(z_0)$ does not depend on the shape of $\gamma$.

\subparagraph{}
If we consider $z_0 \xeof \C$ as a linear operator
$Z_0 \: \C \xto \C \xeof \BO{\C}$ defined by
$Z_0 \,\zeta := z_0 \cdot \zeta$ for $\zeta \xeof \C$,
then the spectrum of $\Z_0$ consists of the single point $z_0$
and the resolvent adopts the form
$R(Z_0,\zeta) \eq (\zeta-z_0)^{-1}$.
Using \eqref{VI_CAUCHY_INTEGRAL_FORMULA_1}, 
we obtain
\begin{equation}\label{VI_CAUCHY_INTEGRAL_FORMULA_2}
f(z_0) \=
\frac{1}{2 \pi i} \, \int_{\,\gamma} \, f(\zeta) \, R(Z_0,\zeta) \, d\zeta.
\end{equation}

\subparagraph{} %- Cauchy's integral formula
The formula \eqref{VI_CAUCHY_INTEGRAL_FORMULA_2} serves as a pattern
for the \textit{holomorphic} or \textit{analytic functional calculus}
in unital and complex Banach algebras like the operator algebra $\BO{\banX}$.
Let $\linA$ be a bounded linear operator acting in $\banX$
(so the spectrum $\sigma(\linA)$ is non-emtpy and bounded),
$\gamma$ be a closed and rectificable path with winding number $1$
containing $\sigma(\linA)$
and $f \: \G \xto \C$ be a bounded and analytic function
defined on an open set $\G \supset \sigma(\linA)$.
Then the linear operator $f(\linA)$ is defined
by the \textit{Riesz-Dunford integral}
\begin{equation}\label{VI_RIESZ-DUNFORD_FORMULA}
f(\linA) \eqdef
\frac{1}{2 \pi i} \, \int_{\,\gamma} \, f(\zeta) R(\linA,\zeta) \, d\zeta.
\end{equation}
The right hand side is a vector-valued Riemann-Stieltjes integral,
and convergence of the intermediate sums to the element
$f(\linA)$ is understood in the sense of the norm topology in $\BO{\banX}$.
The value $f(\linA)$ is well-defined,
since $f$ is bounded and analytic.

\subparagraph{}
Recall
that the set $H^{\infty}(\openO)$ of bounded and analytic functions
$f \: \openO \xto \C$ defined in an open neighbourhood $\openO$
of $\sigma(\linA)$ forms a Banach algebra.
Thus the mapping
\[
\begin{cases}[1.2]
\quad \Phi_{\linA} \: H^{\infty}(\openO) \xto \BO{\banB},\\
\quad
f \mapsto f(\linA) \; \mbox{as defined in \eqref{VI_RIESZ-DUNFORD_FORMULA}} \\
\end{cases}
\]
forms a \textit{Banach algebra homomorphism}
from $H^{\infty}(\openO)$ to $\BO{\banB}$,
that is,
the mapping $\Phi_{\linA}$ is linear, norm-preserving
and respects the multiplicative operations in both algebras.
But $\Phi_{\linA}$ also maps polyomials $p \xeof H^{\infty}(\openO)$
to polynomials $p(\linA) \xeof \BO{\banB}$.
Due to these properties,
$\Phi_{\linA}$ is called a \textit{functional calculus},
which in addition deserves the appendix "analytic",
since the space $H^{\infty}(\openO)$ consists of analytic functions.

\subparagraph{}
It should be mentioned 
that for this kind of functional calculus the
\textit{spectral mapping theorem} is valid,
that is,
$\sigma(f(\linA)) \eq f(\sigma(\linA))$
for $\linA \xeof \BO{\banB}$ and $f \xeof H^{\infty}(\openO)$.
The resulting operator $f(\linA) \xeof \BO{\banX}$
defined by \eqref{VI_RIESZ-DUNFORD_FORMULA}
is called the \textsc{Riesz-Dunford} \textit{operator}
associated with $\linA$ and $f \xeof H^{\infty}(\G)$.

\paragraph{} %- projections in a Banach space
In the following we suppose that
$\lambda$ is an \textit{isolated} point of the spectrum $\sigma(\linA)$,
so $\sigma(\linA) \cap \U_{\lambda} \eq \{\lambda\}$ holds
for every sufficiently small neighbourhood $\U_{\lambda}$ of $\lambda$.
Let $\gamma_{\lambda}$ be a closed and rectificable path,
which surrounds only $\lambda$,
but not any other element of $\sigma(\linA)$.
Then for $x \xeof \banX$ the contour integral
\[
\pi_{\lambda}(x) \eqdef
\frac{1}{2 \pi i} \, \int_{\,\gamma_{\lambda}} \, R(\linA,\zeta) x \, d\zeta
\]
gives rise to a further linear mapping $\pi_{\lambda}$,
called the \textsc{Riesz} \textit{projection} of $\linA$
{\wrt} the spectral value $\lambda$.

\paragraph{} %- properties of the Riesz projection
To the end,
we collect the properties of the mapping $\pi_{\lambda}$ to the following list:

\begin{itemize}[leftmargin=12pt]
\item[1.]
$\pi_{\lambda}$ is linear;

\item[2.]
$\pi_{\lambda}$ owns the \textit{projection property}
$\pi_{\lambda} \circ \pi_{\lambda} \eq \pi_{\lambda}$;

\item[3.]
The spectrum of this mapping only consists of the value $\lambda_0$ itself;

\item[4.]
$\pi_{\lambda}$ commutes with $\linA$,
so $\linA \circ \pi_{\lambda} \eq \pi_{\lambda} \circ \linA$ holds;

\item[5.]
The range $\Pi_{\lambda}$ of $\pi_{\lambda}$ is a closed subspace in $\banX$;

\item[6.]
There is a closed subspace $\Pi_{\lambda}^{\bot} \subset \banX$
such that
$\banX \eq \Pi_{\lambda} \, \oplus \, \Pi_{\lambda}^{\bot}$,
where $\Pi_{\lambda}^{\bot}$ is defined by
$\Pi_{\lambda}^{\bot} := \1_H- \Pi_{\lambda}$;

\item[7.]
Both subspaces $\Pi_{\lambda}$ and $\Pi_{\lambda}^{\bot}$
are $\linA$-invariant, so
$\linA \, \Pi_{\lambda} \subset \Pi_{\lambda}$ and
$\linA \, \Pi_{\lambda}^{\bot} \subset \Pi_{\lambda}^{\bot}$.
One also says that the subspaces 
$\Pi_{\lambda}$ and $\Pi_{\lambda}^{\bot}$
\textit{reduces} the entire space $\banX$;

\item[8.]
The mapping $\pi_{\lambda}^{\bot} := \1_H \pi_{\lambda}$ is a projection again,
and both projections
$\pi_{\lambda}$ and
$\pi_{\lambda}^{\bot} := \1_H -\pi_{\lambda}$
are \textit{orthogonal} to each other,
that is,
$\pi_{\lambda}        \circ \pi_{\lambda}^{\bot} \eq 
 \pi_{\lambda}^{\bot} \circ \pi_{\lambda}        \eq \0$.
\end{itemize}

\subparagraph{}
The concept of {Riesz projection}  may be extended into two directions.
First,
the linear operator $\linA$ may be assumed to be closed
instead of being continuous,
and second
the subset $\{\lambda\} \subset \sigma(\linA)$ may be replaced
by any connected component $\sigma_0(\linA)$ of the set $\sigma(\linA)$.

\section{Essential and Discrete Spectrum}
%------------------------------------------------------------------------------%
The concept of semi-Fredholm operator (see section IV)
can be used to characterize the \textit{essential spectrum}
$\sigma_{ess,1}(\linA)$ of a closed operator $\linA$ acting
in a complex Banach space $\banX$:
a number $\lambda \xeof \C$ belongs to $\sigma_{ess,0}(\linA)$
\iff $\linA - \lambda$ is \textit{not} semi-Fredholm.
But there are further approaches to define the notion of essential spectrum:

\begin{itemize}[leftmargin=6mm]
\item[\,] 
$\sigma_{ess,1}(\linA) := \C \setminus \{\lambda \xeof \C:
\linA-\lambda \mbox{ is upper semi-Fredholm} \}$,

\item[\,] 
$\sigma_{ess,2}(\linA) := \C \setminus \{\lambda \xeof \C:
\linA-\lambda \mbox{ is lower semi-Fredholm} \}$,

\item[\,] 
$\sigma_{ess,3}(\linA) := \C \setminus \{\lambda \xeof \C:
\linA-\lambda \mbox{ is semi-Fredholm} \}$,

\item[\,] 
$\sigma_{ess,4}(\linA) := \C \setminus \{\lambda \xeof \C:
\linA-\lambda \mbox{ is Fredholm} \}$,

\item[\,] 
$\sigma_{ess,5}(\linA) := \C \setminus \{\lambda \xeof \C:
\linA-\lambda \mbox{ is Fredholm with } \chi(\linA-\lambda) \eq 0\}$,

\item[\,] 
$\sigma_{ess,6}(\linA) := \{\lambda \xeof \sigma_{ess,5}(\linA) \:
\mbox{ the scalars near } \lambda \mbox{ are in } \rho(\linA)\}$.
\end{itemize}

\subparagraph{} %- properties of the essential spectrum
The sets $\sigma_{ess,4}(\linA)$ and $\sigma_{ess,5}(\linA)$
are called the \textit{Wolf spectrum} and the \textit{Schechter spectrum},
respectively,
whereas the set $\sigma_{ess,6}(\linA)$ is denoted
the \textit{Browder spectrum}.

\subparagraph{}
All sets $\sigma_{ess,k}(\linA)$ are closed subsets of $\sigma(\linA)$,
and by definition we have the chain of inclusions 
\[
\sigma_{ess,3}(\linA) \xsubset
\sigma_{ess,1}(\linA) \cup
\sigma_{ess,2}(\linA) \xsubset
\sigma_{ess,4}(\linA) \xsubset
\sigma_{ess,5}(\linA) \xsubset
\sigma_{ess,6}(\linA).
\]

\subparagraph{}
In the Hilbert space setting,
however
all sets $\sigma_{ess,k}(\linA)$ coincide and form the largest subset
$\sigma_{ess}(\linA) \xsubset \sigma(\linA)$ such that
for every compact linear operator $\compA \xeof \BO{\banX}$ 
\[
\sigma_{ess}(\linA + \compA) \= \sigma_{ess}(\linA).
\]

\paragraph{} %- discrete spectrum
%------------------------------------------------------------------------------%
On the other hand,
the \textit{discrete spectrum} $\sigma_d(\linA)$ is defined
to be the relative complement of $\sigma_{ess,5}(\linA)$ {\wrt} $\sigma(\linA)$,
hence we have a disjoint decomposition of the spectrum of $\linA$ according to
\[
\sigma(\linA) \eq \sigma_d(\linA) \, \sqcup \, \sigma_{ess,6}(\linA).
\]
The set $\sigma_d(\linA)$ contains those points $\lambda \xeof \sigma(\linA)$,
which are isolated in $\sigma(\linA)$
and for which the rank of the \textit{Riesz projection} $P_{\lambda}$ is finite,
where $P_{\lambda}$ is defined by the \textit{Dunford integral}
\begin{equation}\label{DUNFORD_FORMULA}
P_{\lambda}(x) \eqdef
\frac{1}{2 \pi i} \, \int_{\,\gamma_{\lambda}} \, R(\linA,\zeta) x \, d\zeta
\quad (x \xeof \banX)
\end{equation}
and $\gamma_{\lambda}$ is a closed continuous path enclosing $\lambda$
but no other values in $\sigma(\linA)$.
Elements of the set $\sigma_d(\linA)$ are called \textit{normal eigenvalues}.
One may also introduce the set $\sigma_d(\linA)$ as the collection
of isolated points $\lambda \xeof \sigma(\linA)$
having finite algebraical multiplicity
(this magnitude is the Hamel dimension of the linear space
$E_{gen}(\linA,\lambda) := \cup_{n \in \N} E(\linA^n,\lambda)$
of generalized eigenvalues of $\linA$).

% BIBLIOGRAPHY
%------------------------------------------------------------------------------%
\vspace{5mm}
\begin{thebibliography}{99}
\begin{small}

\bibitem{VI_EN} K.J.\,Engel, R.\,Nagel:
\textit{A Short Course on Operator Semigroups},
Universitext, Springer Verlag, 2006.

\bibitem{VI_Ka} T.\,Kato:
\textit{{Perturbation theory of linear operators}}.
Springer, Grund\-lehren der mathematischen Wissenschaften, Band 132, 1966.

\bibitem{VI_RS1} M.\,Reed, B.\,Simon:
\textit{Modern Methods of Mathematical Physics I: {\FA}},
Academic Press Inc., 1980.

\bibitem{VI_Yo} K.\,Yosida:
\textit{Functional Analysis}.
Springer, 1980.

\end{small}
\end{thebibliography}

%------------------------------------------------------------------------------%
% VII. BANACH ALGEBRAS AND FOURIER TRANSFORMATION ---
%------------------------------------------------------------------------------%
%  §1. Normed Algebras ---
%  §2. Resolvent and Spectrum of Banach Algebra Elements ---
%  §3. The Gelfand Representation ---
%  §4. Involutional Algebras and C*-Algebras ---
%  §5. The Gelfand-Naimark Theorem ---
%  §6. Locally Compact Abelian Groups ---
%  §7. Abstract Fourier Transformation ---
%  §8. Examples of Fourier Transformation ---
%  §9. Fourier-Plancherel Transformation ---
% §10. The Fourier Transform in Distributional Spaces ---
%  ---

\newpage
\thispagestyle{empty}
\def\TitleVII{Commutative Banach Algebras}
\fancyhead[C] {\small{\textsc{VII. \TitleVII}}}
\part{\TitleVII}
%------------------------------------------------------------------------------%
The aim in this part
is to review commutative Banach algebras and
their applications to other topics in functional analysis and operator theory.
They are of own interest due to the celebrated \textit{Gelfand transformation}.
This presence of this fundamental mapping lead us
on the one hand
  to the \textit{Fourier transformation} in certain Lebesgue spaces and
on the other hand
  to the continuous functional calculus for \textit{normal} linear operators
  acting in Hilbert space.

\section{Normed Algebras}
%------------------------------------------------------------------------------%
By a \textit{normed algebra} (rsp. \textit{Banach algebra}) is meant
a normed linear space (rsp. Banach space)
together with a multiplicative operation,
which is compliant with the norm in some sense.

\subparagraph{}
More specifically,
let $(\banB,+,\circ)$ be a $\F$-algebra,
that simultaneously constitutes a normed space $(\banB,\norm{\cdot})$.
Then $\banB$ is called a \textit{normed algebra},
if the following conditions are fulfilled:

\begin{itemize}[leftmargin=16mm,itemsep=2pt]
\item[(A-1)]
$\norm{x \circ y} \, \leq \, \norm{x} \cdot \norm{y}$
for each pair of elements $x,y \xeof \banB$.

\item[(A-2)]
If $\banB$ has an identity $\1 \xeof \banB$, then $\norm{\1} \eq 1$.
\end{itemize}

\subparagraph{}
Condition (A-1) implies
that the map $(x,y) \mapsto x \circ y$ is continuous.
If (A-2) is valid,
then $\banB$ is called a \textit{unital} algebra.
Finally,
$\banB$ is said to be \textit{commutative},
if
\begin{itemize}[leftmargin=16mm,itemsep=2pt]
\item[(A-3)]
\;$x \circ y \eq y \circ x$ \; for all $x,y \xeof \banB$.
\end{itemize}
As for normed linear spaces,
completeness of normed algebras is of interest,
thus a normed algebra is denoted a \textit{Banach algebra},
if it is complete as a normed linear space.

\subparagraph{} %- the unital Banach algebra of bounded linear operators
The most prominent example of a non-commutative Banach algebra is given by
the $\C$-algebra $\BO{\banX}$ of bounded linear operators
$\linA \:\banX \xto \banX$,
where $\banX$ is a real or complex Banach space.
It forms a unital Banach algebra {\wrt} composition of mappings
and the operator norm $\norm{\linA}$ of $\linA \xeof \BO{\banX}$.

\paragraph{} %- Gelfand-Mazur theorem
One of the key results in the theory of complex Banach algebras
is the \textit{Gelfand-Mazur theorem}
(see also \cite{VII_Ka}, p.\,14):

\begin{theorem1}\label{GELFAND_MAZUR_THEOREM}
If each element $x \neq 0$ in a complex Banach algebra $\banB$ is invertible,
then $\banB \cong \C$.
\end{theorem1}

\subparagraph{}
In other words, 
theorem \ref{GELFAND_MAZUR_THEOREM} quotes that
the field of complex numbers is essentially the only complex Banach algebra,
in which each element $x \neq 0$ is invertible.

\paragraph{} %- Neumann series
Let $\banB$ be a complex and unital Banach algebra.
By the \textit{Neumann series} of $x \xeof \banB$ is meant the infinite sum
$
s_x := \sum_{n=0}^{\infty} \, x^n.
$
If $s_x$ is convergent in $\banB$,
then the element $\1-x$ is invertible,
and a short computation proves $(\1-x)^{-1} \eq s_x$.
Like in the case $\banB \eq \C$,
the {Neumann series} is convergent for all $x \xeof \banB$
fulfilling $\norm{x} < 1$.

\subparagraph{}
If $x \xeof \banB$ is invertible,
then any other element $x' \xeof \banB$ having the property
\[
\norm{ x - x'} \, < \, \frac{1}{\snorm{x^{-1}}}
\]
proves to be invertible, too.
Thus all elements in a sufficiently small neighbourhood
of an invertible element have convergent Neumann series
and are invertible,
which implies that
the set $\GL{\banB} $ of invertible elements in $\banB$ is open.
It is called the \textit{group of units} in $\banB$,
since $x \circ y$ and $x^{-1}$ are invertible for $x,y \xeof \GL{\banB} $.
Both operations $(x,y) \mapsto x \circ y$
and $x \mapsto x^{-1}$ are continuous, 
so $\GL{\banB} $ forms a topological group,
when considered in the subset topology of $\banB$.

\paragraph{} %- examples of commutative Banach algebras
In this section we are especially interested
in \textit{commutative} Banach algebras.
Here are some important examples {\wrt} applications:

\begin{itemize}[leftmargin=12pt,itemsep=4pt]
\item[1.]
The field $\C$ of complex numbers may be regarded
as a complex unital Banach algebra in its own right.

\item[2.]
For any topological space $\topX$
the function spaces $C_b(\topX)$ and $C_0(\topX)$,
both endowed with the supremum norm $\norm{f} := \sup_{x} \abs{f(x)} $.

\subparagraph{}
If $\topX$ is \textit{compact},
then the function space $C(\topX)$ consisting of continuous functions
$f\: \topX \xto \C$ forms a Banach algebra,
which in addition is unital with identity $\1(x) := 1$ for $x \xeof \topX$.

\item[3.]
The algebra $L^1(\topG,\Borel{\topG},\mu)$ of Lebesgue-integrable functions
{\wrt} the Haar measure $\mu$ on a locally compact topological group $\topG$,
where multiplication is defined by \textit{convolution}
\[
(f \ast g)(x) \eqdef \int_{\R} f(y) \, g(x-y) \, d\mu
\]
for $f,g \xeof L^1(\topG)$.
Important examples for $\topG$ are the classical groups $\Rd$, $\Z$
and $\mathbb{T} \eq \{z \xeof \C: \norm{z} \eq 1\}$.

\item[4.]
The Banach algebra $L^{\infty}(X)$ of essentially bounded functions
$f \: X \xto \R$, 
where $(X,\Borel{X},\mu)$ may be any measure space.
Addition and multiplication in this algebra are defined pointwise.

\item[5.]
The algebra $\cal{M}^r(\topG)$ of regular complex Borel measures
defined on a locally compact topological group $\topG$.
The multiplication is defined by convolution of measures:
\[
(\mu \ast \nu)(B) \eqdef
(\mu \times \nu) (\{(x,y) \in \topG \xtimes \topG:x+y \in B\})
\]
for any Borel set $B \subset \topG$.
This Banach algebra is even unital, 
since the so-called \textit{Dirac measure} $\mu_0$ owns the property
$\mu_0 \ast \nu \eq \nu \ast \mu_0 \eq \nu$
for all $\nu \xeof \cal{M}^r(\topG)$.
\end{itemize}

\section{Resolvent and Spectrum of Banach Algebra Elements}
%------------------------------------------------------------------------------%
The basic concepts of spectral theory,
as introduced for bounded linear operators acting in a Banach space,
are also well-defined for elements in a unital complex Banach Algebra $\banB$.
Let $\1 \xeof \banB$ be the identity and $x \xeof \banB$ be arbitrary,
then

\begin{itemize}[leftmargin=12pt]
\item[1.]
any number $\zeta \xeof \C$ is called \textit{regular} {\wrt} $x$,
if $x - \zeta \1$ is invertible;

\item[2.]
the set of regular numbers
$\rho(x) := \{ \zeta \xeof \C: (x - \zeta \1)^{-1} \exists \}$
is said to be the \textit{resolvent set} of $x$;

\item[3.]
the mapping $R_x \: \zeta \xeof \rho(x) \mapsto (x - \zeta \1)^{-1}$
is denoted the \textit{resolvent function} and
$R_x(\zeta)$ is called the \textit{resolvent} of $x$
{\wrt} $\zeta \xeof \rho(x)$;

\item[4.]
the complementary set $\sigma(x) := \C \setminus \rho(x)$
is called the \textit{spectrum} of $x$,
and any element $\lambda \xeof \sigma(x)$ is called a
\textit{spectral value} of $x$.
\end{itemize}

\paragraph{} %- properties of the resolvent set and the resolvent
We just reformulate some known results for resolvent sets and 
spectra of bounded linear operators,
which are also true for elements in a unital complex Banach algebra $\banB$.
Because the set $\GL{\banB} $ of invertible elements is open,
the resolvent set $\rho(x)$ of $x \xeof \banB$ forms an open subset in $\C$.
Moreover,
the resolvent function $R(x,\cdot) \: \rho(x) \xto \banB$
is continuous and even analytic,
that is,
$R(x,\cdot)$ can locally be represented by a convergent power series
with coefficients in $\banB$.

\paragraph{} %- properties of the spectrum
On the other hand,
the spectrum $\sigma(x)$ is closed and \textit{non-empty},
which is a generalization of the fact
that $\sigma(\linA) \neq \emptyset$ for any bounded linear operator
$\linA \:\banX \xto \banX$.

\section{Gelfand Representation}
%------------------------------------------------------------------------------%
Given the Banach algebras $\banB_1, \banB_2$,
a linear mapping $\Phi\:\banB_1 \xto \banB_2$
is denoted a \textit{Banach algebra homomorphism},
if it preserves multiplication according to 
\[
\Phi(x \circ y) \= \Phi(x) \circ \Phi(y)
\fa x,y \in \banB_1.
\]
Moreover,
$\banB_1$ and $\banB_2$ are said
to be \textit{isomorphic} to each other, 
if there is a bijective homomorphism $\Phi \:\banB_1 \xto \banB_2$ fulfilling
$\norm{\Phi(x)}_{\banB_2} \eq \norm{x}_{\banB_1}$ for all $x \xeof \banB_1$,
which in this case is called a \textit{Banach algebra isomorphism}
(notation: $\banB_1 \cong \banB_2$).

\paragraph{} %- Gelfand representation
The \textit{Gelfand representation} is a powerful mapping
in the setting of commutative Banach algebras
and provides a Banach algebra homomorphism
between an arbitrary commutative Banach algebra $\banB$ and
the algebra $C_0(\topX)$ of continuous functions vanishing at infinity,
where $\topX$ is assumed to be a locally compact Hausdorff space.

\subparagraph{} %- characters on Banach algebras
A \textit{character} on a complex Banach algebra $\banB$
is a \textit{non-zero} Banach algebra homomorphism $\chi \:\banB \xto \C$.
Characters (often also called \textit{multiplicative functionals}) are
of great importance in the theory of unital commutative Banach algebras,
as we will see in the following.

\subparagraph{}
By the multiplication property of a Banach algebra homomorphism,
each character $\chi$ on $\banB$ satisfyies the inequality
$\norm{\chi} \leq 1$,
so $\chi$ is a continuous mapping and
thus an element of the dual space $\banB'$.
The set $\Gamma$ of all characters is called the \textit{structure space}
or the \textit{spectrum} of the Banach algebra $\banB$.

\subparagraph{} %- Gelfand topology
Due to the inequality $\norm{\chi} \leq 1$ for all $\chi \xeof \Gamma$,
the set $\Gamma$ is a subset of the closed unit ball $B_1(0) \subset \banB'$.
Thus we may endow $\Gamma$ with the subset topology on $B_1(0)$
inherited from the weak-*-topology on $\banB'$,
which is called the \textit{Gelfand topology} on the structure space.
Now,
by the \textit{Alaoglu-Bourbaki theorem}
the weak-*-topology on $\banB'$ is weak*-compact,
implying
that the Gelfand topology must be locally compact and Hausdorff.
Hence the structure space $\Gamma$ becomes into
a locally compact Hausdorff space.
Moreover, 
if $\banB$ is unital,
it can be shown that the set $\Gamma$ is even compact.

\paragraph{} %- Gelfand representation
From now on
let $\banB$ be a \textit{commutative} Banach algebra.
We consider the natural mapping $\gamma \:\banB \xto C(\Gamma)$,
where $\gamma$ assigns to each $x \xeof \banB$ the function
$\hat{x} \xeof C(\Gamma)$ by way of $\hat{x}(\chi) := \chi(x)$,
that is,
the function values of $\gamma$ are given by evaluating
the characters $\chi$ at the element $x$.
The mapping $\gamma$ is called the \textit{Gelfand representation}
of $\banB$.

\paragraph{} %- properties of the Gelfand representation
We collect the main properties of the Gelfand representation
for a commutative unital Banach algebra $\banB$ to the following list:

\begin{itemize}[leftmargin=12mm]
\item[1.]
$\gamma$ is a Banach algebra homomorphism
from $\banB$ to $C(\Gamma)$ or $C_0(\Gamma)$,
depending whether $\banB$ is unital or not.

\item[2.]
If $\banB$ is unital,
then $\hat{\1}$ is the constant function $1$ on $\Gamma$.

\item[3.]
If $\banB$ is unital,
then $x \xeof \banB$ is invertible
iff $\hat{x}$ nowhere vanishes.

\item[4.]
$\norm{ \hat{x} }_{\infty} \leq \norm{x}$ for each $x \xeof \banB$.
\end{itemize} 

\subparagraph{}
The Gelfand representation is injective
iff $\banB$ is a \textit{semisimple} Banach algebra,
that is,
if the intersection of all maximal ideals in $\banB$,
called the \textit{Jacobsen radical} of $\banB$,
is equal to the set $\{0\}$.

\section{Involutional Algebras and C*-Algebras}
%------------------------------------------------------------------------------%
We now will have a look at the theory of C*-algebras
and especially at the Banach algebra $\BO{\hilH}$
of bounded linear operators in a complex Hilbert space $\hilH$.
This algebra contains an important subclass of linear operators,
which are called \textit{normal}
and are defined by satisfying a certain commutativity relation.

\paragraph{} %- involutions
We first define the concept of an \textit{involution}
for an arbitrary complex Banach algebra $\algA$ as follows:
a bijective algebra-homomorphism $* \:\algA \xto \algA$,
$x \mapsto x^*$ is called an \textit{involution},
if for all $\lambda,\mu \xeof \C$ and $x,y \xeof \algA$
\[
\begin{cases}[1.3]
\quad (\lambda x + \mu y)^* \=
\overline{\lambda} \cdot x^* + \overline{\mu} \cdot y^*\,, \\
\quad (x^*)^* = x\,, \\
\quad (xy)^* = y^* x^*.\\
\end{cases}
\]

\subparagraph{} %- properties of involutions
An involution acting on a normed algebra $\algA$
let the norms of the algebra elements unchanged,
that is,
$\norm{x} \eq \norm{x^*}$ holds for all $x \xeof \algA$,
which implies that involutions are always isometric.
Furthermore,
an algebra homomorphism $\Phi \:\algA_1 \xto \algA_2$
between involutional algebras $\algA_1, \algA_2$
is called a \textit{*-homomorphism},
if it is compliant with the involution operations in both algebras,
that is,
$\Phi(x^*) \eq \Phi(x)^*$ holds for all $x \xeof \algA_1$.

\paragraph{}
If a normed algebra $\algA$ is furnished with an involution,
one may define some additional properties for the elements of $\algA$.
In particular,
an element $x \xeof \algA$ is called

\begin{itemize}[leftmargin=30pt]
\item[1.] \textit{normal},       if $x x^* \eq x^* x$.
\item[2.] \textit{self-adjoint}, if $x^* \eq x$.
\item[3.] \textit{skew-adjoint}, if $x^* \eq -x$.
\item[4.] \textit{projection},   if $x x \eq x \eq x^*$.
\item[5.] \textit{unitary},      if $\algA$ is unital,
                                    $x$ is invertible and $x^* \eq x^{-1}$.
\end{itemize}

\subparagraph{}
It can easily be verified
that
all unitary and self-adjoint elements in $\algA$ are normal.
Furthermore,
each element $x \xeof \algA$ may be written as $x \eq x_1 + i x_2$,
where $x_1, x_2$ are self-adjoint elements in $\algA$ defined by the formulas
\[
x_1 \eqdef \frac{1}{2}(x+x^*),
\quad \quad
x_2 \eqdef \frac{1}{2i}(x-x^*).
\]

\paragraph{} %- C*-property
The following concept is crucial in the theory of involutional Banach algebras:
an involution $*:\algA \xto \algA$ of a normed algebra $\algA$ 
is said to have the \textit{C*-property},
if
\[
\norm{x^* x} = \norm{x}^2
\quad \, \mbox{for all }
x \in \algA.
\]
This property leads to the notion of \textit{C*-algebra},
which is defined to be an
{involutional Banach algebra satisfying the C*-property}.

\subparagraph{} %- examples of C*-algebras
Examples of C*-algebras are the function algebras
$C_b(\topX)$ and $C_0(\topX)$, where $\topX$ is an ordinary topological space,
and the algebra $L^{\infty}(X)$ for a measure space $(X,\salgA,\mu)$.
In all three cases the involution is given by conjugation
$f \mapsto \overline{f}$ of functions.
But for the convolution algebra $L^1(\R,+,*)$ 
the involution does not satisfy the C*-property,
hence this algebra cannot be a C*-algebra.

\paragraph{} %- properties of C*-algebras
There are a couple of important properties,
which differ a C*-algebra from an arbitrary involutional algebra.
The first one is given by the fact
that the spectrum of any element $x \xeof \banC$ does not increase
if it is considered as an element
of a proper C*-subalgebra $\cal{U}$ of $\banC$.
In other words,
we have
\[
\sigma_{\cal{U}}(x)
\=
\sigma_{\banC}(x) \, \mbox{ for all } x \in \banC. 
\]
So the spectrum $\sigma(x)$ of $x \xeof \banC$ is unique,
anyway to which specific subalgebra of $\banC$ it is computed.
Second,
for a normal element $x$ in a C*-algebra
the spectral radius can easily be computed by the formula
$r(x) \eq \norm{x}.$
Third, 
self-adjoint and unitary elements of a C*-algebra
can be characterized by properties of their spectra:

\begin{itemize}[leftmargin=30pt]
\item[1.]
$x$ is self-adjoint
iff $\sigma(x) \subseteq \R$. 

\item[2.]
$x$ is unitary
iff $\sigma(x) \subseteq \T \eq \{z \xeof \C: \abs{z} \eq 1\}$.
\end{itemize}

\section{The Gelfand-Naimark Theorem}
%------------------------------------------------------------------------------%
In this paragraph we want to elaborate the Gelfand-Naimark theorem
and the continuous functional calculus based on this theorem.

\subparagraph{} %- C*-subalgebras and generators
Preliminary,
two further concepts are meaningful when dealing with C*-algebras:

\begin{itemize}[leftmargin=12pt]
\item[1.]
A subset $\cal{U}$ of a C*-algebra is called a C*-subalgebra,
if $\cal{U}$ is a \textit{closed} subalgebra,
which is also closed {\wrt} the involution,
that is,
$x^* \xeof \cal{U}$ for all $x \xeof \cal{U}$.

\item[2.]
Let $\setS$ be a non-empty subset of the C*-algebra $\banC$.
Then the intersection of all C*-subalgebras containing $\setS$
is called the C*-algebra \textit{generated} by $\setS$ (notation: $[\setS]$),
and $S$ itself is said to be a \textit{generator} of $[\setS]$.
\end{itemize}

\paragraph{} %- Gelfand-Naimark theorem
Recall the Gelfand transformation $\gamma$ for a commutative Banach algebra,
which is a Banach algebra homomorphism from an arbitrary commutative
Banach algebra $\banB$ to one of the Banach algebras $C(\Gamma)$
and $C_0(\Gamma)$.
In this context, the topological space $\Gamma$
is the structure space of $\banB$,
which turned out to be a locally compact Hausdorff space.
It was also quoted
that the Gelfand transformation is injective for each semisimple Banach algebra,
and this statement is especially true in case of a C*-algebra.
Note
that in general the Gelfand representation
of a semisimple commutative Banach algebra is not surjective,
so the range of the structure space $\gamma$ may be
only a subalgebra of $C(\Gamma)$.

\subparagraph{}
Now,
the question arises
under which conditions the Gelfand representation is bijective
and therefore is a Banach algebra isomorphism
between the Banach algebra $\banB$ and the function algebras 
$C_0(\Gamma)$ and $C(\Gamma)$, respectively.
As already mentioned,
a necessary condition is that of $\banB$ being semisimple,
and a sufficient condition is given by the fact
that $\banB$ is a C*-algebra,
which legitimates the meaning of this class of algebras.

\paragraph{} %- Gelfand-Naimark theorem
We collect these results to the \textit{Gelfand-Naimark theorem}:

\begin{theorem1}\label{GELFAND_NAIMARK_THEOREM}
Let $(\banC,+,\circ)$ be a commutative C*-algebra.

\begin{itemize}[leftmargin=15pt]
\item[(a)]
The Gelfand representation $\gamma$ is surjective.

\item[(b)]
If $\banC$ is unital,
then $\banC$ is isometrically isomorphic to the function algebra $C(\Gamma)$.

\item[(c)]
If $\banC$ is \textit{not} unital,
then $\banC$ is isometrically isomorphic to the function algebra $C_0(\Gamma)$.
\end{itemize}
\end{theorem1}

\subparagraph{}
Thus, 
by this theorem,
each element of a commutative C*-algebra can be considered
as a continuous function defined on a locally compact Hausdorff space.
The theorem emphasizes the meaning of the Banach algebras $C_0(\topX)$
and $C(\topX)$ in case of $\topX$ being a locally compact Hausdorff space.

\subparagraph{}
Conversely,
for a fixed normal element $x \xeof \banC$,
one may send each function $f \xeof C(\sigma(x))$ to a C*-algebra element $f(x)$
by the inverse mapping of the Gelfand representation.
Given a normal element $x \xeof \banC$,
the family of mappings $f \mapsto f(x)$
is called a \textit{continuous functional calculus} for $x$.

\paragraph{} %- spectral mapping theorem for C*-algebras
The following \textit{spectral mapping theorem} establishes a link
between the spectrum $\sigma(x)$ of a normal element $x$ of a C*-algebra
and the range of the continuous function $f$
associated by the Gelfand representation:

\begin{theorem1}
Let $f$ be a continuous complex-valued function,
which is defined on the spectrum $\Gamma$ of the C*-algebra $\banC$.
Then for any normal element $x \xeof \banC$
the equation $\sigma(f(x)) \eq f(\sigma(x))$ holds.
\end{theorem1}

\subparagraph{}
Thus,
if $\banC'$ is the C*-algebra generated
by a normal element $x$ and the identity,
that is, $\banC' := [\{x,\1\}]$,
which is of course commutative,
then $\banC'$ is *-isomorphic to the commutative C*-algebra $C(\sigma(x))$.
In this case any isomorphism maps $\1$ to $\chi_{\sigma(x)}$
and $x$ to the identity function $id \: z \mapsto z$.

\paragraph{} %- continuous functional calculus
We shall apply the Gelfand-Neumark theorem to the commutative C*-algebra $x]$
generated by a normal element $x \xeof \banC$,
its adjoint element $x^*$ and the identity element $\1$.

\subparagraph{}
Recall that
if $\Gamma_x$ is the structure space of $[x]$,
then the Gelfand representation $\gamma_x$ is a C*-isomorphism
between the C*-algebras $[x]$ and $C(\Gamma_x)$,
that is,
one can assign to each function $f \xeof C(\Gamma_x)$
a unique element $f(x) \xeof [x]$
by applying the inverse Gelfand representation $\Phi_x := \gamma_x^{-1}$.

\subparagraph{}
Hence it is possible to define a bijective and continuous mapping
$h_x \: \Gamma_x \xto \sigma(x)$,
where $\sigma(x)$ is the spectrum of $x$.
Since $\Gamma_x$ is compact and the set $\sigma(x)$ is Hausdorff
{\wrt} its metric topology induced by that in $\C$,
a well-known theorem from point-set topology claims
that $\Gamma_x$ and $\sigma(x)$ are homeomorphic to each other.
As a consequence,
there is also a homeomorphism
\[
g_x \: C(\Gamma_{x}) \xto C(\sigma(x))
\]
between the associated function spaces $C(\Gamma_{x})$ and $C(\sigma(x))$,
and finally there is a C*-isomorphism
\[
\begin{cases}[1.3]
\quad \Phi_{x} \: C(\sigma(x)) \xto [x]\,,\\
\quad f \mapsto (\gamma_x^{-1} \circ g_x^{-1})(f)\,,\\
\end{cases}
\]
called the \textit{continuous functional calculus} of the normal element $x$.

\paragraph{}
The following diagram shows the isomorphism in detail:

\begin{center}
\scalebox{0.9}
{
\begin{tikzpicture}[node distance=2.2cm]
\node              (v) at (0,0){$[{x}] \subset \banC$};
\node [right of=v] (w)         {$C(\Gamma_{[x]})$};
\node [below of=w] (u)         {$C(\sigma(x))$};
\draw [-stealth] (v) -- node [midway, above] {$\gamma_{x}$}(w);
\draw [-stealth] (w) -- node [midway, right] {$g_{x}$}     (u);
\draw [-stealth] (u) -- node [midway, below] {$\Phi_{x}$}  (v);
\end{tikzpicture}
}

\vspace{3mm}
\footnotesize{Figure \thesection.1:
continuous functional calculus $\Phi_{x}$
for a normal element $x \xeof \banC$.}
\end{center}

\paragraph{} %- algebra B(H)
The importance of C*-algebras is justified by the fact 
that for each Hilbert space $\hilH$
the associated Banach algebra $\BO{\hilH}$ of bounded linear operators 
forms a C*-algebra, where the involution mapping is given by
\[
* \: \BO{\hilH} \xto \BO{\hilH}, \quad \linA \mapsto \linA^*.
\]

\paragraph{}
If $\hilH$ is a Hilbert space and $\linA \xeof \BO{\hilH}$ normal,
then for each continuous function $f \: \sigma(\linA) \xto \C$
the operator $f(\linA) \xeof \BO{\hilH}$ is well defined.
This is especially true for the operators 
$\abs{\linS} $ and $\sqrt{\abs{\linS} }$.

\section{Locally Compact Abelian Groups}
%------------------------------------------------------------------------------%
Before discussing abstract Fourier transformations,
it is useful to review some concepts coming from general topology,
especially that of a topological group, 
which plays a fundamental role in the theory of Fourier transforms.

\subparagraph{} %- topological groups
There is a subclass within that of general topological spaces,
which is of great importance in higher analysis
and whose members are called \textit{topological groups}.
They are mathematical objects owning the algebraic structure of a group
and a topological structure together.
Both have to be compatible to each other in the following sense:
let $(G,\circ)$ be an ordinary group and a topological space $(G,\topT)$
at once.
Then $G$ is called a \textit{topological group},
if the group operation $(x,y) \mapsto x \circ y^{-1}$ is continuous.

\subparagraph{}
Instead to require continuity of the mapping $(x,y) \mapsto x \circ y^{-1}$,
one can equivalently require
that the mappings $(x,y) \mapsto x \circ y$ and $x \mapsto x^{-1}$
have to be continuous.

\paragraph{} %- examples of topological groups
Well-known examples of topological groups are:

\begin{itemize}[leftmargin=12pt]
\item[1.]
each group $\grpG$ with its discrete topology $2^{\grpG}$,

\item[2.]
the group $\R^{\times} := \R - \{0\}$,

\item[3.]
the group $\C^{\times} := \C - \{(0,0)\}$,

\item[4.]
the group $\R^+ := \{x \xeof \R: x >0\}$,

\item[5.]
the additive groups $(\R,+)$ and $(\C,+)$,

\item[6.]
the \textit{circle group} $\mathbb{T} := \{z \xeof \C: \abs{z} \eq 1 \}$,

\item[7.]
each subgroup $\topU$ of a topological group $\topG$
endowed with the subspace topology
and its closure $\overline{\topU}$ are topological groups.
\end{itemize}

\paragraph{}
Like for purely algebraic groups,
a \textit{homomorphism} between topological groups
is defined to be a \textit{continuous} group homomorphism,
and such a homomorphism $\iota$ is called an \textit{isomorphism},
if it is a group isomorphism and $\iota^{-1}$ is continuous,
too
(we use the notation $\topG \cong \topH$ in this case).

\subparagraph{} %- (R,+) is isomorphic to (R+,.)
A well-known example of a group isomorphism is
the \textit{exponential function} $t \mapsto e^{t}$,
which maps the topological group $(\R,+)$ to the group $(\R^+,\cdot)$.
Rather similar,
the function
$t \xeof [0,1) \mapsto e^{2 \pi it}$
forms a group isomorphism between the topological groups
$\R / \Z$ and $\mathbb{T}$.

\paragraph{} %- automorphisms between topological groups
Because in a topological group $\topG$
the left  translation $x \mapsto g \circ x$ and
the right translation $x \mapsto x \circ g$
are continuous mappings,
they are both examples for \textit{automorphisms},
that is,
isomorphisms from the group $\topG$ to itself.
Hence the topological properties of a topological group
are completely determined by the topology in a neighbourhood
of the identity element $\1 \xeof \topG$,
which gives rise to call a topological group to be \textit{homogeneous}
in this case.

\subparagraph{}
Moreover,
the kind,
how the identity $\1$ of a topological group $\topG$ lies
in a given subgroup $\topU$,
determines the topology of that subgroup.
In particular,

\begin{itemize}[leftmargin=12pt]
\item[1.]
$\topU \subset \topG$ is open iff $\1$ is an inner point of $\topU$.

\item[2.]
$\topU$ is discrete iff $\1$ is an isolated point of $\topU$.
\end{itemize}

\paragraph{} %- circle group
Remember that the {circle group} $\mathbb{T}$ is the group
consisting of all complex numbers of modulus $1$ with group operation
\[
e^{i \alpha} \circ e^{i \beta} \eqdef
e^{i \alpha} \cdot e^{i \beta} \eqdef
e^{i(\alpha + \beta)},
\quad \alpha,\beta \xeof \R.
\]
Endowed with the subset topology inherited
from the standard topology on the field $\C$,
the unit circle $\mathbb{T}$ forms a topological group.
Consider the set $C(\topG,\mathbb{T})$ consisting of all
continuous functions $f \:\topG \xto \mathbb{T}$.
By furnishing this set with the compact-open topology,
it becomes into a Hausdorff space.

\subparagraph{} %- characters in the setting of topological groups
A \textit{character} $\chi$ is a (topological) group homomorphism from a
topological group $\topG$ to the circle group $\mathbb{T}$
and hence a continuous mapping.
We denote $\topG^*$ the set of characters $\chi \:\topG \xto \mathbb{T}$,
which forms a subset of the Hausdorff space $C(\topG,\mathbb{T})$.
We can make $\topG^*$ into a topological space by taking the subset topology
{\wrt} the compact-open topology on $C(\topG,\mathbb{T})$.

\paragraph{} %- dual groups
Performing multiplication of characters $\chi_1,\chi_2 \xeof \topG^*$ by
\[
\chi_1 \circ \chi_2\: g \mapsto \chi_1(g) \cdot \chi_2(g),
\]
we obtain a further character $\chi_1 \circ \chi_2 \xeof \topG^*$.
It can be shown that this operation has all the properties of a group,
and since it is commutative,
$\topG^*$ has the structure of an Abelian group,
called the \textit{character group} or the \textit{dual group} of $\topG$.
Moreover,
the group operation in $\topG^*$ is continuous,
hence the dual group of any topological group forms a topological group
as well.

\paragraph{}
In the following table we list three classical topological Abelian groups
together with their dual groups and characters:

\vspace{3mm}
\renewcommand{\arraystretch}{1.2}\normalsize
\centerline
{
\begin{tabular}{|c|c|c|}
\hline
$\topG$ & $\topG^*$       & characters                           \\
\hline
$\R$    & ${\R}^* \cong \R$ & $\chi_x(t) := e^{itx}, \, x \in \R$  \\
$\Z$    & ${\Z}^* \cong \mathbb{T}$ & $\chi_z(n) := e^{inz},
                 z \in \mathbb{T}$                                \\
$\mathbb{T}$    & ${\mathbb{T}}^* \cong \Z$ & $\chi_n(z) := z^n, 
                 \, n \in \Z$                                     \\
\hline
\end{tabular}
}

\vspace{3mm}
\centerline{\small Table 1: classical topological groups and dual groups.}
\vspace{2mm}

\paragraph{}
Concerning cartesian products of topological groups,
for each topological group $\topG$ and $\topH$
we have
\[
(\topG \times \topH)^*
\, \cong \,
\topG^* \times \topH^*.
\]
As a consequence,
the group isomorphisms ${\Rd}^* \cong \Rd$ hold for any integer $d$.

\subsection{}
%------------------------------------------------------------------------------%
A topological space is called \textit{locally compact},
if every subset owns a compact neighbourhood.
The abbreviation \textit{LCA group} usually stands
for a \textit{locally compact Abelian group}.
The meaning of locally compact groups is based on the fact
that these groups have non-trivial characters,
that is,
we have $\topG^* \neq \{\chi_1\}$ for such a group ${\topG}$,
where $\chi_1$ is the character with $\chi_1(x) \eq 1$ for all $x \xeof \topG$.

\subparagraph{}
Suppose that the Abelian group ${\topG}$ is \textit{locally compact}.
Then it can be shown (for example by the \textit{{\ARZELA}-Ascoli theorem})
that the dual group $\topG^*$ is also locally compact
{\wrt} its compact-open topology.
Let $\topG^{**}$ denote the \textit{bidual group} 
of a topological group $\topG$,
which is defined to be the dual group of its character group.

\paragraph{} %- Pontryagin's duality theorem
One of the main results in the theory of LCA groups is the famous
\textit{Pontryagin duality theorem}:

\begin{theorem1}\label{PONTRYAGIN_DUALITY_THEOREM}
If $\topG$ is a locally compact Abelian group,
then $\topG$ and $\topG^{**}$ are isomorphic to each other.
\end{theorem1}

\subparagraph{}
Notice that
each group element $g \xeof \topG$ gives rise to a character $\chi_x$
on the dual group $\topG^*$ by
$\chi_x(\xi) := \xi(x)$ for all $\xi \xeof \topG^*$.
Thus the mapping $x \mapsto \chi_x \xeof \topG^{**}$ is an injective group
homomorphism,
and Pontryagin duality means
that the canonical embedding $\iota \: \topG \xto \topG^{**}$,
$x \mapsto \chi_x$ is also surjective.

\subparagraph{}
It is easy to see
that the groups listed in table 1 are examples of Pontryagin duality.
This kind of duality also implies remarkable relationships between
the topological properties of a LCA group $\topG$ and its dual group $\topG^*$:

\begin{itemize}[leftmargin=15pt]
\item[1.]
$\topG$ is compact \iff $\topG^*$ is discrete.

\item[2.]
$\topG$ is discrete \iff $\topG^*$ is compact.
\end{itemize}

\paragraph{}
Finally,
let us mention
that it is possible to endow the convolution algebra $L^1(\topG,+,\ast)$
of any LCA group $\topG$
with an involutional mapping $* \:L^1(\topG) \xto L^1(\topG)$,
defined for $f \xeof L^1(\topG)$ by
\[
f^*(x) \eqdef \overline{f(-x)}, \quad \quad (x \in \topG).
\]

\section{Abstract Fourier Transformation}
%------------------------------------------------------------------------------%
Since $L^1(\topG,+,\ast)$ forms a commutative Banach algebra for every
LCA group $\topG$,
we may study its Gelfand transformation in detail,
which in this case is denoted \textit{Fourier transformation}.
If $\Gamma_{\topG}$ is the structure space of the algebra $L^1(\topG,+,\ast)$,
then Fourier transformation maps each function $f \xeof L^1(\topG)$
to a function $\hat{f} \xeof C_0(\Gamma_{\topG})$.
It constitutes a Banach algebra homomorphism
between $L^1(\topG)$ and $C_0(\Gamma_{\topG})$,
where convolution in $L^1(\topG)$ corresponds to multiplication
in $C_0(\Gamma_{\topG})$.

\paragraph{}
Moreover,
each character $\chi \xeof \topG^*$ gives rise to
a continuous multiplicative functional $\gamma \xeof \Gamma_{\topG}$,
that is defined according to
\[
\gamma(\chi) \eqdef \int_{\topG} \, f \, \overline{\chi} \, d\mu,
\]
and reversely each continuous multiplicative functional
$\gamma \xeof \Gamma_{\topG}$ can be represented in this way,
Thus,
we identify the structure space of the convolution algebra $L^1(\topG,+,\ast)$
with the dual group $\topG^*$,
and Fourier transformation appears as the Gelfand transformation
up to this identification:

\begin{itemize}[leftmargin=12pt]
\item[1.]
For $f \xeof L^1(\topG)$
the mapping $\hat{f} \: \widehat{\topG} \mapsto \C$ according to
\[
\hat{f}(\chi) := \int_{\topG} \, f \, \overline{\chi} \, d\mu,
\quad
\chi \in \widehat{\topG}
\]
is called the \textit{Fourier transform} of $f$.

\item[2.]
The mapping
$\FT \: f \xeof L^1(\topG) \xto \hat{f} \xeof C_0(\topG^*)$
is called \textit{Fourier transformation}.
\end{itemize}

\paragraph{} %- properties of Fourier transformation
Let us repeat some properties of Fourier transformation
that are already known from the general Gelfand transformation:

\begin{itemize}[leftmargin=12pt]
\item[1.]
For every $f \xeof L^1(\topG)$ the mapping $\hat{f}$ is continuous on $\topG^*$.

\item[2.]
the function $\hat{f}$ is vanishing at infinity
(\textit{Riemann-Lebesgue theorem}),
that is,
it maps absolutely integrable functions
to elements of the function space $C_0(\topG^*)$.

\item[3.]
$\FT$ is injective,
that is,
$\FT(f) \neq \FT(g)$ for $f \neq g$.

\item[4.]
$\FT$ is norm-preserving,
that is,
$\norm{\FT(f)}_{\infty} \leq \norm{f}_1$.

\item[5.]
$\FT(\alpha f + \beta g) \eq \alpha \FT(f) + \beta \FT(g)$,
    $\alpha,\beta \xeof \R$.

\item[6.]
$\FT(f \ast g) \eq \sqrt{2\pi} \, \FT(f) \cdot \FT(g)$.
\end{itemize}

\subparagraph{}
In addition,
the image of $L^1(\topG)$ under $\FT$ forms a separating subalgebra
of $C_(\Gamma_{\topG})$,
which lies even dense in $C_0(\Gamma_{\topG})$. 

\paragraph{} %- inverse Fourier transformation
By \textit{inverse Fourier transformation} is meant the procedure,
how certain functions in the algebra $L^1(\topG)$ can be retrieved
from the function values of the corresponding Fourier transforms.

\paragraph{}
Since the adjoint group $\topG^*$ of a LCA group $\topG$
with Haar measure $\mu$ is LCA again,
it owns a Haar measure as well,
hence the convolution algebra $L^1(\topG^*)$ is well-defined.
Let $f$ be an element of $L^1(\topG)$
fulfilling the additional condition $\hat{f} \xeof L^1(\topG^*)$.
Then there exists the Fourier transform of $\hat{f}$,
\[
\hat{\hat{f}}(x) \= 
\int_{\topG^*} \hat{f}(\gamma) \overline{x(\gamma)} \, d\nu,
\]
which forms a complex-valued mapping
defined on the biadjoint group $\topG^{**}$.
Due to Pontryagin duality,
$\topG^{**}$ is isomorphic to $\topG$,
and we may define the \textit{inverse Fourier transform} $F$
according to
\[
F(x) \= \int_{\topG^*} \hat{f}(\gamma) \, \overline{\gamma(x)} \, d \mu,
\quad (x \in \topG).
\]
It can be shown that the functions $f$ and $F$ are $\mu$-nearly everywhere
identical on $\topG$.

\paragraph{} %- the algebra of finite Borel measures
Let $\topX$ be any topological space and $\mu$ be a measure
define on the Borel $\sigma$-algebra $\Borel{\topX}$.
A Borel set $B \subseteq \topX$ is called \textit{inner regular} if
\[
\mu(B) \= \sup \,
          \{\mu(\compK): \compK \subset B, \, \compK \mbox{ kompact}\}.
\]
Then a measure $\mu$ is denoted as \textit{inner regular},
if all Borel sets are so.
Hence inner regular measures defined on $\topX$ are already uniquely
defined by their values for compact subsets of $\topX$.
Each inner regular measure on a topological space is said to be \textit{Radon}
if it adopts finite values for compact sets.

\subparagraph{}
Rather similar,
a measure $\mu$ is called \textit{outer regular},
if for each Borel set $B \subseteq \topX$ the equality
\[
\mu(B) \= \inf \,
          \{\mu(\openO): B \subset \openO, \, \openO \mbox{ open} \}.
\]
holds.
Finally,
a measure is called \textit{regular},
if it is inner and outer regular at the same time.

\paragraph{} %- convolution algebras of measures
Let $\MCR{\topG}$ be the Banach space of finite regular Borel measures,
where $\topG$ is a LCA group.
Let $\mu_1, \mu_2$ be measures in $\MCR{\topG}$ and let $\mu_1 \xtimes \mu_2$
be the product measure defined on the group $\topG \xtimes \topG$.
If $B \xeof \Borel{\topG}$ is measurable,
then the set $B'$ is defined by
\[
B' \eqdef \{(x,y) \in \topG \times \topG: x+y \in B\}.
\]
This set is measurable {\wrt} the $\sigma$-Algebra
in $\topG \xtimes \topG$.
Moreover,
the \textit{convolution} operation $\mu_1 \ast \mu_2$
of two measures $\mu_1$ und $\mu_2$ is defined to be the measure
\[
(\mu_1 \ast \mu_2)(B) \eqdef (\mu_1 \times \mu_2)(B')
\quad (B \in \Borel{\topG}),
\]
which owns the following properties:

\begin{itemize}[leftmargin=12pt]
\item[1.]
$\mu_1 \ast \mu_2 \xeof \MCR{\topG}$,

\item[2.]
$(\mu_1,\mu_2) \mapsto \mu_1 \ast \mu_2$ is commutative and associative,

\item[3.]
the inequality
$\norm{\mu_1 \ast \mu_2} \, \leq \, \norm{\mu_1} \cdot \norm{\mu_2}$ holds.
\end{itemize}

\subparagraph{} %- Dirac measure
Hence the Banach algebra $\MCR{\topG,+,\ast}$ is unital and commutative.
Its unit element is the \textit{Dirac measure} $\delta_{\0}$,
which distinguishes from other measures in this algebra by the property
that it is concentrated to the element $\0 \xeof \topG$.

\paragraph{}
In case of $\topG=(\Rd,+)$,
we may regard the algebra $\MCR{\Rd,+,\ast}$ as the Banach algebra,
that is norm-isomorphic to the completion
of the non-unital convolution algebra $L^1(\Rd,+,\ast)$.

\paragraph{} %- Fourier-Stieltjes transformation
\textit{Fourier-Stieltjes transformation} is another phrase for
Gelfand transformation on the unital and commutative Banach algebra
$\MCR{\topG,+,*}$ of finite and regular complex Borel measures
on a LCA group $\topG$.
In fact,
if $\mu \xeof \MCR{\topG,+,*}$,
then the mapping $\hat{\mu}$ defined on the dual group $\topG^*$ by
\[
\hat{\mu}(\gamma) \eqdef
\int_{\topG} -\gamma(x) \, d\, \mu(x), \quad \quad \gamma \in \topG^*
\]
is called \textit{Fourier-Stieltjes transform} of  $\mu$.

\paragraph{} %- properties of Fourier-Stieltjes transformation:
Let $B(\topG^*)$ be the set of functions $\hat{\mu}$
with $\mu \xeof \MCR{\topG,+,*}$.
Then the following properties are valid:

\begin{itemize}[leftmargin=12pt]
\item[1.]
Each function $\hat{\mu} \xeof B(\topG^*)$ is bounded and uniformly continuous.

\item[2.]
If $\sigma \eq \mu \ast \nu$,
then $\hat{\sigma} \eq \hat{\mu} \cdot \hat{\nu}$ holds,
that is,
the mapping $\mu \mapsto \hat{\mu}(\gamma)$ 
is a Banach algebra endomorphism on $\MCR{\topG,+,*}$
for every $\gamma \xeof \topX^*$.

\item[3.]
The set $B(\topG^*)$ is invariant under translation, multiplication and
complex conjugation.
\end{itemize}

\section{Examples of Fourier Transformation}
%------------------------------------------------------------------------------%
It is nearby to consider Fourier transformation
on $L^1(\topG,+,\ast)$ for the classical groups
$\Rd$, $\Z$ and $\mathbb{T}$:

\begin{itemize}[leftmargin=12pt]
\item[1.]
If $\topG \eq \Rd$,
then the character group of $(\Rd,+)$ is the group $(\Rd,+)$ itself.
Hence a function $f \xeof L^1(\Rd)$ owns the Fourier transform
$\hat{f} \: \Rd \xto \C$,
whose function values are defined as follows:
\[
\hat{f}(x) \eqdef
\frac{1}{\sqrt{2\pi}} \int_{- \infty}^{\infty} f(t) \, e^{-itx} \, dt,
\quad (x \xeof \Rd).
\]

\subparagraph{} %- convolution in L^1
Recall that for elements $f,g \xeof L^1(\Rd)$ 
a multiplication $f \ast g$ is defined by
\[
(f \ast g) \, (x) \eqdef \int_{\,\Rd} f(t) \cdot g(x-y) \, dy
\]
and denoted the \textit{convolution} in $L^1(\Rd)$.
This mapping is commutative, associative and fulfills
for all $f,g \xeof L^1(\Rd)$ the estimate
\[
\norm{f \ast g}_1 \, \leq \, \norm{f} \cdot \norm{g}.
\]
Since $L^1(\Rd)$ is complete,
the algebra $(L^1(\Rd),+,\ast)$
constitutes a commutative but non-unital Banach algebra.

\paragraph{} %- Fourier transformation in L1
L$^1$-Fourier transformation in the Euclidean space $\Rd$
is nothing but the Gelfand mapping
\[
\FT \: f \xeof L^1(\Rd) \xto \hat{f} \xeof C_0(\Rd),
\]
which sends each function $f \xeof L^1(\Rd)$
to the function $\hat{f} \xeof C_0(\Rd)$,
called the \textit{Fourier transform} of $f$ and defined by
\begin{equation}\label{VII_L1_FOURIER_TRANSFORM}
\hat{f}(\xi) \eqdef
\frac{1}{(2 \pi)^{d/2}}
\int_{\,\Rd} \, f(x) \, e^{-2 \pi i \scalar{\xi}{x}} \, dx
\quad \quad (\xi \xeof \Rd),
\end{equation}
where $\scalar{\xi}{x}$ is the inner product of $\xi,x \xeof \Rd$.

\paragraph{} %- properties of the Fourier transformation in L1
Let us again collect the remarkable properties of the mapping $\FT$
in the special case $\topG \eq \topG^* \eq \R$:

\begin{itemize}[leftmargin=12pt,itemsep=5pt]
\item[1.]
For each $f \xeof L^1(\Rd)$ the function $\hat{f}$ is continuous on $\Rd$,
so $\hat{f} \xeof C(\Rd)$.

\item[2.]
$\hat{f}$ is vanishing at infinity,
so $\hat{f} \xeof C_0(\Rd)$.

\item[3.]
$\FT$ is injective,
that is,
$\FT(f) \neq \FT(g)$ for $f \neq g$.

\item[4.]
$\norm{\FT(f)}_{\infty} \leq \norm{f}_1$ for $f \xeof L^1(\Rd)$.

\item[5.]
$\FT(\alpha f + \beta g) \eq \alpha \, \FT(f) + \beta \, \FT(g)$
for $\alpha,\beta \xeof \C$,
thus $\FT$ is linear.

\item[6.]
$\FT(f \ast g) \eq \FT(f) \cdot \FT(g)$,
so $\FT$ is an algebra homomorphism between the Banach algebras
$L^1(\Rd,+,\ast)$ and $C_0(\Rd,+,\cdot)$.

\item[7.]
If $f \xeof L^1(\Rd)$ and
$g \: x \xeof \Rd \mapsto \abs{x} f(x) \in L^1(\Rd)$,
then $\hat{f} \xeof C^1(\Rd)$ and
\[
\partial_j f(x) \= \widehat{-i \, \xi_j f(\xi)}
\quad (j=1,\ldots,d).
\]
\end{itemize}

\subparagraph{}
Especially for the LCA group $(\Rd,+)$,
where the dual group of $\Rd$ may be identified with $\Rd$ itself,
inverse Fourier transform of $\hat{f} \xeof L^1(\Rd)$
is obtained according to
\[
F(x) \= \int_Rd \hat{f}(t) \, e^{itx} \, dt, \quad (x \xeof \Rd).
\]

\item[2.]
Let $\topG \eq \Z$.
The normed Haar measure on this group is nothing but the counting measure,
and the character group $(\Z,+)$ is isomorphic
to the multiplicative group $(\mathbb{T},\cdot)$.
Hence the Fourier transform of a sequence $\{f_n\}_{n \in \N}$
is the mapping $\hat{f} \: \mathbb{T} \xto \C$,
which is defined pointwise by the \textit{Fourier series}
\[
\hat{f}(\omega) \eqdef
 \frac{1}{\sqrt{2\pi}} \sum_{n= - \infty}^{\infty} f_n \, e^{-i n \omega},
 \quad (\omega \in \mathbb{T}).
\]

\item[3.]
The character group of the group $(\mathbb{T},\cdot)$
is isomorphic to the group $(\Z,+)$,
hence the Fourier transform of a function
$f \: \mathbb{T} \xto \R$ is the sequence $\hat{f} \: \Z \xto \C$ with members
\[
\hat{f}_n \eqdef
\frac{1}{\sqrt{2\pi}} \int_{\mathbb{T}} f(z) \, e^{-inz} \, dz,
\quad (n \in \Z).
\]
\end{itemize}

\section{Fourier-Plancherel Transformation}
%------------------------------------------------------------------------------%
We want to extend the Fourier transformation to the Hilbert space $L^2(\Rd)$.
Recall that
this space consists of the measurable functions $u \: \Rd \xto \C$ fulfilling
\[
\int_{\,\Rd} \, \abs{u} ^2 \, dx \, < \, \infty,
\]
and the inner product in $L^2(\Rd)$ is given by
\[
\scalar{u}{v} \= \int_{\,\Rd} \, u\,\overline{v} \, dx.
\]
If the functions $u,v \: \Rd \xto \C$ are Lebesgue-integrable
then by definition they belong to the space $L^1(\Rd)$,
and their \textit{Fourier transforms} are well-defined
by formula \eqref{VII_L1_FOURIER_TRANSFORM}.

\subparagraph{}
However, 
if is not only a member of $L^1(\Rd)$, 
but also $u \xeof L^2(\Rd)$,
we define Fourier transformation a little bit different:
\begin{equation}\label{VII_L2_FOURIER_TRANSFORM}
\hat{u}(\xi) \eqdef
\frac{1}{(2 \pi)^{d/2}}
\int_{\,\Rd} \, e^{-2 \pi i \scalar{\xi}{x}} \, u(x) \, dx.
\end{equation}
Then for functions $u,v \xeof L^2(\Rd)$ 
the inner product between $u$ and $v$
and the inner product of their Fourier transforms $\hat{u}$ and $\hat{v}$
are related by \textsc{Plancherel}\textit{'s formula}
\[
\scalar{u}{v} \= \scalar{ \hat{u} }{\hat{v} }.
\]
Hence \textit{Fourier transformation} $\FT \: u \mapsto \hat{u} $
constitutes a partial isometry from the space
$L^{1,2}(\Rd) := L^1(\Rd) \, \cap \, L^2(\Rd)$
to the space $L^2(\Rd)$.

\subparagraph{}
It is a crucial fact now
that $L^{1,2}(\Rd)$ proves to be dense in $L^2(\Rd)$.
Since the mapping $\FT$ is linear and continuous on $L^{1,2}(\Rd)$ ,
it is even \textit{uniformly continuous} on that space
and can be extended to an \textit{unitary},
that is,
a bijective and isometric linear operator
$\FT_P \: L^2(\Rd) \xto L^2(\Rd)$,
called the \textsc{Fourier-Plancherel} \textit{operator}
or the \textit{$L^2$-Fourier transformation} in $L^2(\Rd)$.

\subparagraph{}
But in general a function $u \xeof L^2(\Rd)$ is not an element of $L^1(\Rd)$.
Thus it is not possible to define the $L^2$-Fourier transform $\hat{u}$
of every $u \xeof L^2(\Rd)$
by means of formula \eqref{VII_L2_FOURIER_TRANSFORM}.
However,
for $u \xeof L^2(\Rd)$ and $R > 0$ the function
\[
\hat{u} _R(\xi) \eqdef
\frac{1}{(2 \pi)^{d/2}}
\int_{\,\abs{x} \leq R}
e^{-2 \pi i \scalar{x}{\xi}} \, u(x) \, dx
\]
is well-defined,
and one may obtain the $L^2$-Fourier transform $\hat{u}$ 
as the limit of the functions $\hat{u} _R$ for $R \xto \infty$
{\wrt} the $2$-norm,
that is,
$\lim_{\,R \to \infty} \, \norm{ \hat{u} - \hat{u} _R}_2 \= 0$.

\paragraph{} %- Bessel potential spaces
The Fourier-Plancherel operator can be used
to introduce the class of \textit{Bessel potential spaces}.
First recall
that the mapping $\FT_P$ is well-defined in Hilbert space $L^2(\Rd)$,
therefore it has also a meaning in the linear subspaces $H^k(\Rd)$
for $k \xeof \N$.
We define for each $k \xeof \N$ the real-valued mapping
\begin{equation}\label{VII_DEFINITION_OF_FUNCTION_P}
p_k:x \in \Rd \mapsto (1 + \abs{x} ^{2k})^{1/2}
\end{equation}
and the space
\[
\cal{H}^k(\Rd) \eqdef
\{ f \in L^2(\Rd): p_k \, \hat{f} \in L^2(\Rd) \},
\]
which is Hilbertian with inner product
$\scalar{f}{g}_{p_k} := \scalar{p_k \hat{f}}{p_k \hat{g}}$
for $f,g \xeof \cal{H}^k(\Rd)$ and derived norm
\[
\norm{f}_{p_k}^2 \eqdef \scalar{f}{f}_{p_k} \=
\int_{\,\Rd} \, (1+\abs{x} ^{2k}) \, \abs{\hat{f}} ^2 \, dx.
\]

\subparagraph{}
It turns out
that the so-called \textit{Bessel potential space}
$\cal{H}^k(\Rd)$ is the range of
the \textsc{Fourier-Plancherel} operator $\FT_P$
when restricted to $H^k(\Rd) \subset L^2(\Rd)$.
Hence $\FT_P$ is a bijective and norm-preserving mapping
from $H^k(\Rd)$ to $\cal{H}^k(\Rd)$, {\ie}
\[
\abs{f} _k \= \lVert \hat{f} \rVert _{p_k}
\quad \, \qfa f \in H^k(\Rd).
\]

\subparagraph{}
In order to extend the definition of the spaces $H^k(\Rd)$
to indices $s \xeof \R^+$,
let $p_s$ be defined as in \eqref{VII_DEFINITION_OF_FUNCTION_P}
but with $s > 0$ instead of $k \xeof \N$.
We define
\[
\cal{H}^s(\Rd) \eqdef
\{ \, f \in L^2(\Rd): p_s \hat{f} \in L^2(\Rd) \, \},
\]
then $\cal{H}^s(\Rd)$ given this way
coincides for $s \eq k \xeof \N$ with the classical Sobolev space $H^k(\Rd)$,
which motivates to define $H^s(\Rd) := \cal{H}^s(\Rd)$ for $s \notin \N$.
For $s<0$,
however,
we introduce $H^s(\Rd) := H^{-s}(\Rd)^*$.

\section{The Fourier Transform in Distributional Spaces}
%------------------------------------------------------------------------------%
We finally elaborate that Fourier transformation can be extended
to the space $\TD$ of tempered distributions.

\subparagraph{}
First notice 
that the Fourier transform $\hat{\phi}$
of a Schwartz function $\phi \xeof \SRd$ is well defined,
since $\SRd$ is a linear subspace of the convolution algebra $L^1(\Rd,+,\ast)$.
Moreover, 
Fourier transformation restricted to the Schwartz space $\SRd$
forms a linear mapping $\FT \:\SRd \xto \SRd$,
which sends each function $\phi \xeof \SRd$ to a further function
$\hat{\phi} \xeof \SRd$ according to the well-known formula
\begin{equation}\label{VII_FOURIER_TRANSFORM_OF_SCHWARTZ_FUNCTION}
\hat{\phi}(x) \= \FT(\phi)(x) :=
\frac{1}{(2\pi)^{d/2}}
\int_{\,\Rd} e^{-i <\xi,x>} \phi(\xi) \, d \xi
\quad (x \xeof \Rd).
\end{equation}

\subparagraph{}
The extraordinary meaning of Schwartz functions is founded by the fact that
the Fourier transformation has excellent properties on the space $\SRd$.
More precisely,
besides the well-known features of Fourier transformation in $L^1(\Rd,+,\ast)$,
the mapping $\FT$ is \textit{bijective} on $\SRd$.
Hence the \textit{inverse Fourier transform} exists for all $\phi \xeof \SRd$
and is described by the \textit{Fourier inversion formula}
\[
\check{\phi}(\xi) \= 
\FT^{-1}(\phi)(\xi) \eqdef
\frac{1}{(2\pi)^{d/2}}
\int_{\,\Rd} e^{i \scalar{x}{\xi}} \, \phi(x) \, dx
\quad \quad (\xi \xeof \Rd).
\]

\subparagraph{}
We next extend Fourier transformation in $\SRd$ to the space $\TD$
of tempered distributions by transposition,
that is
\[
(\FT \disT) \, \phi \, := \, \disT(\FT \phi),
\quad \quad \; \phi \xeof \SRd, \; \disT \xeof \TD.
\]
This mapping is bijective in the space $\TD$
and even continuous {\wrt} the weak*-topology on $\TD$.
The inverse Fourier transform of a tempered distribution $\disT$
is represented by the formula
\[
\FT^{-1}(\disT)(\phi) \, := \, \disT(\FT^{-1} \phi)
\quad \quad (\phi \xeof \SRd).
\]
For example,
the Fourier transform $\FT(\delta_0)$ of the Dirac distribution
$\delta_0 \xeof \cal{S}'(\R)$ is the continuous linear functional
\[
\FT(\delta_0)(\phi) \, := \,
\delta_0(\FT(\phi)) \= \hat{\phi}(0) \=
\frac{1}{\sqrt{2\pi}} \, \int_{\,\R} \phi(x) \, dx,
\quad \phi \xeof \SR1\,,
\]
so $\FT(\delta_0)$ is nothing but the distribution induced by
the constant function $x \mapsto \frac{1}{\sqrt{2\pi}}$.

% BIBLIOGRAPHY
%------------------------------------------------------------------------------%
\vspace{5mm}
\begin{thebibliography}{99}
\begin{small}

\bibitem{VII_Ka} E.\,Kaniuth:
\textit{A Course in Commutative Banach Algebras},
Springer GTM 246, 2009.

\bibitem{VII_Ru} W.\,Rudin:
\textit{Fourier Analysis on Groups},
Wiley-Interscience, 1962.

\bibitem{VII_SW} E.\,Stein, G.\,Weiss:
\textit{Introduction to Fourier Analysis on Euclidean Spaces},
Princeton Mathematical Series 32, Princeton University Press, 1971.

\end{small}
\end{thebibliography}

%------------------------------------------------------------------------------%
%VIII. LINEAR OPERATORS IN HILBERT SPACE ---
%------------------------------------------------------------------------------%
%  §1. Numerical Range and External Field of Regularity ---
%  §2. Unitary Operators and Isometries ---
%  §3. Symmetric Linear Operators ---
%  §4. Self-Adjoint Operators ---
%  §5. Essentially Self-Adjointness ---
%  §6. The von Neumann Extension Theory ---
%  §7. Perturbations of Self-Adjoint Operators ---
%  ---

\newpage
\thispagestyle{empty}
\def\TitleVIII{Linear Operators in Hilbert Space}
\fancyhead[C] {\small{\textsc{VIII. \TitleVIII}}}
\part{\TitleVIII}
%------------------------------------------------------------------------------%
The presence of the inner product in a Hilbert space can be used to
analyze linear operators with respect to this additional feature.
For example,
a linear operator $\linA$ with domain of definition $\domain{\linA}$ acting
in a Hilbert space $\hilH$ is called \textit{normal},
if $\linA \linA^* \eq \linA^* \linA$ holds,
where $\linA^*$ is the Hilbert adjoint, an operator acting in $\hilH$ as well,
which is well-defined whenever $\linA$ is densely defined.

\subparagraph{}
Important classes of normal operators are the unitary, symmetric and
self-adjoint ones.
Especially the symmetric and self-adjoint linear operators has been
extensively studied and give rise to a very complete theory.

\section{Numerical Range and External Field of Regularity}
%------------------------------------------------------------------------------%
Besides the spectrum $\sigma(\linA)$,
the \textit{numerical range} $\Theta(\linA)$ plays an important role
in the general theory of a Hilbert space operator $\linA$ 
and is defined by
$\Theta(\linA) :=
\{ \scalar{\linA}{x}: x \xeof \domain{\linA}, \norm{x} \eq 1 \}$.
It is a well-known fact,
given by the \textit{Toeplitz-Hausdorff theorem},
that $\Theta(\linA)$ is a convex set within the complex plane.
Moreover,
let
\[
\Delta(\linA) \eqdef \C \setminus \overline{\Theta(\linA)}
\]
be the \textit{external field of regularity} of $\linA$.
If $\Delta(\linA) \neq \emptyset$,
then it consists of one or two connected components.
To be precise,
if $\Delta(\linA)$ has two connected components,
then the set $\Theta(\linA)$ is either a strip or a ray,
otherwise it forms a half-strip or a half-ray.

\subparagraph{}
For $\zeta \xeof \Delta(\linA)$ we of course have
$\scalar{\linA x}{x} \neq \zeta$
for $x \xeof \domain{\linA}$ and $\norm{x} \eq 1$,
and the following chain of inequalities
\[
0 \, < \, \mbox{dist}(\zeta,\linA) \, \leq \,
\abs{\scalar{\linA x}{x} - \zeta} \=
\abs{\scalar{(\linA-\zeta)x}{x}} \, \leq \,
\norm{(\linA -\zeta) x}
\]
is valid,
where $\mbox{dist}(\zeta,\linA)$ denotes the distance
between the singleton $\{\zeta\} \xsubset \Delta(\linA)$
and the set $\overline{\Theta(\linA)}$,
which in fact is positive,
since these sets are closed and have empty intersection.
Hence the operator $\linA -\zeta$ is bounded below,
which implies
that $\linA-\zeta$ is injective and has closed range.

\subparagraph{} %- properties of symmetric operators
In other words,
the set $\Delta(\linA)$ is a subset of the \textit{field of regularity}
$\Pi(\linA) := \{\zeta \xeof \C: \linA-\zeta \mbox{ is bounded below}\}$.

\section{Unitary Operators and Isometries}
%------------------------------------------------------------------------------%
Let us first consider the class of \textit{unitary} operators.
A bijective linear operator $\linU:\hilH \to \hilH$ acting in a Hilbert space
$\hilH$ is called \textit{unitary},
if
\[
\scalar{\linU x}{\linU y} \= \scalar{x}{y}
\quad \mbox{for all } x,y \in \hilH.
\]
Another definiton for this notion is given by the condition
$\linU^* \eq \linU^{-1}$.

\subparagraph{} %- unitary equivalence of linear operators
The meaning of unitary operators is motivated by the fact
that one can define \textit{unitary equivalence} between operators
in a Hilbert space.
In particular,
two operators $\linA_1$ and $\linA_2$ are called \textit{unitary equivalent},
if there is an unitary operator $U$ such that
\[
\linA_1 \= U \circ \linA_2 \circ U^*.
\]

\subparagraph{}
Each eigenvalue of a unitary operator $\linU$ is of modulus $1$
and the operator norm $\norm{\linU}$ is equal to $1$.
The set of unitary operators $\BU{\hilH}$ acting in $\hilH$ forms a group,
called the \textit{unitary group} of $\hilH$.

\subparagraph{} %- isometric operators
More generally,
an \textit{isometry} is a linear mapping $\isoJ:\hilH_1 \to \hilH_2$
between arbitrary Hilbert spaces $\hilH_1$ and $\hilH_2$,
which preserves the inner products,
that is
\[
\scalar{\isoJ x}{\isoJ y} \= \scalar{x}{y} \quad (x,y \in \hilH_1).
\]
This condition implies that an isometry is always injective.
If $\hilH_1 \eq \hilH_2 \eq \hilH$,
the mapping $\isoJ$ is called a \textit{partial isometry}
(which must not necessary be a surjective mapping).
Finally,
a partial isometry $\isoJ$ is an unitary operator
if and only if
$\isoJ$ is surjective,
so we regard an unitary operator as a special case of an isometry.

\section{Symmetric Linear Operators}
%------------------------------------------------------------------------------%
Let $\hilH$ be a real or complex Hilbert space.
A linear operator acting in $\hilH$ and with domain of definition
$\domain{\linA} \xsubset \hilH$ is called \textit{symmetric}, if
\begin{equation}\label{VIII_SYMMETRY_PROPERTY}
\scalar{\linA x}{y} \= \scalar{x}{\linA y} \qfa x,y \in \domain{\linA}.
\end{equation}

\paragraph{} %- symmetric operators which are continuous
We first consider the symmetry property of operators,
which are continuous and everywhere defined.
Let $\BO{\hilH}$ be the Banach space of \textit{continuous} linear operators
$\linA \: \hilH \xto \hilH$ furnished with the {operator norm}.
Then, if $\linA$ is symmetric,
we have
\[
- \norm{\linA} \, \leq \, \scalar{\linA x}{x} \, \leq \, \norm{\linA},
\]
for $x \in \hilH$ with $\norm{x}=1$,
or equivalently,
\[
\Theta(\linA) \subset [-\norm{\linA}, \norm{\linA}\,].
\]
Hence one may introduce in the class of bounded symmetric operators
acting in $\hilH$ the partial order
\[
\linA_1 \, \geq \linA_2
\;\; \mbox{\iff} \;\;
\scalar{(\linA_1 - \linA_2) \, x}{x} \, \geq \, 0
\quad (x \in \hilH),
\]
which implies that
each symmetric linear operator $\linA \in \BO{\hilH}$ can be enclosed
by the multiplication operators $\linA_{-} := - \norm{\linA}\,\1_H$
and $\linA_+ := \norm{\linA}\,\1_H$.

\subparagraph{}
On the other hand,
the \textit{Hellinger-Toeplitz} \textit{theorem} quotes
that if $\linA$ is symmetric and $\domain{\linA}$ is the whole of $\hilH$,
then $\linA$ must be bounded.
In fact,
let $\linA$ be symmetric with $\domain{\linA} \eq \hilH$,
then for each convergent sequence $(x_n)$ tending to $0$
and with $(\linA x_n)$ tending to $x \in \hilH$ 
we have
\[ 
                    \scalar{x}        {x}       \= 
\lim_{n \xto \infty} \scalar{\linA x_n}{x}       \=
\lim_{n \xto \infty} \scalar{x_n}      {\linA x} \= 
                    \scalar{0}        {\linA x} \= 0,
\]
which implies $x \eq 0$, so $\linA$ is closed 
and even continuous by the closed graph theorem.
Hence we may always assume
that the domain of definition $\domain{\linA}$
of a non-continuous symmetric linear operator $\linA$
is a \textit{proper} subspace of $\hilH$.

\paragraph{}
The symmetry of a linear operator with domain of definition dense in $\hilH$
can be characterized by other properties than \eqref{VIII_SYMMETRY_PROPERTY}:

\begin{theorem1}
Let $\linA$ be a densely defined linear operator with domain of definition
$\domain{\linA}$ in Hilbert space $\hilH$.
Then the following statements are equivalent:

\begin{itemize}
\item[1.]
$\linA$ is symmetric;
\item[2.]
$\scalar{\linA x}{x} \in \R$ for all $x \in \domain{\linA}$;
\item[3.]
$\linA \prec \linA^*$.
\end{itemize}
\end{theorem1}

\subparagraph{}
The last property is of course also valid for symmetric operators
not necessarily being densely defined in the entire space.

\begin{comment}
\paragraph{}
We want to prove the equivalence of 1. and 2.

\subparagraph{}
If $\linA$ is symmetric,
then $\scalar{\linA x}{x} \eq \scalar{x}{\linA x} \eq \overline{\scalar{\linA x}{x}}$,
that is, this expression is real-valued.
Reversely,
if $\linA$ is densely defined in $\hilH$ and
$\scalar{\linA x}{x} \in \R$ for all $x \in \domain{\linA}$ is true,
then $\linA$ must be symmetric.
This can be shown as follows: let $x,y \in \domain{\linA}$, then
\[
\begin{cases}[1.3]
\quad
\R \ni \scalar{\linA (x+y)}{x+y}   & = \; \scalar{\linA x}{x}   +   \scalar{\linA x}{y} +    \scalar{\linA y}{x} + \scalar{\linA y}{y},\\
\quad
\R \ni \scalar{\linA (ix+y)}{ix+y} & = \; - \scalar{\linA x}{x} + i \scalar{\linA x}{y} + -i \scalar{\linA y}{x} + \scalar{\linA y}{y},\\
\end{cases}
\]
which implies $\im \scalar{\linA x}{y} \eq - \im \scalar{\linA y}{x}$
and $\re \scalar{\linA x}{y} \eq  \re \scalar{\linA y}{x}$.
Then
\[
\scalar{\linA x}{y} = \re \scalar{\linA x}{y} + i \im \scalar{\linA x}{y} =
\re \scalar{\linA y}{x} - i \im \scalar{\linA y}{x} =
\overline{\scalar{\linA y}{x}} = \scalar{x}{\linA y},
\]
which shows that $\linA$ is symmetric.
\end{comment}

\paragraph{}
Since the adjoint $\linA^*$ of a symmetric operator $\linA$
forms a closed extension of $\linA$,
the operator $\linA$ itself is closable.

\subparagraph{}
Moreover,
the closure $\overline{\linA}$ proves to be symmetric.
In fact,
to show that $\overline{\linA}$ is symmetric,
let $(x_n)$ and $(y_n)$ be sequences in $\domain{\linA}$
with $x_n \xto x \in \domain{\overline{\linA}}$,
$\linA x_n \xto \overline{\linA} x$
and $y_n \xto y \in \domain{\overline{\linA}}$,
$\linA y_n \xto \overline{\linA} y$.
Since $\scalar{\linA x_n}{y_n} \eq \scalar{x_n}{\linA y_n}$
for all $n \in \N$,
taking the limit for $n \xto \infty$ shows
that the closure of $\linA$ must be symmetric, too.

\paragraph{} %- eigenvalues of symmetric operators
Recall that a number $\lambda \in \C$ is called an \textit{eigenvalue}
of a linear operator $\linA$,
if there is $x \in \domain{\linA}$,
called the \textit{eigenvector} of $\linA$ {\wrt} $\lambda$,
such that $\linA x \eq \lambda x$.
The following properties of eigenvalues and eigenvectors
of a {symmetric} linear operator $\linA$ are easy to prove:

\begin{itemize}[leftmargin=12pt]
\item[1.]
each eigenvalue $\lambda$ of $\linA$ is real;

\item[2.]
eigenvectors $x_1,x_2$
belonging to different eigenvalues $\lambda_1 \neq \lambda_2$
are \textit{orthogonal} to each other,
that is,
$\scalar{x_1}{x_2} \eq 0$.
\end{itemize}

\section{Self-Adjoint Operators}
%------------------------------------------------------------------------------%
If a linear operator $\linA$ acting in $\hilH$ is symmetric,
then every other symmetric linear operator $\linS$ having the property
$\linA \prec \linS$ is called a \textit{symmetric extension} of $\linA$.
Notice 
that $\linA \prec \linS$ always implies $\linS^* \prec \linA^*$,
which is also true for non-symmetric operators $\linA$ and $\linS$.
Furthermore,
let $\linS_1,\ldots,\linS_k$ be a finite sequence
of symmetric extensions of $\linA$,
where $\linS_{i+1}$ extends $\linS_i$ for $i=1,\ldots,k-1$.
Since $\linS_k \prec \linS_k^*$ and
$\linS_{i+1}^* \prec \linS_i^*$ for $i=1,\ldots,k-1$,
we obtain by concatenating all these extensions the chain
\[
\linA \, \prec \, \linS_1
\, \prec \,
\linS_2
\, \prec \,
\ldots
\, \prec \,
\linS_k
\, \prec \,
\linS_k^*
\, \prec \,
\ldots
\, \prec \,
\linS_2^*
\, \prec \,
\linS_1^*
\, \prec \,
\linA^*.
\]
Thus one may hope to find a \textit{maximal} symmetric extension
$\linS_{sa}$ of $\linA$ such that
\[
\linA \, \prec \, \linS_{sa} \= \linS_{sa}^* \, \prec \, \linA^*
\]
holds,
and in this case
the operator $\linS_{sa}$ cannot be extended to a larger symmetric operator.
This expectation suggests to introduce the notion of \textit{self-adjointness}:
a symmetric linear operator $\linA$ is called \textit{self-adjoint},
if $\linA^* \prec \linA$,
or equivalently,
if $\linA \eq \linA^*$.
For an everywhere defined symmetric and linear operator $\linA$ in $\hilH$,
however,
we have $\linA^* \eq \linA$,
so there is no difference between symmetry and self-adjointness in this case,
and the \textsc{Hellinger-Toeplitz} theorem claims that $\linA$ is bounded,
which is a consequence of the {closed range theorem}.
In fact,
since a self-adjoint operator $\linA$ is the adjoint of itself,
it is closed with closed domain of definition $\domain{\linA}$,
hence it is continuous respectively bounded by the closed range theorem.

\paragraph{} %- range condition
To ensure self-adjointess (rsp. essentially self-adjointness)
of a symmetric linear operator,
the \textit{range condition} serves as the default tool
to verify this desired property:

\begin{theorem1}\label{RANGE_CONDITION}
Every densely defined, closed and symmetric operator $\linA$
acting in $\hilH$ is self-adjoint iff
\[
\range{\linA + i} \= \range{\linA - i} \= \hilH.
\]
\end{theorem1}

\begin{comment}
\Proof:
Since $\linA \prec \linA^*$,
we have only to show $\linA^* \prec \linA$, or equivalently,
$y \xeof \domain{\linA^*}$ implies $y \xeof \domain{\linA}$.
Now,
since $(\linA-i)^*x \eq (\linA^* +i)x$ for $x \xeof \domain{\linA}$,
we have
\[
\scalar{\,(\linA-i)x}{y} = \scalar{\linA x}{y} -i\scalar{x}{y} =
\scalar{x}{\linA^*y} -\scalar{x}{-iy} = \scalar{x}{\,(\linA^* +i)y}.
\]
But since $\range{\linA+i} \eq \hilH$,
there exists $z \xeof \domain{\linA}$ such that
$(\linA + i)z \eq (\linA^* +i)y$.
Inserting this in the right-hand side of the last equation,
we obtain
\[
\scalar{\,(\linA-i)x}{y} = \scalar{x}{(\linA + i)z} = \scalar{\,(A-i)x}{z}.
\]
But since in addition $\range{\linA-i} \eq \hilH$, we obtain
$\scalar{u}{y-z} \eq 0$ for every vector $u := (\linA-i)x \xeof \hilH$,
which implies $y=z$.
\qed
\end{comment}

\section{Essentially Self-Adjointness}
%------------------------------------------------------------------------------%
The closure of a symmetric linear operator is not necessarily self-adjoint.
Hence a closable and symmetric operator $\linA$ is denoted
\textit{essentially self-adjoint},
if $\overline{\linA}$ proves to be self-adjoint.
There is a comparable criterion for essentially self-adjointness available:

\begin{theorem1}
A densely defined and symmetric linear operator $\linA$
is essentially self-adjoint \iff
\[
\overline{ \range{\linA + i} } \= \overline{ \range{\linA - i} } \= \hilH.
\]
\end{theorem1}

\subparagraph{} % second criterion for essentially self-adjointness
Essentially self-adjointness can also be characterized by means of the adjoint:
a densely defined and symmetric linear operator $\linA$ in $\hilH$
is essentially self-adjoint
iff the adjoint $\linA^*$ is symmetric,
too.

\section{The von Neumann Extension Theory}
%------------------------------------------------------------------------------%
We want to sketch in this paragraph the principles
of the \textit{extension theory} for symmetric linear operators.
The question is,
whether a given symmetric operator owns one or more self-adjoint extensions.

\subparagraph{}
First notice that the closure of a symmetric linear operator is not
necessarily self-adjoint.
To ensure self-adjointess of a symmetric operator $\linA$,
the kernels of both operators $\linA^* \pm i$ are of special interest.
In fact,
if these are trivial,
then $\linA \pm i$ are surjective and
$\pm i$ are both regular values of $\linA$.

\subparagraph{}
Based on this observation,
\textsc{F.\,Riesz} and \textsc{J.\,v.\,Neumann} 
have developed an extension theory,
which answers the question
under which conditions a symmetric and closed linear operator
can be extended to a self-adjoint operator and how such an extension
can be constructed.
Moreover,
the possible uniqueness of such an extension if of fundamental interest.

\subparagraph{}
Recall that
for a densely defined and closed symmetric operator $\linA$
the relation $\domain{\linA} \subset \domain{\linA^*}$ holds.
But more can be said,
since  the domain $\domain{\linA^*}$ of the adjoint operator
can be represented as sum of three complementary subspaces:
\begin{equation}\label{VIII_NEUMANN}
\domain{\linA^*} \= \domain{\linA} \, \oplus \, \kernel{\linA^*+i}
                                   \, \oplus \, \kernel{\linA^*-i}.
\end{equation}

\subparagraph{} %- defect indices
Thus the linear operators $\linA^*-i$ and $\linA^*+i$
are of special interest.
The numbers
\[
n_+ \eqdef \dim \kernel{\linA^*+i}\,,\quad \quad
n_- \eqdef \dim \kernel{\linA^*-i}
\]
are called the \textit{defect indices} of $\linA$,
where the values $n_{\pm} \eq \infty$ are admissible, too.
One easily derives from \eqref{VIII_NEUMANN}
that a symmetric and closed linear operator $\linA$
with defect indices $(n_+,n_-)$ can be extended to a self-adjoint operator
iff $n_+ \eq n_-$,
and that the operator $\linA$ itself is self-adjoint
iff in addition $n_+ \eq n_- \eq 0$ holds. 
On the other hand,
$\linA$ is called \textit{maximal symmetric},
if at least one of the values $n_+$ and $n_-$ is zero,
which is equivalent to claim that $\linA$ has no proper symmetric extension.

\subparagraph{}
To extend in case of $n_+ \eq n_-$ the symmetric linear operator $\linA$
to a self-adjoint one,
it suffices to have a isometric and linear mapping
\begin{equation}\label{VIII_ISOMETRIC_MAPPING}
U\: \kernel{\linA^*+i} \xto \kernel{\linA^*-i}.
\end{equation}
Then,
if $U$ is given by \eqref{VIII_ISOMETRIC_MAPPING},
we define $\linA_s \,x := \linA x$ for $x \xeof \domain{\linA}$ and
\[
\linA_s \,z \eqdef iz -i \, U z
\]
for $z \xeof \kernel{\linA^* + i}$,
which implies that $\linA_s$ is also symmetric
and hence forms a symmetric extension of $\linA$.
One can show that $\linA_s \eq \linA^*$,
so $\linA_s$ is even self-adjoint.
We emphasize that
the set of all self-adjoint extensions is given
by the totality of isometric linear operators $U$ as defined above.

\section{Perturbations of Self-Adjoint Operators}
%------------------------------------------------------------------------------%
This subject deals with operators of the form $\linA + \linB$,
where $\linA$ is assumed to be self-adjoint and $\linB$ to be symmetric
with domain of definition $\domain{\linB}$ being a superset of $\domain{\linA}$,
so $\linA + \linB$ is well defined
with domain of definition $\domain{\linA + \linB} \eq \domain{\linA}$.
It is our aim to decide
under which conditions on the operator $\linB$
the sum $\linA + \linB$ proves to be self-adjoint.

\subparagraph{} %- A-boundedness
Let $\hilH$ be a complex Hilbert space and the operator
$\linA \: \domain{\linA} \xto \hilH$ be self-adjoint.
Then the following concept is useful:
a symmetric operator $\linB\:\domain{\linB} \xto \hilH$
with $\domain{\linA} \subset \domain{\linB}$
is called $\linA$\textit{-bounded},
if there are constants $\theta, c \geq 0$ such that
\begin{equation}\label{VIII_A_BOUNDEDNESS}
\norm{\linB \,x} \, \leq \, \theta \norm{\linA \,x} \, + \, c \norm{x}
\qfa x \xeof \domain{\linA}.
\end{equation}
In order to ensure $\linA$-boundedness of $\linB$,
it is sufficient to ensure \eqref{VIII_A_BOUNDEDNESS}
only for elements $x$ in a core $\linA_0$ of $\linA$.

\subparagraph{} %- A-bound
The infimum of all numbers $\theta$,
for which there is $c$ such that \eqref{VIII_A_BOUNDEDNESS} holds,
is called the $\linA$\textit{-bound} of $\linB$.
Notice that if $\linB$ is a bounded operator,
then it is $\linA$-bounded with $\linA$-bound $\theta <1$ for any $\linA$.
In addition,
the validness of the estimate
\[
\norm{\linB \,x}^2 \, \leq \,
\theta^2 \norm{\linA \,x}^2 \, + \, c^2 \norm{x}^2
\qfa x \xeof \domain{\linA}.
\]
is a sufficient condition for \eqref{VIII_A_BOUNDEDNESS} to be true.

\paragraph{} %- Kato-Rellich theorem
Let $\linA$ be self-adjoint in Hilbert space $\hilH$ and $\linB$ 
a further symmetric linear operator with domain of definition
$\domain{\linB} \supset \domain{\linA}$.
The \textsc{Kato-Rellich} \textit{theorem}
provides a criterion on self-adjointness of the operator $\linA + \linB$,
which has been independently found by \textsc{T.\,Kato}
and \textsc{F.\,Rellich}:

\begin{theorem1}\label{KATO-RELLICH_THEOREM_1}
Let the linear operator $\linA$ be self-adjoint and 
the linear operator $\linB$ with $\domain{\linB} \supset \domain{\linA}$ 
be symmetric and $\linA$-bounded with $\theta < 1$.
Then $\linA + \linB$ with domain of definition 
$\domain{\linA + \linB} \eq \domain{\linA}$ is self-adjoint.
\end{theorem1}

\subparagraph{} %- Kato-Rellich theorem for essentially self-adjoint operators
A similar result is obtained for the problem,
whether \textit{essentially self-adjointness} of a linear operator $\linA$
is preserved by a symmetric and $\linA$-bounded perturbation $\linB$:

\begin{theorem1}\label{KATO-RELLICH_THEOREM_2}
Let $\linA\:\domain{\linA} \xto \hilH$ be essentially self-adjoint
and assume that $\linB$ is symmetric and $\linA$-bounded,
that is,
$\domain{\linA} \subset \domain{\linB}$ and \eqref{VIII_A_BOUNDEDNESS} holds
with bound $\theta < 1$.
Then the operator $\linA + \linB$
is essentially self-adjoint on $\domain{\linA}$,
where
$
\overline{\linA + \linB} \= \overline{\linA} + \overline{\linB}.
$
\end{theorem1}

% BIBLIOGRAPHY
%------------------------------------------------------------------------------%
\vspace{5mm}
\begin{thebibliography}{99}
\begin{small}

\bibitem{VIII_EE} D.\,E.\,Edmunds, W.\,D.\,Evans:
\textit{Spectral theory and differential operators}.
Oxford University Press, 1987.

\bibitem{VIII_We} J.\,Weidmann:
\textit{Linear Operators in Hilbert Spaces}.
Springer Verlag, 1980.

\end{small}
\end{thebibliography}

%------------------------------------------------------------------------------%
%  IX. SPECTRAL THEORY OF SELF-ADJOINT OPERATORS ---
%------------------------------------------------------------------------------%
%  §1. Spectral Properties of Normal Operators ---
%  §2. Resolvent and Spectrum of Self-Adjoint Operators ---
%  §3. The Discrete Spectral Theorem ---
%  §4. Projections and Projection-Valued Measures ---
%  §5. Functional Calculus and Spectral Operators ---
%  §6. The Spectral Theorem for Self-Adjoint Operators ---
%  §7. Positive Self-Adjoint Operators and Cauchy Problems ---
%  ---

\newpage
\thispagestyle{empty}
\def\TitleIX{Spectral Theory of Self-Adjoint Operators}
\fancyhead[C] {\small{\textsc{IX. \TitleIX}}}
\part{\TitleIX}
%------------------------------------------------------------------------------%
In this section we give an outline on the spectral properties of self-adjoint
linear operators.
A linear operator $\linA$ of this type is also suitable to be decomposed
into sums of projection operators,
where the sum may be mostly countable or be given by an integral representation.
Having such a decomposition available,
it is possible to define \textit{spectral operators} of the form $f(\linA)$,
where the function $f\:\sigma(\linA) \xto \C$ is at least measurable.
The mapping $f \mapsto f(\linA)$ is called a \textit{functional calculus}
for $\linA$ and provides a useful tool to solve evolution equations
governed by a \textit{positive} and self-adjoint linear operator $\linA$.

\section{Spectral Properties of Normal Operators}
%------------------------------------------------------------------------------%
Our next topic will be the spectral theory
of symmetric and self-adjoint operators.
Both kinds of operators can be put together 
by the notion of \textit{normal} operator,
distinguished by the property $\linA \circ \linA^* \eq \linA^* \circ \linA$.
Symmetric and selfadjoint operators are extremely admissable
for the study of their spectral properties.

\paragraph{} %- the operator \linA-\lambda is semi-Fredholm 
In order to analyze the spectrum of a symmetric linear operator,
we start with a consideration of operators of the form $\linA - \lambda$,
where $\linA$ is closed, densely defined and symmetric and $\lambda \xeof \C$
is arbitrary.
Basically,
let $\hilH$ be a complex Hilbert space.

\subparagraph{}
First remember that a number $\lambda \xeof \C$ with $\im \lambda \neq 0$
does not belong to the numerical range $\Theta(\linA)$ of $\linA$,
hence $\lambda$ is contained
in the field of external regularity $\Delta(\linA)$.
This means that $\linA - \lambda$ is bounded below,
and since we assume $\linA$ to be closed,
$\linA - \lambda$ is injective and normally solvable
by the closed range theorem,
thus upper semi-Fredholm with $\alpha(\linA-\lambda) \eq 0$.
By \textsc{Gohberg}'s punctured neighbourhood theorem,
the index $\chi(\linA - \lambda)$
is locally constant for each fixed $\lambda \notin \R$.
In addition,
it can be shown that this magnitude is even globally constant
in the open upper and lower complex half-plane
$\mathbb{H}^+$ and $\mathbb{H}^+$,
respectively.

\subparagraph{}
Before proceeding with the study of the range of the operator $\linA-\lambda$,
notice that
due to the orthogonality relations
the linear subspaces $\range{\linA-\lambda}$
and $\kernel{\linA^*-\overline{\lambda}}$ are coupled by
\[
\hilH \=
\range{\linA - \lambda}
\, \oplus \,
\kernel{\linA^* - \overline{\lambda}}.
\]

\subparagraph{}
In the following
we distinguish between the two cases {\wrt} surjectivity of $\linA - \lambda$,
or equivalently,
to the injectivity of $\linA^* - \overline{\lambda}$:

\begin{itemize}[leftmargin=12pt]
\item[1.]
either $\range{\linA - \lambda} \eq \hilH$,
which implies $\kernel{\linA^* - \overline{\lambda}} \eq \{\0\}$,
then $\linA - \lambda$ is invertible,
$\lambda$ is a regular value of $\linA$,
and $\mathbb{H}^+$ and $\mathbb{H}^-$ belongs to $\rho(\linA)$,
respectively,

\item[2.]
or $\range{\linA - \lambda} \neq \hilH$,
then $\kernel{\linA^* - \overline{\lambda}} \neq \{\0\}$
and $\lambda$ belongs to the residual spectrum $\sigma_r(\linA)$,
so $\mathbb{H}^+$ and $\mathbb{H}^-$ are subsets of $\sigma_r(\linA)$,
respectively.
\end{itemize}

\section{Properties of the Resolvent and the Spectrum}
%------------------------------------------------------------------------------%
As it has already been indicated,
both cases may separately happen for the complex half planes
$\mathbb{H}^+$ and $\mathbb{H}^-$,
so in conclusion
there are four possibilities
how the spectrum $\sigma(\linA)$ of a densely defined and closed
symmetric operator $\linA$ looks like:

\begin{itemize}[leftmargin=20pt,itemsep=2pt]
\item[1.]
$\sigma(\linA) \eq \C$
in case of $\mathbb{H}^+ \cup \mathbb{H}^- \subset \sigma_r(\linA)$,
\item[2.]
$\sigma(\linA) \eq \overline{\mathbb{H}^+}$
in case of $\mathbb{H}^+ \subset \sigma_r(\linA)$ and
$\mathbb{H}^- \subset \rho(\linA)$,
\item[3.]
$\sigma(\linA) \eq \overline{\mathbb{H}^-}$
in case of $\mathbb{H}^- \subset \sigma_r(\linA)$ and
$\mathbb{H}^+ \subset \rho(\linA)$,
\item[4.]
$\sigma(\linA) \subset \R$
in case of $\mathbb{H}^+ \cup \mathbb{H}^- \subset \rho(\linA)$.
\end{itemize}

\subparagraph{}
But if the resolvent set $\rho(\linA)$ contains at least
one element $\zeta \xeof \R$,
then only the last case in the previous enumeration can occur.
In fact,
the property $\sigma(\linA) \subset \R$ leads us again
to the notion of \textit{self-adjointness},
due to the following theorem:

\begin{theorem1}\label{CRITERION_FOR_SELF-ADJOINTNESS}
A closed linear operator $\linA$ is self-adjoint \iff
\,$\sigma(\linA) \subseteq \R$,
and in this case we have for $\im \zeta \neq 0$ the resolvent estimation
\[
\norm{R(\linA,\zeta)} \leq \frac{1}{\abs{\im \zeta} }.
\]
\end{theorem1}

\begin{comment}
\Proof:
We first show that $\range{\linA - \zeta}$ lies dense in $\hilH$,
or equivalently,
that $\kernel{\linA^* - \overline{\zeta}} \eq \{\0\}$.
We compute for $y \xeof \kernel{\linA^* - \overline{\zeta}}$
\[
\im \overline{\zeta} \norm{y}^2 \eq
\im (\overline{\zeta} \scalar{y}{y}) \eq
\im \scalar{\overline{\zeta}y}{y} \eq
\im \scalar{\linA^* y}{y} \eq
\im \scalar{\linA y}{y} \eq 0,
\]
since $\scalar{\linA y}{y}$ is real.
To the end,
remember
that $\linA-\zeta$ has closed range for $\im \zeta \neq 0$,
which implies the assertion.
The resolvent estimation is a consequence of the inequality
$\norm{R(\linA,\zeta)} \leq d(\zeta,\sigma(\linA))^{-1}$,
where $d(\zeta,\sigma(\linA))$ denotes the positive distance between 
the closed sets $\{\zeta\}$ and $\sigma(\linA)$.
\qed
\end{comment}

\subparagraph{}
Indeed,
$\sigma(\linA)$ being a subset of the real line means that
$\mathbb{H}^+$ and $\mathbb{H}^-$ belong to the resolvent set of $\linA$,
which is true \iff the range condition is fulfilled,
that is, 
the operators $\linA \pm i$ are both surjective.
The resolvent estimation is a consequence of the inequality
$\norm{R(\linA,\zeta)} \leq \delta(\zeta,\sigma(\linA))^{-1}$,
where $\delta(\zeta,\sigma(\linA)$ denotes the positive distance between 
$\zeta \xeof \rho(\linA)$ and the spectrum of $\linA$.

\paragraph{} %- residual spectrum of self-adjoint operators
Remember the subsets of the spectrum $\sigma(\linA)$
of a linear operator $\linA$.
They have been introduced by the criterion
how the perturbed operator $\linA - \lambda$ fails to be bijective.
This approach has led us to the notion of
\textit{point spectrum} $\sigma_p(\linA)$,
\textit{continuous spectrum} $\sigma_c(\linA)$ and
\textit{residual spectrum} $\sigma_r(\linA)$,
respectively.

\subparagraph{}
The next result shows
that the division of the spectrum for self-adjoint operators
is simpler than in the arbitrary case:

\begin{theorem1}\label{RESIDUAL_SPECTRUM_OF_SELF-ADJOINT_OPERATOR_IS_EMPTY}
If $\linA$ is self-adjoint,
then the residual spectrum $\sigma_r(\linA)$ is empty,
thus $\sigma(\linA) \eq \sigma_p(\linA) \cup \sigma_c(\linA)$.
\end{theorem1}

\subparagraph{}
This assertion is proved by the relations
$\kernel{\linA} \eq \range{A^*}^{\bot} \eq \range{\linA}^{\bot}$.
By definition,
if $\lambda \xeof \sigma_r(\linA) \subset \R$,
then $\linA-\lambda$ is injective
and consequently $\range{\linA-\lambda} \eq \hilH$.

\paragraph{} %- pure point spectrum and pure continuous spectrum
When dealing with a self-adjoint linear operator $\linA$,
it is advantageous to breakup the set $\sigma(\linA)$
into other subsets as for arbitrary linear operators.

\subparagraph{}
First,
the \textit{pure point spectrum} $\sigma_{pp}(\linA)$ is the set
\[
\sigma_{pp}(\linA) \eqdef
\{ \lambda \in \C:
\range{\linA - \lambda} \= \overline{\range{\linA - \lambda} }
\, \neq \, \hilH \}.
\]

\subparagraph{}
Second,
by \textit{points embedded in continuum} are meant the elements of the set
\[
\sigma_{pec}(\linA) \eqdef
\{ \lambda \in \C:
\range{\linA - \lambda} \, \neq \, \overline{\range{\linA - \lambda} }
\, \neq \, \hilH
\}.
\]
Both definitions do not require $\linA-\lambda$ to be injective,
thus $\sigma_{pp}(\linA)$ and $\sigma_{pec}(\linA)$ 
possibly contain eigenvalues of $\linA$.
We join both subsets together and redefine the \textit{point spectrum}
$\sigma_p(\linA) := \sigma_{pp}(\linA) \sqcup \sigma_{pec}(\linA)$,
that is,
$\sigma_p(\linA)$ for $\linA$ being self-adjoint
consists of those $\lambda \xeof \C$,
for which $\overline{\range{\linA - \lambda}} \neq \hilH$.
This redefinition leads to the following

\begin{theorem1}\label{POINT_SPECTRUM_IS_SET_OF_EIGENVALUES}
The point spectrum $\sigma_p(\linA)$ is exactly
the set of eigenvalues of $\linA$.
\end{theorem1}

\subparagraph{}
Third,
the \textit{purely continuous spectrum} is defined to be the set
\[
\sigma_{pc}(\linA) \eqdef
\{ \lambda \in \C:
\range{\linA - \lambda} \neq \overline{\range{\linA - \lambda} }
\= \hilH \}.
\]
We alternatively introduce the \textit{continuous spectrum}
according to
$\sigma_c(\linA) := \sigma_{pec}(\linA) \, {\sqcup} \, \sigma_{pc}(\linA)$,
which of course is nothing but 
\[
\{ \lambda \in \C:
\range{\linA - \lambda} \, \neq \, \overline{\range{\linA - \lambda} } \}.
\]
In opposite to the originally given definition of $\sigma_c(\linA)$
(see paragraph VI.2),
the set $\sigma_c(\linA)$ may also contain those elements $\lambda \xeof \C$,
for which $\linA-\lambda$ is not injective.

\paragraph{}
For a self-adjoint operator $\linA$
it is also customary to denote $\sigma_{pp}(\linA)$ 
the \textit{discrete spectrum} 
and its complement relatively to
$\sigma(\linA)$ the \textit{essential spectrum}, {\ie}
\[
\sigma_{ess}(\linA) \eqdef
\sigma(\linA) \setminus \sigma_{pp}(\linA),
\]
so we have $\sigma_c(\linA) \subset \sigma_{ess}(\linA)$.

\subparagraph{}
The \textit{discrete spectrum} of $\linA$ consists
of those values $\lambda \xeof \C$,
for which there is a terminating sequence $(e_n)$ in $\hilH$
consisting of pairwise disjoint elements $e_n \xeof S_1(\0)$.
The essential spectrum of $\linA$,
however,
consists of numbers $\lambda \xeof \C$,
for which there is a non-terminating sequence $(e_n)$ in $\hilH$
consisting of pairwise disjoint elements $e_n \xeof S_1(\0)$
such that
\begin{equation}\label{IX_WEYL_SEQUENCE}
\lim_{n \xto \infty} \, \norm{\linA e_n - \lambda e_n} \= 0.
\end{equation}

\subparagraph{}
Each sequence having the property \eqref{IX_WEYL_SEQUENCE}
for some $\lambda \xeof \sigma(\linA)$
is called a \textit{Weyl sequence}.
In summary,
any number $\lambda \xeof \C$ belongs
to the spectrum of a self-adjoint linear operator $\linA$,
if it owns a (possibly terminating) Weyl sequence.

\section{The Discrete Spectral Theorem}
%------------------------------------------------------------------------------%
We next consider linear operators $\compA$
acting in a \textit{separable} Hilbert space $\hilH$,
which are \textit{compact} and \textit{symmetric} at the same time,
so $\compA$ fulfills $\scalar{\compA x}{y} \eq \scalar{x}{\compA y}$
for all $x,y \xeof \hilH$.
We shall call $\compA$ \textit{compact-symmetric} in this case.

\subparagraph{}
Of course,
if $\compA$ is compact-symmetric,
then the properties of compact and and symmetric operators complement,
so we may quote
that the spectrum $\sigma(\compA)$ consists
of a sequence $\{\lambda_n\}_{n \in \N}$
of \textit{real} eigenvalues of $\compA$,
which either terminates or converges to $0 \xeof \C$.
Due to the compactness of $\compA$,
the related eigenspaces $E_{\lambda_n}$ are finite-dimensional.
By the \textit{Gram-Schmidt orthogonalization method},
one can define an orthogonal system $\orthO$ of eigenvectors of $\compA$
such that $\orthO$ spans a dense subspace in $\hilH$
and consequently forms an \textit{orthonormal base} in $\hilH$.

\paragraph{} %- discrete spectral theorem
The \textit{Hilbert-Schmidt spectral theorem},
also often called the \textit{discrete spectral theorem},
claims that
a compact-symmetric operator $\compA$
acting in a separable Hilbert space $\hilH$
can be represented as a mostly \textit{countable} sum
of linear operators $\compA_{\lambda_n}$,
each of them being a simple multiplication operator
$\compA_{\lambda_n} := \lambda_n \1_H$
and acting on a finite-dimensional subspace of $\hilH$.
In other words,
the theorem provides an orthogonal decomposition of $\hilH$
into a sequence of $\compA$-invariant subspaces:

\begin{theorem1}\label{DISCRETE_SPECTRAL_THEOREM}
Let the Hilbert space $\hilH$ be separable.
Given a compact-symmetric linear operator $\compA\:\hilH \xto \hilH$,
there is a null sequence $\{\lambda_n\}_{n \in \N}$
consisting of real eigenvalues and an orthonormal base in $\hilH$
consisting of eigenvectors $\{x_n\}_{n \in \N}$ of $\compA$
such that for all $x \xeof \hilH$ the following representation holds:
\begin{equation}\label{IX_HILBERT_SCHMIDT_REPRESENTATION}
\compA \,x \=
\sum_{n=1}^{\infty} \, \lambda_n \, \scalar{x}{x_n} \, x_n.
\end{equation}
\end{theorem1}

\subparagraph{}
This suggests to order the set of eigenvalues of $\linA$ by their moduli.
Hence we obtain a sequence $\{\lambda_n\}_{n \in \N}$
of eigenvalues $\lambda_n$
such that
\[
\abs{\lambda_1} \, \leq \, \abs{\lambda_1} \, \leq \, \abs{\lambda_3} \, \dots
\, \xto \infty.
\]

\paragraph{} %- resolvent of a compact-symmetric operator
In the following 
let $\linA$ be a {compact-symmetric} operator
acting in a separable Hilbert space $\hilH$.
Then the resolvent operators $R(\linA,\zeta)$ of $\linA$ are symmetric-compact,
too,
and by the discrete spectral theorem \ref{DISCRETE_SPECTRAL_THEOREM} 
the set of eigenvectors of each $R(\linA,\zeta), (\zeta \xeof \rho(\linA)$
forms an orthonormal base in $\hilH$.
By the spectral mapping theorem,
the eigenvectors of $R(\linA,\zeta)$ are determined
by the eigenvectors of $\linA$.
In fact,
let $\xi_1,\xi_2,\ldots$ be the sequence of eigenvalues of $R(\linA,\zeta)$.
Then $(\lambda - \zeta)^{-1}x$ must be of the form
$(\lambda - \zeta)^{-1} x  \eq \xi_n \, x$ for some $n \xeof \N$,
which implies that $\sigma(\linA)$ is a subset of the set
$\{\zeta + {\xi_n}^{-1}: n \in \N \}$.

\section{Projections and Projection-Valued Measures}
%------------------------------------------------------------------------------%
As the projection theorem in Hilbert space shows,
projection mappings
are in one-to-one relationship to closed linear subspaces in $\hilH$.
So it is nearby
to call two projections $\pi_1,\pi_2$ \textit{orthogonal} to each other,
if their related closed subspaces have this property,
and in this case the identity
$\pi_1 \circ \pi_2 \eq \pi_2 \circ \pi_1 \eq \0$ holds,
where $\0$ is the projection related
with the trivial subspace $\{\0\} \subset \hilH$.

\paragraph{} %- properties of projections
In general,
a projection mapping $\pi$ has the following properties:

\begin{itemize}[leftmargin=12pt,itemsep=5pt]
\item[1.]
$\pi$ is idempotent: $\pi \circ \pi \eq \pi$.

\item[2.]
$\pi$ is linear on $\hilH$.

\item[3.]
$\pi$ is bounded with $\norm{\pi} \leq 1$,
hence $\pi \xeof \BO{\hilH}$.\\
If $\setA_{\pi} := \pi(\hilH) \eq \{\0\}$, then $\norm{\pi} \eq 0$,
else $\norm{\pi}=1$.
The only possible eigenvalues of $\pi$ are the real numbers $0$ and $1$.

\item[4.]
Let $\setA_1,\setA_2$ be closed subspaces of $\hilH$
with $\setA_1 \subseteq \setA_2$,
then for the associated projections the equation
$\pi_{\setA_1} \circ \pi_{\setA_2} \eq
 \pi_{\setA_2} \circ \pi_{\setA_1} \eq \pi_{\setA_1}$
holds.
In this case,
the operator $\pi_{\setA_2} - \pi_{\setA_1}$ is also a projection,
and related subspace is given by the set
$\setA_2 \ominus \setA_1 := \setA_2 \cap \setA_1^{\bot}$.

\item[5.]
$\pi \eq \pi^{*}$, that is, $\pi$ is symmetric.
Conversely,
each symmetric and idempotent linear operator
is a projection.

\item[6.]
If $\setA$ is a \textit{finite}-dimensional subspace of $\hilH$,
then $\setA$ is closed,
and the projection $\pi_{\setA}$ can be represented
by means of a set of base vectors $e_1,\ldots,e_n$ in $\setA$, namely
\[
\pi_{\setA} x \= \sum_{i=1}^n \scalar{x}{e_i} \, e_i.
\]
\end{itemize}

\subparagraph{}
Finally,
let $\mathscr{P}(\hilH)$ be the set of projections in $\hilH$,
then there is a partial order relation $\geq$ on $\mathscr{P}(\hilH)$,
defined by means of inclusion between the related closed subspaces:
\[
\pi_{\setA_1} \geq \pi_{\setA_2} \quad
: \Longleftrightarrow \quad \setA_2 \subseteq \setA_1.
\]

\paragraph{} %- projection-valued measures
A \textit{projection-valued measure}
is a mapping $\linP\:\Borel{\R} \xto \BO{\hilH}$,
where $\Borel{\R}$ is the ordinary Borel $\sigma$-algebra in $\R$
and $\linP$ fulfills the following properties:

\begin{itemize}[leftmargin=12mm]
\item[(P-1)]
$\linP(\Omega)$ is a projection operator for each $\Omega \xeof \Borel{\R}$;

\item[(P-2)]
$\linP(\R) \eq \1_{\hilH}$, the identity operator in $\hilH$;

\item[(P-3)]
$\linP$ is \textit{$\sigma$-additiv},
that is,
if $\Omega_1, \Omega_2,\ldots$ are measurable sets in the real line
and $\Omega := \bigsqcup_{\,n \in \N} \Omega_n$,
then
\[
\linP(\Omega) \= \lim_{n \xto \infty} \, \sum_{k=1}^n \, \linP(\Omega_k),
\]
where this kind of convergence has to be understood
in the sense of the strong operator topology in the algebra $\BO{\hilH}$.
\end{itemize}

\subparagraph{}
The projection-valued measure $\linP$ is also called
a \textit{spectral measure} or a \textit{resolution of the identity}.
Notice that
intervals of the form $(-\infty,\lambda]$ are Borel-measurable,
so each projection $\linP((-\infty,\lambda])$ is well-defined.
In fact,
the one-parameter family of these projections constitutes
a projection-valued measure $\linP$ in $\hilH$.

\section{Functional Calculus and Spectral Operators}
%------------------------------------------------------------------------------%
In the sequel
it is our aim to introduce a \textit{functional calculus}
for measurable functions $f\:\R \xto \C$
{\wrt} a projection-valued measure $\linP$.

\subparagraph{}
Let $B(\R)$ be the linear space of \textit{bounded} measurable functions
$f\:\R \xto \C$ endowed with the norm
\[
\norm{f}_{\infty} \eqdef \sup_{x \in \R} \, \abs{f(x) } < \infty,
\]
which in fact is a Banach space.
Then a function $e \xeof B(\R)$ is called \textit{simple},
if it adopts only finitely many values $e_1,\ldots, e_N$.
Moreover, the set
\[
S(\R) := \{\,f\:\R \xto \C: f \mbox{ is simple}\,\}
\]
forms dense subspace of $B(\R)$,
and the supremum norm of $e \xeof S(\R)$ is given by
$\norm{e}_{\infty} \eq \max \, \{\abs{e_1} ,\ldots, \abs{e_N} \}$.

\paragraph{}
After these preliminaries,
the setup of the spectral integral will be carried out
by the following five steps:

\begin{itemize}[leftmargin=12pt,itemsep=2mm]
\item[1.]
First of all,
each spectral measure $\linP$ gives rise
to a complex and real-valued Borel measure on the real line.
For arbitrary $x,y \xeof \hilH$
let the complex measure $\mu_{x,y}$ be defined by
\[
\begin{cases}[1.2]
\quad \mu_{x,y}\: \Borel{\R} \xto \C,\\
\quad \Omega \mapsto \mu_{x,y}(\Omega) := \scalar{x}{\linP(\Omega)y}.\\
\end{cases}
\]
Due to the symmetry of $\linP(\Omega)$,
the derived measure
$\mu_x(\Omega) := \mu_{x,x}(\Omega)$
for $\Omega \xeof \Borel{\R}$ is real-valued.

\item[2.]
We next make sense of the expression $\int_{\R} f \, dP$
for measurable functions $f\:\R \xto \C$.
Let us start with the case, where $f$ is \textit{simple},
that is,
$f$ takes values in the finite set $\{f_1,\ldots,f_N\}$.
Define
\[
\int_{\R} f \, d\linP \eqdef \sum_{n=1}^N f_n \linP(\chi^{-1}(f_n)),
\]
where $\chi^{-1}(f_n)$ is the preimage of $f_n$ in $\Borel{\R}$.
We further mention
that
\[
I_{0,\linP}\:f \xto \int_{\R} f \, d\linP
\]
is a continuous linear mapping from the normed space
$S(\R)$ consisting of simple functions
to the Banach space $\BO{\hilH}$ with operator norm $\norm{I_{0,\linP}}=1$.

\item[3.]
We shall continue with the integration
of \textit{bounded measurable functions} {\wrt} $\linP$.
Due to the B.L.T. theorem,
there exists a unique extension $I_P$ from $I_{0,\linP}$
to the closure $B(\R) \eq \overline{S(\R)}$,
where $\norm{I_P} \eq \norm{I_{0,P}}$.
The mapping $I_P$ is also linear and fulfills
\[
\begin{cases}[1.2]
\quad I_P(\1_{\R}) = \1_{\hilH}\\
\quad I_P(fg) = I_P(f) \circ I_P(g)\\
\quad I_P(\overline{f}) = I_P(f)^*\\
\end{cases}
\quad \quad f,g \in B(\R),
\]
where $\1_{\R}$ is the identity in $\R$,
thus $I_P$ even forms a C*-algebra homomorphism between
the function algebra $B(\R)$ and the operator algebra $\BO{\hilH}$.

\item[4.]
We will also define $I_P$
for an \textit{arbitrary measurable function} $f\:\R \xto \C$.
But this is only possible for those elements $x \xeof \hilH$ fulfilling
\begin{equation}\label{IX_PVM_OPERATOR}
\int_{\R} \abs{f} ^2 \, d \mu_x < \infty.
\end{equation}
The set $D_f$ of elements $x \xeof \hilH$ satisfying \eqref{IX_PVM_OPERATOR}
forms a linear subspace lying dense in $\hilH$.
For $f$ we introduce an approximating sequence $\{f_n\}_{n \in \N}$
of functions $f_n \in B(\R)$ given by
\[
f_n \eqdef \chi_{ \{t \in \R: \abs{f(t)} \leq n\}} \cdot f
\quad (n \in \N).
\]
One can show that $f_n$ is Cauchy in $L^2(\R)$,
hence $I_P(f_n)$ is Cauchy in $\hilH$,
where $I_P$ is the previously defined operator acting on the space $B(\R)$.

\item[5.]
To the end,
for a fixed measurable function $f$
we may define an operator $S(\linP,f)$,
called the \textit{spectral operator} {\wrt} $\linP$,
by virtue of
\begin{equation}\label{IX_FUNCTIONAL_CALCULUS}
\int_{\R} f \, d\linP\: D_f \mapsto \hilH\,,
\quad \quad
x \in D_f \mapsto (\int_{\R} f \, d\linP\,) \, x \eqdef y\,,
\end{equation}
which enjoys the property
\[
(\int_{\R} f \, d\linP)^* \= \int_{\R} \overline{f} \, d\linP.
\]
Especially for $f \eq \1_{\R}$ 
the self-adjoint operator $\linS_P := \linS(\linP,\1)$ is defined by
\[
\domain{\linS_P} := D_{\1_{\R}},
\quad \linS_P x \eqdef (\int_{\R} \1_{\R} \, d\linP) \, x.
\]
\end{itemize}

\section{The Spectral Theorem for Self-Adjoint Operators}
%------------------------------------------------------------------------------%
The spectral theorem is definitely the most important assertion
{\wrt} self-adjoint operators.
Another way to claim this theorem reads as follows:
each self-adjoint linear operator $\linA$ is a \textit{spectral operator},
whose underlying \textit{projection-valued measure} $\linP$
can be retrieved from $\linA$.
In the following we want to make this assertion precise
by going on step by step.

\subparagraph{} %- main axis theorem
In linear algebra
the \textit{main axis theorem} quotes
that each hermitean matrix is diagonalizable
and has purely real-valued eigenvalues.
The generalization of this assertion,
substituting hermitean matrices by densely defined self-adjoint operators
acting in abstract Hilbert space,
leads to the so-called \textit{spectral theorem}.

\paragraph{} %- multiplication operators
Remember from integration theory
that for a given measure space $(X,\salgB,\mu)$
the linear space $L^2(X)$ of measurable and square-integrable functions
$f\:X \xto \C$ forms a Hilbert space with inner product
\[
\scalar{f}{g} \eqdef \int_X f \, \overline{g} \, d\mu.
\]
If $m\:X \xto \C$ is a measurable function,
then the linear mapping $\linM_m$ defined by
\[
\begin{cases}[1.2]
\quad \domain{\linM_m} & := \,
\{\, f \in L^2(X): m \cdot f \in L^2(X)\,\}, \\
\quad (\linM_m f)(x) & := \, m(x) \cdot f(x) \;
\mbox{ for } x \in X \mbox{ and } f \in \domain{\linM_m} \\
\end{cases}
\]
is called a \textit{multiplication operator} with \textit{generator} $m$.
It turns out that the mapping $\linM_m$ is symmetric
iff
$m$ is real-valued,
and since the domains of definition
of the operators $\linM_m$ and $\linM_m^*$ coincide,
each multiplication operator with real-valued generator is self-adjoint.

\paragraph{}
The \textit{spectral theorem} claims that the reverse assertion is also true:

\begin{theorem1}\label{SPECTRAL_THEOREM_VERSION_1}
To each self-adjoint operator
$\linA\:\domain{\linA} \subset \hilH \xto \hilH$
there is assigned

\begin{itemize}[leftmargin=9mm]
\item[1.]
a measure space $(X,\salgB,\mu)$,

\item[2.]
a multiplication operator
$\linM_m$ in $L^2(X)$ with generator function $m\:X \xto \R$,

\item[3.]
a bijective and isometric linear operator $\linU:\hilH \xto L^2(X)$
\end{itemize}

such that $\linA \eq \linU^{-1} \circ \linM_m \circ \linU$ holds.
\end{theorem1}

\paragraph{}
The existence of the operator $\linU$ gives rise to speak of
\textit{unitary equivalence} between the operators $\linA$ and $\linM_m$.
If some operators $\linA_1$ and $\linA_2$ are unitary equivalent,
then their spectra $\sigma(\linA_1)$ and $\sigma(\linA_2)$ are identical.

\subparagraph{}
In summary, 
theorem \ref{SPECTRAL_THEOREM_VERSION_1} implies
that an operator $\linA$ is self-adjoint in Hilbert space $\hilH$
iff
it is unitary equivalent to a real-valued multiplication operator $\linM_m$
on $L^2(X)$,
where $X$ is the base set of a suitable chosen measure space $(X,\salgB,\mu)$.

\paragraph{} %- projection measure variant of the spectral theorem
Using the notion of projection-valued measure,
we are now able to quote a further variant of
the \textit{spectral theorem for self-adjoint operators}:

\begin{theorem1}\label{SPECTRAL_THEOREM_VERSION_2}
Given a self-adjoint linear operator $\linA$,
there is a uniquely defined projection-valued measure $P$
such that $\linA$ is the spectral operator related to $\linP$:
\[
\linA \= \int_{\,\R} \1 \, d\linP.
\]
\end{theorem1}

\paragraph{} %- functional calculus for self-adjoint operators
This version of the spectral theorem is extremely useful
to define operator-valued mappings $f\:\sigma(\linA) \xto f(\linA)$ by
\[
f(\linA) \eqdef \int_{\,\R} f \, d\linP,
\]
where $f(\linA)$ will become into a densely defined operator
acting in $\hilH$.
Indeed,
passing from a self-adjoint linear operator
to its family of projection-valued measures
and then to the spectral operators defined by \eqref{IX_FUNCTIONAL_CALCULUS}
is nothing but a \textit{functional calculus}
for the class of self-adjoint operators.
For example, 
if $\linA$ is self-adjoint and \textit{positive} at once,
that is,
$\scalar{\linA x}{x} \geq 0$ for all $x \xeof \domain{\linA}$,
then we may define operators like
$\abs{\linA} $,
$\sqrt{\linA}$, 
$\linA^n$,
$\exp(\linA)$,
$\sin(\linA)$,
$\cos(\linA)$ and others.

\paragraph{} %- spectral resolutions
It is also possible to formulate the spectral theorem by means of so-called
\textit{spectral resolutions} instead of projection-valued measures.
Preliminary,
to each projection-valued measure $\linP$ in Hilbert space $\hilH$
there is related the concept of \textit{spectral resolution},
which is nothing but a one-parameter family
$\{\linP_{\lambda}\}_{\lambda \xeof \R}$
of linear operators $\linP_{\lambda}\:\hilH \xto \hilH$
such that

\begin{itemize}[leftmargin=22mm]
\item[(R-1)]
$\linP_{\lambda}$ is a projection for each $\lambda \xeof \R$.
\item[(R-2)]
For each $x \xeof \hilH$
the mapping $\lambda \mapsto \linP_{\lambda} x$ is continuous from the right.
\item[(R-3)]
$\lambda < \mu$ implies $\linP_{\lambda} \leq \linP_{\mu}$,
that is,
$\scalar{\linP_{\lambda} x}{x} \leq \scalar{\linP_{\mu} x}{x}$
for all $x \xeof \hilH$.
\item[(R-4)]
$\lim_{\lambda \xto -\infty} \linP_{\lambda} \, x \eq 0$
and
$\lim_{\lambda \xto  \infty} \linP_{\lambda}  \, x= x$
for all $x \xeof \hilH$.
\end{itemize}

\paragraph{}
Each projection $\linP_{\lambda}$ of a given spectral resolution
is related to a closed linear subspace $\setA_{\lambda} \subset \hilH$
and the difference $\linP_{\mu} - \linP_{\lambda}$
for $\mu \geq \lambda$ also forms a projection related to the closed subspace
$\setA_{\mu} \ominus \setA_{\lambda}$.

\paragraph{}
Now,
using spectral resolutions instead of projection-valued measures,
the spectral theorem claims
that each self-adjoint operator $\linA$
owns a spectral resolution $\{\linP_{\lambda}\}_{\lambda \in \R}$
such that
\[
\domain{\linA} \=
\{\, x \in \hilH:
    \int_{\,\R}
\lambda^2 \, d \scalar{\linP_{\lambda} \,x}{x} < \infty\,\}\,,
\]
\[
\scalar{\linA \, x}{y} \=
\int_{\,\R} \lambda \, d\scalar{\linP_{\lambda} \,x}{y}
\quad \mbox{for } x \xeof \domain{\linA} \mbox{ and } y \xeof \hilH.
\]
The latter formula is often written down in the symbolic form
\[
\linA x \= \int_{\,\R} \lambda \, d\linP_{\lambda}x.
\]

\paragraph{} %- relationship between spectrum and spectral resolution
To the end of this paragraph
we want to investigate the relationship between the shape of the spectrum
and the spectral resolution of a self-adjoint operator.

\subparagraph{}
In fact,
the following theorem shows
that the resolvent set and the different kinds of spectra
of a self-adjoint linear operator $\linA$
are determined by its spectral resolution:

\begin{theorem1}\label{SPECTRUM_RESOLUTION_THEOREM}
Let $\linA$ be self-adjoint in $\hilH$
and $\setA_{\mu}$ be the closed subspace in $\hilH$
associated with the projection $\linP_{\mu}$ for some $\mu \xeof \R$.
Then

\begin{itemize}[leftmargin=9mm]
\item[(a)]
$\mu \xeof \R$ belongs to the \textit{point spectrum} $\sigma_p(\linA)$
iff
\[
0 < \dim \; (\setA_{\mu+0} \ominus \setA_{\mu}) < \infty.
\]
%Besides for some number $\epsilon > 0$
%the subspace
%$\setA_{\mu+\epsilon} \ominus \setA_{\mu-\epsilon}$ is finite-dimensional.
%Hence $\mu$ is not only an eigenvalue of $\linA$,
%but also belongs to the set $\sigma_d(\linA)$.

\item[(b)]
$\mu \xeof \R$ belongs to the \textit{continuous spectrum} $\sigma_c(\linA)$
iff
for each $\epsilon > 0$ the subspace
$\setA_{\mu+\epsilon} \ominus \setA_{\mu-\epsilon}$ is infinite-dimensional.

\item[(c)]
$\mu \xeof \R$ is a \textit{regular value} of $\linA$
iff
there is an open interval $(\alpha,\beta) \subset \R$
such that $\mu \xeof (\alpha,\beta)$
and $\linP_{\mu}$ is constant in $(\alpha,\beta)$.
\end{itemize}
\end{theorem1}

\paragraph{} %- representation of the resolvent
Especially the resolvents $R(\linA,\zeta)$ for $\zeta \xeof \rho(\linA)$ 
of a self-adjoint linear operator $\linA$
can be represented by means of the spectral resolution:
\[
R(\linA,\zeta) \, x \= \int_{-\infty}^{\infty}
\frac{1}{\lambda-\zeta} \, dP_{\lambda} x,
\quad (x \in \hilH, \, \im \zeta \neq 0).
\]

\subparagraph{} %- Stone's formula
On the other hand,
\textsc{Stone}\textit{'s formula} provides a constructive method
to describe the spectral resolution of $\linA$ in an interval $[a,b]$
by means of the resolvents $R(\linA,\zeta)$.
Let $x \xeof \hilH$ be arbitrary, then
\[
\frac{1}{2} \{P[a,b] + P(a,b)\} \, x
=
\lim_{\epsilon \downarrow 0} \, \frac{1}{2 \pi i}
\int_a^b
\{R(\linA,t + i \epsilon)\,x -
  R(\linA,t - i \epsilon)\,x \} \,
  dt.
\]

\section{Positive Self-Adjoint Operators and Cauchy Problems}
%------------------------------------------------------------------------------%
In the Hilbert space setting
the analysis of linear operator-differential equations
can be done wth the aid of spectral theory,
where the differential equation $u' + \linA u \eq 0$
is governed by a positive self-adjoint operator $\linA$.
The following existence and uniqueness theorem purely relies
on the spectral theorem and the functional calculus
for unbounded self-adjoint operators.

\begin{itemize}[leftmargin=12pt,itemsep=5pt]
\item[1.] %- first order Cauchy problems
Remember that a linear operator $\linA$ acting in a Hilbert space $\hilH$
is called \textit{positive},
if it is symmetric and accretive.
In this case any solution $x\:[0,\T) \xto \hilH$
of the homogeneous Cauchy problem
\begin{equation}\label{IX_HOMOGENEOUS_CAUCHY_PROBLEM_1}
x'' + \linA x \= \0, \quad x(0) \= x_0
\end{equation}
satisfies the estimate
\begin{equation}\label{IX_ESTIMATION_FOR_NONNEGATIVE_OPERATOR}
\norm{x(t)} \, \leq \, \norm{x(0)}
\qfa t > 0\,,
\end{equation}
which is a consequence of the inequality
$\norm{x(t)}^2 \eq - 2 \, \scalar{\linA x}{x} \leq 0$
and therefore causes the solution of the homogeneous Cauchy problem
to be unique.

\subparagraph{}
To verify the existence of a solution
for \eqref{IX_HOMOGENEOUS_CAUCHY_PROBLEM_1},
let $(P_{\lambda})_{\lambda \in \R}$ be the spectral resolution
of the self-adjoint operator $\linA$ given by the spectral theorem.
We notice that $P_{\lambda} \eq \0$ for $\lambda < 0$,
since $\linA$ is positive.
Then the functional calculus applied to the continuous function
$t \mapsto e^{-t}$ gives rise to the spectral operator
\[
\linL_t(x)
\eqdef
e^{-t \linA} x
\=
\int_{\,0}^{\,\infty} e^{-\lambda t} \, dP_{\lambda} \, x
\quad \quad (t \geq 0, \, x \in \hilH).
\]
We apparently have $\linL_0 \eq \1_X$,
and \eqref{IX_ESTIMATION_FOR_NONNEGATIVE_OPERATOR} implies
$\norm{\linL_t} \leq 1$ for $t > 0$.

\subparagraph{}
In conclusion,
the existence of the solution of \eqref{IX_HOMOGENEOUS_CAUCHY_PROBLEM_1}
is ensured by the following

\begin{theorem1}\label{EXISTENCE_THEOREM_FOR_ABSTRACT_HEAT_EQUATION}
The function $x\:t \mapsto \linL_t(x_0)$ is continuous in $[0,\infty)$,
continuously differentiable in $(0,\infty)$
and forms a solution of problem \eqref{IX_HOMOGENEOUS_CAUCHY_PROBLEM_1}.
\end{theorem1}

\item[2.] %- abstract Schrödinger equation
The functional calculus for self-adjoint linear operators
is also helpful to solve further classes of abstract initial-value problems.
By the \textit{abstract {\SCHROEDINGER} equation}
we mean the linear operator-differential equation
\begin{equation}\label{IX_ABSTRACT_SCHROEDINGER_EQUATION}
x'(t) \, + \, i \linA \, x(t) = \0\,,
\end{equation}
where $\linA$ is a self-adjoint operator in Hilbert space $\hilH$.
The operator $i \linA$ is called \textit{skew-adjoint} in this case.
Each solution $x$ of \eqref{IX_ABSTRACT_SCHROEDINGER_EQUATION}
satisfies the \textit{energy equation}
\begin{equation}\label{IX_ENERGY_CONSERVATION_LAW}
\norm{x(t)} = \norm{x(0)}
\end{equation}
for $t \xeof \R$,
which immediately follows from
\[
\frac{d}{dt} \, \norm{x(t)}^2
\=
\scalar{-i\linA x(t)}{x(t)} + \scalar{x(t)}{-i \linA x(t)} \= 0.
\]
If an initial value $x(0) \eq x_0$ is given,
then \eqref{IX_ENERGY_CONSERVATION_LAW} implies
that the solution of the initial-value problem 
for \eqref{IX_ABSTRACT_SCHROEDINGER_EQUATION} is unique.

\subparagraph{}
A solution of \eqref{IX_ABSTRACT_SCHROEDINGER_EQUATION} is also provided
by semigroup theory and especially by \textsc{Stone}'s theorem
(see paragraph X.5),
since the operator $i\linA$ is the infinitesimal generator of a unitary group
of bounded operators.
To construct a solution for problem \eqref{IX_ABSTRACT_SCHROEDINGER_EQUATION}
by the functional calculus method for unbounded self-adjoint operators,
we define the spectral operator
\[
L_t(x) 
\eqdef
e^{-it \linA} x
\=
\int_{\,-\infty}^{\,\infty} e^{-i \lambda t} dP_{\lambda} \, x
\quad \quad (t > 0),
\]
where $P_{\lambda}$ again is the spectral resolution of $\linA$
having support in the whole real line.
Then the family of operators $L_t$ yields the solution
of the abstract {\SCHROEDINGER} problem:

\begin{theorem1}\label{EXISTENCE_THEOREM_FOR_ABSTRACT_SCHROEDINGER_EQUATION}
The mapping $t \mapsto \linL_t(x_0)$ is the unique solution
of the Cauchy problem for \eqref{IX_ABSTRACT_SCHROEDINGER_EQUATION},
where for each $t > 0$ the operator $\linL_t$ is unitary.
\end{theorem1}

\item[3.] %- second order Cauchy problems
The functional calculus for self-adjoint linear operators may also be applied
to the second-order homogeneous Cauchy problem
\begin{equation}\label{IX_HOMOGENEOUS_CAUCHY_PROBLEM_2}
x'' + \linA x \= \0, \quad x(0) \= x_0, \quad x'(0) \= x_1,
\end{equation}
since the square root operator $\sqrt{\linA}$ is well-defined,
due to the fact that $\linA$ is positive.

\subparagraph{}
Each solution $x$ of problem \eqref{IX_HOMOGENEOUS_CAUCHY_PROBLEM_2}
satisfies the energy conservation law
\[
\norm{x'(t)}^2 \, + \, \snorm{\sqrt{\linA} \, x(t)}^2
\=
\norm{x'(0)}^2 \, + \, \snorm{\sqrt{\linA} \, x(0)}^2
\quad (0 < t < \T),
\]
which implies that the problem owns mostly one solution.

\subparagraph{}
Let $(P_{\lambda})$ again be the spectral resolution of $\linA$.
We define a function $x\:[0,\T) \xto \hilH$
by means of the spectral operator $\linL_t$ induced by the functions
$t \mapsto \cos(\sqrt{t})$ and
$t \mapsto \sin(\sqrt{t})/\sqrt{t}$.
In fact,
the operator $x \mapsto \linL_t x$
is built as the sum of two single spectral operators for each $0 \leq t < \T$,
namely
\[
x(t)
:=
\linL_t(x_0,x_1)
:=
\int_{\,0}^{\,\infty} \cos{\sqrt{\lambda}t} \, dP_{\lambda} \, x_0
\, + \,
\int_{\,0}^{\,\infty} \frac{\sin{\sqrt{\lambda}t}}{\sqrt{\lambda}}
\, dP_{\lambda} \, x_1.
\]

\subparagraph{}
Having defined the family of linear operators $\{\linL_t(x_0,x_1)\}$ this way,
we obtain the following existence and uniqueness theorem
for our second order problem:

\begin{theorem1}\label{EXISTENCE_THEOREM_FOR_ABSTRACT_WAVE_EQUATION}
The mapping $t \mapsto \linL_t(x_0,x_1)$ is continuous in $[0,\T)$,
twice continuously differentiable in $(0,\T)$ and
is the unique solution of problem \eqref{IX_HOMOGENEOUS_CAUCHY_PROBLEM_2}.
\end{theorem1}
\end{itemize}

% BIBLIOGRAPHY
%------------------------------------------------------------------------------%
\vspace{5mm}
\begin{thebibliography}{99}
\begin{small}

\bibitem{IX_RN} F.\,Riesz, B.\,Sz.-Nagy:
\textit{{Functional Analysis}}.
Springer, 1956.

\bibitem{IX_Yo} K.\,Yosida:
\textit{Functional Analysis}.
Springer, 1980.

\end{small}
\end{thebibliography}

%------------------------------------------------------------------------------%
%   X. SEMIGROUP THEORY AND EVOLUTION EQUATIONS ---
%------------------------------------------------------------------------------%
%  §1. Strongly Continuous Semigroups ---
%  §2. Infinitesimal Generators of Semigroups ---
%  §3. Accretive Operators and Contractive Semigroups ---
%  §4. Sectorial Operators and Analytic Semigroups ---
%  §5. Skew-Adjoint Operators and Unitary Groups ---
%  §6. Well-Posedness of Second-Order Cauchy Problems ---
%  ---

\newpage
\thispagestyle{empty}
\def\TitleX{Semigroups of Bounded Linear Operators}
\fancyhead[C] {\small{\textsc{X. \TitleX}}}
\part{\TitleX}
%------------------------------------------------------------------------------%
Our aim in this part
is to sketch the basic notions of semigroup theory,
which forms one of the most important topics within functional analysis
{\wrt} applications to ordinary and partial differential equations.

\subparagraph{}
The notion of \textit{semigroup} has been introduced
by \textsc{J.\,A.\,S\'{e}guier} in 1904.
Later some applications to partial differential equations were detected,
but the first systematic and comprehensive treatement of semigroup theory
was published in the pioneering monograph \cite{X_HP}
by \textsc{E.\,Hille} and \textsc{R.\,Phillips}.
Nowadays semigroup theory forms a powerful method
in solving linear evolution problems
with a multiplicity of applications to partial differential equations
governed by linear-elliptic operators.

\section{Strongly Continuous Semigroups}
%------------------------------------------------------------------------------%
Given a Banach space $\banX$,
a \textit{C$_0$-semigroup} or
a \textit{strongly continuous semigroup of bounded linear operators}
is an operator-valued mapping
\[
\linS:[0,\infty) \xto \BO{\banX},
\]
shortly denoted $\{\linS(t)\}_{t \geq 0}$ and distinguishes by the properties

\begin{itemize}[leftmargin=12mm]
\item[(G-1)]
$\linS(t+s) \eq \linS(t) \circ \linS(s)$ for all $t,s \geq 0$
(the \textit{semigroup property});

\item[(G-2)]
$\lim_{\,t \downarrow 0} \, \linS(t)\,x \eq x$ for all $x \xeof X$.
\end{itemize}

\subparagraph{}
It follows from (G-2),
which is nothing but strong continuity from the right-hand side in $t \eq 0$,
that $\linS(0)$ is the identity mapping in $\banX$.
The conditions (G-1) and (G-2) together
imply continuity from the right of the semigroup even in every $t > 0$.
In case of $\banX \eq \R$,
semigroups share both properties
with the family $\{\cal{E}(t)\}_{\, t \geq 0}$ of exponential functions
$\cal{E}(t): \lambda \in \R \mapsto e^{\lambda t}$,
so $\{\cal{E}(t)\}_{\, t \geq 0}$ serves as an elementary example
of a strongly continuous semigroup.

\subparagraph{} %- classification of strongly continuous semigroups
Notice that the totality of C$_0$-semigroups may be classified as follows:
for each such semigroup $\{\linS(t)\}_{t \geq 0}$
there are constants $M \geq 0$ and $\omega \xeof \R$
such that
\begin{equation}\label{X_SEMIGROUP_GROWTH}
\norm{\linS(t)} \, \leq \, M e^{\omega t}
\quad \for t > 0.
\end{equation}

\begin{comment}
In fact,
let $M \geq 1$ such that $\norm{S(t)} \leq M$ for all $s \xeof [0,1]$.
For $t > 0$ there is $s \xeof [0,1)$ and $n \xeof N$ such that $t \eq n+s$.
Then $\norm{\linS(n)}\leq M \norm{S(n-1)}$ for $n > 1$ and
\[
\norm{\linS(t)} =
\norm{\linS(n+s)} = \norm{\linS(n) \circ \linS(s)} \leq
\norm{\linS(n)} \cdot \norm{\linS(s)} \leq
M^{n+1} = M e^{n \cdot \ln{M}}.
\]
\end{comment}

\paragraph{} %- orbit maps
By the \textit{orbit map} or \textit{trajectory} $\eta_x$
of a C$_0$-semigroup $\{\linS(t)\}_{t \geq 0}$
and associated with an initial element $x \xeof \banX$
we mean the vector-valued mapping $\eta_x:[\,0,\infty) \xto \banX$
given by $\eta_x(t) := S(t) \, x$.
Due to condition (G-1),
strong continuity of $\{\linS(t)\}_{t \geq 0}$
implies continuity of the function $\eta_x$.
But $\eta_x$ is even differentiable under a very mild assumption:
if $\{\linS(t)\}_{t \geq 0}$ is a C$_0$-semigroup,
then for each $x \xeof \banX$
the mapping $\eta_x$ is differentiable in $(0,\infty)$
\iff
$\eta_x$ is right differentiable in $t \eq 0$.

\section{Infinitesimal Generators of Semigroups}
%------------------------------------------------------------------------------%
If $\{\linS(t)\}_{t \geq 0}$ is a strongly continuous semigroup,
its \textit{infinitesimal generator}
(briefly denoted the \textit{generator} of the semigroup)
is defined as a mapping $\linA$ in $\banX$ by means of the graph
\[
\Gamma_{\linA}
\eqdef
\{\,(x,y) \in X \times X: \;
y \= \lim_{\,h \downarrow 0} \, \frac{1}{h}(\linS(h)x-x) \exists \}.
\]
Hence $\domain{\linA}$ of $\linA$ consists of those elements $x \xeof \banX$,
for which the orbit map $\eta_x$ is differentiable from the right in $t \eq 0$,
that is,
\[
\domain{\linA} \eq \{x \xeof \banX: \eta_x \mbox{ is differentiable}\}.
\]
It is easy to show that
$\linA$ is linear,
so it constitutes a linear operator acting in $\banX$
with domain of definition $\domain{\linA} \subset \banX$.

\paragraph{} %- the main theorem on infinitesimal generators}
%------------------------------------------------------------------------------%
The meaning of strongly continuous semigroups is confirmed
by the solution theory for homogeneous Cauchy problems
governed by linear operators.
Indeed,
the \textit{main theorem of semigroup theory} reads as follows:

\begin{theorem1}\label{MAIN_THEOREM_OF_SEMIGROUP_THEORY}
Let $\linA$ be a densely defined and closed linear operator
acting in a Banach space $\banX$.
Then the following assertions are equivalent to each other:

\begin{itemize}[leftmargin=15pt]
\item[(a)]
The operator $-\linA$ generates a strongly continuous semigroup
of type $(\omega,M)$.
\item[(b)]
The resolvent set of $\linA$ fulfills $(-\infty,\omega) \subset \rho(\linA)$
and the resolvent function $R(\linA,\zeta)$ satisfies
\begin{equation}\label{X_HILLE-YOSIDA_CONDITION}
\norm{R(\linA,\zeta)^n}
\, \leq \, M (\omega - \zeta)^{-n} \quad (n \in \N, \zeta < \omega).
\end{equation}
\end{itemize}
\end{theorem1}

\subparagraph{}
The equivalence of these assertions is the content
of the celebrated \textsc{Feller-Miyadera-Hille-Yosida} \textit{theorem}.
However,
since the number of conditions on the operator $\linA$ is infinite,
this theorem has limited worth {\wrt} applications.
But the original theorem for generators of \textit{contraction semigroups}
due to \textsc{Hille} and \textsc{Yosida} in 1948 quotes
that in this case only the first inequality
($n=1$, with $\omega \eq 0$ and $M \eq 1$)
of \eqref{X_HILLE-YOSIDA_CONDITION} has to be verified:

\begin{theorem1}\label{HILLE-YOSIDA_THEOREM}
The linear operator $-\linA$ acting in a Banach space $\banX$
is the infinitesimal generator of a C$_0$-semigroup of \textit{contractions}
\iff

\begin{itemize}[leftmargin=12pt]
\item[(a)]
$\linA$ is densely defined and closed,
\item[(b)]
the interval $(-\infty,0)$ belongs to the resolvent set $\rho(\linA)$,
\item[(c)]
$\norm{\zeta \, R(\linA,\zeta) } \leq 1$ for all $\zeta < 0$.
\end{itemize}
\end{theorem1}

\subparagraph{}
If $-\linA$ owns the properties mentioned in theorem \ref{HILLE-YOSIDA_THEOREM},
then the semigroup $\{\linS(t)\}_{t \geq 0}$ generated by $-\linA$
may be constructed by
\[
\linS(t) \, x \eqdef
\lim_{\,n \xto \infty} (\1_X + \frac{t}{n} \linA)^{-n} \,x 
\fe t > 0 \and x \in \banX.
\]

\section{Accretive Operators and Contractive Semigroups}
%------------------------------------------------------------------------------%
In the following let $\hilH$ be a complex Hilbert space.
Then a linear operator $\linA$ acting in $\hilH$ and with domain of definition
$\domain{\linA} \subset \hilH$ is called \textit{accretive},
if
\begin{equation}\label{X_ACCRETIVE_PROPERTY_1}
\re \, \scalar{\linA \,x}{x} \, \geq 0
\qfa x \in \domain{\linA},
\end{equation}
or equivalently, 
if
$\Theta(\linA) \subset \{\zeta \xeof \C: \re \zeta \geq 0\}$.
Otherwise,
a linear operator $\linA$ is said to be \textit{dissipative}, 
if $-\linA$ is accretive.

\subparagraph{}
We first notice that if $\linA$ is accretive (dissipative),
then $\zeta \linA$ is so for each $\zeta > 0$.
Furthermore,
since 
\[
\norm{ \linA \,x + \zeta x}^2 - \norm{\zeta x}^2
\, \geq \,
2 \zeta \; \re \, \scalar{\linA \,x}{x}
\qfa x \in \domain{\linA},
\]
the operator $\linA$ is accretive \iff
\begin{equation}\label{X_ACCRETIVE_PROPERTY_2}
\norm{(\linA + \zeta) x} \, \geq \, \norm{\zeta x}
\qfa \zeta \geq 0 \and x \in \domain{\linA}
\end{equation}
or eqivalently
\begin{equation}\label{X_ACCRETIVE_PROPERTY_3}
\norm{(\linA + \zeta) x} \, \geq \, \re \zeta \cdot \norm{x}
\qfa \re \zeta \geq 0 \and x \in \domain{\linA}.
\end{equation}

These two properties can be used to define accretivity and dissipativity also
for linear operators in Banach spaces,
where no inner product is present.

\subparagraph{}
Inequality \eqref{X_ACCRETIVE_PROPERTY_1} quotes that
for $\zeta > 0$
the operator $\linA + \zeta$ is bounded below and thus injective.
In addition,
if $\linA$ is closed,
then the bounded below theorem implies that
$\linA$ has closed range,
and by the \textit{closed graph theorem} the inverse operator
\[
(\linA + \zeta)^{-1}: \range{\linA + \zeta} \xto \domain{\linA}
\]
is bounded with norm
$\norm{(\linA + \zeta)^{-1}} \leq \zeta^{-1}$,
so it is even a contraction for $\zeta \geq 1$.

\paragraph{} %- m-accretive operators
Let us now assume that
for the accretive operator $\linA$
there is $\zeta > 0$
such that $\linA + \zeta$ is surjective.
A linear operator with this property is said
to fulfill the \textit{range condition}.
In this case there is an open interval $\J$ containing $\zeta$
such that $\linA + \zeta$ is surjective for all $\zeta \xeof \J$.
Finally,
it turns out that
$\linA + \zeta$ is surjective for every $\zeta > 0$.

\subparagraph{}
An accretive linear operator $\linA$ is called \textit{maximal accretive}
or shortly \textit{m-accretive}
if the range condition is fulfilled.
The following enumeration contains a couple of properties
of the m-accretive operator $\linA$:

\begin{itemize}[leftmargin=12mm]
\item[(a)] $\linA$ is densely defined in $\hilH$;
\item[(b)] $\linA$ is closed;
\item[(c)] $\C^- := \{\zeta \xeof \C: \re \zeta < 0\} \subset \rho(\linA)$;
\item[(d)] $\linA + \zeta$ is either surjective for all $\zeta \xeof \C^+$
          or \textit{not} surjective for all $\zeta \xeof \C^+$.
\end{itemize}

\paragraph{} %- characterization of m-accretive operators
The main theorem of this paragraph
characterizes m-accretive linear operators by other conditions:

\begin{theorem1}
If $\linA$ is accretive,
the following assertions are equivalent to each other:

\begin{itemize}[leftmargin=17pt]
\item[(a)]
$\linA$ is m-accretive;

\item[(b)]
$\linA$ is accretive and does not have a proper accretive extension;

\item[(c)]
$\linA + i\alpha$ is m-accretive for all $\alpha \xeof \R$;

\item[(d)]
$\linA + \alpha$ is m-accretive for all $\alpha > 0$;

\item[(e)]
$\C^- \subset \rho(\linA)$ and $\norm{R(\linA,\zeta)} \leq (\re \zeta)^{-1}$
for all $\zeta \xeof \C^-$;

\item[(f)]
$(-\infty,0) \subset \rho(\linA)$ and
$\sup_{\,t<0} \norm{t\,R(\linA,t)} \leq 1$;

\item[(g)]
$\linA$ is densely defined and closed,
and the adjoint $\linA^*$ of $\linA$ is m-accretive,
too.
\end{itemize}
\end{theorem1}

\subparagraph{} %- kernels of accretive operators
Besides,
if $\linA$ is densely defined in $\hilH$ and m-accretive,
then $\kernel{\linA} \eq \kernel{\linA^*}$ and consequently
$\kernel{\linA} \subset \domain{\linA} \cap \domain{\linA^*}$.

\paragraph{} %- coercive operators
We shall denote a linear operator $\linA$ \textit{coercive},
if $\re \scalar{\linA x}{x} \geq \theta \norm{x}^2$
for some $\theta > 0$ and all $x \xeof \domain{\linA}$.
It should be clear that
every coercive operator is accretive.

In the following we want to characterize the spectral properties (b) and (c)
in theorem \ref{HILLE-YOSIDA_THEOREM}
for a linear operator $\linA$ by equivalent conditions.

\subparagraph{}
By means of the concept of m-accretivity,
the \textit{Lumer-Phillips} \textit{theorem} classifies those linear operators
being infinitesimal generators of a contractive semigroup
without the usage of spectral theoretical concepts:

\begin{theorem1}\label{LUMER-PHILLIPS_THEOREM}
Let $\linA$ be a linear operator in Hilbert space.
Then the operator $-\linA$ generates a contractive C$_0$-semigroup
\iff
$\linA$ is m-accretive.
\end{theorem1}

\begin{comment}
\Proof:
\qed
\end{comment}

\section{Sectorial Operators and Analytic Semigroups}
%------------------------------------------------------------------------------%
Regarding the definition of sectorial operator,
we shall follow close to \cite{X_Ha} and denote
a linear operator $\linA$ in a Banach space $\banX$
and with domain of definition $\domain{\linA}$
(not necessarily dense in $\banX$)
to be \textit{sectorial} of angle $\psi \xeof [0,\pi)$
and vertex $\omega \xeof \R$,
if the following conditions are true:

\begin{itemize}[leftmargin=12mm]
\item[(S-1)']
$\sigma(\linA) \subset \overline{\Sigma_{\,\psi,\omega} } $,
where
$\Sigma_{\,\psi,\omega} :=
\{\zeta \xeof \C \: \zeta \neq \omega, arg(\zeta-\omega) < \psi\}$
is the sector with vertex $\omega$ and semi-angle $\psi$.
This implies
$\C \setminus \overline{ \Sigma_{\,\psi,\omega}} \subset \rho(\linA)$;

\item[(S-2)']
$\sup
\,
\{\, \norm{\zeta R(\linA,\zeta)}:
     \zeta \in \C \setminus \overline{ \Sigma_{\,\phi,\omega}} \,\}
\, < \, \infty$
\, for all $\psi < \phi < \pi$.
\end{itemize}

\subparagraph{}
The operator $\linA$ is called \textit{sectorial},
if it is so for some pair $(\psi,\omega) \xeof [0,\pi) \xtimes \R$.
Notice that
if $\linA$ is sectorial {\wrt} the angle $\psi$,
then there is $0 < \epsilon < \psi$
such that $\linA$ is also sectorial {\wrt} the smaller angle $\psi - \epsilon$.
The infimum $\psi_0$ of all these angles is called
the \textit{angle of sectoriality} of $\linA$.

%\subparagraph*{}
%Furthermore,
%if a sectorial operator $\linA$ is acting in a Hilbert space $\hilH$,
%it is called \textit{m-sectorial}, 
%if it fulfills a \textit{range condition},
%that is,
%there is $\lambda > 0$ such that $\linA + \lambda$ is surjective.
%By this definition,
%the vertex point $\omega \xeof \C$ of $\linA$ 
%always belongs to the resolvent set $\rho(\linA)$
%and $\linA$ is closed due to $\rho(\linA) \neq \emptyset$.
%Since resolvent sets are open,
%the distance between the spectrum $\sigma(\linA)$
%and the vertex point $\omega$ of $\linA$ is strictly positive.

\paragraph{} %- criteria for sectorialness
The following theorem,
whose proof can be found in \cite{X_Lu}, p.\,43,
provides useful criterions to ensure sectorialness of $\linA$:

\begin{theorem1}\label{CRITERION_FOR_SECTORIALNESS}
Let  $\linA$ be a linear operator in $\banX$ with domain of definition 
$\domain{\linA} \subset \banX$ having the following properties:

\begin{itemize}[leftmargin=9mm]
\item[(a)]
the resolvent set $\rho(\linA)$ contains
$\C_{\omega} := \{\zeta \xeof \C: \re \zeta \leq \omega\}$;

\item[(b)]
there is $M > 0$ such that
$\norm{\zeta \, R(\linA,\zeta)} \, \leq \, M$ for all $\zeta \xeof \C_{\omega}$.
\end{itemize}
Then $\linA$ is sectorial.
\end{theorem1}

\paragraph{} %- criterion for sectorialness
A further theorem shows that
in the Hilbert space setting
the notation of sectorial operator is close to \textsc{Kato}'s concept
of \textit{m-sectorial operator}:

\begin{theorem1}
Let $\hilH$ be a Hilbert space and $\linA$ be a closed linear operator
in $\hilH$ with numerical range
$\Theta(\linA) \subset \overline{\Sigma_{\,\psi,0}}$ for $\psi \leq \pi/2$
and $\linA + \1_H$ being surjective.
Then $\linA$ is sectorial of angle $\phi$.
\end{theorem1}

\subparagraph{}
This theorem especially implies that a self-adjoint and positive
linear operator in Hilbert space is sectorial of angle $0$.

\paragraph{} %- properties of sectorial operators in Banach spaces
In the following we collect further features
of a sectorial linear operator $\linA$ acting in ordinary Banach spaces:

\begin{itemize}
\item[1.]
$x \xeof \overline{\domain{\linA}}$ \iff
$\lim_{\,t \xto \infty} t(t+\linA)^{-1}x \eq x$;
\item[2.]
$x \xeof \overline{\range{\linA}}$ \iff
$\lim_{\,t \xto 0} t(t+\linA)^{-1}x \eq 0$;
\item[3.]
If $\banX$ is reflexive, then $\linA$ is densely defined and
$\banX \eq \kernel{\linA} \oplus \overline{\range{\linA}}$.
\end{itemize}

\paragraph{} %- analytic semigroups
Semigroups of bounded linear operators are usually parametrized
on the half line $[0,\infty)$,
but it is also customary to study possible extensions to a symmetric sector
\[
\Sigma_{\,\theta} \eqdef
\{z=re^{i\alpha} \in \C: r > 0\,, \abs{\alpha} < \theta\}
\]
of angle $2 \theta$ for some $\theta \xeof (0,\frac{\pi}{2}]$.

\subparagraph{} %- definition of analytic semigroup
A semigroup $\{\linS(t)\}_{t \geq 0}$ is called \textit{analytic},
if there is an angle $\theta \xeof (0,\frac{\pi}{2})$
and an analytic extension $\linS':\Sigma_{\,\theta} \xto \BO{X}$ of $\linS$,
which is \textit{locally bounded},
that is,
\[
\sup \,
\{ \, \norm{\linS'(z)}: z \in \Sigma_{\,\theta}, \, \norm{z} \leq 1 \}
\, < \, \infty.
\]

\paragraph{} %- functional calculus for sectorial operators
Next we want to sketch a functional calculus for sectorial operators $\linA$,
which especially enables us to define bounded linear operators
of the form $e^{-t \linA}$.

\subparagraph{}
If $\banX$ is a Banach space
and $\linA$ is a sectorial operator acting in $\banX$,
then for each $t > 0$
it is possible to introduce the bounded linear operator
$e^{-t \linA} \xeof \BO{\banX}$ by means of the contour integral
\begin{equation}\label{X_CONTOUR_INTEGRAL}
e^{-t \linA} x \eqdef
\frac{1}{2 \pi i}
\int_{\gamma} \, e^{-t \zeta} \, R(\linA,\zeta)x \, d\zeta
\quad (x \in \banX),
\end{equation}
which generalizes the well-known \textsc{Riesz-Dunford} \textit{formula}.
Here $\gamma$ constitutes a closed curve in the complex plane
that drifts within the sector $S_{\phi,\omega}$
and includes the spectrum of the operator $-t \linA$.
The family of operators $\{\linS(t)\}_{t \geq 0}$
defined by \eqref{X_CONTOUR_INTEGRAL} forms a semigroup.
In addition,
if $\linA$ is densely defined in $\banX$,
then the semigroup is even strongly continuous.

\subparagraph{}
Due to the property of $\linA$ being sectorial,
the domain of the parameter $t$ is not restricted
to the interval $(0,\infty)$,
but can be extended to a certain sector $\Sigma_{\,\theta} \subset \C$,
in which the semigroup property
$\linS(t+s) \eq \linS(t)\,\linS(s)$ for $s,t \xeof \Sigma_{\,\theta}$ holds.
Thus we obtain the following

\begin{theorem1}\label{X_SEMIGROUP_CREATED_BY_SECTORIAL_OPERATOR_IS_ANALYTIC}
If $\linA$ is a sectorial operator,
then the semigroup $\{\linS(t)\}_{t \geq 0}$ of bounded linear operators
generated by $- \linA$ is {analytic}.
\end{theorem1}

\paragraph{} %- sectorially contractive analytic semigroups
We finally turn to semigroups being analytic and contractive at the same time.
Let $\{\linS(t)\}_{t \geq 0}$ be an analytic semigroup.
If in addition $\norm{\linS(\zeta)} \leq 1$
for all $\zeta \xeof \Sigma_{\,\theta}$,
then it is called a \textit{sectorially contractive analytic}
C$_0$-semigroup of angle $\theta$.

\paragraph{}
The generation theorem for this subclass of analytic semigroups
reads as follows:

\begin{theorem1}\label{CRITERION_FOR_SEMIGROUP_TO_BE_SECTORIALLY_CONTRACTIVE}
Let $\linA$ be a densely defined operator on a Banach space $\banX$
and let $\theta \xeof (0,\frac{\pi}{2}]$.
The following assertions are equivalent:

\begin{itemize}[leftmargin=15pt]
\item[(a)]
The operator $-\linA$ generates a sectorially contractive
analytic $C_0$-semigroup of angle $\theta$.

\item[(b)]
$\rho(\linA) \subset  \C \setminus \Sigma_{\,\theta}$ and
$\norm{\zeta(A + \zeta)^{-1}} \leq 1$ for $\zeta \xeof \Sigma_{\,\theta}$.
\end{itemize}
\end{theorem1}

\section{Skew-Adjoint Operators and Unitary Groups}
%------------------------------------------------------------------------------%
Recall that a bijective linear operator $\linU$ in Hilbert space
is called \textit{unitary},
if $\domain{\linU} \eq \hilH$ and $\linU^{-1} \eq \linU^*$.
The spectrum $\sigma(\linU)$ of such an operator $\linU$ forms a subset
of the unit circle in the complex plane.

\subparagraph{} %- unitary groups
A semigroup $\{\linU(t)\}_{t \geq 0}$ of bounded linear operators
acting in a Hilbert space $\hilH$ is called \textit{unitary},
if each operator $U(t)$ is unitary.
Furthermore,
let $\linU_{-t} := \linU_t$ for $t > 0$,
then a \textit{group} $(\linU_t)_{t \in \R}$ of unitary operators is defined,
which fulfills the group property
\[
\linU_{s+t} \= \linU_s \circ \linU_t \quad \quad (t,s \in \R).
\]

\subparagraph{} %- Stone's theorem
Reversely,
\textsc{Stone}'s theorem claims
that the infinitesimal generator of a unitary group is skew-adjoint:

\begin{theorem1}\label{STONE_THEOREM_FOR_UNITARY_GROUPS}
Each unitary group $(\linU_t)_{t \in \R}$ in a Hilbert space
is infinitesimally generated by a linear operator of type $i \linA$,
where $\linA$ is self-adjoint.
\end{theorem1}

It should be mentioned that
if a self-adjoint linear operator $\linA$ is given,
then the unitary group $(\linU_t)_{t \in \R}$ can be constructed
by means of the spectral decomposition of $\linA$ and the functional calculus
for spectral operators with respect to the function $t \mapsto e^{it}$,
that is,
we have for each $t > 0$
\[
\linU_t x \eqdef \int_{\,\R} e^{i \lambda t} \, dP_{\lambda} \,x
\quad (x \in \hilH).
\]


\section{Well-Posedness of Second-Order Cauchy Problems}
%------------------------------------------------------------------------------%
Similar to the first-order Cauchy problem,
a function $x$ is called
a \textit{classical solution} of the second-order homogeneous Cauchy problem
\begin{equation}\label{X_HOMOGENEOUS_CAUCHY_PROBLEM_2}
x'' + \linA x \= \0, \quad x(0) \= x_0
\end{equation}
if $x \xeof C([0,\T),\domain{\linA}) \, \cap \, C^2((0,\T),\banX)$
and satisfies both initial conditions as well as
the differential equation in each point $t \xeof (0,\T)$.

\paragraph{} %- the semigroup approach to second order Cauchy problems
To apply results from semigroup theory to the second order Cauchy problem,
let $U := (x,x')$.
We suppose that $\linA$ is a \textit{positive} and
\textit{self-adjoint} linear operator in a Hilbert space $\hilH$.
Then the differential equation $x''(t) + \linA \, x(t) \eq \0$
in \eqref{X_HOMOGENEOUS_CAUCHY_PROBLEM_2}
can be written as first order problem
\begin{equation}\label{X_TRANSFORMED_SECOND_ORDER_EQUATION}
\begin{cases}[1.2]
\quad U'(t) \= \linB \, U(t) \\
\quad U(0) = (x_0,x_1),\\
\end{cases}
\end{equation}
where the operator $\linB$ is defined by
\begin{equation}\label{X_BLOCK_MATRIX}
\linB \eqdef \begin{pmatrix} \0 & \1 \\ \sqrt{\linA} & \0 \end{pmatrix}.
\end{equation}

\subparagraph{}
The main idea in proving existence and uniqueness of a solution is to show
that
the operator $\linB$ defined on a suitable chosen domain $\domain{\linB}$
is the infinitesimal generator of a semigroup of contractions.
The set $\domain{\linB}$ is defined according to
\[
\domain{\linB} \eqdef
\domain{\linA} \times \domain{\sqrt{\linA} }
\, \subset \,
\domain{\sqrt{\linA} }_{\norm{\cdot}_{\linA}} \times \hilH,
\]
where $\domain{\sqrt{\linA} }_{\norm{\cdot}_{\linA}}$
is the Hilbert space $\domain{\sqrt{\linA} }$
furnished with the graph norm.

\paragraph{} %- Lumer-Phillips theorem for second order problems
The following theorem is for example proved in \cite{X_Go}, p.\,111:

\begin{theorem1}\label{LUMER-PHILLIPS_THEOREM_2}
Let the linear operator $\linA$ be positive self-adjoint
and the operator $\linB$ be defined by \eqref{X_BLOCK_MATRIX}.
Then $\linB$ is the infinitesimal generator of a C$_0$-semigroup 
of contractions,
whose orbit $\eta$ determined by $\eta(0) \eq U(0)$ solves
\eqref{X_TRANSFORMED_SECOND_ORDER_EQUATION}.
\end{theorem1}

% BIBLIOGRAPHY
%------------------------------------------------------------------------------%
\vspace{5mm}
\begin{thebibliography}{99}
\begin{small}

\bibitem{X_EN}  K.-J.\,Engel, R.\,Nagel:
\textit{A Short Course on Operator Semigroups}.
Universitext, Springer, 2006.

\bibitem{X_Go} J.\,Goldstein:
\textit{Semigroups of Linear Operators and Applications}.
Oxford University Press, 1985.

\bibitem{X_Ha} M.\,Haase:
\textit{The Functional Calculus for Sectorial Operators}.
Birk\-h\"{a}user, 2006.

\bibitem{X_HP} E.\,Hille, R.\,S.\,Phillips:
\textit{Functional Analysis and Semigroups}.
AMS Vol.\,31, 1948.

\bibitem{X_Lu} A.\,Lunardi:
\textit{Analytic Semigroups and Optimal Regularity in Parabolic Problems}.
Birk\-h\"{a}user, 1995.

\bibitem{X_Pa} A.\,Pazy:
\textit{Semigroups of Linear Operators
        and Applications to Partial Differential Equations}.
Applied Mathematical Sciences Vol.\,44, Springer, 1983.

\end{small}
\end{thebibliography}

%------------------------------------------------------------------------------%
%  XI. COERCIVE SESQUILINEAR FORMS ---
%------------------------------------------------------------------------------%
%  §1. Generalities on Sesquilinear Forms ---
%  §2. Coercive Forms and Variational Equations ---
%  §3. Triples of Hilbert Spaces ---
%  §4. Linear Operators Associated with a H-Elliptic Form ---
%  §5. Sectorial Forms and m-Sectorial Operators ---
%  §6. The Friedrichs Extension Theorem ---
%  §7. Variational Evolution Equations ---
%  ---

\newpage
\thispagestyle{empty}
\def\TitleXI{Coercive Sesquilinear Forms}
\fancyhead[C] {\small{\textsc{XI. \TitleXI}}}
\part{\TitleXI}
%------------------------------------------------------------------------------%
Our aim in this part
is to present the general theory of coercive sesquilinear forms
and its applications to stationary and evolutional linear equations.
Sesquilinear forms are mostly efficient
when studying abstract problems in the Hilbert space setting.

A further goal is to indicate the correspondence of
sesquilinear forms and linear operators.
In fact,
each sesquilinear form gives rise to a couple of linear operators,
and reversely,
by taking the inner products $\scalar{\linA x}{y}$,
every linear operator $\linA$ defines a sesquilinear form
on its domain of definition $\domain{\linA}$.

\subparagraph{}
The totality of applications of the presented theory
to the boundary and inital-boundary value problems
of linear partial differential equations
are often denoted as \textit{form methods} or as \textit{Hilbert space methods}.

\section{Generalities on Sesquilinear Forms}
%------------------------------------------------------------------------------%
By a \textit{sesquilinear form} defined in complex linear space $\vecV$
is meant a mapping $\sql\:\vecV \xtimes \vecV \xto \C$,
which is linear in the first and conjugate linear in the second argument,
that is,
for all $x,y,z \xeof \vecV$ and $\lambda,\mu \xeof \C$ the computation rules
\[
\begin{cases}[1.2]
\quad \sql(x+y,z) \= \sql(x,z) \, + \, \sql(y,z) \\
\quad \sql(x,y+z) \= \sql(x,y) \, + \, \sql(x,z) \\
\quad \sql(\lambda x,\mu y) \= \lambda \overline{\mu} \, \sql(x,y) \\
\end{cases}
\]
are valid.
If the linear space $\vecV$ is real, 
however,
then $\sql$ is assumed to be linear also in the second argument,
so it becomes into a \textit{bilinear form} in this case.

\subparagraph{} %- quadratic form
The mapping $x \xeof \vecV \mapsto \sql[x] := \sql(x,x) \xeof \C$ is called
the \textit{quadratic form} of $\sql$.
Notice that all values $\sql(x,y)$ can be recovered
from the values $\sql[x]$ and $\sql[y]$ by certain polarization formulas.

\paragraph{} %- representation operators
A sesquilinear form $\sql$ may also be regarded
as a linear mapping $\repA\:\vecV \xto \vecV_a^{\#}$,
where $\vecV_a^{\#}$ is the (algebraic) conjugate space of $\vecV$,
that is,
the linear space of conjugate linear functionals $\phi\:\vecV \xto \C$.
More precisely,
given a sesquilinear form $\sql$,
each element $x \xeof \vecV$ gives rise to a conjugate linear functional
$\phi_x\:\vecV \xto \C$,
defined by $\phi_x(y) := \sql(x,y)$ for $y \xeof \vecV$.
Let $\vecV_a^{\#}$ again be the linear space
of conjugate linear functionals on $\vecV$.
Then the mapping
\[
\repA: \;
\begin{cases}[1.2]
\quad \repA\:V \xto V^{\#},\\
\quad \repA \, x := \phi_x \quad (x \in \vecV)
\end{cases}
\]
proves to be linear and is said to be the \textit{form operator}
or the \textit{representation operator} associated with $\sql$.

\subparagraph{}
Reversely,
for any linear operator $\cal{A}\:\vecV \xto \vecV_a^{\#}$,
there is a sesqui\-linear form $\sql\:\vecV \xtimes \vecV \xto \F$
such that $\sql(x,y) \eq (\cal{A} \, x) y$ for $x,y \xeof \vecV$.
Hence sesquilinear forms acting on a complex linear space $\vecV$
appear in a one-to-one correspondence to linear mappings
acting between the spaces $\vecV$ and $\vecV_a^{\#}$.
This correspondence remains valid,
if $\vecV$ is a \textit{real} linear space,
$\sql$ is a \textit{bilinear} form
and $\cal{A}$ is the associated mapping
from $\vecV$ to the adjoint space $\vecV^{\#}$.

\paragraph{} %- continuous forms
From now on we assume that $\vecV$ is a \textit{normed} linear space
with underlying field $\C$.
Then any sesquilinear form $\sql\:\vecV \xtimes \vecV \xto \C$
is called \textit{bounded},
if there is $\Theta > 0$ such that
\[
\abs{\sql(x,y)} \, \leq \Theta \norm{x} \cdot \norm{y}
\quad \quad (x,y \in \vecV).
\]
Due to the \textit{uniform boundedness principle},
this property is equivalent to say
that $\sql$ is continuous on the space $\vecV \xtimes \vecV$.

\subparagraph{} %- representation of continuous forms
The prototypical example of a continuous form is
the inner product in a complex Hilbert space $\hilH$,
where boundedness is ensured by the \textsc{Cauchy-Schwarz} inequality
$\abs{ \scalar{x}{y} } \leq \norm{x} \cdot \norm{y}$.

\subparagraph{}
Moreover,
for each continuous sesquilinear form $\sql$ on $\hilH$
there is a bounded linear operator $\linB_{\sql}\:\hilH \xto \hilH$
such that
\begin{equation}\label{XI_FORM_REPRESENTATION}
\sql(x,y) \= \scalar{\linB_{\sql} x}{y} \qfa x,y \in \hilH,
\end{equation}
which is a representation for the form $\sql$.
In fact,
if $x \xeof \hilH$ is fixed,
then the mapping $\sql(x,\cdot)\:\hilH \xto \C$
causes a continuous antilinear functional $\phi_x \xeof \hilH_a^*$,
and by \textsc{Riesz}'s representation theorem
there is $\tilde{x} \xeof \hilH$ such that
$\phi_x(y) \eq \scalar{\tilde{x}}{y}$ for all $y \xeof \hilH$.
This suggests to define $\linB_{\sql} x := \tilde{x}$ for every $x \xeof \hilH$.
If \eqref{XI_FORM_REPRESENTATION} holds
for some operator $\linB \xeof \BO{\hilH}$,
then the form $\sql$ is said to be represented by $\linB$.

\section{Coercive Forms and Variational Equations}
%------------------------------------------------------------------------------%
Any continuous sesquilinear form $\sql\:\hilH \xtimes \hilH \xto \C$
is called \textit{coercive}
if there is $\theta > 0$ such that
\begin{equation}\label{XI_COERCIVITY_PROPERTY}
\re \, \sql(x,x) \, \geq \theta \norm{x}^2
\qfa x \xeof \hilH.
\end{equation}
The estimate \eqref{XI_COERCIVITY_PROPERTY} especially implies the validness of
\[
\norm{\linB_{\sql} x}\norm{x}    \, \geq \,
\abs{\scalar{\linB_{\sql} x}{x}} \, \eq  \,
\abs{\sql(x,x)}                  \, \geq \,
\re \sql(x,x)                    \, \geq \,
\theta \norm{x}^2
\]
for all $x \xeof \hilH$.
Hence the representation operator $\linB_{\sql}$ is bounded below
with lower bound $\theta$,
and we may conclude by the fundamental \textit{bounded-below theorem}
that it is injective and the range $\range{\linB_{\sql}}$
forms a closed linear subspace in $\hilH$.

\paragraph{} %- Lax-Milgram theorem
Coercivity of a sesquilinear form $\sql$ is of great interest {\wrt}
the solvability of the so-called \textit{variational problem},
which may be formulated as follows:
given a continuous and conjugate linear functional $\phi \xeof \vecV_a^*$,
it is to find $x \xeof \vecV$ such that
\begin{equation}\label{XI_VARIATIONAL_EQUATION}
\sql(x,y) \= \phi(y) \qfa y \in \vecV.
\end{equation}

\subparagraph{}
The \textsc{Lax-Milgram} \textit{theorem} is of outstanding meaning
in applied functional analysis
and ensures well-posedness of the variational problem
for any right-hand side $\phi \xeof \vecV_a^*$:

\begin{theorem1}\label{LAX-MILGRAM_THEOREM}
If $\sql\:\vecV \xtimes \vecV \xto \C$ is bounded and coercive,
then for each $\phi \xeof \vecV_a^*$
the equation \eqref{XI_VARIATIONAL_EQUATION} owns exactly
one solution $x_{\phi} \xeof \vecV$.
In addition,
$\norm{\,x_{\phi}\,} \leq \theta^{-1} \norm{\phi}$ holds.
\end{theorem1}

\subparagraph{}
In other words,
if $\sql$ is bounded and coercive,
then the form operator $\repA\: \vecV \xto \vecV_a^*$ is bjective.
The bounded inverse theorem implies
that the inverse operator $\repA^{-1}$ is bounded,
too,
so the solution $x_{\phi}$ of \eqref{XI_VARIATIONAL_EQUATION}
depends continuously on $\phi$.

\begin{comment}
\subparagraph{}
The proof of theorem \ref{LAX-MILGRAM_THEOREM} is based on the fact
that the representation operator $\linB_{\sql}$ of $\sql$ is bounded below.

\Proof:
Since $\linB_{\sql}$ is bounded below,
it is injective and $\range{\linB_{\sql}}$ is closed.
It is still to show that $\linB_{\sql}$ is surjective.
If $z \in \range{\linB_{\sql}}^{\bot}$,
then for $\linB_{\sql} z \xeof \range{\linB_{\sql}} $ we have
\[
0 \=
\scalar{\linB_{\sql} z}{z}
\=
\sql(z,z)
\, \geq \,
\theta \norm{z}^2
\, \geq \, 0,
\]
which implies $z \eq 0$ and therefore $\overline{\range{\linB_{\sql}} } \eq \vecV$.
But $\range{\linB_{\sql}}$ is closed,
so $\range{\linB_{\sql}} \eq \vecV$.
Finally,
the inequality $\norm{x_f} \leq \theta^{-1} \norm{f}$
follows from the fact that $\norm{\linB_{\sql}^{-1}} \leq \theta^{-1}$.
\qed
\end{comment}

\paragraph{} %- operator version of the \textsc{Lax-Milgram} theorem
There is also an \textit{operator version} of theorem \ref{LAX-MILGRAM_THEOREM}.
Let $\linA\:\vecV \xto \vecV$ be linear and bounded.
We shall denote $\linA$ \textit{coercive} if
\[
\re \scalar{\linA x}{x} \, > \, \theta \norm{x}^2
\]
holds for some $\theta > 0$ and all $x \xeof \vecV$.
It should be clear
that the mapping $(x,y) \mapsto \scalar{\linA x}{y}$
defines a coercive form on $\vecV$.
Then the \textsc{Lax-Milgram} theorem implies
invertibility of $\linA$ and boundedness of $\linA^{-1}$,
where $\norm{\linA^{-1}} \leq \theta^{-1}$ holds.

\section{Triples of Hilbert Spaces}
%------------------------------------------------------------------------------%
It is often indicated to study abstract linear problems
in the setting of so-called \textit{pairs} $(\vecV,\hilH)$ of Hilbert spaces,
where the space $\vecV$ is continuously and densely embedded
into a second space $\hilH$.
Let us assume that $\vecV,\,\hilH$ are Hilbert spaces and

\begin{itemize}[leftmargin=15mm]
\item[(T-1)]
there is a linear mapping $\iota_1\:\vecV \xto \hilH$
such that $\overline{\range{\iota_1} }  \eq \hilH$;

\item[(T-2)]
there is $c>0$ such that
$\norm{\iota_1 v}_{\hilH} \leq c \norm{v}_{\vecV}$ for all $v \xeof \vecV$.
\end{itemize}

By means of the embedding $\iota_1$,
one can also think of $\vecV$ as a dense linear subspace in the space $\hilH$.

\subparagraph{} %- Gelfand triples
Each Hilbert pair $(\vecV,\hilH)$ with embedding $\iota\:\vecV \xto \hilH$
gives rise to introduce the \textit{conjugate pair} $(\hilH_a^*,\vecV_a^*)$,
where the conjugate space $\hilH_a^*$
can be seen as continuously and densely embedded
into the conjugate space $\vecV_a^*$
by a further embedding $\iota_2\:\hilH_a^* \xto \vecV_a^*$.
Then,
by identifying the spaces $\hilH$ and $\hilH_a^*$ via the \textsc{Riesz} map,
combining both Hilbert pairs leads to the triple
\begin{equation}\label{XI_GELFAND_TRIPLE}
\vecV \, \overset{\iota_1}{\hookrightarrow} \,
\hilH \, \overset{\iota_2}{\hookrightarrow} \, \vecV_a^*,
\end{equation}
sometimes called a \textit{Gelfand triple},
an \textit{evolution triple} or a \textit{rigged Hilbert space}.
The center space $\hilH$ is often denoted
the \textit{pivot space} of \eqref{XI_GELFAND_TRIPLE}.

\subparagraph{} %- property of Gelfand triples
Gelfand triples own a remarkable property.
Let $(\,\cdot\,,\,\cdot\,)$ and $\scalar{\,\cdot\,}{\,\cdot\,}$
be the inner products in $\vecV$ and $\hilH$,
respectively,
then for fixed chosen $x \xeof \hilH$
the action of the antilinear linear functional
$f_x := \iota_2 \, x \xeof \vecV_a^*$ to any element $v \xeof \vecV$
is represented by the inner product of $x$ and $v$ in $\hilH$,
{\ie} $f_x(v) \eq \scalar{x}{\iota_1 \, v}$ holds for all $v \xeof \vecV$.

\section{Linear Operators Associated with a H-Elliptic Form}
%------------------------------------------------------------------------------%
The presented theory of coercive sesquilinear forms in a Hilbert space $\vecV$
given by \S\,2 is often not sufficent for applications.
Hence it is quite natural to examine sesquilinear forms
also {\wrt} pairs $(\vecV,\hilH)$ of Hilbert spaces.
In applications it is often the case
that sesquilinear forms are merely continuous but not coercive.
In the setting of a Hilbert pair $(\vecV,\hilH)$,
however,
we try to make a form $\sql$ coercive
by adding an additional term,
which is just a multiplicity of the inner product in $\hilH$.

\subparagraph{}
Indeed,
if $\sql\:\vecV \xtimes \vecV \xto \C$ is continuous
and the slightly modified form
\[
\sql_{\lambda}(u,v) \eqdef
\sql(u,v) \, + \, \lambda \scalar{\iota_1 u}{\iota_1 v}\, \quad
(u,v \in \vecV)
\]
is coercive for some $\lambda \xeof \R$,
that is,
there is $\theta > 0$ such that
\begin{equation}\label{XI_QUASI_COERCIVE}
\re \sql(u,u)  \, + \,
\lambda \norm{\iota_1u}_{\hilH}^2 \, \geq \, \theta \norm{u}_{\vecV}^2
\quad (u \in \vecV),
\end{equation}
then $\sql$ is called
\textit{quasi-coercive},
\textit{$\hilH$-elliptic}
or simply a \textit{\GARDING} \textit{form}.
Clearly,
if $\sql_{\lambda_0}$ is coercive for some $\lambda_0 \xeof \R$,
then $\sql_{\lambda}$ is so for $\lambda > \lambda_0$.

\paragraph{} %- the operator associated with a quasi-coercive form
Quasi-coercive sesquilinear forms are often used
to define densely defined and closed linear operators
acting in the pivot space $\hilH$ of a given Gelfand triple.

\subparagraph{}
Let $\sql$ be a sesquilinear form defined on a Hilbert space $\vecV$,
which is the left hand side of a Gelfand triple $(\vecV,\hilH,\vecV_a^*)$
with continuous embeddings
$\iota_1\:\vecV \xto \hilH$ and
$\iota_2\:\hilH \xto \vecV_a^*$.
It is our aim to introduce a linear operator $\linA_{\sql}$
with domain space $\hilH$ and range space $\hilH$.
This will be done by restricting the form operator
$\repA\:\vecV \xto \vecV_a^*$ associated with $\sql$
to those elements $u \xeof \vecV$,
for which $\repA \, u \xeof \hilH$.
Indeed,
the domain $\domain{\linA_{\sql}}$ is defined to be the subspace
of those elements $u \xeof \vecV$,
for which $\repA \, u \xeof \iota_2(\hilH) \subset \vecV_a^*$ holds,
briefly denoted
\[
\domain{\linA_{\sql}} \eqdef \{ \, u \in \vecV: \repA \, u \in \hilH \, \}.
\]

\subparagraph{}
Clearly,
the set $\domain{\linA_{\sql}}$ can be regarded as a linear subspace
of $\hilH$ by means of the embedding $\iota_1$.
Then for $u \xeof \domain{A}$ there exists $y \xeof \hilH$
such that $\iota_2 y \eq \repA \, u$.
We set $x := \iota_1 u$ and $\linA_{\sql} x := y$
for $x \xeof \domain{\linA_{\sql}}$
and denote $\linA_{\sql}$
the \textit{operator in $\hilH$ associated with $(\vecV,\sql)$},
or shortly,
the \textit{variational operator} derived from $\sql$.
Another phrase is to say that
the operator $\linA_{\sql}$ is the \textit{part of $\repA$ in $\hilH$}.

\paragraph{} %- main theorem on quasi-coercive forms
The following theorem collects a couple of nice properties 
of the operator $\linA_{\sql}$
associated with a continuous and quasi-coercive form $\sql$.
Besides the \textsc{Lax-Milgram} theorem,
is is our second important theorem:

\begin{theorem1}\label{XI_QUASI-COERCIVE_FORMS_AND_ASSOCIATED_OPERATORS}
Let $(\vecV,\hilH,\vecV_a^*)$ be a triple of Hilbert spaces
and the sesquilinear form $\sql$ be continuous and quasi-coercive
on $\vecV$ with parameter $\lambda_0 \xeof \R$.
Then the following assertions are true:

\begin{itemize}[leftmargin=17pt]
\item[(a)]
$\linA_{\sql}$ is closed.

\item[(b)]
$\linA_{\sql}$ is densely defined in $\hilH$.

\item[(c)]
If the embedding of $\vecV$ into $\hilH$ is \textit{compact},
then the operator $\linA_{\sql} + \lambda$ has a compact inverse
$(\linA_{\sql} +\lambda)^{-1}$ for each $\lambda > \lambda_0$.

\item[(d)]
The operator $\linA_{\sql} + \lambda$ is m-accretive.

\item[(e)]
The operator $\linA_{\sql} + \lambda$ is sectorial.
\end{itemize}
\end{theorem1}

\begin{comment}
\subparagraph{}
To prove that $\linA_{\sql}$ is closed,
let $\{x_n\}_{n \in \N}$ be a sequence in $\domain{\linA_{\sql}}$
with limit element $x \xeof \hilH$
and $(\linA_{\sql} x_n$ be the sequence mapped by $\linA_{\sql}$
with limit element $y \xeof \hilH$.
It is to show that $x \xeof \domain{\linA_{\sql}}$ and $\linA_{\sql} x=y$.
Let $\abs{x} $ and $\norm{x}$ be the norms in $\vecV$ and $\hilH$,
respectively.
Clearly,
we have
\[
\norm{(\linA_{\sql}+\lambda) x_n - y - \lambda x} \xto 0 
\quad \for n \xto \infty.
\]
In the following we shall use the fact
that for $\lambda > \lambda_0$ the operators
$\linA_{\sql}+\lambda$ and $(\linA_{\sql}+\lambda)^{-1}$
forms linear homeomorphism between the normed spaces
$(\hilH,\norm{\,\cdot\,})$
and $(\domain{\linA_{\sql}},\abs{\,\cdot\,} )$,
so the latter one forms a closed subspace
within the normed space $(\vecV,\abs{\,\cdot\,} )$.
By this reason,
\[
x_n = (\linA_{\sql}+\lambda)^{-1} (\linA_{\sql}+\lambda) x_n
\to 
(\linA_{\sql}+\lambda)^{-1} (y + \lambda x)
\quad \for n \xto \infty
\]
{\wrt} the norm $\abs{\,\cdot\,} $ in $\domain{\linA_{\sql}}$.
Since $x$ belongs to $\hilH$ and $\domain{\linA_{\sql}}$ is closed,
the element $(\linA_{\sql}+\lambda)^{-1} x$
belongs to $\domain{\linA_{\sql}}$ with $\linA_{\sql} x \eq y$.

\subparagraph{}
To prove that $\domain{\linA_{\sql}}$ is dense in $\hilH$,
suppose that $y \xeof \hilH$ fulfills $\scalar{x}{y} \eq 0$
for all $x \xeof \domain{\linA_{\sql}}$.
It is to show that $y \eq 0$ in this case.
Since $\linA_{\sql}+\lambda$ is surjective,
there must exist $x_0 \xeof \domain{\linA_{\sql}}$ such that
$(\linA_{\sql}+\lambda) x_0 \eq y$.
Then
\[
0 \=
\scalar{\linA_{\sql} x_0}{x} \, + \, \scalar{\lambda x_0}{x} \=
\sql(x_0,x) \, + \, \lambda \scalar{x_0}{x}.
\]
Especially for $x \eq x_0$ the last equation leads to
\[
0 \=
\scalar{\linA_{\sql} x_0}{x_0} \, + \, \scalar{\lambda x_0}{x_0} \,\geq\,
\theta \norm{x_0}^2
\]
for some $\theta > 0$,
due to to quasi-coercivity of $\sql$,
hence $x_0 \eq 0$ and finally $y \eq \0$.

\subparagraph{}
Finally,
if $\iota_1 \: \vecV \xto \hilH$ is compact,
then $\iota_1 \circ (\linA_{\sql}+ \lambda)^{-1} \: \hilH \xto \hilH$
is compact,
too,
since the composition of a compact linear operator
and a continuous linear operator proves to be compact.
\qed.
\end{comment}

\paragraph{} %- operators associated with adjoint and symmetric forms
It is often helpful to consider the \textit{adjoint form} $\sql^*$
of a given sesquilinear form $\sql$ defined by
\[
\sql^* (u,v) \eqdef \overline{\sql(v,u)}
\quad \quad (u,v \in \vecV).
\]
Notice that if $\sql$ is continuous,
then $\sql^*$ is so.
Moreover,
the form $\sql$ is called \textit{symmetric},
if $\sql \eq \sql^*$ holds\footnote{
\,A well-known example of a symmetric form is the inner product
$\scalar{\,\cdot\,}{\,\cdot\,}_{\hilH}$ in a Hilbert space $\hilH$,
which differs from arbitrary symmetric sesquilinear forms
only by positive definiteness.}.

\subparagraph{} %- symmetric forms imply self-adjoint variational operators
With respect to the variational operators $\linA_{\sql}$ and $\linA_{\sql^*}$,
the following theorem is valid:

\begin{theorem1}\label{XI_QUASI-COERCIVE_FORMS_AND_ASSOCIATED_OPERATORS_2}
If $\sql^*$ is the adjoint form of $\sql$,
then ${\linA_{\sql}}^* \eq \linA_{\sql^*}$.
Furthermore,
if $\sql$ is symmetric,
then $\linA_{\sql}$ is self-adjoint.
\end{theorem1}

\paragraph{} %- generating operators
Given a Hilbert pair $(\vecV,\hilH)$,
the theorem \ref{XI_QUASI-COERCIVE_FORMS_AND_ASSOCIATED_OPERATORS_2}
is very helpful to construct self-adjoint linear operators in $\hilH$
by starting with symmetric forms in $\vecV$.
This approach can especially be performed,
if the form is the inner product in $\vecV$.
The operator $\linA$ associated with the inner product
$\scalar{\,\cdot\,}{\,\cdot\,}_{\,\vecV}$,
called the \textit{generating operator} of $(\vecV,\hilH)$,
is coercive and even self-adjoint in $\hilH$
by theorem \ref{XI_QUASI-COERCIVE_FORMS_AND_ASSOCIATED_OPERATORS_2}.

\paragraph{} %- construction of generating operators
In the following we derive the generating operator
of a pair $(\vecV,\hilH)$ of Hilbert spaces in a straightforward way.
First of all,
each fixed $y \xeof \hilH$
gives rise to a linear functional on $\vecV$ according to
\[
f_y(v) := \scalar{\iota v}{y}_{\hilH}, \quad \quad (v \in \vecV).
\]
The mapping $f_y$ is also bounded,
{\ie} $f_y \xeof \vecV_a^*$.
Then,
by the \textsc{Riesz} representation theorem,
there is $v_y \xeof \vecV$ such that
\[
\scalar{v}{v_y} \= f_y(v) \qfa v \in \vecV.
\]
We therefore obtain a mapping $\linB\:\hilH \xto \vecV$,
by $\linB y := v_y$ for $y \xeof \hilH$,
which proves to be linear and continuous,
and the inner product in $\hilH$ is obtained from that in $\vecV$ by
\[
\scalar{\iota x}{y}_{\hilH} \= \scalar{x}{\linB y}_{\,\vecV}
\quad \quad (x,y \in \hilH).
\]
But the operator $\linB$ is also symmetric and strictly positive.
These properties imply the existence of the inverse operator
\[
\linB^{-1}\:\domain{\linB^{-1} } = \range{\linB} \subset \vecV \xto \hilH,
\]
which indeed is not necessarily bounded,
but uniquely defined by the pair $(\vecV,\hilH)$.

\subparagraph{}
The operator $\linB^{-1}$ is symmetric and strictly positive,
too,
hence there exists the square root operator $\linB^{-1/2}$
having the property
\begin{equation}\label{XI_GENERATING_OPERATOR}
\scalar{\linB^{-1/2} x}{\linB^{-1/2} y}_{\hilH} \= \scalar{x}{y}_{\,\vecV}
\quad \quad (x,y \in \hilH).
\end{equation}
It turns out that $\domain{\linB^{-1/2}} \eq \vecV$,
which can be shown by contradiction.
Especially for $y=x$ in \eqref{XI_GENERATING_OPERATOR} we obtain the equality
\[
\snorm{\linB^{-1/2}x}_{\hilH} \= \norm{x}_{\,\vecV}
\quad \quad (x \in \vecV).
\]

\section{Sectorial Forms and m-Sectorial Operators}
%------------------------------------------------------------------------------%
In this paragraph we introduce
the notations of \textit{closed sectorial form} and
\textit{m-sectorial operator}
and show the relationships between them.

\subparagraph{} %- sectorial forms
Let $\hilH$ be a complex Hilbert space
and $\sql\:\domain{\sql} \xtimes \domain{\sql} \xto \C$ be a sesquilinear form,
but defined on a linear subspace $\domain{\sql}$ of $\hilH$.
Like for linear operators,
the \textit{numerical range} of $\sql$ is defined by
\[
\Theta(\sql) \eqdef
\{\, \sql(x,x): x \in \domain{\sql}, \, \norm{x}=1 \,\}.
\]
We then denote the form $\sql$ to be \textit{sectorial},
if $\Theta(\sql)$ is contained in a sector $\Sigma_{\,\psi,\omega}$,
that is,
\begin{equation}\label{XI_SECTOR_CONDITION}
\Theta(\sql) \, \subset \,
\Sigma_{\,\psi,\omega} \eqdef
\{\zeta \in \C: \zeta \neq \omega, arg(\zeta-\omega) < \psi \},
\end{equation}
where $\omega \xeof \R$ is the vertex
and $0 \leq \psi < \pi/2$ the semi-angle of the sector.
Notice that sectors are symmetric {\wrt} the real line.
Especially for $\psi \eq 0$ let $\Sigma_{\,0,\omega} := (\omega,\infty)$,
and for $\psi \eq \pi/2$ we have $\Sigma_{\,\pi/2,0} \eq \C_+$.
We also say that $\linA$ is $(\psi,\omega)$-sectorial,
if \eqref{XI_SECTOR_CONDITION} holds.

\subparagraph{} %- half-bounded forms
Now,
if $\sql$ is symmetric and $\sql \eq \sql^*$,
respectively,
then $\Theta(\sql) \subset \R$ and $\omega \leq \inf \Theta(\sql)$, 
so in this case we denote $\sql$ to be 
\textit{half-bounded} or \textit{bounded below}
and call $\omega$ the lower bound of $\sql$.
Moreover,
if $\sql$ is half-bounded and $(\psi,\omega)$-sectorial,
then the estimates
\[
\begin{cases}[1.2]
\quad \re \sql(x,x) \, \geq \, \omega \norm{x}^2 \\
\quad \im \sql(x,x) \, \leq \,
      \tan \psi \cdot (\re \sql(x,x) - \omega \norm{x}^2)\\
\end{cases}
\]
hold for all $x \xeof \domain{\sql}$.

\paragraph{} %- closeness of sectorial forms
Let $\re \sql := (\sql + \sql^*)/2$ be the \textit{real part} of $\sql$.
Then for every $\alpha > 0$ the slightly modified form
$[\cdot,\cdot]_{\alpha} := \re \sql - \omega + \alpha$
constitutes an inner product on the linear space $\domain{\sql}$.
It turns out
that the norms $\abs{\cdot} _{\alpha}$
derived from these inner products $[\cdot,\cdot]_{\alpha}$
are equivalent to each other.
We especially take $\alpha=1$,
pay attention to the inner product
\[
[x,y]_1 \= \re \sql(x,y) - \omega \scalar{x}{y} + \scalar{x}{y},
\qfa x,y \in \domain{\sql}
\]
and distinguish between the two cases {\wrt} completeness of the norm
\[
\abs{x} _1 \eqdef {[x,x]_1}^{1/2}
\quad \quad (x \in \domain{\sql}).
\]
In both cases it is our goal to obtain a Hilbert space $\hilH_{\sql}$
with inner product $[x,y]_1$ and norm $\abs{x} _1$
for $x,y \xeof \hilH_{\sql}$,
together with a linear embedding $\iota\:\hilH_{\sql} \xto \hilH$,
which should be continuous if possible,
so $(\hilH_{\sql},\hilH)$ hopefully becomes into a pair of Hilbert spaces.

\begin{itemize}[leftmargin=12pt]
\item[1.]
The normed linear space $(\domain{\sql},\abs{\,\cdot\,} _1$)
is \textit{complete}.\\
We denote the form $\sql$ \textit{closed} in this case,
and $\domain{\sql}$ becomes into a Hilbert space in its own right,
hence we may define $\hilH_{\sql} := \domain{\sql}$ in this case. 
Notice that the norm $\abs{\,\cdot\,} _1$ on $\hilH_{\sql}$
is stronger than the norm $\norm{\,\cdot\,}$ in $\hilH$
restricted to $\hilH_{\sql}$,
so the natural inclusion mapping $\iota\: \hilH_{\sql} \xto \hilH$
proves to be continuous.

\item[2.]
The normed linear space $(\domain{\sql},\abs{\,\cdot\,} _1)$
is \textit{not} complete rsp. $\sql$ is not closed.\\
We define the Hilbert space $\hilH_{\sql}$ as the completion
of $\domain{\sql}$ {\wrt} the norm $\abs{\,\cdot\,} _1$.
Moreover,
Cauchy sequences in $\hilH_{\sql}$ are also Cauchy sequences in $\hilH$,
hence for every $x \xeof \hilH_{\sql}$
there is $x' \xeof \hilH$ such that $x' \eq \iota x$
for a linear map $\iota\:\hilH_{\sql} \xto \hilH$.

\subparagraph{}
We want to extend the form $\sql$ from $\domain{\sql}$ to $\hilH_{\sql}$
as well,
and this will be done by defining
\[
\overline{\sql}(x,y) \eqdef \lim_{\,n \xto \infty} \sql(x_n,y_n)
\]
for $x,y \xeof \hilH_{\sql}$ and arbitrary sequences
$\{x_n\}_{n \in \N}, \{y_n\}_{n \in \N} \subset \domain{\sql}$
with $x_n \xto x$ and $y_n \xto y$.
It turns out that
the definition of $\overline{\sql}(x,y)$ is independent of
the sequences $\{x_n\}_{n \in \N}$ and $\{y_n\}_{n \in \N}$,
so this value is in fact well-defined for $x,y \xeof \hilH_{\sql}$.
\end{itemize}

\paragraph{} %- closable sectorial forms
In summary,
we obtain in both cases a Hilbert space $\hilH_{\sql}$ with inner product
$[\cdot,\cdot]_1$ and norm $\abs{\,\cdot\,} _1$ together with a linear embedding
$\iota\: \hilH_{\sql} \xto \hilH$,
but which in the second case is not necessarily continuous.

\subparagraph{} 
Regarding the second case,
the form $\sql$ is said to be \textit{closable},
if the mapping $\iota$ is continuous.
This definition especially means that every closed form is closable.
If $\sql$ is closable,
then the final result in both cases will be
that there exists a closed form,
namely $\sql$ itself in the first and $\overline{\sql}$ in the second case,
together with a continuous embedding of $\hilH_{\sql}$ into $\hilH$.

\paragraph{} %- sectorialness
Let us continue with the concept of \textit{sectorialness}
of a closed linear operator in a complex Banach or Hilbert space.
This class of operators has been introduced and first studied
by \textsc{T.\,Kato}.
Let again $\Sigma_{\,\psi,\omega}$
be an open sector with vertex $\omega \xeof \R$
and semi-angle $0 < \psi \leq \pi$.
In his book \cite{XI_Ka},
\textsc{Kato} denotes a linear operator $\linA$ in a Hilbert space
to be \textit{sectorial},
if the following conditions are fulfilled:

\begin{itemize}[labelsep=4mm,leftmargin=20pt]
\item[(S-1)]
$\linA$ is {quasi-accretive},
{\ie} $\scalar{(\linA + \omega) x}{x} \geq 0$
for some $\omega \xeof \R$ and $x \xeof \domain{\linA}$,

\item[(S-2)]
the numerical range $\Theta(\linA)$ is a subset of the sector
$\Sigma_{\,\psi,\omega}$.
\end{itemize}

Furthermore,
$\linA$ is called \textit{m-sectorial},
if it is at once sectorial and m-quasi-accretive,
that is,
$\linA + \omega$ is m-accretive for some $\omega \xeof \R$.

\paragraph{} %- correspondence between closable sectorial forms and operators
We finally want to elaborate the correspondence
between closed sectorial forms and m-sectorial operators.
It should be clear
that if $\linA$ is a m-sectorial operator in $\hilH$
with domain of definition $\domain{\linA}$,
then the sesquilinear form
\[
\sql_{\linA}: \;
\begin{cases}[1.2]
\quad \domain{\sql_{\linA}} \eqdef \domain{\linA}, \\
\quad \sql_{\linA}(x,y) \eqdef \scalar{\linA \,x}{y}
\, \for x,y \in \domain{\sql_{\linA}} \\
\end{cases}
\]
is sectorial and closed.

\subparagraph{}
Reversely,
let $\sql_0$ be a closable sectorial form in $\hilH$
with domain $\domain{\sql}$.
We consider the closed completion $\sql$ on the space $\hilH_{\sql}$.
A densely defined linear subspace $D \subset \domain{\sql}$ is denoted
a \textit{core} of $\sql$,
or equivalently,
the set $D$ is a core of $\sql$,
if $\sql$ is the closure of the restriction of $\sql$ to $D$.

\subparagraph{}
\textsc{Kato}\textit{'s first representation theorem} claims that
a densely defined and closed sectorial form gives rise
to a m-sectorial linear operator:

\begin{theorem1}\label{FIRST_REPRESENTATION_THEOREM}
Let $\sql$ be a densely defined and closed sectorial form in $\hilH$
with vertex $\omega$, semi-angle $\psi$ and domain of definition $\hilH_{\sql}$,
and let $\iota$ be the continuous embedding from $\hilH_{\sql}$ into $\hilH$.
Then there is a m-sectorial operator $\linA_{\sql}$ in $\hilH$,
whose domain of definition $\domain{\linA_{\sql}}$ is defined as follows:
Any element $x \xeof \hilH$ belongs to $\domain{\linA_{\sql}}$,
if $x \eq \iota x_0$ for some $x_0 \xeof \hilH_{\sql}$ and
$\sql(x_0,z) \eq \scalar{y}{\iota z}$ holds
for some $y \xeof \hilH$ and all $z \xeof \hilH_{\sql}$,
which suggests to define $\linA_{\sql} \, x := y$.

\paragraph{}
The operator $\linA_{\sql}$ owns a couple of properties:

\begin{itemize}[leftmargin=12pt]
\item[(a)]
$\domain{\linA_{\sql}}$ is dense in $\hilH$
and the preimage $\iota^{-1} \domain{\linA_{\sql}}$
is dense in $\hilH_{\sql}$.

\item[(b)]
If $x \xeof \hilH_{\sql}$, $y \xeof \hilH$
and $\sql(x,y) \eq \scalar{\iota x}{y}$ for all $x$ in a core of $\sql$,
then $\iota x \xeof \domain{\linA_{\sql}}$ and $y \eq \linA (\iota x)$.

\item[(c)]
$\Theta(\linA_{\sql})$ is a dense subset of $\Theta(\sql)$.

\item[(d)]
The adjoint form $\sql^*$ is also closed and sectorial
with vertex $\omega$ and semi-angle $\psi$.

\item[(e)]
The adjoint ${\linA_{\sql}}^*$ of $\linA_{\sql}$
is the $m$-sectorial operator associated with the adjoint form $\sql^*$.
\end{itemize}
\end{theorem1}

\subparagraph{}
Thus m-sectorial linear operators are in one-to-one correspondence
with closed sectorial forms.

\section{The Friedrichs Extension Theorem}
%------------------------------------------------------------------------------%
We want to apply the results of the last paragraph
to the following class of operators:
a \textit{symmetric} linear operator $\linA$ is called \textit{half-bounded},
if there is $\theta \xeof \R$ such that
\[
\scalar{\linA \,x}{x} \, \geq \, \theta \norm{x}^2
\qfa x \in \domain{\linA},
\]
or equivalently,
if $\Theta(\linA) \subset [\theta,\infty)$ for some $\theta \xeof \R$.
The operator $\linA$ is called \textit{positive} (\textit{strictly positive}),
if $\theta \eq 0$ ($\theta > 0$) can be chosen.
In case of $\theta > 0$,
the {Cauchy-Schwarz} inequality implies
that $\linA$ is also bounded below,
hence $\linA$ is surely injective.
In addition,
if $\linA$ is closed,
then the range $\range{\linA}$ is closed in $\hilH$,
which follows from the bounded-below theorem,
but $\range{\linA}$ may be a proper subspace of $\hilH$.

\paragraph{} %- Friedrichs extension
The main result in our topic is given by the fact
that each densely defined and half-bounded linear operator $\linA$
acting in a Hilbert space owns at least one \textit{self-adjoint extension}.

\subparagraph{}
First notice that if $\linA$ is half-bounded with $\theta \xeof \R$,
then it is $(0,\theta)$-sectorial,
and the sesquilinear form $\sql(x,y) := \scalar{\linA x}{y}$
is sectorial and densely defined.
But $\sql$ is also closable,
so we may switch to the sesquilinear form $\overline{\sql}$ defined on the
Hilbert space $\hilH_{\sql}$,
called the \textit{energetic space} of the operator $\linA$.
Then by theorem \ref{FIRST_REPRESENTATION_THEOREM}
there is a m-sectorial and even self-adjoint linear operator $\linA_{\sql}$,
that forms an extension of the half-bounded operator $\linA$,
called the \textit{Friedrichs extension} of $\linA$.

\subparagraph{}
The domain of definition $\domain{\linA_{\sql}}$ of the operator $\linA_{\sql}$
coincides with the linear subspace $\hilH_{\sql} \, \cap \, \domain{\linA^*}$.
Moreover,
the operator $\linA_{\sql}$ owns a useful property:
if one adds to the half-bounded operator $\linA$ the multiplication operator 
$x \xeof \hilH \mapsto \theta x$ with $\theta \xeof \R$,
then the {Friedrichs} extensions of $\linA$ and $\linA + \theta$
are related to each other by
$(\linA + \theta)_{\sql} \eq \linA_{\sql} + \theta$.

\section{Variational Evolution Equations}
%------------------------------------------------------------------------------%
Having the concept of Hilbert-Sobolev space available,
we want to review the variational approach to first and second order
evolution equations,
which has been developed by \textsc{J.-L.\,Lions}
in the midth of the last century.
It can be seen as the straightforward extension of
the \textsc{Lax-Milgram} theory from stationary equations to evolutional ones.

\subparagraph{}
Let $\T > 0$ and denote $\lambda$ the Lebesgue-Borel measure on $[0,\T)$.
By an \textit{abstract variational problem of parabolic type}
in a real Hilbert space triple $(\vecV,\hilH,\vecV^*)$
we mean a $\lambda$-nearly everywhere in $(0,\T)$ valid
linear operator-differential equation
\begin{equation}\label{XI_ABSTRACT_PARABOLIC_VARIATIONAL_PROBLEM}
\begin{cases}[1.2]
\quad u'(t) \, + \, \repA(t) \, u(t) \= f(t), \\
\quad u(0) = u_0, \\
\end{cases}
\end{equation}
where $u_0 \xeof \hilH$,
$f\:[0,\T) \xto \vecV^*$
and $\repA_{t \xeof [0,\T)}$ is a family
of linear operators $\repA\:\vecV \xto \vecV^*$.

\subparagraph{} %- definition of weak solution I
Let $\sql(t)$ be the form associated with the operator $\repA(t)$.
Then we call $u \xeof H^1(0,\T;\vecV,\vecV^*)$
a \textit{weak solution} of \eqref{XI_ABSTRACT_PARABOLIC_VARIATIONAL_PROBLEM},
if $u(0) \eq u_0$ and
\[
\scalar{u'(t)\,}{v(t)} \, + \, \sql(t,u(t),v(t)) \= f(t)(v)
\]
for all $v \xeof \vecV$ and $\lambda$-nearly everywhere $t \xeof [0,\T)$.

\subparagraph{} %- Lions theorem for parabolic problems
Regarding the existence and uniqueness of a \textit{weak solution}
for the initial value problem \eqref{XI_ABSTRACT_PARABOLIC_VARIATIONAL_PROBLEM},
the following theorem is of fundamental meaning:

\begin{theorem1}\label{XI_LIONS_THEOREM_1}
Let $\{\sql(t)\}_{t \xeof [0,\T)}$ fulfill the following conditions:

\begin{itemize}[leftmargin=12pt]
\item[(a)]
the mapping $t \xeof [0,\T) \mapsto \sql(t,u,v)$ is Lebesgue-measurable
for every fixed pair $(u,v) \xeof \vecV \xtimes \vecV$,

\item[(b)]
$\sql$ is \textit{uniformly} bounded in $[0,\T)$,
that is,
there is a constant $M \geq 0$ such that
\[
\abs{\sql(t,u,v)} \, \leq M \cdot \norm{u}_{\vecV} \norm{v}_{\vecV}
\qfa u,v \in \vecV, \; t \in [0,\T)\,,
\]
\item[(c)]
$\sql(t)$ is quasi-coercive for each $t \xeof [0,\T)$.
\end{itemize}

Then for each $u_0 \xeof \hilH$ and each $f \xeof L^2(0,\T;\vecV^*)$
there is exactly one solution $u \xeof H^1(0,\T;\vecV,\vecV^*)$
of \eqref{XI_ABSTRACT_PARABOLIC_VARIATIONAL_PROBLEM},
which continuously depends on the data $u_0$ and $f$.
\end{theorem1}

\subparagraph{}
The proof of this theorem has been published in \cite{XI_Li}
as \textit{Th\'{e}or\`{e}me 1.1}, page 46.
Both terms $u'$ and $\repA(t) u$
in \eqref{XI_ABSTRACT_PARABOLIC_VARIATIONAL_PROBLEM}
belong to the same space $L^2(0,\T;\vecV^*)$ as $f$,
which gives rise to say
that the problem has \textit{maximal regularity} in $\vecV^*$.

\paragraph{} %- second order equations
The concept of weak solution makes also sense
for linear evolution equations of second order.

\subparagraph{}
An \textit{abstract hyperbolic variational problem}
{\wrt} the Hilbert space triple $(\vecV,\hilH,\vecV^*)$ 
and a family of bilinear forms
$\sql(t)\:[0,\T] \xtimes \vecV \xtimes \vecV \xto \R$ is given by
\begin{equation}\label{XI_ABSTRACT_HYPERBOLIC_VARIATIONAL_PROBLEM}
\begin{cases}[1.2]
\quad \scalar{u''(t)\,}{v(t)} + \sql(t,u(t),v(t)) = \scalar{f(t)}{v} 
\quad (v \in \vecV, \\
\quad u(0)  = u_0 \in \vecV, \\
\quad u'(0) = u_1 \in \hilH, \\
\end{cases}
\end{equation}
where $\scalar{\cdot}{\cdot}$ is the inner product in $\hilH$.

\paragraph{}
The next theorem shows that
problem \ref{XI_ABSTRACT_HYPERBOLIC_VARIATIONAL_PROBLEM}
is \textit{weakly solvable} under a couple of conditions
on the involved bilinear forms:

\begin{theorem1}\label{XI_LIONS_THEOREM_2}
Let $\{\sql(t)\}_{t \xeof [0,\T)}$ be a family
of bilinear forms on $\vecV$ such that

\begin{itemize}[leftmargin=20pt]
\item[(a)]
$t \mapsto \sql(t)$ is \textit{uniformly} bounded in $[0,\T)$,

\item[(b)]
$\sql(t)$ is $\hilH$-elliptic for each $t \xeof [0,\T)$,

\item[(c)]
$\sql(t)$ satisfies the symmetry condition
$\sql(t,u,v) \eq \sql(t,v,u)$
for all $u,v \xeof \vecV, \; t \xeof [0,\T)$,

\item[(d)]
the mapping $t \mapsto \sql(t,u,v)$ is continuously differentiable in $[0,\T)$
for each fixed $(u,v) \xeof \vecV \xtimes \vecV$,

\item[(e)]
there is $M_1 \geq 0$ such that
\[
\abs{ \frac{d}{dt} \sql(t,u,v)} \, \leq \, 
M_1 \cdot \norm{u}_{\vecV} \norm{v}_{\vecV}
\qfa \, t \in [0,\T).
\]
\end{itemize}
Then the initial value problem
\eqref{XI_ABSTRACT_HYPERBOLIC_VARIATIONAL_PROBLEM}
has a unique weak solution $u \xeof H^1(0,\T;\vecV,\vecV^*)$
for all $u_0 \xeof \vecV$,
        $u_1 \xeof \hilH$ and
        $f \xeof L^2(0,\T;\hilH)$.
\end{theorem1}

\subparagraph{}
The symmetry condition (c) in \ref{XI_LIONS_THEOREM_2} is essential
to prove this theorem.
Moreover,
the solution map
\[
R\:L^2(0,\T;\hilH) \xtimes \vecV \xtimes \hilH \xto L^2(0,\T;\hilH)
\]
given by $(f,u_0,u_1) \mapsto u$ is linear in each argument and continuous,
thus problem \eqref{XI_ABSTRACT_HYPERBOLIC_VARIATIONAL_PROBLEM} proves
to be well-posed.

% BIBLIOGRAPHY
%------------------------------------------------------------------------------%
\vspace{5mm}
\begin{thebibliography}{99}
\begin{small}

\bibitem{XI_AT} W.\,Arendt, A.F.M. ter Elst:
\textit{From Forms to Semigroups}.
Operator Theory: Advances and Applications, Vol.\,221, pp.\,47-69.

\bibitem{XI_Ka} T.\,Kato:
\textit{Perturbation Theory for Linear Operators}, Springer Verlag,
Grundlehren der mathematischen Wissenschaften, Band 132, 1995.

\bibitem{XI_Li} J.-L.\,Lions:
\textit{\'{E}quations Diff\'{e}rentielles Op\'{e}rationelles et Probl\`{e}mes aux Limites}.
Springer,
Grund\-lehren der mathematischen Wissenschaften, Band 111, 1961.

\end{small}
\end{thebibliography}

%------------------------------------------------------------------------------%
% XII. APPLICATIONS TO DIFFERENTIAL OPERATORS ---
%------------------------------------------------------------------------------%
% § 1. Free Space and Boundary Value Problems ---
% § 2. Operators of Elliptic, Parabolic and Hyperbolic Type ---
% § 3. Functional Analytic Methods for Differential Operators ---
% $ 4. Realizations in Hölder and Lebesgue Spaces --
% § 5. Closeness and Closability of Differential Operators ---
% § 6. Formally Adoint Differential Expressions ---
% § 7. Gårding's Inequality ---
% § 8. Weak Solutions of Elliptic Boundary Value Problems ---
% § 9. The Fredholm Property of Properly Elliptic Operators ---
% §10. Spectral Theory of Strongly Elliptic Operators ---
% §11. Self-Adjoint Extensions of Half-Bounded Operators ---
% §12. Eigenvalue Problems in Bounded Domains ---
% §13. Linear Differential Operators in Distributional Spaces ---
% §14. Free Space Problems ---
% §15. Fourier Transform Methods ---

\newpage
\thispagestyle{empty}
\def\TitleXII{Applications to Differential Operators}
\fancyhead[C] {\small{\textsc{XII. \TitleXII}}}
\part{\TitleXII}
%------------------------------------------------------------------------------%
This final section of this text deals with applications of the presented theory
to linear differential operators and some standard problems
related with them.

\subparagraph{} %- multi-index
Remember first the notation of \textit{multi-index},
which is nothing but a $d$-tuple $\alpha \eq (\alpha_1,\ldots,\alpha_d)$
with values $\alpha_i \xeof \N_0$
and with modulus $\abs{\alpha} := \alpha_1 + \ldots + \alpha_d$.
Multi-indices are used to shorten multiple partial derivatives
$
\PA f \eqdef \partial_1^{\alpha_1}
             \partial_2^{\alpha_2} \ldots
             \partial_d^{\alpha_d} f
$
of a sufficiently often differentiable and complex-valued function $f$,
but also to denote the multiple powers
$\xi^{\alpha} := \xi_1^{\alpha_1}
\cdot \xi_2^{\alpha_2} \cdot \ldots \cdot \xi_d^{\alpha_d}$
of a vector $\,\xi \eq (\xi_1,\xi_2,\ldots,\xi_d) \xeof \C^d$.

\paragraph{} %- linear partial differential operators
By the usage of multi-indices,
a \textit{linear partial differential operator} $\A$ of even order $2m$
is defined by
\begin{equation}\label{XII_EVEN_ORDER_LPDO}
\A(x,\,\cdot\,) \eqdef
\sum_{\abs{\alpha} \leq 2m} c_{\alpha}(x) \, \PA
\end{equation}
for $x \xeof \clG$,
the closure of a \textit{domain} $\G$,
that is,
an open and connected subset of an Euclidean space $\Rd$.
This domain may be bounded or unbounded and even the whole of $\Rd$.

\subparagraph{} %- essentially real differential operators
We shall call $\A$ \textit{essentially real},
if the coefficients $c_{\alpha}$ of highest order $\abs{\alpha} \eq 2m$
are real-valued.
Moreover,
the expression $\A$ is said to be
an \textit{operator with constant coefficients},
if all functions $c_{\alpha}$ are constant in $\clG$.

\section{Free Space and Boundary Value Problems}
%------------------------------------------------------------------------------%
In case of $\G \eq \Rd$ respectively $\bndG \eq \emptyset$,
where $\bndG := \overline{\G} \setminus \interior{\G}$ is the boundary of $\G$,
we denote the linear partial differential equation $\A \, u \eq f$
with $\A$ given by \eqref{XII_EVEN_ORDER_LPDO} a \textit{free space problem}.
On the other hand,
partial differential equations considered in domains
with non-empty boundary are usually furnished with \textit{boundary conditions},
which fix possible solutions and their normal derivatives up to a certain order
in the boundary $\bndG$.
Thus,
when having in mind linear equations of the form $\A\,u \eq f$
with imposed boundary conditions,
we shall speak of \textit{linear boundary value problems}.

\subparagraph{} %- general form of boundary value problems
The most general form of a boundary value problem
posed in a domain $\G \subsetneq \Rd$ is of the form
\begin{equation}\label{XII_LEBVP}
\begin{cases}[1.2]
\quad \A(x,u)     \= f(x)   & (x \xeof \G),\\
\quad \BDO_k(x,u) \= g_k(x) & (x \xeof \bndG \and k \eq 0,\ldots,m{-}1),\\
\end{cases}
\end{equation}
where $\A$ is a {linear} differential operator of order $2m$
and $\BDO_0,\ldots,\BDO_{m-1}$ are \textit{boundary operators} of order $m_k$,
each of them defined by
\[
\BDO_k(x,u) \=
\sum_{\abs{\alpha} \leq m_k} b_{k,\alpha}(x) \, \PA u(x).
\]

\subparagraph{}
To introduce the boundary operators $\BDO_k$,
it is sufficent
that the real or complex-valued functions $b_{k,\alpha}$
are defined in a neighbourhood of the boundary $\bndG$ only.
Then each operator $\BDO_k$ is a linear differential operator in its own right
of order $m_k {\,\leq\,} 2m{-}1$,
which maps the function $u \: \clG \xto \C$
to a function $\BDO_k u \: \bndG \xto \C$
belonging to a so-called \textit{boundary space} $B_k(\bndG)$.

\paragraph{} %- homogeneous and normal boundary operators
The collection $\{B_k\}_{k=0}^{m-1}$ of boundary operators in \eqref{XII_LEBVP}
is said to be \textit{homogeneous}
if $\BDO_k(x,\cdot) \eq 0$ for $x \xeof \bndG$ and $k \eq 0,\ldots,m{-}1$.
Moreover,
it is denoted \textit{normal} if

\begin{itemize}[leftmargin=15mm,itemsep=2pt]
\item[(B-1)]
$m_i \neq m_j$ for $i \neq j$,

\item[(B-2)]
for all $\xi \neq 0$ being normally positioned in $x \xeof \bndG$
we have
\[
\sum_{\abs{\alpha} = m_k} b_{k,\alpha}(x) \, \xi^{\alpha} \neq 0
\quad \quad (k=0,\ldots,m{-}1).
\]
\end{itemize}

\subparagraph{} %- homogeneous Dirichlet systems
Finally,
our collection of boundary operators is called a \textit{Dirichlet system},
if it is normal and the values $m_k$ form a permutation
of the set $\{0,\ldots,m{-}1\}$,
so one may reorder the operators $\BDO_k$ to obtain $m_k \eq k$
($0 \leq k \leq m{-}1$).
The most regarded set of boundary conditions in applications
is the \textit{homogeneous Dirichlet system} defined by
\begin{equation}\label{XII_HOMOGENEOUS_DIRICHLET_SYSTEM}
\BDO_k(x,u) \eqdef (\partial_{\nu}^k u)(x) \= 0
\end{equation}
for $x \xeof \bndG$ and $k \eq 0,\ldots,m{-}1$,
where $(\partial_{\nu}^k u)(x)$ is the $k$-th derivative of a function $u$
along the outer normal vector $\nu$ in $x \xeof \bndG$.

\subparagraph{} %- classical solutions of BVP endowed with a Dirichlet system
If the set of boundary conditions in \eqref{XII_LEBVP} is a Dirichlet system,
then we denote a solution $u \: \clG \xto \C$
of this problem \textit{classical},
if
\begin{itemize}[leftmargin=9mm]
\item[(a)] $u \xeof C^{2m}(\G) \, \cap \, C^{m-1}(\clG)$,
\item[(b)] $u$ fullfills $\A(x,u) \eq f(x)$ in each $x \xeof G$ and
\item[(c)] the $m{-}1$ equations $\BDO_k(x,u) \eq g_k(x)$ hold
           in every $x \xeof \bndG$.
\end{itemize}

\paragraph{} %- boundary condition for second-order problems
Especially to each \textit{second-order} boundary value problem
there is assigned exactly one boundary operator $\BDO_0$,
hence problem \eqref{XII_LEBVP} becomes into
\begin{equation}\label{XII_SECOND_ORDER_BOUNDARY_VALUE_PROBLEM}
\begin{cases}[1.2]
\quad \A(x,u)     = f(x)   & (x \xeof \G),\\
\quad \BDO_0(x,u) = g_0(x) & (x \xeof \bndG).\\
\end{cases}
\end{equation}

\subparagraph{} %- enumeration of boundary conditions for second-order problems
Boundary value problems of the form \eqref{XII_SECOND_ORDER_BOUNDARY_VALUE_PROBLEM}
and governed by a linear differential operator $\A$ of second order
are obtained by various boundary operators $\BDO_0$.
We denote

\begin{itemize}[leftmargin=15mm,itemsep=2pt]
\item[(BC-1)]
\;$\BDO_0(x,u) \eq u(x)$ for $x \xeof \bndG$ a \textit{Dirichlet condition},

\item[(BC-2)]
\;$\BDO_0(x,u) \eq u_{\nu}(x)$ for $x \xeof \bndG$ a \textit{Neumann condition},

\item[(BC-3)]
\;$\BDO_0(x,u) \eq \alpha(x)u(x)+\beta(x) u_{\nu}(x) \eq 1$ for $x \xeof \bndG$ 
a \textit{Robin condition},
where the functions $\alpha, \beta$ must fulfill 
\[
\alpha(x)^2 + \beta(x)^2 \neq 0 \quad \quad (x \in \bndG).
\]
\end{itemize}

\subparagraph{}
Due to physical reasons,
one often has to split the boundary set $\bndG$
into disjoint subsets $\Gamma_1,\Gamma_2,\Gamma_3$,
on which the conditions (BC-1), (BC-2) and (BC-3) are prescribed.
We then shall speak of \textit{mixed} boundary value problems
for the second-order operator $\A$.

\section{Operators of Elliptic, Parabolic and Hyperbolic Type}
%------------------------------------------------------------------------------%
The concept of \textit{ellipticity} has been shown to be helpful in order
to classify linear partial differential expressions
by a simple algebraic criterion.
Let $\A$ be a linear differential operator of order $2m$ be given
by \eqref{XII_EVEN_ORDER_LPDO} and let
\[
p(x,\xi) \eqdef \sum_{\abs{\alpha} = 2m} c_{\alpha}(x) \, \xi^{\alpha}
\quad \for \xi = (\xi_1,\ldots,\xi_d) \in \Rd
\]
be the \textit{characteristic function} or the \textit{main symbol} of $\A$,
which is built up by the summands of highest order in \eqref{XII_EVEN_ORDER_LPDO}.

\begin{itemize}[leftmargin=12mm]
\item[(E-1)]
The expression $\A$ is called \textit{elliptic} in $x \xeof \clG$,
if
\[
\abs{p(x,\xi)} \, \neq \, 0 \qfa \xi \neq 0,
\]
and $\A$ is said to be \textit{uniformly elliptic} in $\G$ (in $\clG$),
if it is elliptic in \textit{each} point $x \xeof \G$ (in $x \xeof \clG$).
We then call a boundary value problem $\A \,u \eq f$ \textit{elliptic},
if $\A$ is so.

\item[(E-2)]
$\A$ is called \textit{properly elliptic} in $\clG$,
if it is elliptic and 
if for each fixed $x \xeof \bndG$ and all linear independent vectors
$\xi, \eta \xeof \Rd$ the main symbol $P_{\A}(x, \xi + \tau \eta)$ has,
{\wrt} the variable $\tau$,
exactly $m$ roots with positive imaginary part and
exactly $m$ roots with negative imaginary part.
It can be shown that
if $\A$ is \textit{essentially real} or $\G \xsubset \Rd$ with $d \geq 3$,
then any elliptic differential operator is also {properly elliptic}.

\item[(E-3)]
The concept of ellipticity is often not good enough
to obtain satisfying results.
Hence the expression $\A$ of order $2m$
is called \textit{strongly elliptic} in $x \xeof \clG$,
if there is $\theta_x> 0$ such that
\[
\re \, (-1)^m \, p(x,\xi)
\, \geq \,
\theta_x \cdot \abs{\xi} ^{2m}
\quad
\fa \xi \neq 0.
\]

\item[(E-4)]
$\A$ is called \textit{uniformly strongly elliptic} in $\clG$,
if there is $\theta {\,>\,} 0$ such that 
\[
\re \, (-1)^m \, p(x,\xi)
\, \geq \,
\theta \cdot \abs{\xi} ^{2m}
\quad
\fa \xi \neq 0 \and x \xeof \clG.
\]
\end{itemize}

\subparagraph{} %- Laplace operator
The most well-known linear partial differential operator of second order
is the \textit{Laplacian}
\[
\Delta = \partial^2 / \partial {x_1}^2 + \ldots
       + \partial^2 / \partial {x_d}^2,
\]
whose negative version $-\Delta$ constitutes the prototype
of what is meant by an {uniformly strongly elliptic operator}.
One easily computes
$\re -p_{-\Delta}(x,\xi) \eq \abs{\xi} ^2$ for $\xi \xeof \Rd$,
hence the expression $-\Delta$ often serves as a first example
when studying linear-elliptic operators.

\paragraph{} %- parabolic and hyperbolic differential operators
Besides the class of {linear-elliptic} operators
there are two further important classes
of linear partial differential operators.
They are called to be of \textit{parabolic} and \textit{hyperbolic} type.
In order to introduce these concepts,
we consider differential expressions on cylindrical subsets
$\J \xtimes \G$,
where $\J \eq (0,\T) \xsubset \R$ and $\G \xsubset \Rd$ is a domain.
Then a \textit{parabolic} differential operator $\A$ is one of the form
\[
\A(t,x,u) \= \frac{\partial u}{\partial t}(x) \+ \A_1(x,u),
\]
where $\A_1$ is \textit{uniformly elliptic} in $\G$,
and $\A$ is of \textit{hyperbolic} type,
if it is of the form
\[
\A(t,x,u) \= \frac{\partial^2 u}{\partial t^2}(x) \+ \A_1(x,u),
\]
where again $\A_1$ is uniformly elliptic in $\G$.
Simple examples of para\-bolic and hyper\-bolic operators are
the \textit{heat operator} $\partial_t    - \Delta$ and
the \textit{wave operator} $\partial_{tt} - \Delta$,
respectively.
With these operators we associate the linear evolution equations
\[
\frac{\partial u}{\partial t} + \A(\cdot,u) = f
\quad \mbox{ and } \quad
\frac{\partial^2 u}{\partial t^2} + \A(\cdot,u) = f,
\]
both of them called \textit{initial value problems},
when considered with the prescribed start value $u(0) \eq u_0$
(and in addition $\dot{u}(0) \eq u_1$ in the hyperbolic case).

\subparagraph{} %- initial value problems
When considering parabolic or hyperbolic equations,
we speak of so-called \textit{initial-boundary value problems},
if for every time $t \xeof (0,T]$ the solution fulfills
a certain set of boundary conditions in every $x \xeof \bndG$,
where we especially have in mind the homogeneous Dirichlet boundary conditions
\eqref{XII_HOMOGENEOUS_DIRICHLET_SYSTEM}.
However,
in case of $\G \eq \Rd$,
we obtain free space problems or pure initial value problems
for the differential operators of parabolic or hyperbolic type.

\section{Functional Analytic Methods for Differential Operators}
%------------------------------------------------------------------------------%
It is the aim of this part to apply
all the presented functional analytic tools to the just enumerated problems.
Let $\A$ be a linear partial differential operator of order $k \xeof \N$
(not necessarily even) and given by
\begin{equation}\label{XII_ANY_ORDER_LPDO}
\A(x,u) \eqdef
\sum_{\abs{\alpha} \leq k} c_{\alpha}(x) \, \PA u.
\end{equation}
Then by a \textit{realization} of $\A$ in a Banach space $\mathbb{E}(\G)$
of functions $u \: \G \xto \C$ is meant any linear operator
\[
\linA: \domain{\linA} \subset \mathbb{E}(G) \xto \mathbb{E}(G)
\]
with domain of definition $\domain{\linA}$.
The Banach space $\mathbb{E}(\G)$ should contain the space $\TF$ 
of test functions\footnote{
\,The space $\TF$ consists of indefinitely differentiable functions
having \textit{compact} support
\[
\supp{\,\phi} \eqdef \overline{\set{x \in \G}{\phi(x) \neq 0}}\,,
\]
usually denoted the \textit{space of test functions}.
This space plays an out\-standing role in analysis,
since it lies dense in a multiplicity of other functional spaces.}
as a dense linear subspace.
We also require for each $u \xeof \domain{\linA}$ that
first $\A \, u$ is well-defined by
$\linA u \eq \A \, u$ and second
$\A \, u \xeof \mathbb{E}(\G)$ holds.
Clearly,
our first condition means that the set $\domain{\linA}$
may contain only functions $u \: \G \xto \C$ being sufficiently often
differentiable in some sense.
But $\domain{\linA}$ should also include the space $\TF$ of test functions,
so $\domain{\linA}$ lies dense in $\mathbb{E}(\G)$ as well 
and $\linA$ appears as a densely defined operator in $\mathbb{E}$.

\paragraph{} %- content of the section
In this section we want to point out the following facts:

\begin{itemize}[leftmargin=2mm,itemsep=4pt]
\item[-] 
for a linear partial differential expression $\A$
there can be defined a couple of densely defined and closable linear operators
$\linA$ acting either in the {\HOELDER} spaces $C^{k,\gamma}(\G)$
($k \xeof \N_0$, $0 < \gamma \leq 1$)
or in the Lebesgue spaces $\Lp$ ($1 \leq p < \infty$);

\item[-] 
using elliptic estimates,
the realization of a strongly elliptic differentiable operator
acting in Sobolev spaces
under homogeneous Dirichlet boundary conditions is closed;

\item[-] 
the study on realizations of differential operators can be supported
by that on their adjoints,
where those are determined by so-called \textit{adjointness theorems};

\item[-] 
by starting with quasi-coercive sesquilinear forms,
created by strongly elliptic differential expressions,
we can obtain realizations in $\Lh$,
which prove to be densely defined and closed;

\item[-] 
weak solutions of initial-boundary value problems exist and are unique,
supposed the governing sesquilinear form is the Dirichlet form
of a strongly elliptic differential operator;

\item[-] 
the realizations of a regularly elliptic expression $\A$ in $\Lh$ are Fredholm,
so the alternative theorem may be applied to the equation $\A \, u \eq f$;

\item[-] 
strongly elliptic expressions give rise to m-accretive and 
sectorial realizations in $\Lh$,
so the related parabolic and hyperbolic initial value problem are well-posed;

\item[-] 
formally self-adjoint and half-bounded differential expressions
lead to self-adjoint realizations in $\Lh$;
\end{itemize}

\subparagraph{}
This enumeration is of course only a small selection 
within the totality of applications of functional analytic methods
to linear partial differential operators and related equations,
but may give a good impression of all the capabilities.

\section{Realizations in {\HOELDER} and Lebesgue Spaces}
%------------------------------------------------------------------------------%
We basically want to distinguish between two classes of functional spaces,
in which realizations of a linear differential operator $\A$
are usually studied:

\begin{itemize}[leftmargin=4mm]
\item[1.] %- realizations in Hölder spaces
The notation of classical derivative $\PA u$,
applied to a function $u \: \G \xto \C$,
leads us to the usage
of the spaces $C^{\gamma}(\G)$ of \textit{{\HOELDER}-constinuous} functions,
which constitute Banach spaces when endowed with a suitable norm.

\subparagraph{}
A realization of a differential operator $\A$ of order $2m$
acting in {\HOELDER} space $C^{\gamma}(\G)$ is nothing but a mapping
\[ 
\linA \: \domain{\linA} \subset C^{\gamma}(\G) \to C^{\gamma}(\G),
\]
where $\domain{\linA}$ forms a linear subspace of $C^{2m,\gamma}(\G)$
restricted by boundary conditions,
for example by a Dirichlet system.

\subparagraph{}
There is an extensive theory on the realizations
of linear partial differential operators acting in {\HOELDER} spaces,
which provide nice results especially for operators of second order.
We should mention \textsc{J.\,Schauder} as the outstanding pioneer
in the investigation of this subject.
But \textsc{Schauder} uses classical methods for the setup of his theory,
coming from potential theory and complex analysis,
so we do not want to go into the details
(a good reference for this subject is \cite{XII_GT}).

\item[2.] % realizations in Lebesgue spaces
Instead of classical derivatives,
the modern theory of partial differential equations
is based on the concept of \textit{weak derivative},
since this approach leads to larger spaces of differentiable functions
than the already mentioned {\HOELDER} spaces,
usually denoted as \textit{Lebesgue-Sobolev spaces}.

\subparagraph{} %- Hilbert space L^2
Recall that
for any domain $\G \xsubset \Rd$
the Hilbert space $\Lh$ is a complete normed space
consisting of measurable functions $f \: \G \xto \C$ such that
\[
\int_{\,\G} \, \abs{f} ^2 \, dx \, < \, \infty,
\]
where the inner product in $\Lh$ is defined by
\[
\scalar{f}{g} \eqdef \int_{\,\G} \, f {\cdot} \overline{g} \, dx
\quad \for f,g \in \Lh.
\]
This space is a specific member in the scale of the so-called
\textit{Lebesgue spaces} $\Lp$ ($1 \leq p < \infty$),
whose norms are defined by
\[
\abs{f} _p \eqdef ( \int_{\,\G} \, \abs{f} ^p \, dx)^{\frac{1}{p}} \,<\,\infty
\]
for equivalence classes of measurable functions $f \xeof \Lp$.

\subparagraph{}
Now a realization of a linear differential operator $\A$ of order $k$
in Lebesgue space $\Lp$ is any linear mapping
\[
\linA \: \domain{\linA} \subset \Lp \to \Lp,
\]
for which $\domain{\linA}$ forms a linear subspace of $W^{k,p}(\G)$,
for example restricted by homogeneous Dirichlet boundary conditions.
\end{itemize}

\section{The Closeness of Strongly Elliptic Dirichlet Operators}
%------------------------------------------------------------------------------%
%-- the closeness of strongly elliptic Dirichlet operators
One often has to deal with the realization $\linA_D$
of a strongly elliptic differential operator $\A$ in $\Lp$
having order $2m$ and domain of definition
\[
\domain{\linA_D} \eqdef \{ u \xeof \WN{m,p}{\G} \: \A\,u \xeof \Lp \},
\]
so the functions in $\domain{\linA_D}$ satisfy all the $m$ boundary conditions
of the homogeneous Dirichlet system \eqref{XII_HOMOGENEOUS_DIRICHLET_SYSTEM}.
By this reason it is reasonable to denote the realization $\linA_D$
a \textit{Dirichlet operator}.

\paragraph{} %- closeness of strongly elliptic differential operators
We turn to the question of closeness of the realization $\linA_D$.
Let the domain $\G$ and the expression $\A$ fulfill the following conditions:

\begin{itemize}[leftmargin=9mm]
\item[(I)\;\;]
$\G$ is bounded in $\Rd$,

\item[(I')\;\,]
$\G$ satisfies the $m$-th \textit{extension property},
that is,
the linear space $R_0(\G)$ consisting
of those functions $\phi \xeof W^{m,p}(\G)$,
which are restrictions $\phi |_{\G}$
of functions $\phi \xeof C_c^{\infty}(\Rd)$,
is dense in $W^{m,p}(\G)$;

\item[(II)\;]
$\A$ is uniformly strongly elliptic in $\G$;

\item[(III)\,]
the coefficient functions $a_{\alpha \beta}$ are essentially bounded in $\G$,
that is,
$a_{\alpha \beta} \xeof L^{\infty}(\G)$;

\item[(III')]
the coefficient functions $a_{\alpha \beta}$ of highest order
$\abs{\alpha} + \abs{\beta} \eq 2m$ are uniformly continuous in the closure
$\clG$ of $\G$.
\end{itemize}

Then,
by showing the validness of the \textit{elliptic estimate}
\[
\abs{u} _{2m} \, \leq \, c \cdot \{ \,\abs{u} _p + \abs{\linA_D u} _p \, \}
\]
for $u \xeof \domain{\linA_D}$,
where $\abs{u} _p \eq \norm{u}$ is the norm of $u$ in $\Lp$,
theorem \ref{II_BROWDER_THEOREM} implies that the operator $\linA_D$ is closed.
Moreover,
$\linA_D$ is injective and the range of $\linA_D$ is a closed subspace in $\Lp$.

\section{Formally Adjoint Differential Expressions}
%------------------------------------------------------------------------------%
%-- formally adjoint expressions and adjointness theorems
As we have already mentioned,
it is often helpful to combine the study of a linear operator
with that of its formally adjoint.
Supposed the linear partial differential operator $\A$ of order $k$
has coefficients $c_{\alpha} \xeof C^{\abs{\alpha} }(\G)$.
We then find for $u \xeof \TF$ and $\psi \xeof \TF$
by repeated partial integration
\[
\int_{\,\G} \sum_{\abs{\alpha} \leq k}
(c_{\alpha} \, \PA u) \cdot \overline{\psi} \, dx
\=
\int_{\,\G} u \cdot \sum_{\abs{\alpha} \leq k}
(-1)^{\abs{\alpha} } \, \overline{\PA(\overline{c_{\alpha}} \, \psi)} \, dx.
\]
Hence there is a \textit{formally adjoint} differential expression
\begin{equation}\label{XII_FAO}
\A^{\dag}(x,u) \eqdef
\sum_{\abs{\alpha} \leq k} (-1)^{\abs{\alpha} }
\PA(\overline{c_{\alpha}(x)} \, u(x))
\end{equation}
such that the identity
\[
\int_{\,\G} \A \, \phi \, \overline{\psi} \, dx 
\=
\int_{\,\G} \phi \, \overline{\A^{\dag} \, \psi} \, dx 
\]
holds for $\phi,\psi \xeof \TF$,
briefly
$
\scalar{\A\,\phi}{\psi} \eq \scalar{\phi}{\A^{\dag}\,\psi}
$
by means of the inner product in $\TF$.
Using \textsc{Leibniz}'s rule to compute the terms
$\PA(\overline{c_{\alpha}} \, u)$ in \eqref{XII_FAO},
there are functions $c_{\alpha}'\: \G \xto \C$
such that
\[
\A^{\dag}(x,u) \=
\sum_{\abs{\alpha} \leq k} c_{\alpha}'(x) \; \PA u(x).
\]
Thus $\A^{\dag}$ is uniquely determined by $\A$
and constitutes a differential expression of type \eqref{XII_ANY_ORDER_LPDO},
too.
It may happen that $\A^{\dag} \eq \A$ holds,
and in this case
we shall call the expression $\A$ \textit{formally self-adjoint}.

\subparagraph{} %- a criterion for closability of differential operators
A useful criterion to ensure closability of a linear operator $\linA$
acting in Hilbert space $\Lh$
and induced by a differential expression $\A$ of order $k$
is given by the following theorem
(see also \cite{XII_Yo}, p.\,78),
whose proof uses the presence of the formally adjoint expression $\A^{\dag}$:

\begin{theorem1}\label{XII_L2_CLOSABILITY_THEOREM}
Let the expression $\A$ of order $n$ be given by \eqref{XII_ANY_ORDER_LPDO}
with coefficients $c_{\alpha} \xeof C^k(\G)$ for $\abs{\alpha} \leq k$ and
the linear operator
\[
\linA: \;
\begin{cases}[1.2]
\quad \domain{\linA} \eqdef \{u \xeof C^k(\G) \cap \Lh \:\A \, u \xeof \Lh \},\\
\quad \linA u \eqdef \A\,u \; \for u \xeof \domain{\linA}\\
\end{cases}
\]
be a realization of $\A$ in $\Lh$.
Then $\linA$ is closable.
\end{theorem1}

\subparagraph{} %- description of the domain of definition of the closure
Note that the domain of definition $\domain{\linA}$
of the operator $\linA$ considered in theorem \ref{XII_L2_CLOSABILITY_THEOREM}
is not restricted to any boundary conditions.
It should also be mentioned that
rather the closability of the differential operator $\linA$
with domain of definition $\domain{\linA}$
like in theorem \ref{XII_L2_CLOSABILITY_THEOREM} is ensured,
it is not an easy task to give an exact description
for the domain of definition of its closure.
But if the coefficients of $\A$ are constant,
then the set $\domain{\overline{\linA} }$ can be determined
and coincides with the Hilbert-Sobolev space $H^k(\G)$.

\subparagraph{} %- minimal and maximal operators
The closure $\linA_{min} := \overline{\linA_0}$
of a closable linear differential operator $\linA_0$
with domain of definition $\domain{\linA_0} \eq \TF$
is denoted the \textit{minimal realization}
or the \textit{minimal operator} of $\A$.
The \textit{maximal operator} $\linA_{max}$ of $\A$,
however,
is by definition the adjoint of the minimal operator of $\A^{\dag}$
with domain of definition $\domain{{\linA^{\dag}}_0} \eq \TF$,
that is,
$\linA_{max} := ({\linA^{\dag}}_{min})^*$.
One always has $\linA_{min} \, {\prec} \, \linA_{max}$,
and often any operator $\linA$ fulfilling
\[
\linA_{min} \prec \linA \prec \linA_{max}
\]
is called a \textit{realization} of $\A$,
in opposite to our selected definition of this concept.
Using this definition,
we also want to consider the linear operators $\linA_{min}$ and $\linA_{max}$
as realizations of $\A$,
too.

\paragraph{} %- adjointness theorems
To the end,
the question arises whether any realization
of the formally adjoint expression $\A^{\dag}$ is in relationship
with the adjoint of some realization of $\A$.
This problem is examined in \cite{XII_Br2} (section\,3: Adjoints)
even in the context of $L^p$-spaces.
\textsc{Browder}'s result for Hilbert spaces $\Lh$
and the Dirichlet realization $\linA_D$
of a regularly elliptic differential expression $\A$ sounds as follows
(corresponding to theorem 5, see p.\,57):

\begin{theorem1}
Let the $2m$-order expression $\A$ be strongly elliptic
with sufficiently smooth coefficients
and defined on a smooth domain $\G \xsubset \Rd$
such that the formal adjoint $\A^{\dag}$ exists and its coefficients
are uniformly bounded in $\G$.

\subparagraph{} %- adjointness theorem in L^2
If $\linA$ and $\linA^{\dag}$ are the realizations of $\A$ and $\A^{\dag}$
in $\Lh$ with homo\-geneous Dirichlet boundary conditions
and domains of definition
\[
\domain{\linA} \= \domain{{\linA}^{\dag}} \eqdef H^{2m}(\G) \cap \HN{m}{\G},
\]
then $\linA^{\dag}$ is the Hilbert adjoint of ${\linA}$,
that is,
${\linA}^* \eq \linA^{\dag}$.
\end{theorem1}

\subparagraph{} %- adjointness theorem in L^p
A similar result holds in the setting of the spaces $\Lp$ for $1 < p < \infty$.
As it is quoted in \cite{XII_Pa}, lemma 7.3.4,
we have 
${\linA_p}^* \eq {{\linA}^{\dag}}_q$,
where $q \eq p/(p-1)$ respectively $\frac{1}{p} + \frac{1}{q} \eq 1$
and both operators
\[
\begin{cases}[1.2]
\quad \domain{\linA_p} \eqdef W^{2m,p}(\G) \cap \WN{m,p}{\G},\\
\quad \linA_p u \eqdef \A\,u \for u \in \domain{\linA_p},\\
\end{cases}
\]
\[
\begin{cases}[1.2]
\quad \domain{{\linA^{\dag}}_q} \eqdef W^{2m,q}(\G) \cap \WN{m,q}{\G},\\
\quad {\linA^{\dag}}_q u \eqdef \A^{\dag} u
      \for u \in \domain{{\linA^{\dag}}_q}.\\
\end{cases}
\]
are densely defined and closed.

\section{{\GARDING}'s Inequality}
%------------------------------------------------------------------------------%
%-- Gardings's inequality and quasi-coercive Dirichlet forms
We next want to illustrate the usage of abstract sesquilinear
forms for the analysis of linear equations of the form $\A \, u \eq f$,
where $\A$ is a linear partial differential operator of order $2m$
and $u,f$ are real or complex-valued mappings
defined on a domain $\G \xsubset \Rd$.

\subparagraph{} %- differential expressions in divergence form
Preliminary,
a linear differential expression $\A$ of order $2m$
is said to be in \textit{divergence form},
if there are functions $a_{\alpha \beta} \xeof C^{\abs{\alpha} }(\G)$
for $\abs{\alpha} , \abs{\beta} \leq m$
such that $\A$ can be written as
\begin{equation}\label{XII_LPDO_DIV}
\A \eqdef
\sum_{\abs{\alpha} \leq m, \abs{\beta} \leq m} (-1)^{\abs{\alpha} } \,
\PA(\,a_{\alpha \beta}(x) \, D^{\beta}).
\end{equation}

In opposite to this definition let us call an arbitary differential expression
given by \eqref{XII_ANY_ORDER_LPDO} to be \textit{in non-divergence form}.

\subparagraph{} %- relationship between divergence and non-divergence forms
It is a matter of fact
that every expression $\A$ in divergence form
can be transformed into an expression $\A'$ in \textit{non-divergence form}.
Reversely,
a non-divergence form expression $\A'$ given by \eqref{XII_ANY_ORDER_LPDO}
can be transformed into an expression $\A$ in divergence form \eqref{XII_LPDO_DIV},
if $c_{\alpha} \xeof C^{2m - \abs{\alpha} }(\G)$ for $\abs{\alpha} \leq 2m$,
that is,
if the coefficients are sufficiently many differentiable in $\G$.
Comparing the expressions \eqref{XII_ANY_ORDER_LPDO} and \eqref{XII_LPDO_DIV},
it turns out that
$\A'$ is strongly elliptic in $\G$ \iff $\A$ is so.

\subparagraph{}
To show the purpose of a linear differential operator being in divergence form,
let the space $\TF$ of {test functions} be furnished with the inner product
\[
\scalar{\phi}{\psi} \eqdef
\int_{\,\G} \phi \, \overline{\psi} \,dx \qfa \phi,\psi \xeof \TF.
\]
Now,
applying $\A$ given in divergence form to $\phi \xeof \TF$,
then multiplying $\A \, \phi$ with $\overline{\psi} \xeof \TF$ and
finally performing $m$-times partial integration leads to
\[
\DF_{\A}(\phi,\psi) \eqdef \scalar{\A \, \phi}{\psi} \=
\sum_{\abs{\alpha} \leq m, \abs{\beta} \leq m} \,
\int_{\,\G}
a_{\alpha \beta} \; \PA \phi \; D^{\beta} \overline{\psi}
\, dx,
\]
since all the boundary terms vanish.
The mapping
\[
\DF_{\A} \eqdef \TF \xtimes \TF \xto \C
\]
is a sesquilinear form,
called the \textit{Dirichlet form} of $\A$ in $\G$.
If the coefficent functions $a_{\alpha \beta}$ of $\A$
are measurable and essentially bounded,
then the mapping $\DF_{\A}$ is continuous
and thus can continuously be extended to the product space
$\HN{m}{\G} \xtimes \HN{m}{\G}$.

\paragraph{} %- Garding's inequality
We are mainly interested in a $2m$-order linear differential operator $\A$
given in divergence form,
which together with the underlying domain $\G$ fulfills the already 
mentioned conditions (I), (I'), (II), (III) and (III').
\textsc{{\GARDING}}\textit{'s inequality} can be seen as a bridge
from the theory of linear-elliptic differential operators
to abstract functional analysis:

\begin{theorem1}\label{XII_GARDING_INEQUALITY_THEOREM}
Let the domain $\G$ fulfill (I),\,(I') 
and let the linear differential expression $\A$ of order $2m$ fulfill
(II),\,\,(III),\,(III').\\
Then there is $\lambda_0 \xeof \R$ and $a_0 {\,>\,} 0$ such that
for all $u \xeof \HN{m}{\G}$
\begin{equation}\label{XII_GARDING_INEQUALITY}
\DF_{\A}(u,u) + \lambda_0 \norm{u}^2 =
\scalar{\,(\A+\lambda_0)\,u}{u} \geq a_0 \, {\abs{u} _m}^2.
\end{equation}
\end{theorem1}

\subparagraph{}
In other words,
theorem \ref{XII_GARDING_INEQUALITY_THEOREM} quotes that
the Dirichlet form associated with the operator $\A + \lambda_0$ is coercive,
hence the {Lax-Milgram} theorem may be applied.
It should be clear that this assertion also holds
for all $\lambda {\,>\,} \lambda_0$.

\paragraph{} %- Garding's inequality for second order operators
The most important class of linear partial differential operators
in divergence form {\wrt} applications is that of \textit{second order},
that is,
our operator $\A$ is of the form
\begin{equation}\label{XII_GENERAL_SECOND_ORDER_DIFFERENTIAL_OPERATOR}
\A(x,\cdot) = 
-\sum_{i=1}^d \,
\partial_i \, (\sum_{j=1}^d a_{ij}(x)) 
+
\sum_{j=1}^d \, (b_j(x) \, \partial_j)
+
c(x),
\end{equation}
which is nothing but \eqref{XII_LPDO_DIV} for $m \eq 1$.
The Dirichlet form $\DF_{\A}$ belonging to $\A$ is the sesquilinear form
\[
\DF_{\A}(u,v) \= \int_{\,\G} 
(\sum_{i,j=1}^d a_{ij} \, \partial_i u \, \partial_j \overline{v} \+ 
 \sum_{j=1}^d b_j \, \partial_j u \, \overline{v} \+
 c \,u \, \overline{v}
)
\, dx.
\]

\subparagraph{} %- coerciveness of the Dirichlet form of second-order operators
A short computation shows
that if the function $c$ is large enough in comparison
with the sum of partial derivatives of the functions $b_j$,
then $\DF_{\A}$ is coercive.
More precisely,
$\DF_{\A}$ proves to be coercive if
\begin{equation}\label{XII_COERCIVITY_CONDITION}
c(x) \, \geq \, \frac{1}{2} \,
\sum_{k=1}^d \, \frac{\partial b_k}{\partial x_k}(x)
\end{equation}
holds in all points $x \xeof \G$.

\subparagraph{}
On the other hand,
if the matrix $A(x) \eq (a_{ij}(x))_{i,j=1}^d$ of second-order coefficients
in \eqref{XII_GENERAL_SECOND_ORDER_DIFFERENTIAL_OPERATOR}
fails to be positive definite for all $x \xeof \G'$,
where $\G' \xsubset \G$ has positive measure,
then $\DF_{\A}$ cannot be coercive on $\HN{1}{\G}$.
Also,
if $A(x)$ is positive definite in all $x \xeof \G$
and condition \eqref{XII_COERCIVITY_CONDITION} does not hold.
then the functions $b_j$ disturb the coercivity of $\DF_{\A}$.
But after all,
the form $\DF_{\A}$ proves at least to be quasi-coercive,
which is nothing but the assertion of \textsc{\GARDING}'s inequality.

\subparagraph{} %- Garding's inequality for expressions of second order
\textsc{\GARDING}'s inequality remains even true
for the Dirichlet forms of second-order differential operators,
if the basic conditions (I)' and (III') are dropped:

\begin{theorem1}\label{XII_GARDING_INEQUALITY_THEOREM_II}
Let $\G \xsubset \Rd$ fulfill (I) and
the differential operator $\A$ of second order and given in divergence form
fulfill (II) and (III).
Then there is $\lambda_0 \xeof \R$ and $a_0 {\,>\,} 0$ such that
\[
\DF_{\A}(u,u) \+ \lambda_0 \norm{u}^2
\=
\scalar{\,(\A {+} \lambda_0)\,u}{u}
\, \geq \,
a_0 \, {\abs{u} _1}^2
\]
for all $u \xeof \HN{1}{\G}$,
where $a_0$ is the uniform lower bound of the eigenvalues
of the matrices $A(x)$ with $x \xeof \G$.
\end{theorem1}

\section{Weak Solutions of Elliptic Boundary Value Problems}
%------------------------------------------------------------------------------%
%-- weak solutions of ellptic boundary value problems
The variational approach to stationary and evolutional boundary value problems
provides the concept of weak solution,
and in this section we want to apply the abstract theory
boundary value problems governed by strictly elliptic differential operators
by quoting the related existence and uniqueness assertions.

\subparagraph{} %- the variational Dirichlet problem
We basically assume that

\begin{itemize}[leftmargin=15mm]
\item[(DP-1)]
$\A$ is a the differential operator of order $2m$,
strongly elliptic in a \textit{bounded} domain $\G \xsubset \Rd$
and given in divergence form;

\item[(DP-2)]
the coefficients $a_{\alpha \beta}$ belong to $L^{\infty}(\G)$
and the top-order coefficients are uniformly continuous in $\clG$.
\end{itemize}

\paragraph{} %- homogeneous Dirichlet problem
It is our first aim to solve the \textit{homogeneous Dirichlet problem}
\[
\begin{cases}[1.2]
\; (\A\,u)(x)\= f(x)& (x \xeof \G),\\
\; u(x) =
\partial u/ \partial \nu (x) =
\ldots = \partial u^{m{-}1} / \partial \nu^{m{-}1}(x) = 0 & (x \xeof \bndG).\\
\end{cases}
\]
Let $\vecV \eq \HN{m}{\G}$ and $\hilH \eq \Lh$,
hence the pair $(\vecV,\hilH)$ is a Hilbert triple.
By \textsc{{\GARDING}}'s inequality,
the Dirichlet form
\[
\DF_{\A}(u,v) \eqdef \scalar{\A \, u}{v} \quad \quad (u,v \in \vecV)
\]
is quasi-coercive with some $\lambda_0 \xeof \R$.
Then,
by means of the {Lax-Milgram} theorem,
we obtain the following \textit{well-posedness theorem}:

\begin{theorem1}\label{XII_DIRICHLET_EXISTENCE_UNIQUENESS_THEOREM}
For $\lambda \geq \lambda_0$, the variational equation
\[
\DF_{\A}(u,v) \+ \lambda \scalar{u}{v}
\= f(v) \quad (v \xeof \HN{m}{\G})
\]
owns for $f \xeof H^{-m}(\G) := \HN{m}{\G}^*$
a unique solution $u \xeof \HN{m}{\G}$ depending
continuously on the right-hand side $f$ and
called the \textit{weak solution} of the homogeneous Dirichlet problem.
\end{theorem1}

\section{The Fredholm Property of Elliptic Operators}
%------------------------------------------------------------------------------%
%-- the Fredholm property of elliptic operators
This statement is properly not correct,
since it only proves to be true for properly elliptic differential operators
defined in a \textit{bounded} domain $\G \xsubset \Rd$
and endowed with certain boundary conditions,
for example those given by a homogeneous Dirichlet system.

\subparagraph{} %- Rellich's embedding theorem
We first want to present \textsc{Rellich}\textit{'s embedding theorem},
which holds for Sobolev spaces whose underlying domain is \textit{bounded}.
It can be seen as a key theorem in the analysis
of linear-elliptic boundary problems on such domains:

\begin{theorem1}\label{XII_RELLICH_EMBEDDING_THEOREM}
If $\G \xsubset \Rd$ is bounded and owns smooth boundary $\bndG$,
then the continuous embeddings
$\iota_k \: \HN{k}{\G} \hookrightarrow \Lh$ are compact for all $k \xeof \N$.
\end{theorem1}

%\subparagraph{}
%A proof of this theorem can be found for example in \ldots.

\subparagraph{} %- elliptic operators as continuous mappings
The basic approach to investigate the Fredholm property
of a properly elliptic operator $\A$ in divergence form
and of order $2m$
consists of two small changes for the Dirichlet realization $\linA_D$
having domain of definition $\domain{\linA} \eq H^{2m}(\G) \cap \HN{m}{\G}$:
let us consider the linear operator $\linA_D$ as

\begin{itemize}[leftmargin=9mm]
\item[(F-1)]
a mapping $\linA_D \: \HN{m}{\G} \xto \Lh$,
which is continuous due to the essential boundedness
of the coefficients $a_{\alpha,\beta}$ in \eqref{XII_LPDO_DIV};

\item[(F-2)]
the sum $\linA_1 + \linA_2$ of two further linear operators,
where we shall assume that
$\linA_1 \: \HN{m}{\G} \xto \Lh$ is the main part 
having weak derivatives of order $2m$ only,
whereas $\linA_2 \: \HN{m}{\G} \xto \Lh$ is the part of $\linA_D$
consisting of weak derivatives with order less than $2m$.
\end{itemize}

\subparagraph{}
Let $\lambda_0$ be the smallest parameter
for the Dirichlet form $\DF_{\linA+\lambda_0}$ to be coercive.
Since
\[
\linA_D \= (\linA_1 + \lambda_0) \+ (\linA_2 - \lambda_0)
\]
and $\linA_2- \lambda_0$ may be seen as a compact perturbation
of $\linA_1+\lambda_0$,
which maps $H^m(\G)$ one-to-one onto $\Lh$ and consequently is Fredholm 
with vanishing index,
the operator $\linA_D$ is Fredholm again with index $\chi(\linA_D) \eq 0$.

\section{Spectral Theory of Strongly Elliptic Operators}
%------------------------------------------------------------------------------%
%-- spectral theory of strongly elliptic operators
Let us consider again a strongly elliptic differential operator $\A$
given in divergence form and its Dirichlet realization $\linA_D$
in a bounded domain $\G \xsubset \Rd$.
By {\GARDING}'s inequality,
the Dirichlet form $\DF_{\A}$ is quasi-coercive
{\wrt} the parameter $\lambda_0$,
hence $\linA_D + \lambda$ proves to be bijective
for all $\re \lambda {\,>\,} \lambda_0$.
Furthermore,
the homogeneous Dirichlet problem in $\Lh$
\begin{equation}\label{XII_HOMOGENEOUS_DIRICHLET_PROBLEM}
\linA_D \, u  \+ \lambda u \= f \quad \quad (f \in \Lh)
\end{equation}
is well-posed,
which is a direct consequence of
theorem \ref{XII_DIRICHLET_EXISTENCE_UNIQUENESS_THEOREM}.

\subparagraph{}
In the special case,
where $f \xeof \Lh$,
the solutions $u$ of \eqref{XII_HOMOGENEOUS_DIRICHLET_PROBLEM}
belong to the Hilbert-Sobolev space $H^{2m}(\G)$ and are called \textit{strong}.
It should be clear that 
each strong solution of \eqref{XII_HOMOGENEOUS_DIRICHLET_PROBLEM}
is also a weak solution.
But in order to show that these solutions $u$ are even classical,
that is,
$u \xeof C^{2m}(\G) \cap C^m(\clG)$,
methods of \textit{regularity theory} for elliptic problems have to be applied.

\paragraph{} %- spectrum of the Dirichlet operator
Let us turn to the spectrum of $\linA_D$.
Since $-\Delta_D + \lambda$ forms a bijective and closed operator
between the spaces $H^{2m}(\G) \cap \HN{m}{\G}$ and $\Lh$
for each $\lambda \xeof C$ with $\re \lambda \geq \lambda_0$,
we surely have
\[
\{ \zeta \in \C: \re \zeta \geq \lambda_0\} \xsubset \rho(\linA_D)
\and
\sigma(\linA_D) \xsubset \{ \zeta \in \C: \re \zeta < \lambda_0\}.
\]
Furthermore,
if $\G$ is bounded,
then $\linA_D$ is an operator with pure point spectrum.

\paragraph{} %- sectorialness of the Dirichlet operator
We can also make assertions on the accretivity and sectorialness
of the operator $\linA_D$.
As theorem \eqref{XI_QUASI-COERCIVE_FORMS_AND_ASSOCIATED_OPERATORS}
(d),\,(e) shows,
$\linA_D + \lambda_0$ is the generator
of a semigroup of bounded linear operators
that is strongly continuous and analytic at the same time.
Hence the initial value problems for the heat and wave equation
governed by the operator $\linA_D + \lambda_0$ is well-posed.
By perturbation theory of semigroups,
this is also true for both problems governed by the operator $\linA_D$ itself.

\paragraph{} %- Gauss-Weierstrass semigroup
Let us finally consider the heat equation
for the negative Laplacian $-\Delta_D$ in the entire space $\Rd$,
which of course is a free space problem.
Theorem \ref{X_SEMIGROUP_CREATED_BY_SECTORIAL_OPERATOR_IS_ANALYTIC} ensures
that the realization $-\Delta_D$ in Hilbert space $L^2(\Rd)$
with domain of definition
\[
\domain{-\Delta_D} \= H^2(\Rd)
\]
is the infinitesimal generator
of a sectorially contractive analytic semigroup $\{W(t)\}_{t \geq 0}$,
denoted the \textit{Gauss-Weierstrass} \textit{semigroup}.
Hence there is a uniquely defined solution $u \xeof C^1(\J \times \Rd,\R)$
for the homogeneous Cauchy problem
\[
\begin{cases}[1.2]
\quad u_t - \Delta u \= 0,\\
\quad u(x,0)         \= u_0(x) \for x \in \Rd \\
\end{cases}
\]
with initial mapping $u_0 \xeof \HN{2}{\Rd}$.
Then one may represent the semigroup $\{W(t)\}_{t \geq 0}$
explicitly by means of the \textit{heat kernel} function $\Phi_d$,
that is,
for all $t > 0$ and $x \xeof \Rd$
\[
(W(t) \, u_0) \, (x)
=
\frac{1}{(4 \pi \, t)^{d/2}}
\int_{\,\Rd} e^{- \frac{\abs{x-y} ^2}{4t}} u_0(y) \, dy
=
(\Phi_d(t) \ast u_0)\,(x).
\]

\section{Self-Adjoint Extensions of Half-Bounded Operators}
%------------------------------------------------------------------------------%
%-- self-adjoint extensions of half-bounded operators
Remember the definition of the differential expression $\DO^{\dag}$.
We then denote a linear differential expression $\DO$
\textit{formally self-adjoint} in case of $\DO^{\dag} \eq \DO$.
If $\DO$ has constant coefficients,
then $\A$ proves to be formally self-adjoint
iff
$c_{\alpha} \xeof \R$ for $\abs{\alpha} $ even
and $c_{\alpha} \xeof i \, \R$ for $\abs{\alpha} $ odd.

\paragraph{} %- examples of formally self-adjoint differential expressions
In the following let us collect a couple of
formally self-adjoint linear differential expressions
and describe some self-adjoint realizations of them:

\begin{itemize}[leftmargin=17pt]
\item[(a)] \textit{Sturm-Liouville operators.}
Our first example of a formally self-adjoint linear differential operator
(given in divergence form) is the \textit{Sturm-Liouville} \textit{operator}
\[
\A_{sl}(u,x) \eqdef -(p(x)\, u'(x))' \+ q(x)\, u(x),
\quad \quad x \xeof (a,b)
\]
with $p(x) {\,>\,} 0$ and $q(x) \geq 0$ for $x \xeof [a,b]$,
$-\infty {\,<\,} a {\,<\,} b {\,<\,} \infty$.
This operator determines the elongations of a vibrating string
being fixed at the endpoints of $[a,b]$.
Then the corresponding Dirichlet form reads as
\[
\DF_{\A_{sl}}(u,v) \=
\int_a^b \,
(p \, u'\, \overline{v}' \+ q \, u \, \overline{v}) \, dx,
\]
which shows that the expression $\A_{sl}$ is symmetric and half-bounded
on the domain of definition $\TF$.
The energetic space $H_{\A_{sl}}$ is the Hilbert-Sobolev space $\HN{1}{a,b}$
containing the restrictions of 
\textit{absolutely continuous functions} $u\: [a,b] \xto \R$
with boundary conditions $u(a) \eq u(b) \eq 0$ to the interval $(a,b)$.

\subparagraph{}
The self-adjoint extension $\linA_{sl,F}$
of the {Sturm-Liouville} operator in $L^2(a,b)$ has domain of definition
\[
\domain{\linA_{sl,F}} = H^2(a,b) \, \cap \, \HN{1}{a,b}.
\]
The operator $\linA_{sl,F}$ has pure point spectrum,
since the interval $(a,b)$ is assumed to be bounded.
Especially in case of $p(x) \eq 1$ and $q(x) \eq 0$ for $x \xeof [a,b]$, 
the eigen\-values and eigen\-functions
of the operator $\A_{sl,F} \eq - {d^2} / {dx^2}$ can exactly be determined.
In fact,
the eigenvalue distribution of $\A_{sl,F}$ is given by
\[
\lambda_n \= \frac{n^2 \pi^2}{(b-a)^2} \quad \quad (n \xeof \N),
\]
so the asymptotical behaviour of the eigen\-values is $\lambda_n \sim n^2$.
Each eigen\-value $\lambda_n$ owns a one-dimensional eigen\-space,
and the eigen\-function $u_n$ of $\lambda_n$ is given by
\[
u_n(x) \eqdef \sin \, (\frac{n \pi x}{b-a}) \quad \for x \in [a,b].
\]
Finally,
the sequence of the eigen\-functions $u_n$
constitutes after suitable renormations an ortho\-normal base in $L^2(a,b)$.

\item[(b)] \textit{(Negative) Laplacian in various domains $\G \xsubset \Rd$.}
In order to apply the theory of self-adjoint operators
to the investigation of the negative Laplacian,
we start with the realization $\linA_{-\Delta,0}$ of the expression $-\Delta$
in the space $\TF$ of test functions.

\subparagraph{} %- Green's formula for the negative Laplacian
Using \textsc{Green}'s second formula,
and due to the boundary condition $\phi(x) \eq 0$ for $x \xeof \partial \G$
and $\phi \xeof \TF$,
we compute for $\phi,\psi \xeof \TF$
\[
\int_{\G} -\Delta \phi \, \overline{\psi} \, dx
\=
\sum_{i=1}^n \int_{\G}
\frac{\partial^2 \phi}{\partial x_i^2} \, \overline{\psi} \, dx
\]
\[
\=
\sum_{i=1}^n \int_{\G}
u \, \overline{\frac{\partial^2 \phi}{\partial x_i^2}} \, dx 
\=
\int_{\G} \phi \, \overline{\Delta \psi} \, dx,
\]
which implies symmetry of the operator $\linA_{-\Delta,0}$.
Moreover,
we obtain for $\phi \xeof \TF$ the inequality
\[
\scalar{-\Delta \phi}{\phi} \=
\sum_{i=1}^n \int_{\G}
\frac{\partial^2 \phi}{\partial x_i^2} \, \overline{\phi} \, dx
\=
\sum_{i=1}^n \int_{\G} \abs{ \frac{\partial \phi}{\partial x_i}} ^2 \, dx 
\, \geq \, 0,
\]
hence $\linA_{-\Delta,0}$ is half-bounded with lower bound $0$
and fulfills the prerequisites
to apply the {Friedrichs} extension method.
Thus we may quote
that the operator $\linA_{-\Delta,0}$ owns at least one self-adjoint extension,
and this assertion is true actually for arbitrary domains $\G \xsubset \Rd$.

\paragraph{} % realizations of the Laplacian
In the following enumeration we consider various self-adjoint realizations
of the negative Laplacian in a domain $\G$,
all of them being extensions of the operator $\linA_{-\Delta,0}$:

\begin{itemize}[leftmargin=5mm]
\item[(i)]
\textit{(Negative) Laplacian in $\Rd$}.
The analysis of the operator $-\Delta$ in the entire Euclidean space $\Rd$
is easy,
since the differential operator proves to be unitarily equivalent to
the multiplication operator $\linA_m$ generated by the function
\[
m \: (x_1,\ldots,x_n) \mapsto (x_1^2 + \ldots + x_n^2) \= \abs{x} ^2
\quad (x \xeof \Rd).
\]
In fact,
the operator $\linA_m$ is real-valued and self-adjoint,
which is a consequence of
the \textit{Fourier-Plancherel transformation}\footnote{
\,We introduce the \textit{Fourier-Plancherel transformation} in $L^2(\Rd)$
by 
\[
\hat{u}(\xi) \eqdef
\frac{1}{(2 \pi)^{d/2}}
\int_{\,\Rd} \, e^{-2 \pi i \scalar{\xi}{x}} \, u(x) \, dx
\]
for all $u \xeof C_c^{\infty}(\Rd)$,
which forms a partial isometry $\FT \: u \mapsto \hat{u} $
from the space $C_c^{\infty}(\Rd)$ to the space $L^2(\Rd)$.
Then for functions $u,v \xeof L^2(\Rd)$ 
the inner product between $u$ and $v$
and the inner product of their Fourier-Plancherel transforms
$\hat{u}$ and $\hat{v}$ are related by \textsc{Plancherel}\textit{'s formula}
$\scalar{u}{v} \eq \scalar{ \hat{u} }{\hat{v} }$.
Since the mapping $\FT$ is linear and continuous on $C_c^{\infty}(\Rd)$,
which is a dense subspace of $L^2(\Rd)$,
it can be extended to an \textit{unitary},
that is,
a bijective and isometric linear operator
$\FT_P \: L^2(\Rd) \xto L^2(\Rd)$,
called the \textit{Fourier-Plancherel} \textit{operator}
or the \textit{$L^2$-Fourier transformation} in $L^2(\Rd)$.}
.
This implies
that the operator $\linA_{sa}$ with domain of definition
$\domain{\linA_{sa}} := H^2(\Rd)$ must be self-adjoint,
too.
The spectrum of $\linA_{sa}$ is that of the operator $\linA_m$,
that is,
it is purely continuous and $\sigma_c(\linA_{sa}) \eq [0,\infty)$.
Moreover,
since $\linA_{sa}$ is the closure of the operator $\linA_{-\Delta,0}$,
the latter proves to be essentially self-adjoint.

\item[(ii)]
\textit{Dirichlet-Laplacian}.
Let the domain $\G \xsubset \Rd$ be \textit{bounded}
and consider $-\Delta$ with domain of definition $\domain{-\Delta} := \TF$.
The related Dirichlet form $\DF_{-\Delta}$ is given by 
\[
\DF_{-\Delta}(\phi,\psi) \eqdef
\int_{\G} \nabla \phi \, \nabla \psi \, dx
\quad \quad
(\phi,\psi \xeof \TF),
\]
which can continuously be extended to the sesquilinear form
$\overline{\DF_{-\Delta}}$ with domain of definition
$\domain{\overline{\DF_{-\Delta}}} := \HN{1}{\G} \xtimes \HN{1}{\G}$.

\subparagraph{}
Due to theorem \ref{XI_QUASI-COERCIVE_FORMS_AND_ASSOCIATED_OPERATORS},
the operator $\linA_D$ associated with the form $\DF_{-\Delta}$
and defined by
\[
\begin{cases}[1.2]
\quad \domain{\linA_D} \eqdef
\{u \xeof \HN{1}{\G} \: \Delta u \xeof L^2(\G)\} \\
\quad \linA_D u \eqdef - \Delta u \; \for u \xeof \domain{\linA_D} \\
\end{cases}
\] 
is self-adjoint,
called the \textit{negative Laplacian with Dirichlet boundary conditions}
or briefly the \textit{Dirichlet-Laplacian}.

\subparagraph{}
We obviously have
$\HN{1}{\G} \cap H^2(\G) \xsubset \domain{\linA_D}$,
and one can prove that
if $\G$ is bounded and owns a $C^2$-boundary or $\G \eq \R_+^n$,
then the reversed inclusion also holds,
hence $\domain{\linA_D} \eq \HN{1}{\G} \cap H^2(\G)$.
This equality is even true in the case,
where $\G$ is bounded and convex.

\subparagraph{}
Furthermore,
if $\G$ has the $1$-extension property,
then the space $\HN{1}{\G}$ is compactly embedded into the space $\Lh$
and $\linA_D$ is an operator with pure point spectrum.
In this case it is an important property of the Dirichlet-Laplacian
that the number $0$ does \textit{not} belong to the spectrum $\sigma(\linA_D)$,
so $\linA_D$ is invertible and the operator $\linA_D^{-1}$ proves to be bounded.

\item[(iii)]
\textit{Neumann-Laplacian}.
Let $\G \xsubset \Rd$ be a bounded domain with $C^1$-boundary $\partial \G$
and let the realization $\linA_N$ of the negative Laplacian in $\Lh$
be defined by
\[
\begin{cases}[1.2]
\quad \domain{\linA_N} \eqdef
\{ u \xeof H^2(\G) \: u_v(x)= 0, \; x \xeof \partial \G \} \\
\quad \linA_N u \eqdef - \Delta u \; \for u \xeof \domain{\linA_N}. \\
\end{cases}
\]
The realization $\linA_N$ is said to be the
\textit{(negative) Laplacian with Neumann boundary condition}
or briefly the \textit{Neumann-Laplacian}.
It owns $\lambda \eq 0$ as the minimal eigen\-value,
and all constant functions in $\domain{\linA_N}$ are eigen\-functions
of $\lambda \eq 0$.
\end{itemize}

\item[(c)] \textit{{\SCHROEDINGER} operators.}
By the \textit{{\SCHROEDINGER} operator} is meant
the linear partial differential expression given by
\[
\HAM_q(u,x) \eqdef (-\Delta u)(x) \+ q(x),
\]
defined in the entire Euclidean space $\Rd$.
In the special case,
where $q \eq 0$,
self-adjointness of the operator $-\Delta$ acting in $L^2(\Rd)$
and with domain of definition $H^2(\Rd)$ is well-known.

\subparagraph{} % potential functions and multiplication operators
The so-called \textit{potential function} $q \: \Rd \xto \R$
gives rise to a multiplication operator $M_q$ acting in $L^2(\Rd)$,
that is,
\[
\begin{cases}[1.2]
\quad \domain{M_q} \eqdef \{u \in L^2(\Rd) \: q {\cdot} u \in L^2(\Rd)\},\\
\quad M_q u        \eqdef q {\cdot} u \; \for u \in \domain{M_q}   ,\\
\end{cases}
\]
so the operator $\HAM_q$ can be seen as a perturbation
of the negative Laplacian by the operator $M_q$.

\paragraph{} %- Schroedinger equation
By the time-dependent \textit{{\SCHROEDINGER} equation} is meant
the initial value problem
\[
i \, \dot{u} \= -\Delta u \+ M_q u
\]
with prescribed initial value $u(0) \eq u_0 \xeof H^2(\Rd)$,
where it is supposed that $\HAM_q \eq -\Delta + M_q$
is a self-adjoint operator in $L^2(\Rd)$
with domain of definition $\domain{\HAM_q} \eq H^2(\Rd)$.
The existence and uniqueness of a solution for this problem is ensured
by the functional calculus for self-adjoint operators
using the spectral resolution of the operator $\HAM_q$
and the function $f:\lambda \xeof \R \mapsto e^{it \lambda}$.
The solution $u$ of the problem is then provided
by a group of unitary operators $\{U(t)\}_{t \in \R}$,
which implies $u(t) \eq U(t) u_0$ for every $t \xeof \R$
(see also section IV.8 for the treatment
of the abstract {\SCHROEDINGER} equation).

\paragraph{} %- spectral analysis of Schroedinger operators
The study of \textit{time-independent solutions}
of the {\SCHROEDINGER} equation leads to the eigenvalue problem
\begin{equation}\label{XII_SCHROEDINGER_EIGENVALUE_PROBLEM}
\HAM_q u \eqdef -\Delta u + M_q u \= \lambda u,
\end{equation}
which means to find those real numbers $\lambda$,
for which there are non-trivial solutions
of \eqref{XII_SCHROEDINGER_EIGENVALUE_PROBLEM}.
The discrete spectrum $\sigma_d(\HAM_q)$ of the operator $\HAM_q$
is in fact of great interest,
since the eigen\-functions of eigen\-values $\lambda \xeof \sigma_d(\HAM_q)$
describe the stationary states of a quantum mechanical system.

\subparagraph{} %- spectrum of Schroedinger operators
The presence of a potential function $q$
can fairly change the spectrum of a {\SCHROEDINGER} operator 
like the example of the \textit{harmonic oscillator} shows:
the spectrum of the negative Laplacian is the connected set $[0,\infty)$,
but the spectrum of the operator $-\Delta {\,+\,} M_q$ 
with $q \: x \mapsto \frac{1}{2}x^2$ consists of discrete points only.

\subparagraph{} %- lower and upper bounds for the discrete spectrum
It is a further and interesting problem in quantum mechanics
to determine lower and upper bounds for the discrete spectrum
of the operator $\HAM_q$.
In the following we want to present three assertions,
which contain information on the bounds of the set $\sigma_d(\HAM_q)$,
namely
the \textit{Birman-Schwinger} \textit{bound}, 
the \textit{Lieb-Cwikel-Rozenblum} \textit{bound} and
    \textsc{Kato}\textit{'s theorem}:

\begin{itemize}[leftmargin=4mm]
\item[1.]
The \textit{Birman-Schwinger} \textit{bound.}
This assertion presents an upper bound for the set $\sigma_d(\HAM_q)$,
which is valid for certain potentials $q$ in case of $d \eq 3$.
Let $q \xeof L^{\infty}(\R^3)$ and 
\[
B_q \eqdef
\int_{\,\R^3} \int_{\,\R^3} \frac{q(x) \, q(y)}{\abs{x-y} ^2} \ dx \, dy
\, < \, + \infty.
\]
Then the total number $N(\HAM_q)$ of eigen\-values of the operator $\HAM_q$
can be estimated according to
\[
N(\HAM_q) \, \leq \, \frac{1}{16 \pi^2} \, B_q.
\]

\item[2.]
The \textit{Lieb-Cwikel-Rozenblum} \textit{theorem}
quotes an estimate on the number $N_(\HAM_q)$
of negative eigen\-values of the operator $\HAM_q$:
if $d > 3$ and $q \xeof L_{loc}^{\infty}(\Rd)$,
then
\[
N(\HAM_q) \, \leq \, c_d \int_{\,\Rd} \min \{q(x),0\}^{n/2} \, dx,
\]
where $c_d$ is a constant only depending on the dimension $d$.

\item[3.]
Our last result is often denoted as \textit{Kato}\textit{'s theorem}.
It does not make assertions on bounds for the discrete spectrum,
but quotes the absence of positive eigen\-values
under mild assumptions on the function $q$:
if $q \xeof L_{loc}^{\infty}(\Rd)$ and
\[
\lim_{\abs{x} \to \infty} \, \abs{x} \, q(x) \= 0,
\]
then the operator $\HAM_q$ does not own any positive eigen\-value.
\end{itemize}
\end{itemize}

\section{Eigenvalue Problems in Bounded Domains}
%------------------------------------------------------------------------------%
We now turn to the eigenvalue equation $\A \, u  - \lambda u \eq f$
of a strongly elliptic differential operator $\A$ of second order
in a \textit{bounded} domain $\G \subset \Rd$
with homogeneous Dirichlet boundary conditions $u(x) \eq 0$ for $x \xeof \bndG$.

\subparagraph{}
Sine the realization $\linA$ of $\A$ in Hilbert space $\Lh$
is an operator with pure point spectrum,
we may conclude
that the eigenvalue problem is well-posed for every $\lambda \xeof \R$
up to countable many isolated points $(\lambda_k)$
with $\abs{\lambda_k} \xto \infty$ for $k \xto \infty$,
and for these points the space of solutions of the homogeneous equation
$\A \, u \eq \lambda_k u$ is finite-dimensional.

\paragraph{} %- eigenvalue distributions of the Laplacian
The mentioned facts will be illustrated for the negative Laplacian $-\Delta$
in bounded domains and under homogeneous Dirichlet boundary conditions.
Unfortunately,
the eigenvalues and eigenfunctions can exactly be determined only in the case,
where the shape of $\G$ is of elementary type:

\begin{itemize}[leftmargin=12pt,itemsep=5pt]
\item[1.]
If $\G$ is a finite interval $(a,b) \subset \R$,
then the homogeneous boundary conditions are $u(a) \eq u(b) \eq 0$.
In this case the eigenvalue distribution is given by
\[
\lambda_n \= \frac{n^2 \pi^2}{(b-a)^2} \quad \quad (n \in \N),
\]
and the related eigenfunctions are the mappings
\[
u_n(x) \eqdef \sin \, (\frac{n \pi x}{b-a}) \; \for x \in (a,b)\,
\]
forming an orthonormal base in $L^2(a,b)$.
Thus,
the smallest eigenvalue is the number $\pi^2/(b-a)^2$,
and the asymptotical behaviour of the eigenvalues is $\lambda_n \sim n^2$.
Notice that for each eigenvalue $\lambda_n$
the eigenspace of $-\Delta-\lambda_n$ is $1$-dimensional.

\subparagraph*{} %- Sturm-Liouville operator
The Sturm-Liouville operator $\cal{L}\, u \eq -(pu')' + qu$
with domain of definition $H^2(a,b) \, \cap \, \HN{1}{a,b}$
also forms an operator with pure point spectrum,
supposed $(a,b)$ is a bounded subset of the real line.
Especially for $p(x) \eq 1$ and $q(x) \eq 0$ for $x \xeof (a,b)$ 
we obtain the negative Laplacian $- d^2 / dx^2$.

\item[2.]
In case of $\G \eq (0,a) \xtimes (0,b)$ we obtain similar results.
In fact,
the separation approach $u(x,y) \eq u_1(x) \cdot u_2(y)$ 
for the eigenvalue problem $-\Delta u \eq \lambda u$ leads to
separated differential equations for $u_1$ and $u_2$ with solution functions
\[
u_{1,n}(x) \= \sin \, (\frac{n \pi x}{a}),
\quad \quad
u_{2,k}(x) \= \sin \, (\frac{k \pi y}{b}),
\]
eigenvalues
$\lambda_{n,k} \eq \frac{n^2\pi^2}{a^2} + \frac{k^2\pi^2}{b^2}$
and related eigenfunctions $u_{1,n} \cdot u_{2,k}$.

\item[3.]
Finally,
the eigenvalue problem for the $2d$-circle reads as
\[
\G \= \{\, (r,\theta): 0 \leq r < 1, 0 \leq \theta \leq 2\pi\,\}.
\]
It is appropriate to transform the eigenvalue equation
into an equation {\wrt} polar coordinates $(r,\theta)$
\[
v_{rr} \, + \,
\frac{1}{r} v_r \, + \,
\frac{1}{r^2} v_{\theta \theta} \, + \,
\lambda v \= 0\,,
\]
and boundary conditions
$v(1,\theta) \eq 0$ for $0 \leq \theta \leq 2 \pi$.
The separation approach $v(r,\theta) \eq R(r) \cdot \Theta(\theta)$ leads
to encoupled differential equations for $R$ and $\Theta$:
\[
\begin{cases}[1.2]
\quad r^2 R''(r) \, + \, rR'(r) \, + \,
(\lambda r^2-\sigma)R(r) \= 0 \,, \\
\quad \Theta''(\theta) \, + \,
\sigma \, \Theta(\theta) \= 0 \,, \\
\end{cases}
\]
where the second one is solved by
$\Theta_n(\theta) \eq a \cos(n \theta) + b \sin(n \theta)$ for $a,b \xeof \R$.
The index $n := \sqrt{\sigma}$ can only be of integer type,
due to the $2 \pi$-periodicity of the function $\Theta$.
Moreover,
the solutions $R_n(r)$ of the first differential equation
are the \textit{Bessel functions}
\[
J_n(\sqrt{\lambda}r) \eqdef
\frac{\lambda^n r^n}{2^n} \,
\sum_{k=0}^{\infty}
\frac{(-1)^k}{k! (n+k)!} \, \frac{\lambda^k r^{2k}}{4^k}.
\]
The eigenvalues of this problem coincide with the squares
of the infinitely many zero points $\mu_{n,1},\mu_{n,2},\ldots$
of the functions $J_n$ ($n \xeof \N_0$)
having the eigenfunctions
\[
v_{n,m}(r,\theta) \=
J_n(\mu_{n,m} \cdot r) \cdot
\{\,a \cos \, (n \theta) \, + \, b \sin \, (n \theta)\,\}.
\]
\end{itemize}

\section{Linear Differential Operators in Distributional Spaces}
%------------------------------------------------------------------------------%
We next study the action of linear partial differential operators 
in distributional spaces.
First recall
that if $\psi \xeof \SF$ and $\disF \xeof \D'(\G)$,
then $\psi \disF$ defined by
\[
(\psi \disF) \phi \, := \, \disF (\psi \phi), \quad \quad \phi \xeof \TF
\]
also belongs to $\D'(\G)$. 
This implies that the linear differential operators $c_{\alpha} \, \PA$ 
with $c_{\alpha} \xeof \SF$ are acting on the space $\D'(\G)$,
since for $\disF \xeof \D'(\G)$
the mapping $c_{\alpha} \PA \disF$ is a distribution again.
By taking linear combinations of such operators,
we conclude that the action
of a linear partial differential operator
\[
\DO(x,u) \, := \,
\sum_{\abs{\alpha} \leq 2m} c_{\alpha}(x) \; \PA u(x)
\]
having smooth coefficients and acting on the space $\D'(\G)$ is well-defined,
which of course is especially true in case of $\G \eq \Rd$.

\paragraph{} % definition of fundamental solution
In the following we introduce the concept of \textit{fundamental solution}
for a linear partial differential operators $\DO$
having smooth coefficients $c_{\alpha}$.
It is nothing but a distribution $\Phi_{\DO} \xeof \D'(\Rd)$
that forms a solution of the differential equation
$\DO \Phi_{\DO} \eq \delta_0$,
where $\delta_0$ is the Dirac distribution {\wrt} the origin $x \eq 0$.

\paragraph{}
It is useful to know a fundamental solution of $\DO$
as the following computation shows:
let $\disF \xeof \D'(\Rd)$ be any Schwartz distribution and consider the
inhomogeneous differential equation $\DO \, \disU \eq \disF$.
If convolution between $\Phi_{\DO}$ and $\disF$ exists,
then the equalities
\[
\DO(\Phi_{\DO} \ast \disF)  \=
(\DO \Phi_{\DO}) \ast \disF \, = \;
\delta_0 \ast \disF         \= \disF
\]
are valid,
so the convolutional product $\Phi_{\DO} \ast \disF$
solves the distributional differential equation $\DO \, \disU \eq \disF$. 
Notice that
there is a distributional solution $\disU$ of this equation
especially for all distributions $\disF$ having compact support.

\paragraph{}
The fundamental solution of a linear differential operator
need not be unique,
which can be showed by the following example:
let $\disU \xeof \D'(\Rd)$ be a solution of the homogeneous equation
$\DO \, \disU \eq 0$,
then the distribution $\Phi_{\DO} + \disU\,$
is a fundamental solution of $\DO$, too.

\paragraph{} %- Malgrange-Ehrenpreis theorem
Regarding the existence of fundamental solutions for linear 
differential operators,
the \textsc{Malgrange-Ehrenpreis} \textit{theorem}
has been shown to be a milestone in PDE theory:

\begin{theorem1}\label{XII_MALGRANGE-EHRENPREIS_THEOREM}
Each linear differential operator $\DO \neq 0$
with constant coefficients
has a fundamental solution in the distributional space $\D'(\Rd)$.
\end{theorem1}

\subparagraph{}
It should be mentioned 
that there are counterexamples which demonstrate
that this theorem is not valid for differential operators
whose coefficients are non-constant.

\paragraph{} %- examples of fundamental solutions
In the following we describe the fundamental solutions
of three linear differential operators,
namely
the \textit{Laplace operator},
the \textit{heat operator} and
the \textit{wave operator}.

\begin{itemize}[leftmargin=12pt,itemsep=5pt]
\item[1.]
The fundamental solution $\Phi_{\Delta}^d$ of the Laplacian $\Delta$
depends on the number of spatial coordinates.
Since $\Delta$ is invariant {\wrt} translations and rotations,
it is nearby to conjecture $\Phi_{\Delta}^d(x) \eq \phi_d(\abs{x} )$.
Indeed,
the functions 
\[
\Phi_{\Delta}(x) \=
\begin{cases}[1.5]
\quad \frac{1}{2} \abs{x}                         & (d=1)      \\
\quad - \frac{1}{2 \pi} \ln \abs{x}               & (d=2)      \\
\quad \frac{1}{(d-2) \, \omega_d } \abs{x} ^{2-d} & (d \geq 3) \\
\end{cases}
\]
are solutions of the differential equations
$\Delta \Phi_{\Delta}^d \eq \delta_0$ for $d \xeof \N$,
where $\omega_d$ is the Lebesgue-Borel measure of
the $d$-dimensional unit sphere.
Each other fundamental solution of the Laplacian is of the form
$\Phi_{\Delta}^d + P_1$,
where $P_1$ is any polynomial of degree less than two.

\item[2.]
The fundamental solution of the \textit{heat operator} 
$\DO := \partial_t - \Delta u$ in $\Rd$ is given 
by the so-called \textit{heat kernel} function
\[
\Phi_d(x,t) \=
\begin{cases}[1.5]
\quad \frac{1}{(4 \pi t)^{d/2}} \cdot e^{- \frac{\abs{x} ^2}{4t}},
         & \; x \xeof \Rd, \; t > 0,   \\
\quad 0, & \; x \xeof \Rd, \; t \leq 0.\\
\end{cases}
\]
Each other fundamental solution of the heat operator
adopts the form $\Phi_d + c$ ($c \xeof \C$),
so the kernel of the heat operator 
purely consists of constant distributions.

\item[3.]
When investigating the fundamental solutions
of the hyperbolic \textit{wave operator} $\Box := \partial_{tt} - \Delta u$,
one has to distinguish between even and odd spatial dimensions.

\subparagraph{}
Let $H$ be the \textit{Heaviside function},
defined by $H(t) := 0$ for $t<0$ and $H(t) := 1$ for $t \geq 0$.
Then we obtain for $d \leq 3$ 
\[
\Phi_{\Box}^d(x,t) \= \;
\begin{cases}[1.5]
\quad \frac{1}{2} H(t-\abs{x} )                                      & (d=1)\\
\quad \frac{1}{2 \pi} \, (t^2 - \abs{x} ^2)^{-1/2} \, H(t- \abs{x} ) & (d=2)\\
\quad \frac{1}{2 \pi} \, \delta(t^2 - \abs{x} ^2) \, H(t)            & (d=3).\\
\end{cases}
\]

\subparagraph{}
It is a notable fact
that all the higher fundamental solutions of the wave operator
can be derived from the functions $\Phi_{\Box}^1$ and $\Phi_{\Box}^2$.
Using the variable $r := \abs{x} $,
the functions $\Phi_{\Box}^{d+2}$ ($d \xeof \N$) are recursively given by
\[
\Phi_{\Box}^{d+2}(t,r) \=
\frac{-1}{2\pi r} \cdot \frac{\partial}{\partial_r} \Phi_{\Box}^d(t,r),
\quad d \geq 1.
\]
To determine the support of a fundamental solution $\Phi_{\Box}^d$,
we first define for fixed $x_0 \xeof \Rd$ the \textit{light cone}
\[
C^+(x_0) \, := \, \{(t,x): t > 0, \abs{x-x_0} < t \}
\]
and obtain the support of the fundamental solution according to
\[
\supp{\Phi_{\Box}^d} \=
\begin{cases}[1.5]
\quad \partial C^+(0),   \quad     d \mbox{ odd  or } d \geq 3,\\
\quad \overline{C^+(0)}, \quad\;\; d \mbox{ even or } d=1.\\
\end{cases}
\]
\end{itemize}

\paragraph{} % hypoelliptic differential operators
A linear differential operator $\DO$ is called \textit{hypoelliptic},
if {\wrt} the linear differential equation $\DO u \eq f$
the condition $f \xeof \SF$ implies $u \xeof \SF$.
Restricting to differential operators with constant coefficients,
hypoellipticity can be characterized
by means of a feature of the fundamental solution $\Phi_{\DO_c}$:
$\DO_c$ is hypoelliptic
\iff
$\Phi_{\DO_c} \eq \disF_f$,
where $\disF_f$ is a regular distribution 
with generating function $f$ being indefinitely differentiable
in every $x \neq 0$.

\subparagraph{}
We next introduce the concept of \textit{singular support} $\Sigma(\disF)$
for $\disF \xeof \D'(\Rd)$:
it is the smallest subset $S \xeof \Rd$
such that $\disF$ remains regular on $\Rd \setminus S$.
Then hypoellipticity of linear differential operators
with constant coefficients are characterized as follows:

\begin{theorem1}\label{XII_THEOREM_ON_HYPOELLIPTIC_OPERATORS}
A linear partial differential operator $\DO_c$ with constant coefficients
is hypoelliptic
\iff
$\,\Sigma(\DO_c \, \disF) \eq \Sigma(\disF)$ for all $\disF \xeof \D'(\Rd)$.
\end{theorem1}

\subparagraph{}
Knowing the fundamental solutions of the elliptic operator $\Delta$ and
the parabolic heat operator $\partial_t-\Delta$,
it turns out that both operators are hypoelliptic.
But theorem \ref{XII_THEOREM_ON_HYPOELLIPTIC_OPERATORS} implies
that the hyperbolic wave operator $\Box \eq \partial_{tt}-\Delta$
is not hypoelliptic.

\section{Free Space Problems}
%------------------------------------------------------------------------------%
To show the benefit of knowing the fundamental solution
of a linear partial differential operator with constant coefficients,
we consider free space problems ($\G \eq \Rd$)
related to the Laplacian, the heat and the wave operator.

\begin{itemize}[leftmargin=12pt,itemsep=5pt]
\item[1.]
The \textit{Poisson equation} in the entire space $\Rd$ is given by
\begin{equation}\label{XII_CLASSICAL_POISSON_EQUATION}
\Delta \, u \= - 4\pi \,f, \quad f \xeof C(\Rd).
\end{equation}

\subparagraph{} % Poisson equation
A function $u \:\Rd \xto \R$ is called a \textit{classical solution}
of \eqref{XII_CLASSICAL_POISSON_EQUATION},
if it is twice continuously differentiable
and fulfills the differential equation in each $x \xeof \Rd$.
However,
$\disU \xeof \D'(\Rd)$ is said to be a \textit{distributional solution}
of the Poisson equation,
if
\begin{equation}\label{XII_DISTRIBUTIONAL_POISSON_EQUATION}
\Delta \, \disU \= - 4\pi \,\disF, \quad \disF \xeof \D'(\Rd)
\end{equation}
holds in the sense of distributions.

\subparagraph{}
Let $\Phi_{\Delta}$ be the fundamental solution of the Laplacian.
If the convolutional product $\Phi_{\Delta} \ast \disF$ exists,
which is especially true for each regular distribution $\disF_f$
with generating function $f \xeof C_c(\Rd)$,
then the solution of \eqref{XII_DISTRIBUTIONAL_POISSON_EQUATION}
is obtained by convolution,
that is,
$\disU \eq \Phi_{\Delta} \ast \disF$.
For the physically interesting case $d=3$
and for $f \xeof C_c(\R^3)$ we get a classical solution $u_f$,
which is denoted the \textit{Newton potential} for the function $f$:
\[
u_f(x) \=
\int_{\,\R^3} \, \frac{f(y)}{\abs{x-y } }  \, dy, \quad x \xeof \R^3.
\]

\item[2.]
By the \textit{heat equation} equation in $\Rd$ is meant
the inital value problem
\begin{equation}\label{XII_CLASSICAL_HEAT_EQUATION}
\begin{cases}[1.2]
\quad u_t(t,x) - \Delta u(t,x) = f(t,x) & (x \xeof \Rd, \, t > 0) \\
\quad u(0,x) \= 0,                      & (x \xeof \Rd). \\
\end{cases}
\end{equation}

\subparagraph{} % heat equation
We first define the concept of distributional solution
for problem \eqref{XII_CLASSICAL_HEAT_EQUATION}:
let $\disU_0 \xeof \D'(\Rd)$ and $\disF \xeof \D'(\R^{d+1})$ are
Schwartz distributions satisfying 
$\supp{\disF} \xsubset [0,\infty) \xtimes \Rd$,
then $\disU \xeof \D'(\R^{d+1})$ is called a \textit{distributional solution},
if $\supp{\disU} \xsubset [0,\infty) \xtimes \Rd$ and
\begin{equation}\label{XII_DISTRIBUTIONAL_HEAT_EQUATION}
\disU_t \, - \Delta \, \disU \= \disF + \disU_0 \otimes \delta_0,
\end{equation}
where $\delta_0 \xeof \D'(\R)$ is the Dirac distribution.
Here the initial condition of the problem has been transformed
into an inhomogeneous part of the differential equation.

\subparagraph{}
The existence of a solution for \eqref{XII_DISTRIBUTIONAL_HEAT_EQUATION}
is ensured,
if $\disU_0 \xeof \E'(\Rd)$ and $\disF \xeof \E'(\R^{d+1})$.
In this case the solution $\disU \xeof \D'(\Rd)$ is given by
convolution of the fundamental equation $\Phi_d$ 
with the inhomogeneous parts of \eqref{XII_DISTRIBUTIONAL_HEAT_EQUATION},
\[
\disU \=
\Phi_d \ast (\disF + \disU_0 \otimes \delta_0) \=
\Phi_d \ast \disF \, + \, \Phi_d \ast \disU_0.
\]

\subparagraph{}
Given a classical solution $u$ of \eqref{XII_CLASSICAL_HEAT_EQUATION},
the related Schwartz distribution $\disF_u$ is a solution
in the distributional sense.
Conversely,
let $\disF_u$ be regular and solve \eqref{XII_DISTRIBUTIONAL_HEAT_EQUATION},
then $u$ forms a classical solution,
if some regularity conditions are valid:

\begin{itemize}[leftmargin=13pt]
\item[a.]
$\disU, \disU_0$ and $\disF$ are regular distributions
with associated functions $u, u_0, f \xeof \Lc$,

\item[b.]
$u_0$ is continuous on $\Rd$ and
$f$ is continuous on $[0,\infty) \xtimes \Rd$,

\item[c.]
$u$ is twice continuously differentiable on $[0,\infty) \xtimes \Rd$.
\end{itemize}

\item[3.]
The initial-value problem for the \textit{wave equation} in $\Rd$
is given by

\begin{equation}\label{XII_CLASSICAL_WAVE_EQUATION}
\begin{cases}[1.2]
\quad u_{tt}(t,x) - \Delta u(t,x) \= f(t,x) & (x \xeof \Rd, \, t > 0) \\
\quad u(0,x) = u_0(x)                       & (x \xeof \Rd) \\
\quad u_t(0,x) = u_1(x)                     & (x \xeof \Rd. \\
\end{cases}
\end{equation}

\subparagraph{} % wave equation
As for the heat equation,
we define for \eqref{XII_CLASSICAL_WAVE_EQUATION}
the concept of distributional solution,
where again both initial value conditions are transformed
into inhomogeneous parts of the equation by means of tensor products.

\subparagraph{}
Let $\disU_0 \xeof \D'(\Rd)$, $\disU_1 \xeof \D'(\Rd)$
and $\disF \xeof \D'(\R^{d+1})$ be distributions,
then $\disU$ is called a distributional solution
of \eqref{XII_CLASSICAL_WAVE_EQUATION},
if
\begin{equation}\label{XII_DISTRIBUTIONAL_WAVE_EQUATION}
\disU_{tt} - \Delta \disU \=
\disF + \disU_0 \otimes \delta_0' \, + \, \disU_1 \otimes \delta_0 \,,
\end{equation}
where $\delta_0$ is the Dirac distribution for $t \eq 0$ 
and $\delta_0'$ is the distributional derivative of $\delta_0$
{\wrt} the variable $t$.

\subparagraph{}
Since the support of the function
$\disU_0 \, \otimes \, \delta_0' \, + \, \disU_1 \, \otimes \, \delta_0$
belongs to the set $\R^{d+1}$,
there is to each distribution $\disF \xeof \D'(\R^{d+1})$ with support
\[
\supp{\disF} \= \{\, (x,t) \xeof \Rd \xtimes \R: t \geq 0 \,\}
\]
a further distribution
$
\disU_F :=
\Phi_{\Box}^d \ast (\disF + \disU_0 \otimes \delta'_0
+ \disU_1 \otimes \delta_0)
\xeof \D'(\R^{d+1}),
$
which in fact solves the distributional differential equation
\eqref{XII_DISTRIBUTIONAL_WAVE_EQUATION}.

\subparagraph{}
Each classical solution of \eqref{XII_CLASSICAL_WAVE_EQUATION}
is also a distributional solution.
However,
the reverse statement only holds under the conditions 

\begin{itemize}[leftmargin=13pt]
\item[a.]
$\disU, \disU_0, \disU_1$ and $\disF$ are regular
and generated by ordinary functions $u, u_0, u_1, f \xeof \Lc$,

\item[b.]
$u_0 \xeof C^1(\Rd)$ and $u_1 \xeof C(\Rd)$,

\item[c.]
$f \xeof C([0,\infty) \xtimes \Rd)$ and

\item[d.]
$u \xeof C^2((0,\infty) \xtimes \Rd) \cap C([0,\infty) \xtimes \Rd)$.
\end{itemize}
\end{itemize}

\section{Fourier Transform Methods}
%------------------------------------------------------------------------------%
The usability of Fourier transformation mainly relies on the fact
that elementary differential operators $\partial_k$,
applied to Schwartz functions $\phi$,
appear as multiplication operators to the Fourier transform $\FT(\phi)$
and via verce,
that is, 
we always have
\[
\partial_k \, \phi(x) \, \overset{\FT}{\longmapsto}
\, i \, \xi_k \cdot \hat{\phi}(\xi)\,,
\quad \quad
i \, x_k \phi(x) \, \overset{\FT}{\longmapsto}
\, \partial_k \, \hat{\phi}(\xi)\,,
\]
for $k=1,\ldots,d$,
as one can verify by partial integration.
For any linear differential operator $\DO_c$
with constant coefficients $c_{\alpha}$
and symbol $S_{\DO_c}$ we accordingly obtain 
\[
(\DO_c \, \phi)(x) \, \overset{\FT}{\longmapsto} \,
S_{\DO_c}(i \xi) \cdot \hat{\phi}(\xi).
\]
This is nothing but the fact
that all the differential operators $\DO_c$ are Fourier multipliers.
They operate on functions $\phi \xeof \SRd$ like multi\-plications
with the symbols $S_{\DO_c}$ on the Fourier trans\-form $\FT \phi$.
Hence a linear differential operator $\DO_c$ with constant coefficients
is conjugated to its symbol function $S_{\DO_c}$ in the sense of
\[
\DO_c \phi \=
(\FT^{-1} \circ S_{\DO_c} \circ \FT) \, \phi \=
\FT^{-1}(S_{\DO_c} \hat{\phi})
\qfa \phi \xeof \SRd.
\]

\paragraph{} %- the general usage of the Fourier transform in PDE theory
We want incident
how Fourier transformation can be used
to solve linear differential equations in the whole space $\Rd$
and governed by the Laplacian.
The basic idea is to transform the differential equation
either into a purely algebraic equation
or into an ordinary differential equation.
Once these are solved,
inverse Fourier transformation is applied
to get the solution of the original problem.

\subparagraph{}
Two examples for these methods will be presented.
In the first one,
the bijectivity of Fourier transformation on the Schwartz space
is used to solve a specific differential equation,
while in the second one the simplified computation of convolution
products after performing Fourier transformation is helpful
to find fundamental solutions of the underlying differential operators.

\begin{itemize}[leftmargin=12pt,itemsep=5pt]
\item[1.]
The equation $\DO_c \, u \eq f$ with $f \xeof \SRd$
and governed by a differential operator $\DO_c$ having constant coefficents
can be converted by applying Fourier transformation
into the algebraic equation
\begin{equation}\label{XII_ALGEBRAIC_EQUATION_OF_LDE}
S_{\DO_c}(i \xi) \, \hat{u}(\xi) \= \hat{f}(\xi)
\quad \quad (\xi \xeof \Rd).
\end{equation}
Now,
if the symbol function $S_{\DO_c}$ of $\DO_c$ owns the property
\begin{equation}\label{XII_SYMBOL_CONDITION}
S_{\DO_c}(i \xi) \neq 0
\qfa \xi = (\xi_1,\ldots,\xi_d) \xeof \Rd,
\end{equation}
equation \eqref{XII_ALGEBRAIC_EQUATION_OF_LDE}
can be resolved for $\hat{u}(\xi)$,
and the solution $u$ of $\DO_c \,u \eq f$
is obtained by applying inverse Fourier transformation:
\[
\hat{u}(\xi) = \hat{f}(\xi) \, / \, S_{\DO_c}(i \xi)
\]
and
\[
u(x) \= \frac{1}{(2 \pi)^{d/2}} \,
\int_{\,\Rd}
e^{i \xi x} \cdot S_{\DO_c}^{-1}(i \xi) \cdot \hat{f}(\xi) \, d \xi.
\]

\item[2.]
If \eqref{XII_SYMBOL_CONDITION} holds,
the Fourier transform method can also be applied
to solve the distributional differential equation
$\DO_c \, \disU \eq \disF$ with $\disF \xeof \TD$.
Since $\DO_c$ is bijective on $\SRd$ and 
\[
P_{\DO_c}(i \xi) \= P_{\DO_c^*}(-\xi)
\qfa \xi \xeof \Rd
\]
holds,
the formally adjoint operator $\DO_c^*$
is bijective on $\SRd$ and so its adjoint operator
$\DO_c \eq (\DO_c^*)^*$ is bijective on the dual space $\TD$.
Thus,
if \eqref{XII_SYMBOL_CONDITION} is valid,
the equation $\DO_c \, \disU \eq \disF$
is uniquely solvable for each tempered distribution $\disF$.

\subparagraph*{}
Especially for the differential operator 
$\Delta{_\lambda} := -\Delta + \lambda$ and $\lambda > 0$ we obtain
\[
S_{\Delta_{\lambda}}(i \xi) = \abs{\xi} ^2 + \lambda \neq 0
\]
for $\xi \xeof \Rd$,
hence the solution $u$ of $\Delta_{\lambda} u \eq f$ with $f \xeof \SRd$ is
\[
u(x) \=
\frac{1}{(2 \pi)^{d/2}} \,
\int_{\,\Rd} e^{i \xi x}
\cdot \frac{\hat{f}(\xi)}{\abs{\xi} ^2 + \lambda} \, d \xi, \quad (x \xeof \Rd).
\]
\end{itemize}

\paragraph{} %- partial Fourier transformation
Partial Fourier transformation is a method
where the operator $\FT$ is applied only to the spatial variables
of the partial differential equation.
The result is an ordinary differential equation
{\wrt} the time variable $t$.
Then solving this equation and applying inverse Fourier transformation
leads to the solution of the original problem.

\paragraph{} % examples of partial Fourier tranformation
In the following we sketch
how {partial Fourier transformation} is useful to solve
the initial-value problem of the heat and wave equation
in the entire space $\Rd$.

\begin{itemize}[leftmargin=12pt,itemsep=5pt]
\item[1.]
The initial-value problem of the heat equation is given by
\[
\begin{cases}[1.2]
\quad u_t - \Delta u = 0   & (t > 0)     \\
\quad u(0,x) \= f(x)       & (x \xeof \R). \\
\end{cases}
\]
Applying Fourier transformation to both sides of the differential equation
{\wrt} the spatial variables $x_1,\ldots,x_d$,
we obtain
\begin{equation}\label{XII_PARTIAL_FOURIER_TRANSFORMATION_1}
\frac{\partial}{\partial t} \hat{u}(t,\xi) \, + \, \xi^2 \, \hat{u}(t,\xi)
\= 0,
\quad \quad \quad
\hat{u}(0,\xi) = \hat{f}(\xi),
\end{equation}
where $\xi \eq (\xi_1,\ldots,\xi_d)$ are the variables
in the \textit{phase space}.

\subparagraph{}
Then multiplying both sides of the first equation
in \eqref{XII_PARTIAL_FOURIER_TRANSFORMATION_1} by $e^{\abs{\xi} ^2 t}$
leads to
\[
\frac{\partial}{\partial t} \, e^{\abs{\xi} ^2 t} \, \hat{u}(t,\xi) \= 0
\]
with solution
$\hat{u}(t,\xi) \eq \phi(\xi) \, e^{- \abs{\xi} ^2 t}$.
The mapping $\xi \mapsto \phi(\xi)$ is
the integration constant for each fixed point $\xi$.
The second equation in \eqref{XII_PARTIAL_FOURIER_TRANSFORMATION_1}
implies $\phi(\xi) \eq \hat{f}(\xi)$,
hence the Fourier trans\-formed solution $\hat{u}$ 
of the initial-value problem is nothing but
\begin{equation}\label{XII_PARTIAL_FOURIER_TRANSFORMATION_2}
\hat{u}(t,\xi) = \hat{f}(\xi) \cdot e^{- \abs{\xi} ^2 t}.
\end{equation}

\subparagraph{}
The solution of the original equation is obtained
by applying inverse Fourier transformation and the rule
\[
\FT^{-1}(f \cdot g) \= \FT^{-1}(f) \ast \FT^{-1}(g)
\]
to each side of \eqref{XII_PARTIAL_FOURIER_TRANSFORMATION_2}:
\[
u(t,x) \=
\FT^{-1}(\hat{f}(\xi) \cdot e^{- \abs{\xi} ^2 t}) \=
\FT^{-1}(\hat{f}(\xi)) \ast \FT^{-1}(e^{- \abs{\xi} ^2 t})
\]
\[
\= f(x) \ast \FT^{-1}(e^{- \abs{\xi} ^2 t})
\=
\frac{1}{\sqrt{4\pi t}}
\int_{\,\R} e^{-\abs{x-y} ^2/(4t)} \,f(y)\,dy.
\]
The function
\[
H_t(x,y) \, := \,
\frac{1}{\sqrt{4\pi t}} \, e^{- \abs{x-y} ^2/(4t)}
\=
\Phi_{\partial_t - \Delta}(t,x-y)
\]
is called a \textit{heat kernel},
since it constitutes the kernel of the integral operator
$\cal{H}_t:f \mapsto H_t \ast f$.

\item[2.]
The method also enables to solve the initial-value problem
for the wave equation
\[
\begin{cases}[1.2]
\quad u_{tt} - \, \Delta u \= 0  & (t > 0),    \\
\quad u(0,x) \= f(x)             & (x \xeof \Rd),\\
\quad u_t(0,x) \= g(x)           & (x \xeof \Rd).\\
\end{cases}
\]

\subparagraph{}
Partial Fourier transformation let become this equation into
\[
\hat{u}_{tt}(t,\xi) \= \abs{\xi} ^2 \hat{u}(t,\xi),
\quad \quad
\hat{u}(0,\xi) \= \hat{f}(\xi),
\quad \quad
\hat{u_t}(0,\xi) \= \hat{g}(\xi),
\]
and the general solution of the ordinary differential equation
in phase space is 
\[
\hat{u}(t,\xi)
\=
\hat{f}(\xi) \, \cos( t \abs{\xi} )
\, + \,
\hat{g}(\xi) \: \frac{\sin(t \abs{\xi} )}{\abs{\xi} }.
\]

\subparagraph{}
Unfortunately,
performing inverse Fourier transformation to $\hat{u}$
for any $d >1$ is not an easy task.
But for $d=1$ we obtain
\[
u_1(x,t) \=
\FT^{-1}\left(\hat{f}(\xi) \cdot \cos(\abs{\xi} ) \right)
\=
\frac{1}{2}\,[f(x+t) \, + \, f(x-t)],
\]
\[
u_2(x,t) \=
\FT^{-1}\left(\hat{g}(\xi) \cdot
\frac{\sin(\abs{\xi} )}{\abs{\xi} } \right)
\=
\frac{1}{2}\, \int_{-t}^{+t} g(x+s) \, ds,
\]
thus the solution of the wave equation is given by $u \eq u_1+u_2$,
known as the \textsc{d'Alembert} \textit{solution}
\[
u(t,x) \=
\frac{1}{2} \, [f(x+t) \, + \, f(x-t)]
\, + \,
\frac{1}{2t}\, \int_{-t}^{+t} g(x+s) \, ds.
\]
\end{itemize}

% BIBLIOGRAPHY
%------------------------------------------------------------------------------%
\vspace{5mm}
\begin{thebibliography}{99}
\begin{small}

\bibitem{XII_Br2} F.\,Browder:
\textit{On the Spectral Theory of Elliptic Differential Operators I}.
Math. Annalen 142, pp.\,22-130, 1961.

\bibitem{XII_GT}  D.\,Gilbarg, N.S.\,Trudinger:
\textit{Elliptic Partial Differential Operators of Second Order}.
Springer, 2001.

\bibitem{XII_Pa} A.\,Pazy:
\textit{Semigroups of Linear Operators
        and Applications to Partial Differential Equations}.
Applied Mathematical Sciences Vol.\,44, Springer, 1983.

\bibitem{XII_RR} M.\,Renardy, R.\,Rogers:
\textit{An Introduction to Partial Differential Equations}.
Texts in Applied Mathematics 13, Springer, 2004.

\bibitem{XII_Yo} K.\,Yosida:
\textit{Functional Analysis}.
Springer, 1980.

\end{small}
\end{thebibliography}

%------------------------------------------------------------------------------%
%                               END OF DOCUMENT                                %
%------------------------------------------------------------------------------%
\end{document}
